% 本文件由 LaTeX 排版 Agent 根据有效 translation.json 自动生成。
% author=John Chrysostom; work_id=1901; title=雕像讲道集

\chapter{雕像讲道集}

% structure: p0001 source=h1 target=omitted reason=page-title-already-rendered-by-parent

\Emph{讲道一}\newline{} \Emph{讲道二}\newline{} \Emph{讲道三}\newline{} \Emph{讲道四}\newline{} \Emph{讲道五}\newline{} \Emph{讲道六}\newline{} \Emph{讲道七}\newline{} \Emph{讲道八}\newline{} \Emph{讲道九}\newline{} \Emph{讲道十}\newline{} \Emph{讲道十一}\newline{} \Emph{讲道十二}\newline{} \Emph{讲道十三}\newline{} \Emph{讲道十四}\newline{} \Emph{讲道十五}\newline{} \Emph{讲道十六}\newline{} \Emph{讲道十七}\newline{} \Emph{讲道十八}\newline{} \Emph{讲道十九}\newline{} \Emph{讲道二十}\newline{} \Emph{讲道二十一}\par

\SourceRecord{Translated by W.R.W. Stephens.；Nicene and Post-Nicene Fathers, First Series；Vol. 9.；Edited by Philip Schaff.；Buffalo, NY: Christian Literature Publishing Co.,；1889.；<http://www.newadvent.org/fathers/1901.htm>.}

% structure: p0001 source=h1 target=section reason=page-title

\section{论雕像训导第一讲}

\Emph{这篇讲道是在安提约基雅的旧教堂中宣讲的，当时\Person{圣金口若望}还是一位司铎，论及宗徒所说的\BibleRef{弟前 5:23}：\BibleQuote{为了你的胃和你的屡次不适，略用点酒。}}\par

1. 你们听见了宗徒的声音，那来自天上的号角，那属神的竖琴！因为正如号角吹响可怕而战争的音符，既惊吓敌人，又振奋己方沮丧的灵，并用极大的胆量充满他们，使那些留意它的人对魔鬼不可战胜！同样，作为竖琴，它以迷人心魄的旋律温柔地抚慰，使邪念的纷扰安息；并藉着愉悦，将许多益处注入我们内。你们今天听见了宗徒对\Person{弟茂德}谈论各种必要的事！因为他就按手之事写信给他，说：\BibleQuote{不可轻易给人按手，也不可在别人的罪上有分。} \BibleRef{弟前 5:22} 他解释了这种过犯的可怕危险，指出如此行的人将与他人共同承受他人所犯之罪的惩罚，因为他们藉着按手将权能赋予他们的邪恶。紧接着他又说：\BibleQuote{为了你的胃和你的屡次不适，略用点酒。} 今天他又向我们谈论了仆人的服从、守财奴的疯狂、富人的傲慢，以及其他各种事项。\par

2. 既然无法逐一详述，你们愿意我们选择所诵读的部分中的哪一部分作为我们向你们爱德讲话的主题？因为如同在草地中，我在所读的内容中看到许多不同的花朵；大量的玫瑰和紫罗兰，以及不少的百合；到处散布着圣神各种丰盛的果实和浓郁的芬芳。不，不如说，读圣经不仅是草地，而且是乐园；因为这里的花朵不仅芬芳，而且有滋养灵魂的果实。那么，你们希望我们今天把所诵读的哪一部分呈献给你们呢？你们是否希望那看似更微不足道、任何人都容易理解的部分，成为我们现在所处理的？对我来说，这似乎合适，而且我相信你们会同意这个看法。那么，这看似比任何事都更简单的是什么？不就是那似乎很容易、任何人都会说的吗？嗯！那是什么？\BibleQuote{为了你的胃和你的屡次不适，略用点酒。} 那么，让我们将全部讲话用在这个主题上；我们这样做，并非为了爱赞美，也不是因为我们研究展现口才的能力（因为即将说的事不是我们的，而是圣神的恩宠所启示的）；而是为了激励那些过于懈怠的听众，并使他们相信圣经宝藏的伟大；以及匆匆浏览它们既不安全，也非没有危险。因为如果一个如此简单明了的经文，正如这一句，在众人看来似乎没有什么值得坚持的，却为我们提供了丰富财富的手段和通向最高智慧的门径，那么，其他的经文，那些立即显出其本有财富的，更会用其无尽的宝藏来满足那些留意它们的人。因此，我们确实不应匆匆掠过那些被认为浅显的圣经句子；因为它们也出自圣神的恩宠；但这恩宠从不是微小的、低劣的，而是伟大、奇妙而且配得上赐予者的慷慨。\par

3. 因此，让我们不要粗心地听；因为就连那些烘焙矿土的人，当他们将矿土投入熔炉时，不仅拿起金块，也极其小心地收集细小的颗粒。既然我们也同样要烘焙从宗徒矿藏中提取的金子，不是将它投入熔炉，而是将它存于你们灵魂的思想中；不是点燃地上的火焰，而是燃起圣神的火，让我们勤奋地收集这些小颗粒。因为这句谚语虽简短，但其德能却伟大。因为珍珠也有其固有的市场，不是由于物质的尺寸，而是其本性的美丽。同样，读圣经也是如此；因为世俗的学问大量地传播其琐碎之事，并以虚妄的闲聊之流淹没其听众，却让他们空手而归，没有收集到任何或大或小的益处。然而，圣神的恩宠并非如此，相反，藉着简短的句子，它将神圣的智慧植入所有留心的人中，并且一句经文常常为接受它的人提供足够的人生旅途储备。\par

4. 既然它的财富如此巨大，让我们自我警醒，以清醒的头脑接受所说的；因为我准备将我们的讨论潜入极深的深处。这劝告本身无疑被认为无关主题且多余，许多人这样谈论：\Quote{\Person{弟茂德}难道不能自己判断需要使用什么，而要等他的老师来教导他？而且老师不仅给指示，还写在信里，像刻在铜柱上一样刻在他的书信里？当他以公开方式写给他的门徒时，指导这类性质的事，难道他不感到羞耻吗？} 为此，为使你们明白，这劝告非但不过分，反而是必要且极为有益的；并且这事不是出于\Person{保禄}，而是出于圣神的恩宠，即这（我说）不是一条口头的诫命，而是一条写在信里，并藉着书信传给后代所有的人的，我将立即进行证明。\par

5. 因为除了已经提及的课题之外，还有另一个课题，使一些人同样感到困惑，他们内心自问：天主为何容许一个对祂怀有如此信心、其骨骸和遗物能驱逐魔鬼的人，陷入如此虚弱的状况？这不只是因为他在患病，而是持续地、长久地患病；通过这些反复而持久的虚弱，他甚至得不到片刻的喘息。有人会问：\Quote{这如何看得出来呢？}从\Person{保禄}本人的话中就可看出，因为他没有说\Quote{因那疾病}，而是说\Quote{因那些疾病}；不仅说\Quote{疾病}，而且明确提到这些疾病是持续的，当他说\Quote{你屡次的疾病}。因此，无论是什么人，若因长期患病而怨声载道、意志消沉，都当留意这一点。\par

6. 但议题不仅是：作为一个圣人，他患病，且如此持续地患病；同时他还受托管理世界公共事务。因为假如他是那些隐居山巅、独居修院、选择远离一切事务生活的人之一，那么现在探讨的问题便不会有如此困难；但一个被推入人群之中、手中掌管着如此众多教会、监督整个城市和民族、甚至整个世界，并以如此热忱和勤勉的人，却要受疾病之困！这才是最令人困惑之处，若人不恰当地思考此事。因为，即便不是为了自己，至少为了他人，他理应健康。反对者说：\Quote{他是最好的将领。战争不仅由他对抗不信者，而且对抗魔鬼，对抗\Person{撒殚}本身。所有敌人都极为猛烈地争战，分散军队，俘虏囚犯\BibleRef{弟后 2:26}；但这个人却能将无数人带回真理，然而他却生病了！因为，即使这疾病对事业没有其他损害，仅此一点就足以使信徒灰心丧气。如果士兵看到他们的将军卧病在床，就会灰心松懈，更何况当信徒看到那位行了许多奇迹的导师，身体持续患病和受苦时，他们岂不更可能表现出人性软弱？}\par

7. 但不仅如此。这些怀疑论者还提出进一步的问题，问为何\Person{弟茂德}既没有医治自己，也没有被他的导师\Person{保禄}在他处于这种状态时医治。宗徒们可以复活死人、驱逐魔鬼、轻易战胜死亡，却连一个病人的身体都不能复原！虽然对其他人的身体，无论是在他们生前还是死后，他们都显示了如此非凡的能力，却不能恢复一个失去活力的胃！更有甚者，\Person{保禄}在行了那么多伟大的奇迹之后，并不以此为羞耻，也不脸红，在写给\Person{弟茂德}时，却命他求助于饮酒的疗效。不是说饮酒是可耻的。千万不要！因为这样的规则属于异端者；但令人惊奇的是，他认为没有这种帮助就无法修复一个失常的肢体，并不以为耻。然而，他非但不以此为耻，反而将此显明给后世。你看，我们已将主题探讨到如此深度，这个看似微小的问题，却充满了无数疑问。好，让我们进行解答；因为我们如此深入地探讨问题，是为了在吸引你们的注意之后，将解释安全地保存起来。\par

8. 但在解决这些问题之前，请允许我说一说弟茂德的品德和保禄的爱顾。有谁比他更心肠仁慈呢？他身在远方，为诸多挂虑所环绕，却如此关心他弟子胃口的健康，并仔细地写信纠正他的失调。又有谁能比得上弟茂德的品德？他如此轻视奢华，嘲笑丰盛的筵席，以致因过度的刻苦和严苛的禁食而生病。其实他并非天生如此虚弱，而是因禁食和只饮水而损害了胃的力气。你可以听保禄本人小心地阐明这一点。因为他并非简单地说\Quote{少用点酒}，而是在先说了\Quote{不要再只饮水}之后，才提出关于饮酒的劝告。而这\Quote{不要再}一词明确证明，他直到那时一直饮水，因此才变得虚弱。那么，谁不惊叹他的神圣智慧和严谨呢？他抓住了诸天，跃上了德行的最高点。他的老师为此作证说：\Quote{我打发弟茂德到你们那里去，他在主内是我所爱和忠信的儿子}\BibleRef{格前 4:17}；保禄称他为\Quote{儿子}、\Quote{忠信可爱的儿子}，这话足以表明他拥有各种美德。因为圣人们的判断不凭私情或敌意，而是摆脱一切偏见。弟茂德若按血肉是保禄的儿子，也不会如此令人羡慕；而他现在令人敬佩，是因为他虽与保禄无血肉关系，却藉着虔敬的亲情引入了宗徒的义子名分，在一切事上精确地保存了他属神智慧的印记。因为如同一只小牛犊与公牛拴在一起，他也与保禄同负一轭，无论他到世界何处，都不因年轻而减轻负担，反而以热忱之心效仿老师的劳作。保禄本人再次为此作证说：\Quote{不要让人轻视他，因为他如同我一样，做主的工作}\BibleRef{格前 16:10}。你看见他如何见证弟茂德的热忱与他自己完全一致吗？\par

9. 此外，为了不让人认为他说这些话是出于偏爱或客气，他让听众亲自为儿子的美德作证，说：\Quote{你们知道他的贤德，他怎样如同儿子侍奉父亲一样，与我一同为福音服务}\BibleRef{斐 2:22}；也就是说，\Quote{你们已经体验过他的德行和他那被认可的品格。}然而，尽管他达到了这样的善行高度，却并未因此得意自满，反而充满焦虑和恐惧。因此他也严格禁食，不像许多人那样，当他们保持十个月或二十个月的刻苦后，便立刻完全放弃。我所说的这位弟茂德绝无此种心态，也没有对自己说这样的话：\Quote{我还需要禁食做什么？我已克制了自己，战胜了私欲，制伏了身体，吓退了魔鬼，驱逐了恶魔，复活了死人，洁净了癞病人，成了敌对势力的可畏者；我还需要禁食做什么？还须从那里寻求安全吗？}他没有说，也没有想任何这类的话。相反，随着他积累的善行越来越多，他越加恐惧战栗。他从他的导师那里学到了这种属神的智慧。因为就连保禄，在他被提升到第三层天、被提到乐园里、听到不可言传的话、参与如此奥秘、并如带翼者般走遍普世之后，当他写信给格林多人时却说：\Quote{我恐怕我传报给别人，自己反被淘汰}\BibleRef{格前 9:27}。如果保禄在那么多显著善行之后仍然害怕——他能说：\Quote{世界于我已被钉在十字架上，我于世界亦然}\BibleRef{迦 6:14}——那么，我们更应当害怕，而且越是积累了众多善行，就越当如此。因为那时魔鬼变得更凶猛，更残暴，当他看见我们谨慎地规整生活时！当他看见德行的货品已积聚，负载变得沉重时，他便急于制造更惨重的沉船！因为那个微不足道、卑贱的人，即使失足跌倒，也不会给共同的善业带来如此大的损害。但那个如同站在德行的某高点上明显可见、为众人所知、为众人所钦佩的人，若受袭击而跌倒，便会造成巨大的毁灭和损失。这不仅是因为他从这高处坠落，而且使许多仰望他的人变得更为懈怠。如同在身体上，某个肢体毁坏了不一定造成大损伤，但若眼睛失明或头部受重伤，整个身体便无用了；同样，对于圣人和那些已行最高善行者，我们也必须说：当他们被熄灭、染上任何污点时，他们给整个身体的其余部分带来普遍的、难以忍受的伤害！\par

10. \Person{弟茂德}既知此事，就处处加强自己；因为他知道青年时期是困难的年纪，不稳定，易受骗，极易滑倒，需要极强的缰绳。它确实像一堆易燃物，容易从外界引火，很快燃烧；因此他注意处处压制它，尽力减弱火焰。那难以驾驭和倔强的马，他用力勒住，直到驯服它的任性；直到使它温顺；并完全控制，交到作为御者的理性手中。\Quote{让身体软弱，}他说，\Quote{但不要让灵魂软弱；让肉体受约束，但不要让精神向天奔驰的赛跑被阻止。}此外，人们尤其会为此惊奇：此人虽如此患病，与这样的虚弱挣扎，却没有对天主的事业漠不关心，反而比那些身体强健的人飞得更快；有时去\Place{厄弗所}，有时去\Place{格林多}，常去\Place{马其顿}和\Place{意大利}；在陆地和海上，与老师一起出现在各处，在一切事上分享他的挣扎和不断的危险；而他灵魂的神圣智慧并未因身体的软弱而蒙羞。对天主的热情就是这样！它赋予翅膀如此轻盈！因为正如那些体质良好健康的人，如果灵魂卑怯、懒惰、愚蠢，力量也无用；同样，那些极度虚弱的人，如果灵魂高贵而清醒，软弱不会造成伤害。\par

11. 然而，这劝告和忠告，似乎给一些人放纵饮酒的权威。但并非如此。如果我们仔细考察这句话，它反而相当于推荐节制。因为请想一想：\Person{保禄}起初并没有一开始就给出这劝告。但当他看到所有力量都倒下时，才给出了劝告；即使那时也不是简单地，而是有一定的先决限制。他不只是说\Quote{用酒}，而是\Quote{少许}酒；不是因为\Person{弟茂德}需要这劝告和建议，而是因为我们需要。为此，在给他写信时，他为我们规定了饮酒的量和限度；吩咐他饮到刚好纠正紊乱的程度；使身体恢复健康，而不引发另一种疾病。因为过量饮酒产生的身体和灵魂的疾病并不少于大量饮水，反而更多、更严重；它给心灵带来情欲的战争和邪念的风暴，此外还使身体的坚实体质变得松弛软弱。因为长期受过多水分干扰的土地，其分解程度还不如人体力量因不断酗酒而削弱、松弛、疲惫的程度。因此，让我们谨慎避免任何一方过度，同时我们照顾身体的健康，并剪除其奢侈的倾向。因为天主赐给我们酒，不是要我们醉酒，而是要我们清醒；要我们喜乐，而不是自找痛苦。\BibleQuote{酒，}它说，\BibleQuote{使人的心欢喜，}但你们却使它成为悲伤之事；因为醉酒的人极度阴沉，巨大的黑暗笼罩他们的思想。当它有最好的节制引导时，它是最好的药物。眼前的经文也对那些诽谤天主造物的异端者有用；因为如果酒是被禁止的东西，\Person{保禄}就不会允许它，也不会说它应该被使用。不仅针对异端者，也针对我们弟兄中的愚钝者，他们看到有人因醉酒而丢脸时，不责备这些人，反而责备天主赐予的果实，说：\Quote{不要有酒。}那么我们应该这样回答他们：\Quote{不要有醉酒；因为酒是天主的工程，但醉酒是魔鬼的工程。酒不制造醉酒，而是无节制产生醉酒。不要控告天主的造物，而应控告同类的疯狂。但你，在省略责备和纠正罪人的同时，竟藐视你的恩主！}\par

12. 因此，当我们听到人们说这样的话时，我们应该堵住他们的嘴；因为产生醉酒的并非饮酒，而是缺乏节制。醉酒！一切邪恶的根源。酒被赐予是为了恢复身体的软弱，而不是推翻灵魂的力量；为了除去肉体的疾病，而不是摧毁精神健康。那么，不要因过度使用天主的恩赐而给愚昧无耻之人以把柄。因为有什么比醉酒更可悲呢！醉酒者是活着的尸体。醉酒是自选的魔鬼，无法原谅的疾病，不容辩解的堕落，人类共同的耻辱。醉酒者不仅在集会中无用，不仅在公共和私人事务中无用；而且光是看见他就是最令人厌恶的事，他的呼吸是臭气。醉汉的打嗝、打哈欠和言语，既令人不快又冒犯，完全被看见和与之交谈的人所厌恶；而这些邪恶的冠冕是，这种疾病使天堂对醉汉不可及，不让他们赢得永福：因为除了今生遭受此病之人的羞耻，还有严重的惩罚在等待他们！那么，让我们戒除这邪恶的习惯，并听\Person{保禄}说：\Quote{用少许酒。}因为即使是这少许，他也是因他的软弱而允许的；所以，如果软弱没有困扰他，他也不会强迫他的门徒允许自己即使少量，因为我们应该总是根据情况和需要来使用必需的食物和饮料，这些是赐给我们的；决不能超出我们的需要，也不要做任何无意义和无目的的事。\par

13. 既然我们现在已经知道了保禄的温柔关怀和弟茂德的德行，那么，接下来，让我们将讲论转向对这些问题的实际解决。这些问题是什么呢？有必要再次提及它们，以便解决之道更加清晰。那么，天主为何允许这样一位圣人，一位受托管理如此多事务的人，陷入疾病的状态？而且，无论是弟茂德本人还是他的老师，都没有力量纠正这种失调，反而需要藉饮酒来获得帮助？的确，这就是所提出的问题。但是，有必要提出一个精确的解决之道；这样，如果有人不仅陷入类似的疾病和不适，而且陷入贫困、饥饿、锁链、折磨、挫败、诽谤，以及所有属于现世生命的邪恶，即使他们是伟大而奇妙的圣人，你们仍然能够从今天将要陈述的内容中，为他们的情形找到对那些喜欢挑剔之人的精确而清晰的答复。因为你们曾听到许多人提出这样的问题，例如：\Quote{为什么某个温和而谦逊的人，竟被另一个无法无天、邪恶的人每日拖到审判席前，遭受无数恶事，而天主却允许此事？} \Quote{又为何另一个人因诬告而被不公正地处死？} 反对者说：\Quote{某人被淹死了；另一个人被推下悬崖；我们可以举出许多圣人，无论是在我们自己的时代，还是在我们祖先的时代，他们遭受了各种变化多端的磨难。} 因此，为了让我们看到这些事情的原因，也为了我们自己不受困扰，也不忽视那些因此遇到绊脚石的其他人的情况，我们应该认真留意即将提出的理由。\par

14. 因为关于圣人所遭遇的多样而多变的苦难，我要向你们的爱德宣告八种理由。因此，请所有人以最严格的注意力转向我，要知道，如果经过这一切，尽管有如此之多的理由，我们仍然和找不到一个理由时一样困惑和不安，那么我们将不再有原谅或借口来为所发生的事绊倒。\par

第一个理由是天主允许他们遭受恶事，以免他们因伟大善工和奇迹而过于轻易地陷入骄傲自满。\par

第二个理由，是免得别人对他们有超出人性的过高评价，把他们当作神而不是人。\par

第三个理由，是使天主的权能得以显明，即在软弱和受捆绑的人身上，天主的权能得以施展、得胜并推进所宣讲的圣言。\par

第四个理由，是使这些圣人自身的忍耐更加显著，因为他们事奉天主，不是为了报酬；反而在遭受如此多恶事之后，仍然显明对天主不变的好意。\par

第五个理由，是使我们的心智对复活的道理变得明智。因为当你看到一个正义的人，一个充满德行的人，遭受成千上万的恶事，并这样离开现世生命时，你完全被迫（即使不情愿）去思考未来的审判；因为如果人们不允许那些为自己劳苦的人没有工资和报酬就离去，那么天主更不能设计让那些如此劳苦的人空手离开。但如果祂最终不能意图剥夺那些人的劳苦报酬，那么必然有一个时间，在今生结束之后，他们将在那里接受他们现世劳苦的报偿。\par

第六个理由，是使所有陷入逆境的人，通过观看这些人，并记住他们所受的苦难，得到充分的安慰和缓解。\par

第七个理由，是当我们劝勉你们效法这些人的德行，并对你们每个人说：\Quote{效法保禄，效法伯多禄}时，你们不会因为他们善工的卓越性而懒散地退缩，不进行这样的效法，认为他们分享了不同的本性。\par

第八个理由，是当有必要称任何人有福或相反时，我们可以学习该把谁视为有福，把谁视为不幸和可怜。\par

这些就是理由；但有必要从圣经中确立这一切，并精确地表明所有关于这个主题的陈述并非人类推理的发明，而是圣经的判决。因为这样，我们所说的将更可信，并更深地印在你们心中。\par

15. 那时，患难对圣人们是有益的，好让他们操练节制和谦卑，免得因自己的奇迹和善工而自高自大；天主允许患难正是为了这个目的；我们可以听见\Person{达味}先知和\Person{保禄}说出同样的话。前者说：\BibleQuote{上主，为受苦，正是于我有益，好使我能学习你的规矩。}后者在说了\BibleQuote{我被提到第三层天上，}并被接到乐园之后，接着说：\BibleQuote{并且免得我因那高超的启示而过于高傲，就在身体上给了我一根刺，就是撒旦的使者来拳击我。}有什么比这更清楚呢？他说：\BibleQuote{免得我过于高傲，}为此缘故，天主允许\BibleQuote{撒旦的使者来拳击我；}所谓撒旦的使者，他并非指特定的魔鬼，而是指为魔鬼服役的人，即不信者、暴君、外邦人，他们坚持不懈地骚扰他，无休止地折磨他。他所说的正是这样：\Quote{天主本能够压制这些迫害和接连的患难；但既然我被提到第三层天上，并被接到乐园，免得我因这些高超的启示而自高自大、看重自己，他就允许这些迫害，任凭这些撒旦的使者以迫害和痛苦来拳击我，使我不要过分自高。}因为\Person{保禄}、\Person{伯多禄}以及所有像他们一样的人，虽然是圣洁而了不起的人，正如他们确实是，但他们终究是人，需要十分谨慎，免得轻易自高自大；而且他们比其他人更需谨慎。因为没有什么比充满善工的良心和充满信心的灵魂更容易使人自大。因此，为使这些人不受这类事的害，天主允许有试探和患难；这些试探和患难大有能力使他们谦抑，并劝他们在一切事上操练节制。\par

16. 这一点也有助于彰显天主的能力，你们甚至可以从同一位宗徒那里学到，他告诉了我们前面的事。免得你们说（正如不信者所想的那样）：天主在允许这事时，是个软弱的存在，容许这样的人不断受折磨，是因为不能拯救自己的人脱离危险：我说，正是这件事，请观察\Person{保禄}如何通过这些事件来说明，不仅表明这些事件远非指责祂软弱，反而更引人注目地向所有人证明祂的能力。因为他说了：\BibleQuote{在身体上给了我一根刺，就是撒旦的使者来拳击我，}并这样表明了他反复的试探，接着又补充说：\BibleQuote{为了这事，我三次求主使它离开我；祂对我说：我的恩宠够你用的，因为我的能力在软弱中才得以圆满。} \BibleRef{格后 12:8} 他说：\Quote{我的能力}意思就是\Quote{在你软弱时才能看见；然而藉着你这看似软弱的人，所传的道被显扬，并被播撒到四方。}因此，当他被带到监狱，在受了大量鞭打之后，他反而囚禁了狱警。\BibleRef{宗 16:24} 他的脚上了木狗，手上了锁链；而半夜他们在唱赞美诗时，监狱震动。你们看，祂的能力如何是在软弱中得以圆满？如果\Person{保禄}自由自在，并且震动了那建筑，事情就不会如此奇妙。他说：\Quote{为此，}\Quote{你被捆绑；墙壁将四面震动，囚犯将被释放；好让我的能力显得更大，当你被监禁、带着锁链时，所有被捆绑的人都得以释放。}正是这个情况当时使狱警惊愕：他如此被强力拘禁，却仅仅藉着祈祷，就震动了地基，打开了监狱的门，释放了所有囚犯。这并非唯一的事例。在\Person{伯多禄}和\Person{保禄}自己身上，以及在别的门徒身上，我们也可以不断看到这样的事；在迫害中，天主的恩宠始终茂盛，与患难相伴出现，从而宣扬祂的能力。因此祂说：\BibleQuote{我的恩宠够你用的，因为我的能力在软弱中才得以圆满。}\par

17. 但为表明许多人若非看见他们承受如此苦难，就太容易将超乎人性的事想象在他们身上，请听保禄在这一方面是如何畏惧的：\BibleQuote{因为我即使愿意夸耀，也不算是狂妄，但我不这样做，是怕有人把我看得超过他所看见我的，或从我所听到的。} \BibleRef{格后 12:6} 可他所指的是什么？他说，我能谈论更大的奇迹，但我不愿意，恐怕奇迹的宏大在人心中对我产生过高的想法。为此，伯多禄在治好瘸子后，当众人都惊奇他们时，为遏止民众，并说服他们明白他们并非凭自己的力量或本性的能力行出这事，便说：\BibleQuote{你们为什么注视我们，好像凭我们自己的能力和虔诚使这人行走？} \BibleRef{宗 3:12} 又在吕斯特辣，众人不仅惊奇，还牵出公牛，戴上花冠，准备向保禄和巴尔纳伯献祭。请看魔鬼的恶意：主藉以清除世间不敬虔的同一批人，那仇敌却试图藉他们引入不敬虔，再次说服人将人当作神——这正是他在古时所行的。而这尤其引入了偶像崇拜的原则和根源。因为许多人在战争中获胜、树立凯旋柱、建造城市、并给当时的人们带来各种此类益处后，就被民众尊为神，受到庙宇和祭台的尊崇；希腊诸神的全部名单就是由这样的人组成。因此，为使这情况不至于发生在圣人身上，天主许可他们不断被放逐、被鞭打、患病；好让身体的软弱和试炼的众多，能使当时与他们在一起的人确信：行这奇事的是人，他们并未贡献自己的任何能力；而是纯粹的恩宠藉他们行这一切奇迹。因为倘若人们将行过卑贱之事的人当作神，那么当他们行从未有人见过或听过的奇迹时，若他们未曾经历任何属于人性的事，人们就更容易将他们当作神了。因为当他们被鞭打、被推下悬崖、被监禁、被放逐、每日身处危险时，尚且有人陷入这种不敬虔的想法，那么若他们未曾忍受任何属于人性的事，人们岂不更会如此看待他们？\par

18. 这就是受苦的第三个原因；第四个原因是，圣人可能不被认为是因为盼望现世的顺遂而事奉天主。因为许多生活在放荡中的人，当他们常被众人责备，并被邀请去行德行的劳苦时；当他们听到圣人在巨大患难中仍保持喜乐而受称赞时，他们便以此为由攻击圣人的品格；不仅人如此，连魔鬼自己也起了这样的怀疑。因为当约伯拥有巨大财富、享受富足时，那邪恶的魔鬼因约伯的缘故被天主责备，无话可说；当它既不能回答对自己的控告，也不能指责这义人的德行时，它立刻采取这个辩护，这样说：\BibleQuote{约伯敬畏天主岂是无益的吗？你不是四面围上篱笆保护他吗？}\BibleRef{约 1:9}\Quote{那么，}它说，\Quote{那人行善是为了报偿，于是他享受如此多的富足。}那么天主做了什么？祂要表明，祂的圣人侍奉祂并非为了报偿，便剥夺了他所有的富足，将他交给贫穷，并允许他陷入重病。后来祂责备约伯毫无根据地怀疑时，说：\Quote{他仍然持守自己的纯正；你虽激动我攻击他，无故地毁灭他，这也是徒然。}因为对圣人来说，侍奉天主本身就是足够的报偿和补偿；因为对于爱者来说，爱他所爱的对象就是足够的报偿；此外他别无他求，也不以任何事为大于此。如果人对人尚且如此，何况对天主呢？因此，天主为了显明这一点，便赐下了比魔鬼所求的更多；因为魔鬼说：\BibleQuote{伸手击打他的骨和肉，}\BibleRef{约 2:5}但天主并非这样说，而是：\Quote{我把他交在你手里。}因为正如在外面的竞赛中，强壮且身体状态良好的斗士，当他们被浸油的衣服紧紧包裹时，不易被看清；但一旦脱下这衣服，赤裸着被带进竞技场，那时他们肢体的匀称最令四周的观众惊叹，因为再没有什么遮掩他们了；约伯也是如此。当他被所有财富包裹时，众人看不出他是怎样的人。但当他像摔跤手脱下衣服一样抛弃财富，赤裸地来到虔诚的搏斗中时，这样赤身露体，他令所有看见他的人惊讶；\BibleRef{约 1:21}以至天使的剧场看着他的灵魂坚韧而欢呼，为他加冕而鼓掌！因为正如我所说，当他穿着所有财富时，众人并不如他像脱去衣服一样赤裸地展示在剧场中那样清楚地看见他，那时全世界都赞叹他灵魂的力量，这不仅由他被剥夺一切所证明，也由他的争战和他对病痛的忍耐所证明。如我先前所说，天主并没有亲自击打他，免得魔鬼再次说：\Quote{你宽恕了他，并没有施加必要的考验；}而是将牲畜的毁灭和肉身的权柄交给了那仇敌。\Quote{我确信，}祂说，\Quote{这位摔跤手；所以我不禁止你给他加上任何你想要的挣扎。}但正如那些精通摔跤场技巧、并有理信赖自己的技艺和体力的人，往往不正面抓住对手，也不取得平等优势，而是允许对手拦腰抱住自己，以便取得更辉煌的胜利；同样，天主让魔鬼拦腰抱住这位圣人，使他能在如此不利的攻击下得胜，将对手摔倒，从而他的冠冕更加荣耀！\par

19. 这是经过考验的金子！随你愿考验它，随你愿审视它，你必找不到任何渣滓。这不仅向我们显示了他人的坚韧，也带来了更多的安慰；因为基督说：\Quote{几时人为了我而辱骂迫害你们，捏造一切坏话毁谤你们，你们是有福的。你们欢喜踊跃罢！因为你们在天上的赏报是丰厚的，因为你们以前的先知，人也曾这样迫害过他们。}保禄在致马其顿人的书信中安慰他们说：\BibleQuote{弟兄们！你们曾效法了天主在基督耶稣内的各教会，因为你们也由同乡的同胞受了苦，正如他们由犹太人受了苦一样。}\BibleRef{得前 2:14}他又同样地安慰希伯来人，列举了所有曾生活在火窑、坑穴、旷野、山谷、山洞中，忍饥受穷的义人。\BibleRef{希 11:34}因为同受苦楚给受苦的人带来一些安慰。\par

20. 但这同样为复活提出论证，请听同一\Person{保禄}又说：\Quote{我若按人的方式在\Place{厄弗所}与野兽搏斗，对我有何益处呢？如果死人不复活。}\BibleRef{格前 15:32}\newline{}又说：\Quote{我们只在今生有希望，就比众人更可怜。}\BibleRef{格前 16:19}\newline{}他告诉我们，我们在今生遭受无数灾祸；如果没有任何来生可盼望，还有什么比我们的状况更悲惨呢？因此，显然我们的境遇并非局限于当前状态；这从我们的考验中便可显明。因为天主绝不容许那些忍受了如此多且大的灾祸、在无数考验和危险中度尽今生的人，得不到远为更丰厚的赏报；如果他不能容忍这点，那么他必定预备了另一个更美好更光明的生命，在其中他要加冕那些为敬虔而奋斗的人，并在全世界面前表扬他们。因此，当你看到一个义人受困苦、遭患难，在疾病、贫穷以及无数其他灾祸中度过此生时，你要对自己说：如果没有复活和审判，天主就不会允许一个为他忍受如此大患难的人，在未享受任何好处之前就离世；由此明显可见，天主为这样的人预备了另一个更甜美、更可忍受的生命。因为若非如此，他就会容忍许多恶人在今生享乐，而许多义人却留在万般痛苦中；但因为另有生命被预备，他要在其中按各人的行为报应各人——对恶人出于他的邪恶，对义人出于他的德行——因此他容忍，当他看到前者忍受邪恶，后者生活在奢华之中。\par

21. 我也要努力从圣经中提出另一个理由。它是什么呢？就是当人们被劝勉追求同一德行时，我们不至于说他们是另一种本性的分享者，或不是人。为此，有人论及伟大的\Person{厄里亚}说：\Quote{厄里亚是与我们相同情欲的人。}\BibleRef{雅 5:17}\newline{}又说：\Quote{我也是与你们相同情欲的人。}\BibleRef{智 7:1}\newline{}这便保证了本性的共通性。\par

22. 但你也可从中学到，这也教导我们应当称那些我们该称为有福的人为有福，这一点从以下即可明见。因为你听见\Person{保禄}说：\Quote{直到此时此刻，我们还是又饥又渴，赤身露体，受人拳打，居无定所。}\BibleRef{格前 4:11}\newline{}又说：\Quote{主所爱的，他必管教，又鞭打凡所收纳的儿子。}\BibleRef{希 12:6}\newline{}那么显然，我们并非应当称赞那些享受安逸的人，而是那些为天主受苦、在患难中的人，效法那些度德行生活、培养敬虔的人。因为先知这样说：\Quote{他们的右手是邪恶的右手。他们的女儿装扮美丽，装饰犹如圣殿。他们的仓库饱满，溢流而出；他们的羊群繁殖，在街上众多；他们的牛犊肥壮。没有篱笆的倒塌，也没有出口；街上没有喧闹。他们称那情况如此的人民有福。}\newline{}但是先知啊，你说什么？\Quote{有福的，}他说，\Quote{是以主为天主的人民；}不是财富充裕的人民，而是以敬虔为装饰的人民；他说，那人民我视为有福，尽管他们遭受无数艰难！\par

23. 但是如果需要加上第九个理由，我们可以说，这患难使受患难的人更蒙悦纳：\Quote{患难生忍耐，忍耐生老练，老练生盼望，盼望不至于羞耻。}\BibleRef{罗 5:3}\newline{}你看见吗？出于患难的老练，在我们心中坚定了对将来之事的盼望，并且持守于试炼使我们对未来有美好的盼望。因此我说不是轻率的：这些患难本身向我们指明了复活的盼望，并使受试炼的人变得更好；因为经上说：\Quote{如同金子在炉中被炼，可爱的人也在谦卑的炉中被炼。}\BibleRef{德 2:3}\par

24. 此外还有第十个理由需要提及，那是什么？就是我之前经常提到的：即如果我们有任何污点，就可以借此除去。族长使这事清楚明白，对富翁说：\Quote{拉匝禄已经得到了他的恶事，}\BibleRef{路 16:25}\newline{}因此\Quote{他得安慰。}\newline{}此外，我们还可以找到另一个理由，即我们的冠冕和赏报因此增加。因为患难越猛烈，赏报也越加增，甚至远超过：\Quote{现在的苦楚，}经上说，\Quote{不配与将来要显于我们的荣耀相比。}\BibleRef{罗 8:18}\newline{}因此，我们既有许多理由可以提出为圣徒的苦难辩护，就不要因我们的试炼而错怪、忧伤或烦恼；而是要自己操练自己的灵魂，并教导别人也如此行。\par

25. 亲爱的，你若看见一人生活于德行，紧握属神的智慧，中悦天主，却遭受无数苦难，切莫绊倒！你若见一人致力于属神的事务，即将成就一些有益的事，却立刻被取代，也不要灰心！因为我知道许多人常提出这样的问题：\Quote{有这样一个人，}他们说，\Quote{正前往某位殉道者的圣殿朝圣；在运送钱财给穷人时，遭遇船难，丧失一切。另一个人，做同样的事时，落入强盗手中，几乎丧命，赤身露体地离开那地方。}那么，我们该说什么？为什么？在这两种情形中，都不必悲伤。因为即使那人遭遇船难，他仍然拥有完整的义德果实，因为他已经尽了自己的一份。他收集了钱财，妥善保管，随身携带，出发朝圣；但后来的船难并非出于他的意愿。\Quote{但天主为什么允许这事发生？}为要使他成为经得起考验的人。\Quote{可是，}有人说，\Quote{穷人失去了那笔钱财。}你岂不比创造他们的天主更关心穷人吗？因为即使他们失去了这些东西，天主也能从别处为他们提供更大的财富。\par

26. 那么，我们不要质问祂的作为，而要凡事归光荣于祂。因为祂常常允许这样的事发生，并非轻率或毫无目的。除了祂不会忽略那些本应从这些财富中得到安慰的人，反而为他们提供其他生活来源外，祂也使遭遇船难的人更加经得起考验，并为他预备更大的赏报；因为在这样的灾难中感谢天主，远胜过施舍。因为不仅是我们施舍的，连我们被他人剥夺而坚忍承受的，也为我们带来许多果实。为使你知道后者确实是更大的事，我从\Person{约伯}身上所发生的事来证明。他拥有财富时，向穷人敞开家门，将他所有的一切施舍出去；但他敞开家门给穷人时，并不如他听到房屋倒塌而耐心承受时那样显赫。他用羊群羊毛给赤身露体者穿衣时并不显赫，如同他听到火降下烧尽他所有羊群却仍感谢时那样显赫和著名。从前，他是爱人的；如今，他是爱智慧的。从前，他怜悯穷人；如今，他感谢主！他没有对自己说：\Quote{为什么发生这事？羊群被烧尽了，数千穷人曾靠它们维持生计；如果我不配享有如此丰裕，至少祂应为了那些受惠者而顾惜我。}\par

27. \Person{约伯}没有说出这样的话，也没有这样想，因为他知道天主安排万事，为使人得益。为使你知道，他在丧失一切后感谢天主，比他在拥有一切时施舍，给了魔鬼更沉重的一击；请注意，当他拥有时，魔鬼尚能提出某种猜疑，尽管是虚假的，它仍能提出：\Quote{约伯敬畏祢，难道是徒然的吗？}但当它夺走了一切，剥光了他的一切，这人却仍对天主保持同样的善意，从那时起，它无耻的口被堵住了，再也无可指责。因为义人比先前更显赫。因为高贵并感恩地承受丧失一切，远胜于富足时施舍；这事已在义人身上得到证明。从前，他对同仆有极大的仁慈；如今，他对主显出极深的热爱！\par

28. 我这样延长讲论并非没有目的；因为有许多人，常常在从事慈悲工作如供养寡妇时，被剥夺了全部财产。有些人因火灾意外失去一切；有些人遭遇船难；还有些人因诬告和类似的侵害，虽然做了许多施舍，却陷入极端的贫困、疾病和痛苦，且未得到任何人的帮助。为避免我们像许多人那样说：\Quote{没有人知道任何事；}刚才所说的足以消除在这点上的所有困惑。假设有人反对说：\Quote{某某人做了许多施舍后，失去了一切？}失去一切又怎样？如果他因这损失而感谢，他将从天主那里获得更大的恩宠！他不仅会像\Person{约伯}那样得到双倍，而是在来世得到百倍。但若他在这里受苦，他坚忍承受一切的事实，将为他带来更大的宝藏；因为天主允许他从富足落入贫困，正是为以此召唤他进行更频繁的操练和更大的战斗。是否像常发生的那样，火灾吞没了你的房屋，烧毁并吞噬了你所有的财产？要记住\Person{约伯}所发生的事；感谢主，祂本可阻止却没有阻止；你将得到如同你将所有财富存入穷人手中一样大的赏报！但你是否在贫困、饥饿和无数危险中度日？要记住\Person{拉匝禄}，他曾与疾病、贫困、孤苦以及无数其他考验搏斗，而且是在如此崇高的德行之后！要记住宗徒们，他们生活在饥饿、口渴和赤身露体中；记住先知、族长、义人们，你会发现他们并非在富有或奢华者之中，而是在贫穷、受苦和忧伤者之中！\par

29. 你把这些话对自己说，向主感恩，因为他使你成为这一方，并非恨你，而是极其爱你；因为他若没有极其爱那些人，就不会允许他们那样受苦，因为他藉这些苦难使他们更显荣耀。没有比感恩更好的事了；同样，也没有比亵渎更坏的事了。我们不应感到奇怪，当我们专注于属灵之事时，会遭受许多痛苦。因为正如盗贼不会在只有干草、糠秕和稻草的地方凿墙并严密看守，而是在有金银的地方；同样，魔鬼也特别攻击那些从事属灵事务的人。哪里有德行，那里就有许多陷阱！哪里有施舍，那里就有嫉妒！但我们有一件最好的武器，足以击退所有这些攻击，就是在一切事上感谢天主。请问，\Person{亚伯尔}在向天主献初果时，不是死在他兄弟的手中吗？但天主仍然允许了此事，并非恨那尊敬他的人，而是极其爱他；并且除了那美好祭献所得的赏报外，还藉殉道为他预备了另一顶冠冕。\Person{梅瑟}想要保护一个受欺压的人，却陷入极大的危险，并被流放国外。\EditorialNote{出谷纪第二章} 天主也允许了这一切，为叫你学习圣人的忍耐。因为如果我们预先知道不会遭受任何严重的苦，然后才着手宗教工作，我们就似乎不算做了什么大事，因为我们有这样的安全保证。但事实上，那些这样做的人更令人惊叹，正因为此：因为他们虽然预见到危险、惩罚、死亡和无数灾祸，却没有停止那些善工，也没有因对恐怖的预期而减少热忱。\par

30. 因此，正如三个青年所说：\Quote{有一位天主在天上，他能拯救我们；即便不救，大王，愿你知道，我们决不事奉你的神，也不朝拜你所立的金像。}\BibleRef{达 3:17} 你也要在履行对天主的任何职责时，预料种种危险、种种惩罚、种种死亡；如果这些事情发生，不要惊奇，也不要惊慌。因为经上说：\Quote{我儿，你若前来事奉主，要预备你的灵魂接受试探。}\BibleRef{德 2:1} 因为没有人选择战斗，却期望不受伤而赢得冠冕！因此，你既已承担与魔鬼全面战斗，就不要想过一种没有危险、充满安逸的生活！天主没有在此世向你许诺他的赏报和他的应许；而是在来世的一切辉煌！所以，如果你自己行了一件善事却遭受相反的对待，或者你看见别人遭受这事，就应欢喜快乐，因为这对你成为更高赏报的根源！不要沮丧，也不要放弃热忱，更不要变得麻木；反倒要更加急切地前进；因为就连宗徒们宣讲时，虽然被鞭打、被石击，常被囚禁在监狱中，他们不仅在脱离危险之后，而且在那些危险之中，也更加勇敢地宣讲真理的信息。可以看到保禄在监狱中，甚至在锁链中，教训人并施洗；而且在法庭上、在船难中、在暴风雨中、在千万危险中，也做着同样的事。你也要效法这些圣人，只要你能，就不要停止善工；虽然你看见魔鬼千万次阻挠你，也决不要后退！你或许带着你的财富，遭遇了船难；但保禄携带圣言，远胜于一切财富，前往罗马，并遭遇船难，并忍受了无数艰难。他自己也表明了这一点，当他说道：\Quote{我们曾多次想去你们那里，但撒殚阻止了我们。}\BibleRef{得前 2:18} 天主允许了这事，借此更丰富地彰显他的大能，并显示魔鬼所做或阻止的一切，既没有减少也没有中断福音的宣讲。因此，保禄在一切事上感谢天主；并且知道他自己因此更被认可，就在每个场合都表现出极大的热忱，不让任何阻碍阻止他！\par

31. 因此，每当我们受挫于属灵工作时，让我们屡次重新着手去做；不要问：\Quote{天主为什么允许这些阻碍？}因为他允许它们，正是为向他人更多显示你的热切和你的大爱；这正是爱者的特别标志：永不停止做他所爱者所喜悦的事。的确，软弱无力、无精打采的人，会在第一次冲击后退缩；但精力充沛、警觉的人，即使被阻挠一千次，也会更加投身于天主的事；尽己所能完成一切，并在一切事上感恩。那么，让我们这样做吧！感恩是巨大的宝藏；是丰厚的财富；是不能被夺走的美善；是强大的武器！同样，亵渎会增加我们眼前的不幸，并使我们失去比已失去的多得多的东西。你失去了金钱吗？如果你感恩，你就赢得了你的灵魂，并获得了更大的财富，因为你得到了天主恩宠更大的分额。但如果你亵渎，那么除此以外，你还失去了你自己的安全；你也没有重新获得你的财富；是的，连你原本拥有的灵魂，你也牺牲了！\par

32. 既然我们的论述现在转向亵渎的话题，我愿向诸位请求一项恩惠，作为我这番讲话和与你们交流的回报，那就是：你们要替我纠正这城中的亵渎者。若你们在公共街道或广场中央听到有人亵渎天主，就上前去斥责他；若有必要施加击打，不要吝惜。打他的脸，击他的口，以这击打圣化你的手；若有人控告你，拉你到法庭，你就跟从他们去；当法官在审判席上向你问责时，你要大胆地说，这人亵渎了天使的君王！因为若必须惩罚那些亵渎地上君王的人，更何况那些侮辱天主的人呢？这是共同的罪行，公开的伤害；凡愿意的人，都可以提出控告。愿犹太人和希腊人知道，基督徒是这城的拯救者，是它的守护者、保护者和导师。也愿放荡和乖僻的人知道，他们必须敬畏天主的仆人；若他们有时想说出这类话，他们必会四处环顾，甚至因自己的影子而战栗，生怕有基督徒听见了所说的话，会突然扑上来严厉地责罚他们。你们没有听说过\Person{若翰}所做的事吗？他看见一个暴君推翻婚姻的法律，便勇敢地在广场中央宣告：\BibleQuote{你不可娶你兄弟\Person{斐理伯}的妻子。}\BibleRef{谷 6:18} 但我劝你们起来，不是反对君王或法官，不是为被践踏的婚姻法令，也不是为受侮辱的同仆，而是要求你们惩治一个同等的弟兄，因为他向主傲慢。真的，若我对你们说：惩罚和纠正那些违法的君王或法官，你们岂不说我疯了？但\Person{若翰}确实这样做了。所以这对我们也不算过分。现在，至少纠正一个同仆、一个同等的人；即使必要牺牲，也不可退缩不去责罚一个弟兄。这就是你们的殉道，因为\Person{若翰}也是一个殉道者。虽然他没有被命令献祭或崇拜偶像，但为了被藐视的神圣法律，他献出了头颅。你也当这样为真理奋斗至死，天主必为你作战！不要给我这样冷漠的回答：\InnerQuote{这与我何干？我与他毫无关系。}唯独与魔鬼我们毫无关系，但与众人我们有许多共同之处：因为他们与我们同享一样的本性，居住在同一大地，以同样的食物滋养，有同一的主，领受同一的法律，被邀享同样的福份。所以我们不要说与他们毫无关系；这是属魔鬼的话语，是魔鬼般的不人道。因此，不要说出这样的话，而要显出合乎弟兄情谊的温柔关怀！\par

33. 至于我，我以最确凿的把握承诺，并向诸位全体担保：如果你们所有在场的人愿意承担起这城居民的安全，我们必将迅速使全城得到改善。即使这里只有城中最小的一部分——数量上最少，但在虔敬方面却是最主要的——也是如此。让我们承担起弟兄们安全的责任！一人为热忱所燃起，就足以革新整个团体！但不仅是一人、二人、三人，而是如此众多的人能够照顾那些被疏忽的人，那么，若不是因为我们自己的懈怠，而非缺乏力量，大多数人是不会丧亡跌倒的。这岂非荒谬吗？当我们碰巧在广场上看见一场争斗，我们就冲进去，劝解争斗者！但我何必说争斗？若我们偶然看见一头驴跌倒了，我们都急忙伸手扶它起来。而我们却忽略了正在丧亡的弟兄们！亵渎者就是一头驴：他无法承受愤怒的重负，跌倒了。上前扶起他，用言语和行动，用温和与严厉；让医治的方法多样。若我们这样尽自己的本分，为邻人的安全劳苦，我们很快就成为那些受我们纠正的人所渴慕和爱戴的对象；而超乎一切的是，我们将享受那储备的美好事物。愿天主恩赐我们众人获得这一切，藉着我们的主耶稣基督的恩宠和怜悯；愿光荣、权能和尊崇，藉着祂、偕同祂，偕同圣神，归于圣父，自今而永，及于万世。阿们。\par

\SourceRecord{Translated by W.R.W. Stephens.；Nicene and Post-Nicene Fathers, First Series；Vol. 9.；Edited by Philip Schaff.；Buffalo, NY: Christian Literature Publishing Co.,；1889.；<http://www.newadvent.org/fathers/190101.htm>.}

% structure: p0001 source=h1 target=section reason=page-title

\section{论雕像讲道集第二篇}

\Emph{曾在安提约基雅旧堂，当时被称为旧堂，担任司铎时所讲。论及因推倒虔诚伟大的德奥多西皇帝雕像而生的骚乱所降临于城市的灾难。并论宗徒的话：\BibleQuote{命令那富有的，不要心高气傲。} \BibleRef{弟前 6:17}。并论贪财。}\par

1. 我要说什么，或宣讲什么呢？当下的时节是流泪之时，而非言语之际；是哀悼之时，而非讲论之际；是祈祷之时，而非讲道之际。所行的事如此大胆，其危害之重，创伤之深，非人力所能医治，迫切需要上天的援助。昔日的约伯，当他失去一切后，坐在灰堆上；他的朋友们听闻后赶来，远远看见他，便撕裂外衣，撒灰在头上，大大哀哭。\BibleRef{约 2:8}现在，周围所有的城市都应当这样做，来我们城市，以同情的心哀哭我们所遭受的。那时约伯坐在灰堆上；现在，我们的城市坐在巨大的罗网中央。因为正如魔鬼那时猛烈地攻击那义人的羊群、牛群和一切所有，现在它也向整个城市发怒。但那时和现在，天主都允许了；那时，为了藉苦难的伟大使义人更光荣；现在，为藉这患难的极端使我们更清醒。请允许我为当前的状态哀哭。我们已经沉默了七天，正如约伯的朋友们那样。\BibleRef{约 2:13}请允许我今天开口，哀悼这共同的灾难。\par

2. 谁，亲爱的，迷惑了我们？谁嫉妒了我们？这一切变化从何而至？从前没有比我们的城市更尊贵的！现在，从来没有更可怜的！民众如此有序和安静，甚至像一匹驯良温顺的骏马，总是顺从统治者的手，如今却如此突然地与我们脱缰，以至于造成了如此大的恶事，几乎无人敢提及。\par

我现在哀悼悲叹，不是为了那将要来临的愤怒之烈，而是为了所显现的疯狂之极！因为即使皇帝不被激怒，不发怒，即使他既不惩罚，也不报复；请问，我们怎能承受所行的这些事的羞辱？我发现讲道的言语被哀叹打断了；我几乎不能开口、启唇、动舌或发出一个音节！因此，就像缰绳一样，悲伤的重压勒住我的舌头，阻止我想说的话。\par

3. 从前没有比我们的城市更幸福的了；现在也没有比它更忧郁的了。如同蜜蜂绕着蜂巢嗡嗡，从前居民每天在广场上穿梭，所有人都称我们因人数众多而有福。但看哪，如今这蜂巢变得孤单！因为正如烟雾驱散蜜蜂，恐惧也驱散了我们的蜂群；先知哀悼耶路撒冷时所说的话，我们现在可以恰当地说：\Quote{我们的城市成了\BibleQuote{像失去叶子的橡树，又如无水浇灌的园子}。}\BibleRef{依 1:30}因为正如一个园子，当灌溉失败时，露出光秃秃的树木，没有果实；我们的城市现在也是如此。从上而来的帮助离弃了它，它荒凉地站在那里，几乎所有的居民都被剥夺了。\par

4. 没有比自己的家乡更甜美的了；但现在，却成了没有比它更苦涩的了！人人都从那生育之地逃跑，如同逃脱罗网。他们遗弃它如同遗弃地牢；他们从中跳出，如同跳出火坑。正如一所房子被火焰吞没时，不仅住在里面的人，而且所有附近的人，都极其匆忙地逃离，只求拯救赤裸的身体；同样现在，当皇帝的愤怒如从天降的火时，每个人都急着及时出去，在火势蔓延到他们之前挽救赤裸的身体。现在，我们的灾难成了一个谜：一场没有敌人的逃亡；一场没有战争的驱逐；一场没有被俘的俘虏！我们未见过蛮族的火，也未见过敌人的面：但我们却经历着俘虏的痛苦。现在所有人都听到了我们的灾难；因为他们接待我们的流亡者，从他们那里得知了落在我们城市上的打击。\par

5. 然而我并不为此感到羞耻，也不脸红。让所有人了解这城市的苦难，使他们同情自己的母亲，从整个大地齐声向天主呼求；并同心一致为这普世的乳母和父母向天上的君王恳求。近日我们的城市被震动；但现在居民的内心在动摇！那时房屋的地基震动，但现在每颗心灵的地基都在颤抖；我们都亲眼看见死亡每天在我们面前！我们生活在不断的恐惧中，忍受着加音的惩罚；一种比从前那些监狱囚犯更可怜的惩罚；我们现在经历一种新的、奇异的围困，远比普通围困可怕。因为被敌人围困的人只是被关在城墙内；但如今连广场对我们都变得无法通行，每个人都缩在自己家的墙壁内！正如被围困者不能安全地走出城墙，因为敌人在外面扎营；同样，对于这座城市的许多居民来说，也不能安全地出去或公开露面；因为有些人到处搜捕无辜和有罪的人；甚至在广场中央就抓住他们，不加礼节，随意拖到法庭。因此，自由人与他们的家奴被困在屋内；焦急地、细致地询问那些他们可以安全提问的人：\Quote{今天谁被捉走了？谁被带走或惩罚了？怎么回事？怎么发生的？}他们过着比任何死亡更悲惨的生活；被迫每天为别人的灾难哀悼；同时为自己的安全颤抖，与死人无异；因为他们已被恐惧吓死。\par

6. 但如果有人没有这种恐惧和痛苦，选择进入广场，他立刻被凄凉的景象逼回自己的住所；几乎找不到一两个人，而且他们低垂着头，目光沮丧地游荡，而几天前这里的人群比河流的流水还更无休止地涌动。然而所有这些现在都离我们而去！如同在一片茂密的橡树林中，许多树木被砍得东倒西歪，景象令人悲伤，就像一个头上有多处秃斑的头；同样，城市本身，居民减少，只有零星几人出现，现在变得凄凉，给目睹这一切的人投下沉重的悲伤迷雾。不仅是地面，就连空气的特性，甚至太阳光线的圆圈，在我看来都显得哀伤，更加暗淡；这并非元素改变本性，而是我们的眼睛被悲伤的云雾所迷惑，无法清晰或同样愉悦地接受光线。这就是古时先知哀叹时所说的：\BibleQuote{太阳要在午间落下，白昼变为黑暗。}\BibleRef{亚 8:9}他说这话，并非因为晨星会亏缺或白昼消失，而是因为那些忧愁的人，由于痛苦之黑暗，甚至中午的光线也无法察觉；现在正是如此。无论向何处望去，无论是地面还是墙壁；无论是城市的柱廊还是邻居，他似乎都看见黑夜和深深的黑暗；一切都充满忧郁！有一种恐怖的寂静，处处孤独；那人潮熟悉的喧闹被扼杀；仿佛所有人都已沉入地下，沉默无语现在占据了城市；所有人都像石头，灾难像塞住舌头的塞子压迫他们；他们保持最深沉的沉默，是的，这种沉默就像敌人来袭，用火与剑将他们一举消灭！\par

7. 现在是适合说：\BibleQuote{你们要召来哀悼的妇女，让她们前来；又召来精通哀歌的妇女，让她们放声哀哭。让你们的眼睛流下泪水，你们的眼睑涌出泪泉。}\BibleRef{耶 9:17}的时刻。山丘啊，你们要放声痛哭；群山啊，你们要哀号！让我们召唤整个受造界同情我们的苦难。如此伟大的城市，东方天空下的众城之首，正有从文明世界中心被拔除的危险！她曾有过那么多儿女，如今突然成了无子女的，没有谁前来帮助她！因为那被侮辱者在世上没有相称尊严的人；他是一位君主，是尘世一切的高峰和首脑！因此，让我们投奔天上的君王。让我们呼求祂的援助。如果我们得不到上天的恩宠，那么对于我们的遭遇就没有安慰了！\par

8. 我在此本想结束这篇讲辞；因为忧苦中的人心不适于将谈论拖得太长。如同稠密的云层形成后，飞入太阳的光线之下，又将全部光辉反照给太阳一样，忧苦的云层一旦遮蔽我们的灵魂，便拒绝让言语轻易通过，反而将其扼住，强留其中。这不仅发生在讲者身上，也发生在听者身上；因为它既不容许言语从讲者灵魂中自由涌出，也不容许它以其自然的能力沉入听者心思。因此，古时的犹太人，当他们在泥土和砖块中为奴时，也无心听梅瑟屡次向他们宣讲关于未来解救的伟大之事；沮丧使他们的心思无法接纳训言，并封闭了听觉。因此，就我个人而言，我本想在此结束讲辞；但想到云层不但能拦截阳光的前行，也常有相反的情况发生：因为太阳不断以炽热照射云层，将其消融，常常从云层中间穿透而出，骤然照耀，欣然迎向观者的目光。我今天也期望如此；言语不断与你们的心思相伴，住在其中，我希望冲破忧苦的云层，再次以惯常的训诲照亮你们的理智！\par

9. 但请你们留心听！暂且侧耳倾听！摆脱这沮丧吧！让我们恢复往日的习惯；正如我们素来在此聚会欢喜一样，如今也这样行，将一切交托给天主。这也会有助于我们实际脱离灾祸。因为若主看见祂的言语被用心聆听，并看见我们对神圣智慧的热爱经得起如今时日的考验，祂必会迅速再次扶起我们，从眼前的暴风雨中带来平静而幸福的转变。因为基督徒与不信者之间也应以此区分：即高贵地承当一切；并借着对未来的希望，超越人间灾祸的袭击。信徒站立在磐石上；为此，他不会被波涛的冲刷所倾覆。即便诱惑的波浪涌起，也无法触及他的脚。他站立得如此高，任何这样的冲击都无法触及。所以，亲爱的，我们不要沮丧！天主创造了我们，祂关心我们的安全远胜于我们自己。我们对于遭受可怕不幸的忧虑，远不及那赐予我们灵魂、又赐予我们如此多美好事物之主。让我们乘上这些希望的翅膀，以惯常的乐意聆听将要讲论的事。\par

10. 亲爱的，我前些日子向你们作了一次长篇讲论，但我看见你们都紧随其后，无人中途转身。我为你们的乐意感谢你们，并已获得了劳苦的报酬。但除了那专注之外，还有另一个我曾经向你们要求的报酬；或许你们知道并记得它。那报酬是什么呢？就是你们应惩罚和责斥城中的那些亵渎者；你们应制止那些对天主暴戾无礼的人！我不认为当时我是凭自己说这些话；而是天主预见到将发生的事，将这些话注入我的心思；因为如果我们惩罚了那些胆敢如此行的人，现在所发生的事就绝不会发生。若迫于必要，去冒险，甚至为惩戒和纠正这样的人而受苦（这会为我们带来殉道者的荣冠），也比现在因这些人的悖逆而恐惧、战栗、等待死亡要好得多！看，恶行是少数人所为，但责备却落在众人身上！看，因着这些人，我们如今都处于恐惧之中，并正在承受这些人胆敢行的事的惩罚！但如果我们及时抓住他们，将他们驱逐出城，惩戒他们，矫正这有病的肢体，我们就不至于遭受眼前的恐怖。我知道这城的风气自古以来高贵；但某些外乡人和混杂种族的人——受诅咒而有害的品性——对自己的得救无望，行了所行的事。正因如此，我一直提高声音，不断作见证说：让我们惩罚那些亵渎者的狂妄——控制他们的精神，并为他们的得救而操心——即使为此而死，这行为也必为我们带来巨大的益处：让我们不要忽略对我们共同的主人的侮辱；忽略这样的事必会给我们的城带来某种大祸！\par

11. 我曾预言这些事，如今它们确实发生了——我们正为那懈怠付出代价！你们忽略了加于天主的侮辱！——看，祂容许皇帝受侮辱，并使极度的危险临到众人，好让我们以这恐惧为那懈怠付出代价；那么，我预言这些事，并殷勤敦促你们的爱德，难道是徒然无益的吗？但尽管如此，什么事也没做。然而，现在让我们行动起来；借着眼前的灾祸受管教，现在让我们制止这些人的狂乱。让我们封住他们的口，如同封闭毒泉一般；并引导他们转向相反的方向，那么已经抓住这城的恶事必定会止息。教会不是剧院，让我们仅为消遣而听讲。我们应当带着益处离开此地，在离开这地方之前获得一些新的巨大收获。因为如果我们只是暂时被吸引，然后空手回家，内心毫无所获，那么我们的聚会就是徒然且无理性的。\par

12. 我要这些掌声、欢呼和骚动的赞许有何用？我所寻求的赞美，是你们在行为上活出我所说的一切。那时我才是一个可羡慕和幸福的人，不是当你们赞许，而是当你们以全部乐意遵行我所听到的。那么，让各人纠正他的邻人，因为经上说：\BibleQuote{你们要彼此劝勉}\BibleRef{得前 5:11}。若不如此，各人的罪行必将给全城带来某种普遍的、不可容忍的损害。看，我们虽不察觉自己参与此事，却与那些胆大妄为之徒同样惊恐！我们惧怕皇帝的愤怒降临到所有人身上；而我们仅以\Quote{我不在场；我没有同谋，也没有参与这些行为}来辩护是不够的。他可能回答说：\Quote{为此，你将受罚，并受最重的刑罚，因为你当时不在场；你没有制止、没有阻止闹事者，也没有为皇帝的尊荣冒任何风险！你没有参与这些胆大妄为之举？我称赞这点，并视为好事。但你当时没有制止这些事的发生。这是罪责的缘由！}这些话，若我们默许继续对天主的侮辱和冒犯，也将从天主那里听到。因为那将塔冷通埋在地里的人，并非因自己所为的恶事受审——他已将所有交托的全数归还——而是因他没有使之增值；因他没有教导他人；因他没有存入银行家手中，即他没有劝诫、没有忠告、没有责备、没有改正那些无法无天的罪人——他的邻人。为此，他被毫无缓刑地打发到那些不可忍受的刑罚中！但我深信，尽管你们从前没有，如今至少会从事这项纠正的工作，不再忽视对天主的冒犯。因为所发生的事，即使无人警告，也足以使那些即使如此迟钝的人确信，他们必须为自身安全而努力。\par

13. 但现在是我们继续从圣保禄那里为你们摆出惯常的筵席的时候，通过处理今日读经的主题，并把它摆在你们众人面前。那么，今日所读的经文是什么？\BibleQuote{你要嘱咐今世富足的人，不要自高}\BibleRef{弟前 6:17}。当他说\BibleQuote{今世富足的人}时，他表明还有别的富足人，即来世富足的人：比如那拉匝禄，在今世贫穷，在将来却富足；不是富有金银和如此易朽暂存的财富，而是富有那些不可言喻的美善，就是\BibleQuote{眼所未见，耳所未闻，人心所未曾想到的}\BibleRef{格前 2:9}。因为这才是真正的财富和富足，即善而无掺杂，且不受任何变化。那轻视他的富翁却不是这样，他成了人类最贫穷的人。后来至少当他寻求得一滴水时，连那也没有得到，他竟陷入如此极端的贫困。为此，他称他们为\BibleQuote{今世}富足，为教导你：世界的财富与今世一同消灭。它不再前进，也不与迁移的主人一同改变地方，而常常在主人结局之前就离开他们；因此他用这话显示：\BibleQuote{也不要依赖无定的钱财}；因为再没有比财富更不可靠的了。我曾多次说过，并且不会停止说：它是一个逃跑的、忘恩负义的仆人，没有忠信；即使你给它加上一万条锁链，它也会拖着锁链逃走。的确，拥有它的人常常用门闩和门把它关起来，派奴隶四周守护。但它说服了这些仆人，带着它的守卫一同逃走；把他的看守像锁链一样拖着，这监守竟是如此不可靠。那么，还有什么比这更不可靠？还有什么比沉迷于它的人更可怜？当人们以全部热情去收集如此脆弱易逝的东西时，他们不听先知的话：\Quote{祸哉，那些信赖自己力量，并以自己财富众多而夸耀的人。}告诉我为何要宣布这祸？——他说：\Quote{他积蓄财宝，却不知为谁收取}——因为劳作是确定的，而享受却不确定。非常常见的是，你辛劳受苦，却是为了仇敌。你去世后的财富遗产，在许多情况下落到了那些曾用千百种方式伤害你、谋害你的人手中，把罪恶分给你，把享受留给别人！\par

14. 但此处值得探究，为何他不说：\Quote{吩咐今世富足的人，不可富有；吩咐他们成为贫穷；吩咐他们舍弃所有的；}而说：\Quote{吩咐他们不可心高气傲。}因为他知道，骄傲是财富的根源和基础；若有人懂得谦逊，就不会对此事过分在意。请告诉我，你身边带着这么多仆人、食客、谄媚者，以及其他排场，究竟是为何？并非为了需要，而只是为了骄傲，好叫你借此显得比别人尊贵！此外，他知道，若财富用于必需之事，并不禁止。因为如我所说，酒并非坏事，醉酒才是坏事。贪财者是一回事，富人是另一回事。贪财者并不富有；他缺少许多东西，既然需要许多，就永不能富有。贪财者是财富的看守者，而非主人；是奴隶，而非主人。因为他宁愿割下一块肉给人，也不愿给人他那埋藏的金子。仿佛有人命令并强迫他丝毫不得动用这些宝藏，他便极其谨慎地看守保存，戒用自己的财物，如同戒用别人的财物一般。的确，它们不是他的。因为他既不能决意施与他人，也不能分给穷人，纵然要受无穷的惩罚，他怎能认为那是自己的呢？他对于自己没有自由使用和享用的东西，又如何能拥有呢？但除此以外，保禄并不惯于对每个人都一样吩咐，而是迁就听者的软弱，正如基督也曾如此。因为当那富人来到祂面前，问祂关于生命的事，祂并没有直接说：\Quote{去，变卖你所有的}——\BibleRef{玛 19:16}，而是暂且不提，先对他讲其他诫命。后来那人挑战祂说：\Quote{我还缺少什么？}祂并没有简单说：\Quote{变卖你所有的}，而是说：\Quote{你若愿意成全，去变卖你所有的}——\BibleRef{玛 19:21}。\Quote{我由你决定。我给你完全的选择权。我不加给你任何必须。}因此，保禄也没有对富人讲论贫穷，而是讲论谦卑；既因听者的软弱，也因他深知，只要他能引导他们操练节制、远离骄傲，就能很快使他们解脱对富贵的渴求。\par

15. 而且，在给了\Quote{不可心高气傲}的劝诫之后，他还教导他们如何才能避免如此。如何避免呢？就是要思考财富的本性，它是多么不确定、多么不可靠！因此他接着说：\Quote{不要依靠无定的财物。}富人并非拥有很多的人，而是付出很多的人。亚巴郎是富人，但他并不悭吝；因为他没有将自己的念想转向此人的房屋，也没有贪求那人的财富；而是走出去，环顾四周，看看哪里有客旅或穷人，好扶助贫穷，款待行旅。他没有用黄金覆盖屋顶，而是将帐篷支在橡树旁，满足于树叶的荫蔽。然而他的住处如此卓著，以致天使也不羞于与他同住；因为他们寻求的不是居所的华丽，而是灵魂的德行。因此，亲爱的，让我们效法此人，将我们所拥有的施予穷人。那住处虽简陋，却比君王的大殿更卓越。从未有君王接待过天使；但他，住在橡树下，仅支一帐幕，竟配得那荣誉：并非因他居所简陋而受此荣誉，而是因他灵魂的伟大和其中贮藏的财富而享有此恩惠。\par

16. 那么，让我们不要装饰我们的房屋，而应装饰我们的灵魂，胜过装饰房屋。因为用大理石徒然无益地装饰墙壁，却让赤身露体的基督被忽视，岂不令人羞愧？房屋对你有何益处，人啊！你离世时能带走它吗？你离世时不能带走它。但你的灵魂，你离世时必定会带走！看，如今这场大危险已临到我们！让你的房屋站在你身边！让它们将你从威胁的危险中拯救出来！但它们不能！你们自己就是见证，你们正离开它们，匆匆奔向旷野；害怕它们如同害怕陷阱和网罗！让财富现在来帮助吧！但此刻不是它们帮助的时候！如果财富的力量在人的愤怒面前显得不足，那么在神圣而无情的审判台前更将如此！如果只是人被激怒和得罪，即使是黄金也无济于事，那么当天主发怒时，金钱的力量就更彻底无能了，因为祂不需要财富！我们建造房屋是为了有个住所，不是为了炫耀。超出我们所需的，就是多余和无用的。穿上比脚大的鞋，你会无法忍受，因为它妨碍走路。同样，超出需要的房屋，也是你迈向天堂的障碍。你想建造宽敞华丽的房屋？我不禁止，但不要建在地上！你要在天堂为自己建造帐幕，并能容纳他人——永不毁坏的帐幕。你为何为转瞬即逝、必将遗留在世的事物而疯狂？没有比财富更不可靠的了。今日为你，明日对你。它到处激起嫉妒者的眼目。它是敌对的同伴，家中的仇敌；你们这些拥有财富并千方百计将其埋藏隐匿的人就是见证；正如现在，我们的财富本身使危险对我们更难承受！你确实看见穷人随时准备行动，无牵无挂，为一切做好准备；而富人却极度困惑，四处游荡，寻找可埋藏黄金之地，或寻找可存放之人！人啊，你为何寻求你的同仆？基督已准备好接收并为你保管存款；不仅保管，还要增加，并以丰厚利息偿还。无人能从祂手中强行夺走。祂不仅保管存款，还因此使你脱离危险。因为在人当中，接受信托财宝的人认为他们帮了我们一个忙，保管了他们所接收的；但基督却相反；当祂接收你存放的财宝时，祂不说祂施了恩，而是说祂受了恩；并且对于祂对你财富的保管，祂不向你索取报酬，反而给你报酬！\par

17. 那么，当我们忽视那位能够保管、并因信托而感恩、且给予巨大而不可言喻的回报，而我们却将财宝委托给那些无力保管、且自以为施恩于我们、最终只归还所交之物的人时，我们还能有什么辩护或借口呢？你在此世是客旅和朝圣者！你在天上拥有自己的家乡！将一切转移到那里——这样在真正享受之前，你在此已能享受报偿。那以美好希望为养料、对未来充满信心的人，在此已尝到天国的滋味！因为通常没有什么比美好的未来希望更能修复灵魂、使人更好；所以如果你将财富转移到那里，你就能为你的灵魂提供合适的安息。因为那些将所有精力用于装饰住所的人，虽然外在丰富，却忽略了内在，让他们的灵魂荒芜、肮脏、布满蛛网。但如果他们对外在事物漠不关心，而将全部注意力倾注于心灵，从各方面美化它；那么这些人的灵魂将成为基督的安息之所。有了基督作为居住者，还有什么能更蒙福呢？你想富有吗？让天主成为你的朋友，你比所有人更富有！——你想富有吗？不要心高气傲！——这准则不仅适用于未来，也适用于现在。因为没有像富人那样令人嫉妒的了；但当加上骄傲，就形成双重的悬崖；战争四面八方变得更猛烈。但如果你懂得节制，你的谦卑就会削弱嫉妒的暴政；你无论拥有什么，都能安全地拥有。因为美德的本性如此：它不仅在将来对我们有益，在此世也带来现世的赏报。\par

18. 那么，我们不要在财富上高傲，也不要在任何事上高傲；因为即使在属灵的事上，高傲的人也会跌倒和败亡，更何况在属血肉的事上呢。让我们记住我们的本性，回想我们的罪过，明白我们是什么；这将成为我们完全谦卑的充分基础。不要对我说：\Quote{我已经积攒了这么多年的收入；无数的金塔冷；每日增长的收益。} 无论你说多少，你都是徒然，毫无益处。常常在一个时辰内，是啊，在短短的一瞬间，就像轻尘被从上而下的狂风骤雨吹扫出房屋一样，这一切都被一阵风从家中扫除。我们的生命充满了这样的例子，圣经也充满了这类教训。今朝富，明日贫。因此，我常常在阅读遗嘱时微笑，遗嘱说，让某人拥有这些田地或这所房屋的所有权，而另一个人拥有其使用权。因为我们都有使用权，但没有一个人拥有所有权。因为即使财富在我们一生中保持不变，我们最终也不得不，无论愿意与否，将它们转移到他人手中；我们只享受了它们的使用权，而离开此世时，对所有权却赤身露体、一无所有！由此可见，只有那些藐视使用、嘲笑享受的人，才真正拥有财产的所有权。因为那人把财产从他那里抛出去，施舍给穷人，他就按应当的方式使用了它；当他离去时，他便带着这些东西的所有权离开，甚至在死亡中也不会被剥夺，反而在那时全部收回；是啊，远超过这些，在那审判之日，当他最需要它们的保护，当我们所有人都必须为我们所做的事交账时。所以，如果有人希望完全拥有他的财富、使用权和所有权，就让他自己从这一切中解脱；因为，确实，不这样做的人，无论如何都会在死亡时与之分离；而且常常在死前，在危险和无数祸患中失去它们。\par

19. 不仅如此，灾难在于变化来得突然；而且富人没有经验去忍受贫穷。但穷人却不然；因为他信赖的不是金银这些无生命的物质，而是\Quote{天主，他丰厚地赏赐我们一切享用}。所以富人比穷人处于更多的不确定之中，因为他经历着频繁而多样的变化。这是什么意思？\Quote{他将一切丰富地赐给我们享用}。\BibleRef{弟前 6:17} 天主慷慨地赐予那些比财富更必需的一切；例如，空气、水、火、太阳；以及这一类的一切。富人不能说他所享受的阳光比穷人更多；他不能说他所呼吸的空气更丰富：但这些都同样提供给了所有人。也许有人会说，为什么那些更伟大、更必需的福分，以及维持我们生命的东西，天主使之成为共有的；而更小、更不珍贵的（我指的是金钱）却并非如此共有呢？这是为什么呢？为的是我们的生命能够受到操练，为的是我们能有修德的训练场。因为如果这些必需品不是共有的，也许那些富人，习于他们的贪婪，会勒死那些穷人。因为他们若为了金钱就这样做，那么为了那更必需的物品，他们更会如此。再者，如果金钱也是普遍的财产，并且以同样的方式提供给所有人，那么施舍的机会和仁爱的条件就会被剥夺。\par

20. 那么，为了我们能安全地生活，我们存在的源泉被造为共有的。另一方面，为了我们有机会获得荣冠和美誉，财产没有被造为共有的；为的是我们憎恶贪婪，追随正义，并慷慨地把我们的财物分施给穷人，借此我们能为自己的罪过获得某种减轻。天主使你富有，你为什么使自己贫穷？他使你富有，为的是你帮助有需要的人；为的是你通过对他人的慷慨，获得自己罪过的赦免。他给了你金钱，不是要你把它封锁起来以致毁灭，而是要你为你的救恩而倾流而出。为此，他也使财富的占有变得不确定和不稳定，企图以此缓和你们对财富的疯狂热度。因为如果财富的占有者，即使现在对之毫无信心，反而看到由此产生的无数陷阱，尚且如此热切地渴望这些事物；倘若财富再增添安全和稳定的因素，他们会放过谁？他们会不掠取什么？什么寡妇？什么孤儿？什么穷人？\par

21. 因此，我们不要以为财富是极大的善；因为极大的善不在于拥有金钱，而在于拥有敬畏天主和一切虔诚。看，如果这里有一个义人，在天主面前大有勇气，即使他是凡人中最贫穷的，他也足以解除我们当前的祸患！因为他只需向天伸开双手，呼求天主，这乌云就会消散！但现在黄金大量囤积；却比泥土更无用，无法救我们脱离迫在眉睫的灾祸！不仅在这种危难中；如果疾病或死亡或任何此类灾祸降临到我们身上，财富的无能就完全显明，因为它不知所措，在这些事件中不能给我们任何慰藉。\par

22. 财富似乎有一事胜过贫穷，即它生活在每日的奢华之中，并在宴席上获得丰盛的快乐。然而，这同样也可以在穷人的餐桌上看到；而他们在此享受一种比富人更大的快乐。不要对此感到惊奇，也不要以为我所说的是悖论；因为我将通过事实的证据向你们阐明此事。你们当然知道，且都承认，在宴席中，引起快乐的通常不是食物的性质，而是享用者的心情；例如，当一个人饥饿时来到餐桌，即使是最普通的食物，也会比任何佳肴、调味品或千百种精致的菜肴更甜美。但那些不等待需要，或不先等到饥饿（如富人的习惯）就来到餐桌的人，尽管面前摆满了最精美的珍馐，却感觉不到快乐，因为他的食欲没有被事先激发。为使你们知道实际情况就是如此，除了你们都是见证之外，让我们听听圣经告诉我们同样的真理：经上说：\Quote{饱足的人厌弃蜂蜜，饥饿的人苦物也甜。}\BibleRef{箴 27:7} 然而有什么比蜂蜜和蜂巢更甜的呢？但他（指圣经）说，对于不饥饿的人，它并不甜。又有什么比苦物更苦涩的呢？然而对于贫穷的人，它们却是甜的。穷人是带着需要和饥饿来吃饭的，而富人并不等待这一点，我想每个人都知道。因此，他们得不到纯正而完全的快乐。不仅食物如此，任何人都能看出饮料也是如此；如同饥饿是快乐的原因，远胜于食物的品质一样，渴望也通常使饮品变得最甜美，即使饮用的只是水。这正是先知所暗示的，他说：\Quote{他以石中蜂蜜饱飨了他们。} 但我们并未在任何地方读到梅瑟从磐石中带出蜂蜜，而是在整个历史中读到河流、水泉和清凉的溪流。那么这是什么意思呢？因为圣经绝不说谎。正因为他们在干渴和疲惫中，发现这些水流如此清凉，为了显示这种饮品的快乐，他称水为蜂蜜，并非其本质变成了蜂蜜，而是因为饮用者的状况使这些水流比蜂蜜更甜。你看，口渴的状况如何使饮品变甜？是的，许多穷人常常在疲倦、困苦、焦渴时，甚至以我所说的那种快乐享用这样的水流。但富人饮用甜美的、有花香、具有酒所有优点的酒时，却体验不到这种快乐。\par

23. 每个人都能察觉，同样的事也发生在睡眠上。因为柔软的床榻、银饰的床架、整座房屋的安静，以及诸如此类的一切，通常都不如劳动、疲劳、睡眠的缺乏和躺下时的困倦更能使睡眠甜蜜愉快。关于这一点，事实的经验，甚至在事实经验之前，圣经的断言就作证了。曾经生活在奢华中的撒罗满，想要说明此事，就说：\BibleQuote{劳工的睡眠是甘饴的，无论他吃得少或多。}\BibleRef{训 5:12} 为什么他加上\Quote{无论他吃得少或多}？这两样通常都导致失眠，即贫乏与饮食过度：前者使身体干枯，使眼睑僵硬而不能闭合；后者使呼吸紧促压迫，引起许多痛苦。但劳动是如此有说服力，以致即使这两样临到他身上，仆人也能入睡。因为整个白天，他们在各地奔跑，服侍主人，被推搡压迫，几乎没有喘息的时间，但他们在睡眠的快乐中获得了辛劳和苦工足够的回报。这样，由于天主对人的仁慈，这些快乐不是用金银买得的，而是以劳动、辛苦、需要和各种锻炼获得的。富人则不然。相反，他们躺在床上，常常整夜无眠；尽管想出许多办法，却得不到那种快乐。但穷人从每日劳动中解脱，四肢完全疲乏，几乎还未躺下就沉入酣然、甘饴而真实的睡眠，享受着这并非微不足道的报酬，即他一日劳苦的果实。因此，既然穷人比富人睡眠、饮水、进食都更有乐趣，那么财富被剥夺了似乎胜过贫穷的那一点好处后，还有什么价值可言呢？为此，从起初，天主就使人劳动，不是为惩罚或管教，而是为改善和教育。当亚当过着不劳动的生活时，他失落了乐园；但当那位宗徒大量劳作、辛苦努力，说：\BibleQuote{在劳苦辛苦中，昼夜工作。}\BibleRef{得前 2:9} 时，他被提升到乐园，升到了第三层天！\par

24. 所以，我们不要轻视劳动；不要轻视工作；因为在天国之前，我们从劳动中获得最大的报酬，从中得到喜乐；不仅是喜乐，还有比喜乐更大的，最纯净的健康。因为除了缺乏胃口之外，许多疾病也攻击富人；但穷人却免于医生的手；即使有时他们患病，也很快康复，因为他们远离一切娇气，体格强健。对于明智地承受贫穷的人，贫穷是一笔巨大的财富，一个不可夺走的宝藏；最坚固的杖；一条无可阻挡的得利之路；一个安全的住所，无陷阱之虞。穷人受压迫吗？但富人更容易遭受恶意算计。穷人被轻视和侮辱。但富人被人嫉妒。穷人不像富人那样容易被攻击，因为富人在各方面给魔鬼和暗敌提供无数把柄；由于他的事务繁多，他成了众人的仆人。由于需要许多事物，他不得不奉承许多人，并极其奴性地服侍他们。但穷人如果知道如何灵性明智，甚至魔鬼本人也无法攻击他。因此，\Person{约伯}在此之前虽然坚强，但当他失去一切时，他变得更加强大，并从魔鬼那里赢得了辉煌的胜利！\par

25. 除此之外，穷人如果知道如何灵性明智，就不可能受到伤害。现在，关于快乐，我曾经说过，它不在于肴馔的昂贵，而在于进食者的心态；对于侮辱，我也这样说：侮辱的产生或消除，不在于侮辱者的意图，而在于承受者的心态。例如：有人用许多语言侮辱了你，不论适合不适合重复。如果你嘲笑这些侮辱，不把话放在心上，表明自己凌驾于打击之上，那么你就没有被侮辱。正如如果我们有金刚不坏之身，即使一千支箭从四面八方射来，我们也不会受伤，因为箭矢造成伤口，不在于投掷者的手，而在于承受者的身体；同样，在这里，侮辱之所以成为真实且有损名誉的，不在于施加者的愚昧，而在于被侮辱者的软弱。如果我们知道如何真正明智，我们就不会被侮辱，也不会遭受任何严重的伤害。有人侮辱了你，但你并没有感觉到？你没有痛苦。那么你就没有被侮辱，反而给了对方一击，而不是受到了打击！因为当侮辱者发现他的打击没有触及被辱骂者的灵魂时，他自己会更加烦恼；而被辱骂者沉默不语时，侮辱的打击就会反过来，自动反弹回攻击者身上。\par

26. 因此，亲爱的诸位，让我们在所有事情上都灵性明智，这样贫穷就不能伤害我们，反而会大大有益于我们，使我们比最富有的人更显赫、更富有。因为请告诉我，谁比\Person{厄里亚}更贫穷？但正因如此，他超越了所有富人，因为他如此贫穷，而且这种贫穷是他从心灵富足中自愿选择的。因为他认为所有财富的丰裕都不及他的慷慨大度，也不值得他的灵性智慧，所以他接受了这种贫穷；因此，如果他看现世事物有什么价值，他就不会只拥有一件外衣。但他如此藐视今生的虚荣，视所有黄金如街上的泥土，以至于他除了那件衣服外一无所有。因此，国王需要那个穷人，拥有那么多黄金的国王却依赖那个只有一张羊皮的人。因此，羊皮比紫袍更辉煌，义人的山洞比君王的宫殿更灿烂。因此，当他升天时，他留给门徒的只有那张羊皮。\Quote{凭此，}他说，\Quote{我与魔鬼搏斗；拿起这个，你要武装起来对抗他！} 因为贫穷是一件强大的武器，一个不可攻破的避难所，一座不可动摇的堡垒！\Person{厄里叟}继承了羊皮作为最大的遗产；它确实是如此；比所有黄金更宝贵。从此以后，那个\Person{厄里亚}成了双重人物：天上的\Person{厄里亚}和地上的\Person{厄里亚}！我知道你们认为那个义人是有福的，你们每个人都希望成为那个人。如果我告诉你们，我们中间所有领受洗礼的人，所领受的比他还要大得多呢？因为\Person{厄里亚}留给门徒一张羊皮，但天主子升天时留给我们祂自己的身体！\Person{厄里亚}确实在他升天之前脱下了他的外衣；但基督为了我们的缘故留下了它；然而祂升天时还保留着它。所以，我们不要沮丧。不要哀叹，也不要害怕时代的艰难，因为那为所有人倾流祂的血，并让我们再次分享祂的体和血的主，祂会拒绝为我们的安全做什么呢？因此，怀着这些希望，让我们不断恳求祂；让我们在祈祷和恳求中热切；让我们严格地关注一切其他德行；这样我们就能逃脱当前的危险，并获得将来的美善。愿天主恩赐我们所有人都能配得上这些，因我们的主\Person{耶稣基督}的恩宠和慈爱，藉着祂，并同祂，偕同圣神，光荣归于圣父，直到永远。阿们。\par

\SourceRecord{Translated by W.R.W. Stephens.；Nicene and Post-Nicene Fathers, First Series；Vol. 9.；Edited by Philip Schaff.；Buffalo, NY: Christian Literature Publishing Co.,；1889.；<http://www.newadvent.org/fathers/190102.htm>.}

% structure: p0001 source=h1 target=section reason=page-title

\section{雕像讲道集 第三讲}

\Emph{论\Person{Flavian}的离去，他是\Place{安提约基雅}的主教，已为使节前往\Person{德奥多西}皇帝，为城市求情。论司祭职的尊严。什么是真正的斋戒。诽谤比吞噬人体更恶劣。最后，论那些因叛乱而被处死的人，并反驳那些抱怨许多无辜者被逮捕的人。}\par

1. 当我注视那空着的宝座，失去了我们的导师，我同时欢喜又哭泣。我哭泣，因为我看到我们的父亲不在我们中间！但我欢喜，因为他已踏上为我们保全性命的旅程；他要去从皇帝的愤怒中拯救如此众多的群众！这对你们是荣耀，对他则是冠冕！对你们是荣耀，因为你们被分派了这样一位父亲；对他是冠冕，因为他如此疼爱他的子女，并以实际行动证实了基督所说的。因为他得知\Quote{善牧为羊舍命}\BibleRef{若 10:11}，便离开了；为我们众人冒险献出他的生命，尽管有许多事情阻碍他离去，迫使他留下。首先，他的年岁已至暮年的极限；其次，身体的虚弱和一年的节气，以及他必须在神圣节庆时在场；此外，他的独生妹妹此刻已奄奄一息！然而，他轻看了亲族、年迈、虚弱、季节的严酷以及旅途的劳顿；他将你们和你们的安全置于一切之上，冲破了所有这些束缚。并且，如同少年一般，这位老人如今正加速前行，被热忱的翅膀托举！因为基督（他说）为我们舍己，我们既然担负了如此众多民众的职责，若不准备为托付给我们的人的安全做任何事、受任何苦，我们还能有何借口或宽恕呢？因为（他继续说）圣祖雅各伯在管理羊群、牧养无理性的羊，并向人交账时，彻夜不眠，忍受酷暑严寒和一切气候的严酷，唯恐有一头牲畜失丧；何况我们这些管理并非无理性而是属神的羊群，且要向天主而不是向人交账的人，岂能在任何方面懈怠，或畏缩不前，不做任何有益于羊群的事呢？此外，正如后者比前者更尊贵（人比牲畜，天主比人），我们也应当表现出更大、更热切的关心和勤勉。他深知，他现在所关心的不只是一座城市，而是整个东方。因为我们的城市是东方所有城市之首和母亲。为此，他愿面对一切危险，没有任何事能将他留在此地。\par

2. 因此，我相信有美好的希望；因为天主不会不屑于垂视这样的热诚和勤奋，也不会让祂的仆人空手而回。我知道，当他刚一见到我们虔诚的皇帝，并被皇帝见到时，他就能立刻凭他的面容平息皇帝的愤怒。因为圣人们不仅言语，连面容都充满恩宠。他本身也赋有丰富的智慧，并且精通神圣法律，他将像梅瑟对天主说的那样对他说：\Quote{如今，求你赦免他们的罪——不然，就把我连同他们一起杀掉。}\BibleRef{出 32:31} 因为圣人们的胸怀是这样的：他们认为与子女同死比无子女的生存更甜蜜。他还会以这特殊的时节为辩护，并藏身于逾越节的神圣庆节之后；他要提醒皇帝，基督正是在这个时节赦免了全世界的罪。他要劝勉皇帝效法他的主。他还要提醒他那关于一万塔冷通和一百德纳的比喻。我知道我们父亲的胆识，他必会毫不犹豫地用这个比喻警告皇帝，说：\Quote{你要小心，免得在那一天你也听到这样的话：\InnerQuote{你这恶仆，因为你求了我，我就免了你所有的债；你也应当宽免你的同仆。}\BibleRef{玛 18:32}} 你这样做对自己比对他们更有益处，因为通过宽免这些轻微的冒犯，你赢得了对更大过犯的赦免。在这番话中，他还要加上那在神圣奥迹中入门时所教给他的祈祷，并说：\Quote{宽免我们的罪债，如同我们也宽免得罪我们的人。}\par

3. 他还要告诉他，这罪过不是全城共犯，而是某些陌生人和冒险者的行为，这些人行事毫无计划，而是充满各种胆大妄为和无法无天；少数人的混乱行为不应成为铲除如此伟大城市的理由，也不应惩罚那些未曾作恶的人；即使全城都是犯法者，他们也已受了足够的惩罚——那么多天被恐惧吞噬，每天都等待被处死，流亡逃难；活得比已定罪的犯人更悲惨，手中握着生命，毫无逃脱的希望！\Quote{让这惩罚（他会说）足够吧。不要再继续你的愤怒！以你对待同仆的仁慈，使天上的审判者对你宽恕！想想这城市的伟大，现在的问题不是关于一个、两个、三个或十个灵魂，而是无法计数的广大民众！这是关乎全世界首都的问题。这是基督徒最先被称为基督徒的城市。尊崇基督。尊敬那最先宣告这个为所有人所喜爱和甜蜜名字的城市！这城市曾是宗徒的帐幕，义人的居所！而现在这是首次也是唯一一次对抗其统治者的起义；所有过去的时间都将为这座城市的风俗作证。因为如果民众持续叛乱，也许有必要以这种罪恶作为警戒；但如果这在全部时间中只发生一次，那么显然这罪行并非出于城市的习惯，而是那些在不幸时刻纯粹偶然来到此处之人的过犯。}\par

4. 这些事以及更多的事，司铎将更大胆地说出来；皇帝会听从他们；一个仁慈，另一个忠信；所以我们两方面都抱有美好的期望。但与其依靠我们导师的忠信和皇帝的仁慈，我们更依赖天主的仁慈。因为当皇帝被恳求时，司铎在恳求时，祂自己会介入，软化皇帝的心，激发司铎的舌头；使他的说话流利——预备另一方的心思以接受所说的话，并带着极大的宽厚应允恳求。因为我们的城市在基督眼中比所有其他城市更宝贵，既由于我们祖先的美德，也由于你们自己的美德。正如\Person{伯多禄}在宗徒中最先宣讲基督，同样，如我之前所说，这城市是诸城中最先以基督徒的称号装饰自己的，如同一种卓越的冠冕。但如果只有十个义人之处，天主就应允拯救其中所有居民，那么当我们不仅有十个、二十个或两倍之多，而且有更多——他们以全部严谨侍奉天主——时，我们为何不期待一个有利的结果，并确信我们所有人的生命得保全呢？\par

5. 我听见许多人说：\BibleQuote{君王的威吓有如狮子的怒啸}\BibleRef{箴 19:12}，他们充满了沮丧和哀叹。那么我们应该对这些说什么呢？就是那位说过\BibleQuote{豺狼与羔羊将一起牧放，豹子将与小山羔同卧，狮子将像牛一样吃草}\BibleRef{依 11:6}的，有能力将狮子变为温顺的羔羊。让我们因此恳求祂；让我们派遣使者到祂那里；祂必会平息皇帝的愤怒，救我们脱离迫在眉睫的苦难。我们的父亲已奉命前往那里。让我们从这里派遣使者到天上的尊威！让我们用祈祷帮助他！教会的团体能成就大事，如果我们以忧伤的灵魂和痛悔的心情献上祈祷！无需渡海，也无需长途跋涉。我们中的每个男人和女人，无论在教堂中聚会，还是留在家里，都要极其诚恳地呼求天主，祂必会应允这些恳求。\par

这从哪里显得明显呢？因为祂极其渴望我们常常投奔祂，在一切事上向祂祈求；没有祂不行任何事、不说任何话。因为人们，当我们反复用我们的事务打扰他们时，就变得懒惰、回避，并对我们态度不悦；但天主恰恰相反。不是当我们不断为事务祈求祂时，而是当我们不这样做时，祂特别不悦。至少听听祂责备犹太人时所说的话：\BibleQuote{你们商量，却不靠我；你们立约，却不靠我的圣神。}\BibleRef{依 30:1}因为这是爱者的习惯：他们渴望所爱者的一切事都由他们自己完成，并且没有他们就不做任何事、不说任何话。为此，天主不仅在那时，而且在别处也以同样的责备说话：\BibleQuote{他们自立为王却不靠我，他们自订计划却不让我知道。}\BibleRef{欧 8:4}所以我们不要迟缓地不断投奔祂；无论是什么邪恶，它终将找到合适的解决之道。\par

6. 难道有人使你惊惶吗？赶快投奔主，你必不受任何恶害。古人就是这样从灾祸中获得解脱的；不仅男子，妇女也是如此。有一位希伯来妇女，名叫艾斯德尔。这位艾斯德尔正是用这种方法，在犹太人即将被交付毁灭时，拯救了整个犹太民族。因为当波斯王下令要彻底消灭所有犹太人，而无人能阻挡他的愤怒时——这位妇女脱下华丽的衣袍，穿上苦衣并撒上灰烬，恳求仁慈的天主与她一同去见王；她向天主献上祈祷，说了这些话：\BibleQuote{主，请使我的言语蒙悦纳，在我口中放上动听的话。}愿这也是我们为我们的神师向天主献上的祈祷。因为如果一位妇女为犹太人恳求，能平息野蛮人的愤怒，那么我们的神师为如此伟大的城市，并与如此伟大的教会一同恳求，岂不更能说服这位最温和仁慈的皇帝吗？因为如果他有权赦免对天主所犯的罪过，那么他岂不更能除去并涂抹那些冒犯他人的罪过吗？他本人也是一位统治者，并且是比另一位更高贵的统治者。因为神圣的律法甚至将君王的头颅也置于他的手下。每当需要从上而来的任何恩惠时，皇帝惯于投奔司铎；而非司铎投奔皇帝。司铎也有他的护胸甲，即正义的护胸甲。他也有腰带，即真理的腰带，以及更为尊贵的鞋子，即和平福音的鞋子。他也有剑，不是铁剑，而是圣神的剑；他也有冠冕戴在头上。这套全副武装更加辉煌。武器更宏伟，言辞更自由，力量更强大。因此，凭借他权威的分量和他自己灵魂的伟大，更超乎一切的是，凭借他对天主的希望，他将以极大的自由和极大的谨慎对皇帝说话。\par

7. 因此，让我们不要对我们的得救绝望，而是让我们祈祷；让我们呼求；让我们恳求；让我们以许多眼泪作为大使前往天上的君王那里！我们还有这斋戒作为盟友，作为这美好代祷的助手。因此，正如当冬天过去、夏天来临时，水手将船拖入深海；士兵擦亮武器，备好战马待战；农夫磨利镰刀；旅行者勇敢踏上长途旅程；摔跤手脱衣赤身准备比赛。同样，当斋戒到来时，如同一种属灵的夏天，让我们作为士兵擦亮武器；作为农夫磨利镰刀；作为水手约束思想以抵御过度欲望的波浪；作为旅行者踏上通往天堂的旅程；作为摔跤手脱衣准备比赛。因为信徒同时是农夫、水手、士兵、摔跤手和旅行者。因此圣保禄说：\BibleQuote{我们战斗不是对抗血和肉，而是对抗率领者，对抗掌权者。所以要穿上天主的全副武装。}\BibleRef{厄 6:12}你看见摔跤手了吗？你看见士兵了吗？如果你是摔跤手，就必须赤身投入战斗。如果你是士兵，就必须全副武装站在战线上。那么这两件事如何可能：赤身却不赤身？穿着却未穿着？如何？我告诉你。脱去世俗的事务，你就成了摔跤手。穿上属灵的武装，你就成了士兵。脱去世俗的忧虑，因为这是摔跤的季节。穿上属灵的武装，因为我们与魔鬼有激烈的争战。因此，我们需要赤身，以免在摔跤时给魔鬼任何可抓之处；并且要全副武装，以免从任何方面受到致命打击。耕耘你的灵魂。砍去荆棘。播下虔敬的言语。用许多关怀培育并滋养智慧之美的佳苗，你就成了农夫。圣保禄会对你说：\BibleQuote{劳苦的农夫理当先得果实。}他本人也实践了这门技艺。因此，在写给格林多人的书信中，他说：\BibleQuote{我栽种了，阿颇罗浇灌了，然而使之生长的却是天主。}\BibleRef{格前 3:6}磨利你因贪饕而变钝的镰刀——借着斋戒磨利它。踏上通往天堂的道路；尽管崎岖狭窄，仍要踏上它，并继续前行。你怎能做到这些呢？通过克制你的身体，使它服从。因为当道路变窄时，贪饕带来的肥胖是极大的障碍。抑制过度欲望的波浪。驱散恶念的风暴。保全小船；展现技巧，你就成了舵手。但我们将有斋戒作为这一切的基础和导师。\par

8. 我所说的，并非多数人所守的那种斋戒，而是真正的斋戒；不仅是戒肉，更是戒罪。因为斋戒的本质是，除非按相应的法则进行，否则不足以使实行者得救。因为\Quote{角力者}，经上说，\Quote{若不按规矩竞赛，就不能得冠冕。}\BibleRef{弟后 2:5} 因此，为了使我们在经历了斋戒的劳苦之后，不致丧失斋戒的冠冕，我们应当明白，必须如何并以何种方式来进行这事；因为那法利塞人也曾守斋，\BibleRef{路 18:12} 但后来却空手而归，失去了斋戒的果实。税吏并未守斋，却比那守斋的更蒙悦纳；这是为使你们知道，除非其他一切义务随之而行，斋戒是无益的。尼尼微人守了斋，赢得了天主的眷顾。\BibleRef{纳 3:10} 犹太人也曾守斋，却毫无益处，反而带着责罚离去。既然对于不知道应如何守斋的人而言，斋戒的危险如此之大，我们就应当学习这操练的法则，以免我们\Quote{盲目奔跑}，\Quote{打空气}，或在争战时与影子搏斗。斋戒是一剂良药；但这剂药，即使再有益，也常常因使用者的无能而变得无用。因为还必须知道何时应用它，所需的剂量，以及适宜的体质、地区性质、季节、相应的饮食，以及其他许多细节；其中任何一项若被忽略，就会破坏所有其余的部分。如果当身体需要医治时，我们需要如此精确，那么当我们关心灵魂、寻求医治心灵的疾病时，岂不更应以最大的精确度去观察和探求每一个细节吗？\par

9. 那么，让我们看看尼尼微人是如何守斋的，以及他们如何从那义怒中被解救出来——\Quote{无论人畜、牛群、羊群，都不可尝什么，}\BibleRef{纳 3:7}（先知）说。你说什么？告诉我——难道连非理性的受造物也要守斋，马和骡子都要披上麻衣吗？\Quote{也应如此，}他答道。因为正如当某位富翁去世时，亲戚们不仅给男仆女仆披上麻衣，也给马匹披上麻衣，吩咐它们由马夫牵着跟随送葬队伍到坟墓，以此表示灾祸的重大，并邀请众人怜悯；同样，当那城即将被毁灭时，连非理性的受造物也被裹上麻衣，被置于斋戒的轭下。\Quote{这是不可能的，}他说，\Quote{非理性的受造物不能借着理性认识天主的义怒；让他们借着斋戒来学习，使他们知道这是天主的打击。因为如果这城被倾覆，它不仅会成为我们这些住在其中者的共同坟墓，也会成为它们的坟墓。既然它们要分担惩罚，也让它们分担斋戒吧。}但在这行为中，他们还有一个目的，也是先知们常做的。因为这些先知，当他们看到从天上降下可怕的惩罚，而受罚者无言以对，满负羞耻，不配得任何宽恕或借口——不知如何是好，也不知从何处为被定罪者求得辩护时，他们就借助非理性的受造物，以悲剧的方式描述它们的死亡，借它们来转求，以它们可怜可悲的毁灭作为恳求的理由。因此，从前当饥荒降临于犹太人，大旱压迫他们的土地，一切都被消耗时，有一位先知这样说：\Quote{幼小的母牛在栏中跳跃；牛群哭泣，因为没有草场；田野的一切牲畜都仰望你，因为溪水干涸了。}\BibleRef{岳 1:17} 另一位先知哀叹干旱的灾祸，又这样说：\Quote{母鹿在田野产仔，却弃之不顾，因为没有草；野驴站在林中，像龙一样喘气，它们的眼睛失明，因为没有草。}\BibleRef{耶 14:5} 此外，你们今天也听见岳厄尔说：\Quote{新郎从洞房里出来，新娘从内室出来，还有吃奶的婴儿。}请问，他为什么叫这样年幼的年龄来祈求？岂不是明明出于同样的原因？因为所有达到成年人年龄的，都激起了天主的义怒，所以他说，让那没有过犯的年龄向被激怒的天主祈求吧。\par

10. 但是，如前所述，我们来看是什么化解了那不可挽回的义怒。仅仅是斋戒和麻衣吗？我们说不，而是他们整个生命的改变。这从何可见？从先知的话本身。因为他论述了天主的义怒和他们的斋戒，当他论及和好并教导我们和好的原因时，这样说：\Quote{天主看见了他们的作为。}\BibleRef{纳 3:10} 什么作为？是他们守了斋？是他们披上了麻衣？都不是；他略过这些，补充说：\Quote{他们各人离开了自己的恶行，上主就后悔了，不把所说的灾祸降在他们身上。}你们看见了吗？斋戒并没有解救他们脱离这危险，而是生命的改变使天主对这些外邦人仁慈和善待。\par

11. 我说这些话，并非要贬抑禁食，而是要尊崇禁食；因为禁食的尊荣不在于戒除食物，而在于远离罪恶的行为；那将禁食仅限于戒除肉食的人，便是特别贬抑它。你禁食吗？以你的行为来证明！你问是什么样的行为？若你看见穷人，就怜悯他！若你看见仇敌，就与他修好！若你看见朋友得荣耀，不要嫉妒他！若你看见美貌女子，就避开她！因为不仅口要禁食，眼、耳、脚、手以及我们身体的所有肢体也要禁食。让双手禁食，远离掠夺和贪婪。让双脚禁食，停止奔向那些不合法的表演。让双眼禁食，学会绝不粗鲁地注视美貌的面容，或沉迷于陌生的美色。因为观看是眼睛的食物，但若这食物是不合法或被禁止的，它就破坏了禁食，并扰乱了灵魂的全部安全；但若是合法且安全的，它便装饰了禁食。因为最荒唐的事莫过于为了禁食而戒除合法的食物，却用眼睛去触摸被禁止的东西。你不吃肉吗？那就不要通过眼睛滋养淫荡。也要让耳朵禁食。耳朵的禁食在于拒绝接受恶言和诽谤。经上说：\Quote{你不可接受虚假的报告。}\par

12. 也要让口禁戒可耻的言语和辱骂。因为我们戒除鸟和鱼，却咬噬吞吃弟兄，有什么益处呢？毁谤者吃弟兄的肉，咬邻人的身体。因此，保禄说出这句可怕的话：\Quote{你们若互相咬噬吞吃，要小心，免得彼此消灭。}\BibleRef{迦 5:15}你没有把牙齿钉在肉里，却把诽谤钉在灵魂里，造成了恶意猜疑的创伤；你在一千种方式中伤害了自己、他和许多其他人，因为诽谤邻人时，你使听诽谤的人变得更糟：倘若他是个恶人，他找到恶行的同伴就更不在乎；倘若他是个义人，他就被引向傲慢和自大，因他人的罪而自以为了不起。此外，你打击了教会的共同利益；因为所有听到的人不仅指责那所谓的罪人，而且这羞辱也加在基督徒团体上；你听到不信的人不是说：\Quote{某人是奸淫者或放荡者；}而是，代替那犯罪的个人，他们指责所有基督徒。除此之外，你使天主的荣耀受亵渎；因为正如我们有好名声时祂的名受荣耀，同样，当我们犯罪时，它被亵渎和侮辱！\par

13. 第四个理由是，你羞辱了那位被恶名所传的人，并因此使他比以前更无耻，因为你把他置于敌意和对抗之中。第五，你使自己遭受惩罚和报复，因为你卷入了与你毫不相干的事。因为不要有人回答说：\Quote{那么，我说假话时才是毁谤者，但若我说真话，就不是了。}即便你说的是真话，这也是罪过。因为那个法利塞人确实说真话毁谤了税吏；但这对他毫无益处。我问，难道后者不是税吏和罪人吗？每个人都知道他是税吏。但与此同时，因为法利塞人说了他的坏话，他便失去了一切益处，离开了圣殿。你希望纠正一个弟兄吗？哭泣；向天主祈祷；把他带到一边，劝告、规劝、恳求他！保禄也这样做了，他说：\Quote{Lest,}他说，\Quote{等我再来时，我的天主使我在你们面前谦卑，我要为许多从前犯罪、还没有悔改他们所做的不洁、淫乱和放荡而哀哭。}\BibleRef{格后 12:21}向罪人显示你的爱德。说服他，提醒他的罪是出于对他得救的关心和焦虑，而非出于揭露他的愿望。抱住他的脚；拥抱他；你若真心要医治他，就不要羞惭。医生们也常做这类事，当他们的病人难以取悦时，通过拥抱和恳求，他们最终说服病人服用有益的药剂。你也这样做。向司铎展示伤口；这是关心、照顾他，并为他焦虑者的行为。\par

14. 如今我不仅劝戒那些说恶言的人，也劝戒那些听别人被毁谤的人，要堵住耳朵，效法那先知所说：\Quote{暗中毁谤邻人的，我必惩治。}要对你的邻人说：\Quote{你有人可称赞或极力推崇吗？我张开耳朵，接受芬芳的油；但若你有恶言要说，我堵住你话语的入口——因我不接纳粪土和污秽。这给我带来什么益处呢？得知某人是坏人，最大的伤害和损失莫过于此！}要对他说：\Quote{让我们为自己的过错焦虑，思想如何为自己的过犯交账，并对自己的生活表现出这种好奇和多管闲事。我们从不考虑自己的事，却这样好奇地窥探别人的事，那还有什么借口或宽恕可寻呢？}正如窥探别人的房屋，路过时观察里面是什么，是卑鄙至极的；同样，探究别人的生活，也是极度的不雅。但更可笑的是，那些过着这种生活、忽略自己事务的人，当提到任何这些秘密之事时，竟恳求并嘱咐听者不要再向任何人提起，从而表明他们做了该受责备的事。因为你若恳求他不要告诉别人，那你当初就不该先告诉他。那事在你手中本是安全的；如今你泄密之后，却为它的安全担忧。若你希望它不被传到别人那里，就自己不要说。而当你把保管之责交托给别人时，再嘱咐他、叫他发誓为所说之话保密，便是多余无益的了。\par

15. \Quote{但毁谤是甘甜的。}不，不说恶言才是甘甜的。因为说了恶言的人从此好争辩，他猜疑、恐惧、后悔、咬自己的舌头。他胆战心惊，唯恐他所说的话何时传到别人耳中，给说话者带来大危险和毫无益处、不必要的敌意。但守口如瓶的人，必满有平安和喜乐地度日。经上记载：\BibleQuote{你听见了话语，让它与你同死；放心，它不会胀破你。}\BibleRef{德 19:10}这是什么意思？\Quote{让它与你同死？}就是熄灭它，埋葬它，不容它出去，甚至不让它动；最好的办法是小心不要容忍别人行毁谤之事。若你偶受其影响，就埋葬它，消灭所说的话，使之遗忘，好使你如同未曾听见的人，在极大的平安与安稳中度此世生活。若毁谤者知道我们比被他们控告的人更厌恶他们，他们自己便会放弃这恶习，纠正罪过；之后必称赞并宣扬我们是他们的救星和恩人。因为，如同说好话、称赞是友谊的开端，同样，说坏话、毁谤是敌对、仇恨和无数争吵的开端和基础。我们自己的事之所以更被忽略，没有别的原因，正是由于窥探和干预他人之事的习惯；因为一个惯于毁谤、忙于别人生活的人，不可能照顾自己的生命。他把全部心思都花在干预别人的事上，因此他自身的一切必然被抛在一边、无人照管。因为一个人若将自己所有的闲暇都用于焦虑地省察自己的罪、并加以判断，尚且很难进步；而你总是忙于别人的事，何时才会留意自己的恶呢？\par

16. 所以，亲爱的，让我们逃避毁谤！我们知道它就是撒殚的深渊，是他设网罗埋伏的地方。为使我们忽略自己的状况，从而使我们的罪责更沉重，魔鬼就引领我们陷入这习惯。不仅如此，这不仅是一件非常严重的事——我们将来要为所说的话交账——而且我们还会因这些手段使自己的过犯更沉重，剥夺自己的一切借口。因为一个刻薄审视别人行为的人，永远无法为自己所犯的罪获得宽恕。因为天主要按我们过犯的性质，也按你评判别人的判断来定案。因此他告诫说：\BibleQuote{你们不要判断，以免受判断。}\BibleRef{玛 7:1}因为无论何种罪，在那里将不再显得像当初所犯的那样，而是会因你对同你所判的罪而得到重大且不可宽恕的加增。正如仁慈、怜悯、宽恕的人剪除了自己大部分罪过，同样，刻薄、残忍、不宽容的人也大大增加自己过犯的严重性。因此，让我们从口中除去一切毁谤，知道如果我们不戒除它，即使我们吃灰，这刻苦也于我们无益。\BibleQuote{因为不是入口的，而是出口的，才污秽人。}\BibleRef{玛 15:17}假如有人翻搅粪池，当你经过时，你岂不责骂那做这事的人吗？对毁谤者也当如此。因为翻搅的粪池使闻者嗅觉所受的冒犯，远不及翻搅别人的罪过、暴露不洁的生活，对听闻者的灵魂所造成的冒犯和扰乱。因此，让我们戒绝毁谤、污言秽语、亵渎；不要口出恶言攻击邻人，也不要攻击天主！\par

17. 因为我们的许多恶言者已陷入如此疯狂，竟从同仆身上举起舌头攻击他们的主人。但这是多么大的恶事，你可以从我们目前所卷入的事件中得知。一个人受了侮辱，看！我们都恐惧战栗，无论是那些犯下侮辱之罪的人，还是那些全然不知情的人！但天主每天都受侮辱！为什么我说每天？——其实是每时每刻，受富人、穷人、安逸者、受苦者、诽谤者、被诽谤者的侮辱，然而从未有人提及此事！因此，祂容许我们的同仆被侮辱，为的是藉由这次侮辱所发生的危险，让你学到主的慈祥！而且，尽管这是我们第一次且唯一的冒犯，我们并不因此期望获得开脱或宽恕。但我们每天都激怒天主，毫无悔改的迹象，而祂却以极大的忍耐容忍这一切！你看到主的慈祥有多么伟大吗？然而，在这次暴行中，那些做错事的人被抓住投入监狱，并付出了代价；但我们仍然害怕，因为被侮辱的人尚未知晓所发生的事，也未宣判，我们都在颤抖。但天主每天听到对祂的侮辱，却无人理会，尽管天主是如此的慈悲和爱人。在祂那里，只需承认罪过，就可以撤销控告。但在人那里则完全相反。当犯罪者承认时，他们反而受到更重的惩罚；这正是在当前事例中所发生的。有些人死于刀剑，有些人死于火刑；有些人被喂野兽，不仅是成人，还有儿童。无论是年龄的幼小、民众的骚乱、他们是在魔鬼的驱使下做出这些行为、还是苛捐杂税被认为难以忍受、贫穷、与众人一同犯罪、承诺以后绝不再犯，还是其他任何理由，都无法拯救他们；他们被毫无缓刑地押往刑场，武装的士兵在两边押送和看守，以防有人劫走罪犯；母亲们远远跟随，看着自己的孩子被斩首，却不敢哀悼自己的不幸；因为恐惧胜过了悲伤，害怕胜过了天性！就像人们从陆地上看到那些遭遇海难的人，深感痛苦，却无法接近并拯救溺水者一样，在这里，母亲们因害怕士兵而受束缚，如同被许多波浪阻隔，不仅不敢走近孩子并拯救他们免于定罪，甚至不敢流泪。\par

18. 你们确实由此看到天主的慈悲，是何等不可言喻、何等无限、何等超越一切描述！在这里，被侮辱者与侮辱者是同一本性；而且一生中只有一次经历此事；并且不是当面做的，也不是在他目睹或耳闻时发生的；然而，那些行恶者无人得到宽恕。但就天主而言，绝无此类事情可说；因为人与天主之间的差距，大得任何语言都无法表达；而且祂每天都受侮辱，尽管祂就在现场，看见并听到一切：然而祂不发出闪电，也不命令海水泛滥陆地淹没所有人；也不命令大地裂开吞没所有侮慢者；但祂容忍、忍耐，并且仍然愿意宽恕那些侮辱祂的人，只要他们悔改并承诺不再做这些事！现在正是宣告的时刻：\Quote{谁能述说主的大能？谁能宣扬祂的一切赞美？} 有多少人不仅推倒，而且践踏了天主的肖像！因为当你勒紧债务人、剥光他、拖走他时，你就在践踏天主的肖像。请听保禄确实说过，\BibleQuote{男人不应蒙头，因为他是天主的肖像和光荣。}\BibleRef{格前 11:7} 又听天主亲自说：\BibleQuote{让我们照我们的肖像，按我们的模样造人。}\BibleRef{创 1:26} 但若你说人并非与天主同体——那有什么关系？因为铜像也与皇帝不同体；然而，藐视铜像的人仍受到了惩罚。同样，对于人类，即使他们不与天主同体（事实上确是如此），他们仍被称为祂的肖像；且因这名号，理当受到尊敬。但你为了一点金子，就践踏他们、勒紧他们、拖走他们；而至今你仍未受到任何惩罚！\par

19. 但愿此刻迅速有某种有利而吉利的转变！我确实预先警告并作证说，即使这片乌云消散，而我们仍处于同样懈怠的状态，我们将不得不承受比现在所害怕的更重的灾祸；因为我并不害怕皇帝的愤怒，更不如说害怕你们的懈怠。当然，仅仅祈求两三天是不够的作为补偿，而是必须改变我们的整个生活，并且当戒避邪恶时，应持续坚持德行。因为正如体弱的人，若不坚持持久的养生法，则从两三天饮食节制中得不到任何益处；同样，那些在罪恶中的人，若不时刻保持节制，也不会从两三天悔改中得到任何益处。因为正如经上所说，那洗了澡，后来又沾染污泥的人，什么也没得到；同样，那忏悔了三天，又回到先前状态的人，什么也没成就。因此，现在我们不要像从前那样行事。因为许多次，当地震、饥荒和干旱突然降临，我们在变得清醒温和三四天后，却又回到旧路上。正因如此，我们当前的患难发生了。但即使我们以前没有这样做，如今至少让我们所有人坚持同样的虔敬；让我们保持同样的温良，以免我们需要再次受打击。难道天主不能阻止所发生的事吗？祂却容许它发生，为使那些轻视祂的人，因敬畏同仆而更加清醒！\par

20. 但不要让任何人说许多有罪的人逃脱了，许多无辜的人受了惩罚。因为我听到许多人常这样说；不仅在这次叛乱中，而且在许多其他类似情况下。那么我应该对那些作此观察的人如何回答呢？那就是，如果被捉的人在这次叛乱中是无辜的，那么他在此前犯有其他更严重的过犯，由于他事后没有悔改，就在此时受到了惩罚。因为天主对待我们的方式正是如此。当我们犯罪时，祂并不立即惩罚过犯，而是放过它，给我们悔改的时间，好使我们能改正和皈依。但如果我们因为没有受罚，就认为罪行也被抹去了，并轻视它；那么在某个我们意想不到的地方，我们之后必定会受罚。这样做的目的是，当我们犯罪而不受罚时，我们不应无所畏惧，除非我们改正，因为我们知道我们必会在意想不到之处落入惩罚。因此，亲爱的，如果你犯罪而没有受罚，不要变得傲慢，反而因此更要警惕，因为知道天主乐意时再施报是一件容易的事。因此祂没有惩罚你，好使你获得悔改的时间。因此，我们不要说，某人无辜却受罚；另一个有罪却逃脱，因为那受罚的人虽如我所说是清白的，却为别的过犯付了代价；而那位现在逃脱的人，若不悔改，将在另一个网罗中被捉。如果我们这样心怀意念，我们就永远不会忘记自己的罪，反而永远恐惧战栗，生怕必须受罚，从而容易回忆它们。因为没有什么比惩罚和管束更能使罪被记起。这一点由\Person{若瑟}的兄弟们表明。因为当他们出卖了那位义人，十三年过去后，他们怀疑自己陷入了惩罚，并担心性命，就想起了自己的罪，彼此说：\Quote{我们实在有罪，对不起我们的弟弟若瑟。} \BibleRef{创 42:21} 你们看到恐惧如何使他们的罪过被记起吗？当他们犯罪时，他们未曾察觉；但当他们害怕受罚时，他们就记起来了。因此，知道这一切之后，让我们改变并修正我们的生活；让我们在想到脱免迫在眉睫的苦难之前，先想到宗教和德行。\par

21. 同时，我愿将三条诫命铭刻在你们心中，使你们在四旬期内能为我完成它们：即不诽谤任何人，不以任何人为仇敌，并彻底戒绝起誓的恶习。正如当我们听说课征某种银钱税时，每个人回到家中，召集妻子、儿女和仆人，考虑并商议如何缴纳这笔税款；同样，让我们也如此对待这些属神的诫命。让每个人回家后，召集妻子和儿女，对他们说：\Quote{今日课征了一项属神的税，靠着这税将获得某种解脱并除去这些灾祸；这税不会使缴纳者变穷，反而更富足；也就是说，不以任何人为仇敌，不诽谤任何人，并且完全不发誓。}让我们省察，让我们思考，让我们决意如何履行这些诫命。让我们竭尽全力。让我们彼此劝诫。让我们互相纠正，免得我们以负债者的身份进入来世，然后需要向他人借贷，遭受愚拙贞女的命运，并从不朽的救恩中坠落。若我们如此整顿我们的生活，我向你们保证并承诺，从此将摆脱当前的灾难，除去这些可怕的祸患；而比一切更大的，是享受未来美好的事物。因为本该将全部德行托付给你们；但我认为最好的改正方法是，将律法分部分地接受并付诸实践，然后再进行其他部分。因为正如在一块田地中，农夫逐渐地挖遍所有地方，慢慢完成他的工作；同样，若我们为自己定下规矩，无论如何要在本四旬期内正确实践这三条诫命，并将它们安全地交托给良好习惯的保管，我们就能更轻松地进行其余部分；并且借此办法达到神性智慧的顶峰，我们在此生将收获希望之善果；在来生，我们将满怀信心地站在\Person{基督}面前，享受那不可言喻的福乐。愿天主恩赐，使我们众人都堪当获得这福乐，赖\Person{我们的主耶稣基督}的恩宠和慈爱，愿光荣归于\Person{圣父}及\Person{圣神}，至于无穷世之世。阿们。\par

\SourceRecord{Translated by W.R.W. Stephens.；Nicene and Post-Nicene Fathers, First Series；Vol. 9.；Edited by Philip Schaff.；Buffalo, NY: Christian Literature Publishing Co.,；1889.；<http://www.newadvent.org/fathers/190103.htm>.}

% structure: p0001 source=h1 target=section reason=page-title

\section{第四篇关于雕像的讲道}

\Emph{对民众关于坚毅和忍耐的劝勉，以约布和巴比伦三童的事迹为例。讲道以关于戒绝起誓的讲话结束。}\par

1. 愿天主受赞颂！祂安慰了你们忧伤的心灵，安抚了你们激动的情绪！因为你们现在所表现出的聆听渴望和热忱，证明了你们获得了不小的安慰。因为在痛苦中、被沮丧阴云压抑的灵魂，不可能有力量准备好聆听任何讲论。但我看到你们以极大的善意和强烈的热忱留心听我们讲话；并且因着你们对聆听的渴爱，已经摆脱了忧愁的思绪，搁置了当前的苦难感。为此，我与你们一同由衷感谢天主：灾祸并未胜过你们的明智，恐惧未曾削弱你们的活力，患难未能熄灭你们的热忱，危险不曾使你们的热忱干涸，对人的恐惧未能胜过对天主的渴望，时代的艰难未能摧毁你们的勤勉；不但没有摧毁，反而增强了它；不但没有松懈，反而使之更加专注；不但没有熄灭，反而点燃了它。市场虽空荡，教会却充满；前者提供忧伤的材料，后者则是喜乐和精神欢愉的场合！因此，亲爱的，当你前往市场，看到荒凉的景象而哀叹时，就逃回你的母亲（教会）那里，她立刻会以她的众多子女安慰你，向你展示众弟兄的完整合唱团，驱散你所有的沮丧！因为在城中，我们渴望见到人，如同居住在旷野中的人一般；但当我们避难于教会时，我们却因人多而拥挤不堪。正如当大海咆哮、被猛烈风暴激怒时，恐惧迫使所有人从外面逃入港口；同样现在，市场的波涛和城市的暴风，从四面八方将众人驱入教会，并以爱的纽带将肢体紧密相连。\par

2. 让我们即使因这些事也感谢天主，因为我们从患难中收获了如此多的果实，从试炼中获得了如此大的益处。如果没有试炼，就没有冠冕；如果没有角力，就没有奖品；如果没有赛道，就没有荣誉；如果没有患难，就没有安宁；如果没有冬天，就没有夏天。这不仅在人身上可以观察到，甚至在种子上也是如此；因为在那种情况下，我们期待麦穗发芽并生长，就需要大量的雨水、大量的云层聚集和大量的霜冻；播种的季节也是多雨的季节。既然如今冬天，不是元素的冬天，而是灵魂的冬天，已经来临，让我们也在这个冬天播种，以便在夏天收获；让我们播种眼泪，以便收获喜乐。这不是我的话，而是先知许诺的：\BibleQuote{含泪播种的，必含笑收割。}落下的雨水不会像眼泪的倾盆那样使种子发芽生长，它洗净灵魂，浇灌心志，使教义的嫩芽迅速向前生长。为此，也需要深耕。先知在这样说时表明了这一点：\BibleQuote{你们要开垦荒地，不可撒种在荆棘中。}\BibleRef{耶 4:3}因此，如同当农夫将犁置于田地上时，他从深处翻起泥土，提前为种子预备安全的安身之处，使它们不散落在表面，而是隐藏在地的子宫中，安全地扎根；同样，我们也有责任这样做，并利用患难的犁铧开垦心灵的深处。因为另一位先知在这样说时劝诫道：\BibleQuote{你们要撕裂你们的心，而不是你们的衣服。}\BibleRef{岳 2:13}让我们撕裂我们的心，这样，如果其中有任何邪恶的植物、任何诡诈的念头，我们就可以把它连根拔除，为虔诚的种子提供纯净的土壤。因为如果我们现在不开垦荒地；如果我们现在不播种；如果我们现在不在患难和斋戒之时以眼泪浇灌，我们何时才会感到痛悔？当我们在安逸和奢侈中时吗？但这是不可能的。因为安逸和奢侈通常导致懒惰，正如患难又引导我们回到勤勉；它使那曾四处游荡、在多对象中梦想的心志回归自身。\par

3. 我们不要因这种沮丧而悲伤，反而要感谢天主，因为磨难带来巨大的益处。农夫用极大的辛劳收集种子，撒下后，他祈求降雨；无知的人旁观，会对所发生的一切感到惊讶，或许自言自语：\Quote{这个人在做什么？他正在撒播他收集来的东西；不但撒播，还殷勤地把它与泥土混合，以致他将难以再将其聚集起来；此外，与泥土混合后，他还渴望一场大雨，使他所投下的一切都腐烂，变成泥浆。}这样的人看到雷声在云中爆发、闪电向下击打时也会恐惧。但农夫并不如此。他看见大雨时欢喜快乐。因为他所注目的不是眼前，而是期盼未来。他不理会雷声，而是计算他禾捆的数目。他不思想腐烂的种子，而是茁壮的麦穗；不思想恼人的雨水，而是打谷场上令人愉悦的尘土。同样，我们也应如此：不注目我们目前的磨难和其中的痛苦，而是瞩目它所可能带来的益处——它将结出的果实。让我们等待打谷场的禾捆；因为如果我们警醒，就能从当下收取许多果实，充满我们心灵的粮仓。如果我们警醒，就不仅完全不会因这磨难受到任何伤害，反而会收获数不尽的益处。但如果我们懒惰，连平安也会毁灭我们！这两件事都对不留心的人有害，但对生活严谨的人都有益处。正如金子：如果浸在水中，仍显示其固有的美；即使落入熔炉，也会比先前更加光亮；但另一方面，如果泥土或草与水混合，则一个溶解，一个腐烂；如果落入火中，则一个干枯，一个烧尽；义人与罪人的情况也是如此！因为如果义人享受安息，他仍然光辉灿烂，如同浸在水中的金子；即使落入考验，他更加光辉，如同经历火的试炼的金子；但罪人若得享安息，便如同草和泥遇水一样软弱腐烂；若经历考验，便如草和泥遇火一样被烧毁消灭！\par

4. 我们切莫因现时的苦难而灰心；因为如果你尚有残余的罪过，它们将被苦难消去，轻易地被焚毁；但若你拥有德行，你将因此更加显赫卓越。因为你若常保持警醒和节制，便能超然于一切伤害之上。因为导致跌倒的，并非考验的性质，而是受考验者的懈怠。所以，你若渴望喜乐、安逸和享乐，就不要寻求享乐或安逸，而要寻求一个充满忍耐、并能彰显刚毅的灵魂；因为若你没有这忍耐，不仅考验会使你蒙羞，而且安逸也将更彻底地摧毁你。为了证明并非罪恶的攻击，而是心灵的懈怠颠覆了我们的救恩，请听基督的话：\Quote{凡听了我的话而实行的，我要把他比作一个聪明人，他把房屋建在磐石上：雨淋、水冲、风吹，袭击那座房屋，它却不坍塌，因为基础是建在磐石上。} 又说：\Quote{凡听了我的话而不实行的，要把他比作一个愚昧人，他把房屋建在沙土上：雨淋、水冲、风吹，袭击那座房屋，它就坍塌了，且坍塌得很惨。} \BibleRef{玛 7:24} 你看见了吗？造成坍塌的，并非这些考验的冲击，而是建造者的愚昧？因为那里有雨，这里也有雨；那里有洪水，这里也有洪水；这里有大风吹袭，那里同样有。一个人建了房屋，另一个人也建了房屋。建筑相同，考验相同，结局却不相同，因为基础不同。建筑的倒塌是由于建造者的愚昧，而非考验的性质；否则那建在磐石上的房屋就该倒塌，然而它却安然无恙。但你莫以为这些话仅仅是指房屋而言；因为这番话论及的是灵魂，并借其行为证明它是听从或拒绝天主之言的。约伯就是这样建造了他的灵魂。雨淋下来——因为天火降下，烧尽了他所有的羊群；洪水涌来——那接连不断、接踵而至的灾祸使者，向他报告牲畜、骆驼和儿女的毁灭。大风吹袭——他妻子刻薄的话：\Quote{你诅咒天主，死了吧！} \BibleRef{约 2:9} 然而房屋没有倒塌：灵魂没有被颠覆：义人没有亵渎，反而如此感恩说：\Quote{主赐予的，主又收回。愿主的名受赞美。} \BibleRef{约 1:21} 你看见了吗？导致倾覆的，不是考验的性质，而是懒惰者的疏忽？因为苦难使强者更强。谁说的这话？是生活在苦难中的那个人，蒙福的保禄，他这样说：\Quote{苦难生忍耐，忍耐生考验，考验生望德。} \BibleRef{罗 5:3} 正如狂风猛烈冲击坚固的树木，将它们向各处摇摆，非但不能拔起，反而使它们更加坚实稳固；同样，圣洁且度虔敬生活的灵魂，不会因考验和苦难的侵袭而跌倒，反而由此激发更大的忍耐；正如蒙福的约伯，这些苦难使他更加显赫可敬。\par

5. 如今，有一个人向我们发怒，一个和我们同样有情感、同样有灵魂的人，我们就害怕；但在\Person{约伯}身上，却是一个邪恶恶毒的恶魔发怒；不，他不仅是发怒，还发动了各种阴谋，使出一切诡计；然而，即使如此，他仍不能征服那义人的刚毅。但是这里是一个人，他时而发怒，时而又和好；而我们却吓得要死。在那场合，是一个魔鬼在作战，它永不与人性讲和，反而对我们种族发动了一场没有条约的战争、一场没有休战的战斗；然而，那义人却嘲笑它的箭矢。那么，如果我们不能承受人的考验，我们还有什么借口，或什么宽恕可言呢？我们这些在恩宠之下受教于如此属神智慧的人，而这个人，在恩宠之前，在旧约之前，竟如此高尚地忍受了这场最惨烈的战争！因此，亲爱的，我们应当常常彼此谈论这些事，并用这样的话语鼓励自己。因为你们是证人，你们的良心也是证人，我们从这次考验中已获得了多少益处！放荡的人如今变得清醒；胆大的人变得温良；懒惰的人变得积极。那些从未见过圣堂，却经常在戏院消磨时光的人，如今整天留在圣堂里。那么请告诉我，你为此悲伤吗？因为天主通过恐惧使你认真，通过磨难引导你意识到自己的安全？但是你的良心痛苦吗？是的，你的心神每天像被箭刺穿，等待着死亡和最大的愤怒？然而，从那里我们也会在德行上获得巨大进步，如果我们的虔敬因困苦而变得更加恳切。因为天主能够今天就把你从这一切邪恶中解救出来。但是，直到祂看到你被净化，直到祂看到皈依发生，以及坚定不动摇的悔改，祂才会完全除去磨难。金匠直到看到金子精炼好了，才会把它从炉中取出；同样，天主在我们彻底改善之前不会挪去这云。因为允许这场考验的祂自己，知道除去它的时间。弹竖琴的人也是如此；他既不过度拉紧琴弦，以免拉断，也不过分放松，以免破坏和谐的共鸣。天主就是这样行事。祂既不使我们的灵魂处于持续的安逸状态，也不使其处于长期的磨难中，而是按祂的裁夺使用二者；因为祂既不容许我们享受持续的安逸，以免我们变得懈怠，也不让我们处于持续的磨难中，以免我们沉沦其中而绝望。\par

6. 那么，让\Emph{我们}把除去我们罪恶的时间交给祂；让\Emph{我们}只祈祷；让\Emph{我们}度虔敬的生活：因为这是我们的工作，转向德行；而使我们摆脱这些罪恶是天主的工作！因为祂确实比受试炼的你更渴望扑灭这火；但祂正等待着你的救恩。既然磨难来自安息，那么在磨难之后，也当期待安息。因为并非总是冬天，也并非总是夏天；并非总是波浪，也并非总是风平浪静；并非总是黑夜，也并非总是白昼。因此，磨难并非永久，但也会有安息；只是在我们磨难中，要常常感谢天主。因为那三个青年被投入火窑，连这也没有使他们忘记虔敬；火焰也没有吓倒他们，反而比那些坐在内室、毫无惊扰的人更加恳切，他们被火包围，仍向天上发出那些神圣的祈祷——因此，火成了他们的墙壁，火焰成了他们的衣袍；火窑成了泉源；它接收他们时他们是捆绑的，释放他们时他们是自由的。它接收了会死的身体，却像对待不朽的身体一样不下手！它知道它们的本性，却尊重他们的虔敬！暴君捆绑了他们的脚，而他们的脚却捆绑了火的作用！哦，奇妙的事！火焰释放了那些被捆绑的人，后来它自己被那些曾受捆绑的人捆绑；因为青年们的虔敬改变了事物的本性；或者不如说，它并没有改变本性，但更奇妙的是，它中止了它们的作用，尽管它们的本性依然存在。因为它没有熄灭火焰，却使火尽管燃烧也失去效力。这真是奇妙而不可思议，不仅发生在这些圣者的身体上，也发生在他们的衣服和鞋上。正如在宗徒们的情形中，\Person{保禄}的衣服驱逐了疾病和魔鬼\BibleRef{宗 19:12}，\Person{伯多禄}的影子\BibleRef{宗 5:15}赶走了死亡；照样，在这个事例中，这些青年的鞋也熄灭了火的威力。\par

7. 我不知该如何说，因为这奇事超出一切描述！火的力量既被熄灭又未被熄灭：当它接触这些圣人的身体时，它被熄灭；但当需要烧断他们的绑索时，它又未被熄灭；因此它烧断了他们的绑索，却没有触碰他们的脚踝。\BibleRef{达 3:25} 你看它多么接近？然而火并未受骗，不敢侵入绑索之内。暴君捆绑，火焰释放；为叫你同时学到外邦人的残暴与元素的顺从。他为何在要将他们投入火中之前捆绑他们？为使奇迹更大，使记号更难以解释，使你不至于认为所见的是视觉幻象。假若那火不是火，它就不会烧毁绑索；更有甚者，它就不会吞噬炉外的士兵；但事实上，它对炉外之人显示了威力，对炉内之人却显示了顺从。然而，我祈祷你观察，事事可见魔鬼如何以与天主仆人作战的同样手段，推翻自己的权势；这不是出于他的意愿，而是因为天主的智慧与丰富的谋略将他的所有武器和计谋转向他自己的头；那件事上确实如此。当时魔鬼激动暴君，既不许圣人的头被剑砍下，也不许他们被交给野兽，或受任何类似的刑罚；而是将他们投入火中；为使这些圣人的遗骸一点不剩，身体完全烧尽，骨灰与柴灰混在一起。然而天主恰恰利用这一情况来铲除不敬。如何？我将告诉你。火被波斯人视为神明；居住在那国的外邦人至今仍以许多敬礼尊崇它。天主因此渴望根除不敬的材料，允许刑罚采取这种方式，为能在这些拜火者眼前赐予祂的仆人胜利；以明确的事实说服他们：外邦人的神不仅恐惧天主，也恐惧天主的仆人。\par

8. 再思量，这场胜利的冠冕是由敌人编织的，敌人自身也成了这凯旋的见证。因为经上说：\BibleQuote{\Person{拿步高}派人召集总督、省长、队长、法官、司库、检察官和各省的首长，来参加雕像的开光礼，他们都聚集了。}\BibleRef{达 3:2}敌人搭好了戏台，他自己召集观众，布置了竞技场；这戏台不是由偶然的或无名之辈组成，而是由所有尊贵和有权势的人组成，为叫他们的见证在群众中更可信。他们应召而来是为了一个目的，却目睹了另一件事离去。他们来是为了朝拜雕像，离去时却嘲笑雕像，因这些青年身上发生的奇迹而惊叹于天主的权能。请留意，这场展示的场地设在何处。不是城市，也不是精选的围场，而是平坦开阔的原野，为这世界的戏台提供了空间。因为他在城外\Place{杜辣}平原上竖立了雕像，传令官前来呼喊：\BibleQuote{众民、各族、各方言的人哪，有令传与你们：你们一听见角、笛、琴、瑟、筝、鼓和各种乐器的声音，就应俯伏朝拜金像；谁不俯伏朝拜，立刻就被投入烈火窑中。}（因为朝拜偶像实在是堕落）\BibleRef{达 3:4}你看见这斗争被弄得多么艰难，网罗多么无法抗拒，深渊多么深，两边都是悬崖？但不要害怕。敌人增加计谋，他就越彰显青年的勇气。正因此，才有多乐器的交响，才有燃烧的火窑，为叫快乐和恐惧同时围攻在场者的灵魂。其中有顽固不屈之人？\Quote{让各种乐器的旋律，}他说，\Quote{诱惑并软化他。}但他超越这诡计？\Quote{让火焰的景象惊吓他。}于是恐惧和快乐同时在场：一个从耳朵攻入灵魂，另一个从眼睛。但这些青年的高尚品格无论如何也不被征服；正如他们落入火中却战胜了火焰，同样他们也嘲笑了一切欲望和一切恐惧。因为魔鬼为他们预先准备了一切。他对自己的臣民毫不怀疑，极有信心无人会抗拒王的命令。但当众人俯伏屈服后，只有青年被带到中间；为叫胜利更加显赫，他们独自在如此众多的人群中得胜，并被宣告为胜利者。如果他们在最初、还没有人跌倒时就勇敢行动，那可能不会如此令人惊奇。但最伟大、最惊人之处是，众多跌倒的人群既没有吓倒他们，也没有削弱他们。他们没有像许多人常说的那样想：\Quote{如果我们是第一个、也是唯一朝拜雕像的人，那是罪过；但如果我和这千万人一起做，谁会不谅解？谁会认为我们不值得辩护？}他们看见如此众多暴君俯伏在地时，没有说或想这类事。请你也想想那些控告者的邪恶，他们多么恶毒、刻毒地提出控告！他们说：\BibleQuote{有一些犹太人，就是您派管理巴比伦省事务的。}他们不仅提到民族，还提起他们的尊贵地位，为激怒国王；几乎像是说：\Quote{这些奴隶、这些俘虏、无城之人，您却立他们管理我们。但他们轻视这样的尊荣，并侮辱赐予此尊荣的人！}因此他们说：\BibleQuote{您所派管理巴比伦省事务的那些犹太人，不遵守您的命令，不事奉您的神。}\BibleRef{达 3:12}控告成了他们最大的赞美，指控成了颂词；由于是仇敌提出的，这见证无疑。王于是作什么？他命令将他们带到中间，为从各方面恐吓他们。但什么也不能使他们惊慌：既不是王的怒气，也不是孤身处在众人中，也不是火焰的景象，也不是号角声，也不是全体观望者对火焰的注视；他们嘲笑这一切，仿佛要投入清凉的水泉中，进入火窑时说出那有福的话：\BibleQuote{我们不愿事奉您的神，也不朝拜您所立的金像。}\BibleRef{达 3:18}\par

9. 我并非毫无理由地提及这段历史，而是为使你们学习：无论是君王的愤怒，或是士兵的暴力，或是仇敌的嫉妒，或是被掳，或是穷困，或是烈火，或是火窑，或是千万种恐惧，都不能使义人羞愧或恐惧。因为如果在那位不虔敬的君王面前，青年们并未因暴君的愤怒而惊慌，那么，我们有一位仁爱慈悲的皇帝，岂不更应当满怀信心，为这磨难而感谢天主吗？因为我们从刚才所说的话知道，磨难若使人以刚毅承受，便能使人在天主和人面前更加荣耀。诚然，倘若他们未曾为奴，我们就不会知道他们的自由！倘若他们未曾被掳，我们就不会学到他们灵魂的高贵！倘若他们未曾被流放离开地上的祖国，我们就不会知道他们天上国民身份的卓越！倘若地上的君王未曾对他们发怒，我们就不会知道他们蒙受天上君王何等的恩宠！\par

10. 那么，你也一样，如果你有天主作你的朋友，即使你落入火窑，也不要绝望；同样，如果祂发怒，即使你在乐园里，也不要以为安全。因为亚当确实在乐园里，然而，当他激怒了天主，乐园对他毫无益处。这些青年在火窑里，然而，由于他们被认可，火窑丝毫没有伤害他们。亚当在乐园里，但当他懈怠时，他被打倒了！约伯坐在粪堆上，但由于他警醒，他得胜了！乐园比粪堆好多少啊！然而地方的优越丝毫没有帮助那居民，因为他背叛了自己；同样，地方的卑劣也丝毫没有伤害那以德行四面坚固的人。至于我们自己，让我们坚固我们的灵魂；因为如果财富的损失甚至死亡威胁我们，却无人能夺去我们的信仰，我们便是最幸福的人。基督赞扬这件事，祂曾说：\BibleQuote{你们要机警如同蛇。} \BibleRef{玛 10:16} 因为正如蛇暴露全身以保全头部，你也当如此。即使为了保全你的信仰，暴露财富、身体、今生，或一切，都是必要的，也不要灰心！因为如果你带着那信仰离开此世，天主必以更丰盛的光彩恢复你的一切，并以更大的荣耀复活你的身体；代替财富的，将是那超越一切描述的福乐。约伯岂非赤身坐在粪堆上，忍受着比千万死亡更痛苦的生命吗？然而，由于他没有丢弃他的虔敬，他先前所有的一切都更丰富地归还给他：身体的健康和美好、满堂的儿女、财产，而比这一切更重要的是他那忍耐的荣冠。因为正如树木一样，若有人摘去果子又拔掉叶子，甚至砍去所有枝条只留下树根，这树仍会完整地重新生长，且更美丽，我们也是如此。如果虔敬的根还在，纵使财富被夺去，身体被毁坏，一切事物都会以比先前更大的荣耀归还给我们。因此，让我们抛弃一切焦虑和多余的牵挂，回归自我；用德行的装饰来装扮身体和灵魂，将我们的肢体转化为正义的器具，而非罪恶的器具。\par

11. 首先，让我们训练我们的舌头成为圣神恩宠的仆人，从口中驱除一切恶毒和狠毒，以及使用可耻言辞的习惯。因为我们能够使我们的每一个肢体成为邪恶或正义的工具。那么，请听人如何使舌头成为工具：有些人成为罪恶的工具，另一些人成为正义的工具！\BibleQuote{他们的舌头是锋利的剑。} 但另一位论及自己的舌头说：\BibleQuote{我的舌头是敏捷文士的笔。} 前者带来毁灭，后者书写了神圣法律。因此，一个是剑，另一个是笔，这不是出于其本性，而是出于使用者的选择。因为这一个舌头和那一个舌头的本性相同，但作用不同。同样，关于口舌，我们也看到同样的事。因为这些人满口污秽和邪恶，因此论及他们时，有这样指控的话：\BibleQuote{他们的口充满咒骂与苦毒。} 而他的口却不同，正如：\BibleQuote{我的口要宣讲智慧，我的心默思明达。} 又有些人，他们的手充满不义，指责这些人时，他说：\BibleQuote{不义在他们的手中，他们的右手充满了贿赂。} 但他自己的手只习惯于向天伸展。因此他也论及它们说：\BibleQuote{我举手（愿它成为）晚祭。} 关于心，也可以看到同样的事；因为他们的心确实是愚昧的，但此人的心是真诚的；因此他论及他们说：\BibleQuote{他们的心是虚妄的；} 但论及自己的心说：\BibleQuote{我的心正在酝酿一篇美好的言论。} 关于耳朵，也可以看到同样的情况；因为有些人的听觉如同野兽，不能被感动或同情；指责这样的人时，圣咏作者说：\BibleQuote{他们像聋蝮蛇，塞住耳朵。} 但他的耳朵是神圣言语的容器，这他又表明出来，当他说：\BibleQuote{我要侧耳听比喻，我要在竖琴上揭示我的隐语。}\par

12. 知道这些事，就让我们用德性从各方面坚固自己，这样我们必将避免天主的义怒；让我们使身体的肢体成为正义的器具；并让我们管教眼睛、口、手、脚、心、舌和整个身体，只用于德性的事奉。让我们记住那三条诫命，我曾对你们的爱德讲论过，劝勉你们不要看任何人作仇敌，也不要说任何亏负你们的人的坏话；并从你们的口中除去发誓的恶习。至于前两条诫命，我们将在另一机会向你们讲论；但在整个本周内，我们将对你们谈论关于发誓的事；这样从较容易的诫命开始。因为克服发誓的习惯根本不是什么劳苦，只要我们稍加努力，彼此提醒、劝告、观察，并要求那些忘记自己的人交账和受罚。因为如果我们不除去灵魂的恶习，禁食又有什么益处呢？看，我们整天禁食；到了晚上，我们将摆设宴席，不像昨夜那样，而是改变了的、更庄重的宴席。我们当中有人能说自己今天也改变了生活，改变了他的恶习，如同改变了他的食物一样吗？真的，我想没有！那么我们的禁食有什么益处呢？因此我劝勉，并将不停地劝勉，你们应分别承担每一条诫命，花两三天时间去实行它；正如有些人在禁食上彼此竞争，表现出奇妙的竞赛（有些人整两天不进食；另一些人不但从宴席上弃绝酒和油，而且弃绝一切菜肴，只取面包和水，在整个四旬期内坚持这种做法）；同样，我们也要彼此竞赛，消除频繁发誓的恶习。因为这比任何禁食更有益；这比任何苦修更有利。我们为禁食所表现出来的同样关心，也应在戒除发誓上表现出来；因为我们若不顾那些禁止的事，却把一切关心都花费在无关紧要的事上，就难免被视为极其愚昧；因为吃并不被禁止，但发誓是被禁止的；然而，我们竟放弃那些允许的事，胆敢去犯那些禁止的事！因此，我恳求你们的爱德做些改变，并从今天开始显出改变。因为如果我们以这样的热心度过整个四旬期，在本周内达到完全不发誓的实践；在下一周内熄灭愤怒；再下一周内连根拔除说坏话；然后再改正剩余的事；这样步步前行，我们将渐渐地到达德性的顶峰；我们必将逃脱眼前的危险；并使天主仁慈；众百姓必将再次回到我们的城市；我们将教导逃亡者，我们的安全希望既不是在于地方的保障，也不是在于逃跑和隐居，而是在于灵魂的虔敬和品德的善行。这样，我们将获得今世和来世的美善；愿天主恩赐我们众人都有份于这些美善！因我们的主耶稣基督的恩宠和慈爱，借着祂、偕同祂，圣父偕同圣神，光荣归于祂，从今直到永远。阿们。\par

\SourceRecord{Translated by W.R.W. Stephens.；Nicene and Post-Nicene Fathers, First Series；Vol. 9.；Edited by Philip Schaff.；Buffalo, NY: Christian Literature Publishing Co.,；1889.；<http://www.newadvent.org/fathers/190104.htm>.}

% structure: p0001 source=h1 target=section reason=page-title

\section{论雕像第五篇讲道}

\Emph{本篇继续前篇讲道的劝勉。劝勉民众勇敢承受皇帝即将来临的义怒。以\Person{约伯}和\Person{尼尼微人}为例。说明人不应怕死，而应怕罪。解释何为悲惨的死亡；并以严肃的劝阻结束，反对滥用誓词。}\par

1. 关于三位青年和\Place{巴比伦}火窑的讲论，昨日似乎给了你们的爱德不小的安慰；而\Person{约伯}的榜样，以及那比任何王座更可敬的粪堆，更是如此。因为观看王座，对观看者毫无益处，只有暂时的快乐，毫无裨益；但从观看约伯的粪堆，却能获得各种益处，甚至许多神圣的智慧和安慰，以培养忍耐。因此，直到今日，许多人进行长途朝圣，甚至漂洋过海，从地极赶到远至\Place{阿拉伯}，为要看见那粪堆；看见后，就亲吻那土地，那地曾容纳了如此得胜者的角力场，并接受了比所有金子更宝贵的鲜血！因为紫袍的光芒，远不及那身体——它被染红，不是用别人的血，而是用自己的血！就连那些伤口，也比各种珠宝更宝贵！因为珍珠的本性对我们的生活毫无帮助；也不能满足拥有者任何必要的需求。但那些伤口，却是对所有悲伤的安慰；为使你们知道这是真理，设想有人失去一个心爱的独子。你给他看一万颗珍珠，也不能安慰他的悲伤，减轻他的痛苦；但若使他想起\Person{约伯}的伤口，你就能轻易地通过这样说而施予安慰：\Quote{你为什么悲伤，人啊？你失去一个儿子；但那蒙福的人，在丧失全部子女之后，又在自己身上遭受灾病，赤身坐在粪堆上，全身血流如注，肉体逐渐腐烂；他原是正义、诚实、如此虔诚的人，远离一切恶行，甚至天主亲自为他的德行作证。}这样说，你就能熄灭受苦者的一切悲伤，消除他的一切痛苦。如此，义人的伤口变得比珍珠更有用！\par

2. 那么，请你们想象这位角力者；想象你们看见那粪堆，以及他坐在其中！那金像！镶嵌着宝石！我不知道如何表达：因为我找不到任何珍贵的材料可与那染血的身体相比！那肉体的本性远超过一切物质，无论多么昂贵，都无可比拟地更珍贵；那些伤口比阳光更灿烂；因为阳光照亮肉体的眼睛，而伤口照亮心灵的眼睛！这些伤口使魔鬼完全盲目！因此，在那打击之后，他退却并不再出现。亲爱的，你也由此学习磨难有何益处！因为当义人富有、享受安逸时，魔鬼有指控他的理由。尽管是虚假的，但他仍有权说：\Quote{约伯岂是无缘无故事奉你？} 但在剥夺他一切、使他贫穷之后，他甚至连口也不敢再开。当他富有时，魔鬼准备与他角力，并威胁要打倒他；但当他使他贫穷，夺去他所有，并把他投入最深苦难之后，他便退却了。当他的身体健全时，魔鬼举手攻击他；但当他击打他的肉体后，他便逃跑了——败退了！你看，对于警醒的人，贫穷远比财富更好、更有益；疾病和衰弱比健康更好；考验比安宁更好，因为它使角斗者更光荣、更有力！\par

3. 谁曾见过或听过如此惊人的竞赛？世上竞赛的斗士，当他们击碎对手的头颅时，才算得胜，并被加冕！但这位对手，当他击打义人的身体，用各种疮痍将其穿孔，并使之极度虚弱时，却被征服而退却了。即使他刺穿了义人的每一根肋骨，也并未因此获利；因为他未能夺去义人隐藏的宝藏，反而使他更显耀于我们面前；并且通过那次刺穿，他让所有人都有特权窥视义人的内心，并完全辨明他的全部财富！当他期望获胜时，却带着极大的羞辱退却，并且再没有说出一句话！何事发生，啊魔鬼？你为何退却？难道你所选择的一切没有成就吗？你不是夺走了他的羊群、牛群、马群和骡群吗？你不是也摧毁了他的一群儿女吗？并且将他的身体击成了碎片。你为何退却？\Quote{因为，}他说，\Quote{我所选择的一切都已成就，然而我最渴望成就的、并为此做这一切的那件事，却没有成就：他没有亵渎！因为正是为了这个，他接着说，我做了这一切；既然这事没有成就，我并未因夺走他的财富、或因他儿女的毁灭、或因加在他身上的灾祸而获益；相反，我所图谋的结果适得其反：我使我的敌人更加光荣；我增添了其声誉的光彩。}亲爱的，你看出患难的赏报何其大吗？他的身体先前美丽健康，但被这些伤口穿透后，却变得更加可敬！羊毛在染色前固然美丽，但当它变成紫色时，便获得了一种难以形容的美丽和额外的风韵。但是，如果他未曾剥夺他，我们就不会知道胜利者的良好状态；如果他未曾用疮痍刺穿他的身体，内在的光芒就不会闪耀。如果他未曾使他坐在粪堆上，我们就不会知道他的财富。因为一位国王坐在宝座上，并不如这人坐在他的粪堆上那样显赫夺目！因为在王座之后，便是死亡；但在那粪堆之后，便是天国！\par

4. 那么，收集所有这些理由，让我们从压迫我们的沮丧中振作起来。我向你们陈述这些史事，并非要你们称赞所说的，而是要你们效法这些高尚之人的德行和忍耐；好使你们从事实本身学到，在人世间的祸患中，除了罪以外，没有什么是可畏惧的；既不是贫穷，也不是疾病，不是侮辱，不是恶待，不是羞辱，也不是死亡——这被视为万恶中最坏者。对于热爱属灵智慧的人而言，这些都只是灾祸的名称；这些名称没有实质性的实在。真正的灾祸在于得罪天主，并做任何令祂不悦的事。因为请告诉我，死亡中有什么可怕的呢？岂非因为死亡更快地将你带入平安的港湾，以及那无扰的生活吗？即使人不杀你，自然律本身最终也会悄然临到你，使灵魂与身体分离；即便我们恐惧的这件事现在不发生，它不久也会发生。\par

5. 我这样说，并非预料任何可怕或悲惨的事：愿天主禁止！而是因为我为那些害怕死亡的人感到羞耻。请告诉我，一面期待那\Quote{眼所未见，耳所未闻，人心所未曾想到的}如此美好的事物，一面却对这享受犹豫不决，漠不关心且怠惰，不仅怠惰，而且恐惧战栗？你为死亡而苦恼，而\Person{保禄}却为现世生命而叹息，在写给罗马人的书信中说：\BibleQuote{受造物一同叹息，连我们这些已蒙受圣神初果的，也自己叹息。}\BibleRef{罗 8:22}他这样说话，并非谴责现世，而是渴望来世。他说，\Quote{我尝到了恩宠，}他说，\Quote{便不愿忍受这迟延。我有了圣神的初果，就向着全备推进。我升到了第三层天；我看见了那说不出的荣耀；我望见了那光芒四射的宫殿；我得知我在此逗留时所失去的喜乐，因此我叹息。}因为假设有人带领你进入王宫，向你展示墙上处处闪耀的金子以及其余一切辉煌的景象；若他从那里随后又带你回到一个穷人的茅舍，并应许在短时间内带你重返那些宫殿，并在那里赐给你永久的居所；请告诉我，你难道不会因渴望而憔悴，甚至为这几日而感到急躁吗？同样，你要思想天上和地上，并与\Person{保禄}一同叹息，不是因为死亡，而是因为现世的生命！\par

6. 但你说，让我像\Person{保禄}一样，我就永不怕死。请问，人哪，什么阻止你成为像\Person{保禄}一样的人？他不是穷人吗？他不是帐棚匠吗？他不是地位卑微的人吗？因为如果他富有而且出身高贵，穷人被召叫效法他的热忱时，就可以拿贫穷作借口；但现在你无话可说。因为这个人操劳手工，并且靠每日劳作养活自己。而你，从一开始就继承了你祖辈的信仰；从幼年起就在圣经的研习中被养育；而他却曾是\Quote{亵渎者、迫害者和施暴者}\BibleRef{弟前 1:13}，并且蹂躏了教会！然而，他如此顷刻间改变，以至于在热忱的激烈程度上超过了众人，他喊说：\Quote{你们当效法我，如同我效法基督一样}\BibleRef{格前 11:6}。他效法了主；而你这个从起初就在虔敬中受教育的人，难道不肯效法一个同仆吗？一个后来藉皈依而信德的人？难道你不知道，那些在罪恶中的人，活着也是死的；而那些生活在正义中的人，即使死了，也仍然活着？\BibleRef{弟前 5:6} 这不是我的话。这是基督对\Person{玛尔大}说的宣言：\Quote{信从我的人，即使死了，仍要活着}\BibleRef{若 11:5}。我们的道理难道是个神话吗？如果你是基督徒，就相信基督；如果你相信基督，就用你的行为向我证明你的信德\BibleRef{雅 2:18}。但你怎么证明呢？藉着轻视死亡：因为在这方面我们与不信者不同。他们害怕死亡是自然的；因为他们没有复活的希望。但你，正向着更美好的事物前进，并且有机会默想未来的希望，既然确信复活，却同时像那些不信复活的人一样害怕死亡，你还有什么借口呢？\par

7. 但有人说，我不怕死亡，也不怕死的过程，而是怕悲惨的死亡、被斩首。我问，那么\Person{若翰}死得悲惨吗？因为他被斩首。或者\Person{斯德望}死得悲惨吗？因为他被石头砸死；所有殉道者都按这种说法死得悲惨；因为有些人被火烧死；有些人被刀剑杀死；有些人被投入海洋；有些人被推下悬崖；还有些人被投入野兽口中，这样死去。人啊，卑贱的死去不是以暴力结束生命，而是死在罪恶中！至少请听先知就此事所作的道德训诲，他说：\Quote{罪人的死是邪恶的}。这很公正；因为在离开此生之后，有无法忍受的惩罚；不灭的报复；有毒的虫；不熄的火；外边的黑暗；不可解的锁链；切齿的哀号；磨难、痛苦和永恒的正义。\par

8. 既然这样的恶事等待罪人，他们即使在家中床上结束生命，又有什么好处呢？同样，另一方面，对义人来说，如果他们通过刀剑、钢铁或火焰放下今生，当他们将要前往不朽的美善事物时，这也不能伤害他们。的确，\Quote{罪人的死是邪恶的}。那轻视\Person{拉匝禄}的富人就是这样的死。他自然死亡，在家中床上，有亲属环绕，但离世后到另一个世界却经历了火刑；在现世所享受的一切快乐中，他在那里连一点安慰也得不到！但\Person{拉匝禄}却不是这样；因为他躺在石板上，狗来舔他的疮，他遭受了暴力死亡（因为有什么比饥饿更痛苦呢？），但他离世后却享受永恒的福乐，安息在\Person{亚巴郎}的怀抱中！那么，暴力死亡在哪方面伤害了他呢？或者富人不死于暴力，又有什么益处呢？\par

9. 但有人说：\Quote{我们不怕死于暴力，而是怕不义地死去；怕和那些被怀疑犯了罪而我们与此无关的人一样受到惩罚。}你说什么，告诉我？你怕不义地死去，却希望义地死去吗？有谁那么悲惨可怜，当可以选择不义地死去时，反而愿意死于正义呢？因为如果必须怕死，那就在死亡正义地临到我们时怕它；因为那不义地死去的人，恰恰因此与众圣人成为有分。许多蒙天主悦纳和显扬的人都遭受了不义的死亡；首先是\Person{亚伯尔}。因为他没有得罪他的兄弟，也没有伤害\Person{加音}；相反，因为他尊敬了天主，所以被杀害。但天主允许了这事。你想，是因为爱他，还是因为他恨他？很明显，是因为爱他，并想通过那次最不义的谋杀使他的冠冕更加明亮。那么你看到吗，我们不应该怕死于暴力，也不应怕不义地死，而应怕死在罪恶中？\Person{亚伯尔}不义地死了。\Person{加音}活着，呻吟颤抖！那么请问，两人中谁更有福？是那在正义中安息的人，还是那活在罪恶中的人？是不义地死去的人，还是那正义地受罚的人？你愿意我向你的爱德宣告，我们为什么怕死吗？天国的爱没有渗透我们，对将来事物的渴望没有燃烧我们；否则我们会轻视一切现世之物，如蒙福的\Person{保禄}一样。另一方面，我们也不敬畏地狱；因此死亡是可怕的。我们感觉不到那里惩罚的不可忍受性；所以我们用死亡代替罪来害怕；因为如果对死亡的恐惧占据我们的灵魂，对地狱的恐惧就无法进入。\par

10. 我要努力显明这一点，不是从遥远的事物，而是从我们眼前的事，从这些日子在我们当中发生的事。当皇帝的信函到来，下令征收那被认为难以忍受的赋税时，所有人都骚动起来；所有人都抱怨，认为这是严重的苦难，表示愤慨；当他们相遇时，彼此说：\Quote{我们的生命不值得活下去，城市完了——没有人能承受这沉重的负担；}他们痛苦万分，仿佛处于极端的危险之中。此后，当叛乱真的发生，某些卑劣的、彻底卑劣的人，践踏法律，推倒雕像，使所有人陷入极大的危险；而现在，由于皇帝的愤怒，我们为性命而恐惧，这金钱的损失不再刺痛我们。但取而代之的是，我听到所有人说的是另一种话：\Quote{让皇帝拿走我们的财产吧，我们乐意被剥夺田地和所有物，只要有人能确保我们赤裸身体的平安。}因此，如同在死亡的恐惧临到我们之前，财富的损失折磨我们；而在这些无法无天的暴行发生后，随之而来的死亡恐惧驱散了因损失而来的悲伤；同样，如果地狱的恐惧占据了我们的灵魂，死亡的恐惧就不会占据它们。但正如身体一样，当两种痛苦抓住我们时，更强烈的通常掩盖较弱的；现在也是如此，如果对将来惩罚的畏惧停留在灵魂中，它将掩盖所有人类的恐惧。所以，如果有人总是努力记起地狱，他将嘲笑各种死亡；这不仅会将他从当前的困境中解救出来，甚至会将他从将来的火焰中救出。因为总是害怕地狱的人，永远不会落入地狱之火中；因这不断的恐惧而清醒！\par

11. 请允许我在合适的时机对你们说：\BibleQuote{弟兄们，在见识上不要做小孩子，但在邪恶上却要做小孩子。}\BibleRef{格前 14:20}因为，如果我们害怕死亡，却不畏惧罪，这是我们幼稚的恐惧。小孩子也害怕面具，却不害怕火。相反，如果偶然被带到点燃的蜡烛旁，他们会毫无顾忌地把手伸向蜡烛和火焰；而一个如此可鄙的面具却令他们害怕；但他们不畏惧火，那才是真正可怕的东西。同样，我们害怕死亡，这不过是一个可以蔑视的面具；却不畏惧罪，那才是真正可怕的；它像火一样吞噬良心！这通常不是因为事物的本性，而是因为我们自己的愚蠢；所以如果我们一旦思考死亡是什么，我们就永远不会害怕它。那么，请问，死亡是什么？就像脱下一件衣服。因为身体之于灵魂如同衣服；死后暂时脱下，我们将会更加荣耀地重新穿上。死亡最多是什么？是暂时的旅程；是比平常更长的一觉！所以如果你害怕死亡，你也应该害怕睡觉！如果你为垂死的人痛苦，也要为那些吃喝的人悲伤，因为这是自然的，那也是自然的！不要让自然的事令你悲伤；而要让你因邪恶选择而来的事悲伤。不要为垂死的人悲伤；要为活在罪中的人悲伤！\par

12. 你还要我提出另一个我们害怕死亡的原因吗？我们生活不严谨，良心不清白；因为如果是这样，没有什么能惊吓我们，无论是死亡、饥荒、财富的损失，还是任何其他类似的事。因为生活有德的人，不会被这些事中的任何一件伤害，也不会被剥夺内心的喜乐。因为他有美好的希望支持，没有什么能使他沮丧。有人能做什么使高尚的人忧愁呢？夺走他的财富？他还有天上的财富！将他逐出祖国？他将前往那在上面的城市！给他戴上锁链？他仍然良心自由，对外在的锁链无动于衷！杀死他的身体？他还要复活！正如与影子搏斗、击打空气的人，不能击中任何人一样；所以与义人作战的人，只是在击打影子，浪费自己的力气，不能对他造成任何伤害。请允许我确信天国；如果你愿意，就在今天杀死我。我将为这次杀戮感谢你，因为你使我快速获得那些美好事物！\Quote{然而，这正是我们特别哀叹的，}有人说，\Quote{就是被我们众多的罪所阻碍，我们无法进入那国度。}既然如此，就停止哀叹死亡，哀叹你的罪吧，好使你得以脱离它们！悲伤确实存在，不是要我们为财富的损失、死亡或任何其他类似的事悲伤，而是要我们用它来除去我们的罪。我借一个例子来说明这一点。治病的药物只针对它们能去除的疾病，而不是对那些它们完全无用的疾病。例如（我想把事情说得更清楚），一种只能对眼疾有益而对其他疾病无效的药，人们完全可以说它是专为眼睛而造的，不是为胃、手或其他任何器官。现在让我们把这个论证转到悲伤的主题上；我们会发现，在发生在我们身上的所有事情中，它除了纠正罪之外毫无益处；因此显而易见，它只是为了摧毁罪而存在的。现在让我们审视一下降临在我们身上的每一种邪恶，让我们把沮丧当作补救方法，看看它产生什么益处。\par

13. 有人被罚财产：他变得忧愁，但这并不能弥补他的损失。有人失去了儿子：他悲伤，但他不能使死者复活，也不能惠及逝者。有人被鞭打、挨打、受辱；他变得忧愁。这不能撤销侮辱。有人染上疾病，而且是最严重的病；他沮丧。这不能消除他的疾病，反而使之更严重。你看见了吗？在所有这些情况下，忧愁都没有任何益处。假设有人犯了罪，并且忧愁。他抹去了罪；他脱离了过犯。这如何显明呢？因主的宣告；因为，论到某个犯了罪的人，祂说：\Quote{因他的不义，我使他忧愁一时；我看见他悲伤，他步履沉重；我医治了他的道路。}因此\Person{保禄}也说：\BibleQuote{依照天主的旨意而有的忧愁，能产生悔改，以致于得救，不使人后悔。}\BibleRef{格后 7:10}既然我所说的清楚显示，既不能靠忧愁的介入来弥补财富的损失、侮辱、毁谤、鞭打、疾病、死亡或任何其他此类事物，只有忧愁能抹去并消除罪恶，那么显然，这就是忧愁存在的唯一原因。因此，我们不要再为财富的损失而忧愁，而只应在犯罪时忧愁。因为在这种情况下，忧愁带来的收益是巨大的。你被罚款了吗？不要沮丧，因为这样对你毫无益处。你犯罪了吗？那么就应该忧愁：因为这是有益的；并且要思考天主的技能和智慧。罪为我们产生了这两样东西：忧愁和死亡。因为\BibleQuote{你吃的那一天}，祂说，\BibleQuote{你必定死}；并对女人说：\BibleQuote{你将在忧愁中生育儿女。}\BibleRef{创 2:17}而借着这两样，祂除去了罪，并安排母亲被她的后代摧毁。因为死亡和忧愁都除去罪，首先从殉道者的例子明显可见；从\Person{保禄}对犯罪者所说的话也清楚可见，他这样说：\BibleQuote{因此，在你们中间有许多软弱和患病的人，且多有长眠的。}因为，他说，你们既然犯了罪，你们就死了，所以你们借着死亡从罪中得释放。因此他接着说：\BibleQuote{我们若省察自己，就不致于受审判。但我们受审判，乃是主所惩罚，以免我们和世人一同被定罪。}\BibleRef{格前 11:31}正如从木头中生出的虫子，吃掉了木头；飞蛾从羊毛中产生，又吃掉了羊毛；同样，忧愁和死亡生于罪，又吞噬了罪。\par

14. 那么，我们不要怕死，只应怕罪，并为罪而忧愁。我说这些话，并非预料什么可怕的事，绝无此事！而是希望你们在惊惶时，常怀着这样的心情，并确实履行\Person{基督}的法律。因为\Person{基督}说：\BibleQuote{谁不背起自己的十字架跟随我，就不配是我的。}\BibleRef{玛 10:38}祂这样说，不是要我们把木头背在肩上，而是叫我们常将死亡摆在眼前。正如\Person{保禄}天天死去，嘲笑死亡，轻视今世的生命。因为你确实是个士兵，时常武装待命；但一个怕死的士兵，永远不能做出高贵的行动。同样，一个基督徒如果害怕危险，也不会做出任何伟大或可钦佩的事；不仅如此，他还容易被击败。而勇敢高尚的人却不如此。他保持坚不可摧、不可战胜。正如那三个青年，当他们不怕火时，就从火中逃脱；同样，我们若不怕死，也必完全脱离死亡。他们不怕火（因为被烧不是罪），但他们怕罪，因为行亵渎是罪。让我们也效法他们和所有这样的人，不要害怕危险，这样我们就能安全地度过危险。\par

15. 至于我，\BibleQuote{我不是先知，也不是先知的儿子}\BibleRef{亚 7:14}，但我对未来的事却明白得如此清楚，并且大声清楚地宣告：如果我们改变自己，关心自己的灵魂，并停止不义，那么就没有什么会令人不愉快或痛苦。这一点，我从天主对人的爱，以及祂为人、为城市、为国家、为全体人民所做的一切，清楚知道。因为祂曾威胁尼尼微城，说：\Quote{还有三天，尼尼微就要被毁灭。}\EditorialNote{约纳书第三章}我问，那么，尼尼微被毁灭了吗？那城被摧毁了吗？不，恰恰相反；它既复兴了，而且变得更加著名；时间过去那么久，并未抹去它的荣耀，我们至今仍颂扬和钦佩它。因为从那时起，它成了所有犯罪者的一个绝好的避风港，不让他们沉沦于绝望，反而召唤所有人悔改；借它的所作所为，以及从天主那里获得的恩宠，说服人永不绝望于自己的救恩，反而活出最好的生命，并给他们立下美好的希望，使他们确信结局无论如何终将是好的。因为谁听了这样的榜样，即使是最懒惰的人，岂会不被激励呢？\par

16. 天主甚至宁愿自己的预言落空，好使城市不致倾覆。或者更确切地说，预言并未落空。因为如果那些人继续行同样的恶事，而判决没有生效，或许有人会指控所说的预言。但如果当他们改变了，停止了恶行，天主也止息了祂的怒火，谁还能再挑剔预言，或指证所说的为虚谎呢？事实上，天主从起初就制定的同一条律则，藉先知向所有人颁布，在那时被严格遵循。那么，这条律则是什么？祂说：\BibleQuote{我说一判决，关于一国或一邦，要拔除、拆毁、毁灭它；但若那邦人悔改他们的恶行，我也要后悔我所说的要加给他们的怒气。} \BibleRef{耶 18:7} 遵守这条律则，祂拯救了悔改的人，并从怒火中释放了停止恶行的人。祂知道外邦人的德行，因此祂赶快遣先知到那里。于是，那城市在听到先知的声音时非常震动，但恐惧不但没有伤害它，反而使它受益。因为那恐惧是它得救的原因。威胁实现了从危险中的解救。推翻的判决阻止了推翻。啊，多么奇异惊人！威胁死亡的判决却带来了生命！判决公布后竟被取消，与世间的判官完全相反！因为在他们那里，宣布判决使它生效，完全确立；但相反，在天主那里，公布判决却使它被取消。因为如果判决没有被公布，犯罪的人就不会听到；如果他们没有听到，就不会悔改；如果没有悔改，就不能避免惩罚，也不能获得那惊人的解救。这何其惊人：判官宣布判决，而被判罪的人以自己的悔改化解了判决！他们确实没有像我们现在这样逃离城市，而是留在城中，使城市屹立。城市本是网罗，他们却使它成为堡垒！它是深渊和峭壁，他们却把它变成安全的高塔！他们听说房屋将要倒塌，却没有逃离房屋，而是逃离了自己的罪恶。他们没有像我们现在这样各自离家出走，而是各自离开了恶道；因为他们说：\Quote{我们为何以为城墙带来了怒火？我们是创伤的根源，因此我们应该提供药物。} 因此，他们信赖的不是更换住所，而是改变习惯。\par

17. 外邦人尚且如此！我们岂不羞愧，岂不掩面？因为他们改变习惯，我们却只更换住所，暗中搬移财物，行醉酒之人的事！我们的主向我们发怒，我们却忽略平息祂的怒火，只把家具从一处搬到另一处，东奔西跑，寻找可存放财产的地方；我们更应该寻求的是在哪里安全地存放我们的灵魂；或者更确切地说，我们不应寻求，而应将灵魂的安全托付给德行和正直的生活。因为当我们对仆人发怒不满时，如果他不是为自己的不悦辩护，而是下楼到他的房间，收拾衣服，捆好所有动产，打算逃跑，我们岂能容忍这种轻慢？那么，让我们停止这不合适的努力，让我们各人对天主说：\Quote{我往哪里去躲避祢的神？我往哪里逃，躲避祢的面？} 让我们效法外邦人的神性智慧。他们甚至在不确定的情况下悔改！因为判决没有这样的条款：\Quote{你们若回转悔改，我就重建这城}，而只是说\Quote{再过三天，尼尼微必倾覆}。那么他们怎么说？\BibleQuote{谁知道天主是否会后悔祂所说的要加给我们的灾祸？} \BibleRef{纳 3:9} 谁知道？他们不知道事情的结果，却并未忽略悔改！他们不熟悉天主显示仁慈的方式，却凭着不确定而改变！因为他们既不能参考其他悔改而得救的尼尼微人；也没有读过先知；没有听过族长；没有享受劝告或接受训诫；也没有说服自己必定能藉悔改平息天主。因为威胁并不暗示这一点；但他们对此心存疑虑、犹豫不决，却仍殷勤悔改。我们还有什么理由推脱呢？那些对结果毫无把握的人却显出了如此巨大的改变；而你有理由信靠天主的仁慈，屡次领受祂眷顾的保证，听过先知、宗徒，并从实际事件中受教，却仍然没有效法他们达到同样的德行！他们的德行确实伟大！但天主的仁慈更为伟大！这从威胁的巨大本身就可以看出。因此，天主没有在宣告中加上\Quote{但若你们悔改，我必宽恕}，为的是通过设定一个无限制的判决来增加恐惧，并在恐惧增加后，迫使它们更快地悔改。\par

18. 先知确实感到羞愧，因为他预见到事情的结果，并猜想他所预言的事将不会实现。然而，天主并不羞愧，祂只渴望一件事，即人的救恩，并纠正了自己的仆人。因为当约纳进入船中时，天主立刻在那里掀起汹涌的海浪；好让你知道，哪里有罪，哪里就有风暴；哪里有悖逆，哪里就有波涛的翻腾。尼尼微城因为尼尼微人的罪而震动；船因为先知的悖逆而震动。于是水手们把约纳抛入深海，船便得救了。那么，让我们溺毙自己的罪，我们的城邦必得安全！逃跑对我们肯定无益；因为对约纳没有益处，反而伤害了他。他确实逃离了陆地，却没有逃离天主的愤怒；他逃离了陆地，却把风暴带到了海上；他非但没有从逃跑中得到任何好处，反而使收留他的人也陷入了极端的危险。当他坐在船中航行时，尽管有水手、舵手和船上一切必要的设备，他仍处于极大的危险中。然而，被抛入深海并通过惩罚除去罪过之后，他被送入那不稳定的容器——我是指鲸腹——却享有极大的安全。这是为了教导你：如同没有船能对生活在罪中的人有益一样，那除去罪过的人，大海不能淹死他，怪兽也不能毁灭他。的确，海浪接收了他，却没有使他窒息。鲸鱼接收了他，却没有毁灭他；动物和元素都把托付给它们的安然无恙地归还给了天主；通过这一切，先知被教导要有人情和慈悲，不要比野兽、无思想的水手或狂暴的波浪更残忍。因为就连水手也不是一开始就立刻交出他，而是经过许多强迫之后；大海和怪兽也以极大的善意守护了他；这一切都在天主的引领之下。\par

19. 因此他回来了；他宣讲；他警告；他劝服；他保存；他惊吓；他改正；他建立；只通过一次，而且是第一次宣讲！他不需要许多日子，也不需要持续的商议；仅仅说了这些简单的话，就使所有人悔改！因此，天主没有直接把他从船里带入城中；而是水手把他交给大海；大海把他交给鲸鱼；鲸鱼把他交给天主；天主把他交给尼尼微人；通过这漫长的迂回，祂带回了这逃犯，为要教导所有人：不可能从天主的手中逃脱；无论一个人走到哪里，拖着罪过，他将遭受无数祸患；即使没有人在场，整个受造物也会从四面八方极其猛烈地起来反对他！那么我们不要通过逃跑，而要通过品格的改变来确保我们的安全。天主向你发怒，是因为你留在城里吗？祂发怒是因为你犯了罪。因此，除去罪过，并从你伤口的根源那里移除邪恶的源头。因为医生也指示我们用相反的方法治疗相反的病。例如，发烧是由丰盛的饮食引起的吗？他们用禁食的疗法治疗疾病。有人因悲伤生病吗？他们说欢笑是合适的药。同样，我们也应该这样对待灵魂的疾病。懈怠引发了怒火吗？让我们用热忱驱散它，并在行为上表现出巨大的改变。我们有斋戒，这是一位非常强大的助手和盟友；除了斋戒，还有迫近的苦难和危险恐惧。那么，现在正是时候，让我们在灵魂上做工；因为我们很容易就能说服它去做我们选择的任何事；因为那惊慌恐惧、摆脱一切奢华、活在恐惧中的人，能够毫不困难地实践道德智慧，并非常乐意地接受美德的种子。\par

20. 因此，让我们说服它首先通过避免发誓来做出这个更好的改变；因为尽管我昨天和前天都对你讲了同一件事，但无论是今天、明天还是后天，我都不会停止就此提出劝告。我为什么说明天和后天呢？直到我看见你们改正了，我才会停止这样做。如果那些违反这条法律的人不感到羞耻，那么我们这些劝他们不要违反的人，更不应该觉得频繁劝告是可耻的事。因为不断提醒人同样的话题，并不是说话者的过错，而是听者的过错，因为他们需要持续教导简单且容易遵守的诫命。有什么比不发誓更容易的呢？这只是一个习惯的好行为。既不费力，也不耗费财富。你渴望学习如何克服这个软弱，如何摆脱这个恶习吗？我将告诉你一个特定的方法，如果你遵循它，你一定能掌握它。如果你看见你自己或任何人，无论是你的仆人、孩子或妻子，陷入这个恶习中；当你不断提醒他们而他们仍不改正时，命令他们不吃完晚饭就去休息；并把这判决加在他人和自己身上——一个不会带来伤害只会带来益处的判决。因为属神行为的本质就是如此：它们带来益处和迅速的改正。舌头经常受惩罚，因口渴而紧缩，因饥饿而痛苦，即使没有人监督，也得到充分的告诫；即使我们是最愚笨的人，当我们在整整一天中被惩罚的严重性所提醒时，我们就不需要其他的建议和劝告了。\par

21. 你们为我所说的鼓掌。但仍要以行动向我显示你们的鼓掌。否则我们在此聚会有何益处？假设一个孩子每天上学，但如果他什么也没学到，难道我们会对他说\Quote{他每天都去那里}而感到满意吗？我们岂不认为这是极大的过失，即天天去，却毫无效果？让我们自己思考这一点，并对自己说：这么长时间我们在教堂聚会，领受最庄严的圣体圣事，其中含有极大的益处；而如果我们回去时仍和来时一样，毫无改正缺点，那我们来这里有何益处？因为大多数行动都不是为了行动本身，而是为了通过它们所达到的效果；例如，播种者播种不是为了播种本身，而是为了他也能收获；因为若不如此，播种就是损失，种子腐烂而无任何益处。商人航行不是为了航行本身，而是为了通过出行增加财富；因为若不达到这一点，就会产生极大的祸害，商人的航行只是亏损。让我们也联系自身思考这一点。我们在教堂聚会，也不是仅仅为了在此消磨时间，而是为了回去时获得巨大的属灵益处。如果我们空手而回，毫无所得，那么这种勤奋反而成了对我们的谴责！为了避免此事和极大的祸害，在离开此地回家后，让朋友之间、父亲与子女、主人与仆人彼此操练，并训练自己完成所指定的功课；这样，当你们再次回来，听到我们就同一主题给予劝告时，你们就不会被控告的良心所羞愧，反而会欢喜快乐，因为你们发现自己已经履行了劝告的大部分。\par

22. 我们不要只在这里谈论道德。因为这种临时的劝告不足以根除全部的邪恶；在家里，丈夫也要从妻子那里听到这些事，妻子从丈夫那里听到。且让所有人之间有一种竞争，努力在履行这条律法上争先；让那已领先并改正了行为的人，责备那仍然落后的人，以便用这些讥讽更激励他。那有缺陷、尚未改正行为的人，让他看着那超过他的人，奋发竞争，赶快赶上他。如果我们对这些事采取忠告，并热切关心它们，我们其他的事务就会迅速调整。你要关心天主的事，祂必关心你的事！不要对我说：\Quote{如果有人强加给我们起誓的必要，怎么办？如果他不相信我们，怎么办？}因为确实，当一条法律被违反时，就不应提及必要；因为只有一种必要是不可免除的，即不冒犯天主！然而，我再说这一点：暂时砍掉多余的誓言，那些无益且无任何必要的誓言；向你的家人、朋友和仆人所发的誓言；如果你除去了这些，其他方面你就不再需要我了。因为那张被良好训练以畏惧和避免频繁发誓的口，即使有人强迫它一千次，它也决不会同意重新滑入同样的习惯。相反，正如现在，我们费了很大力气和极大的恳求，通过警告、威胁、劝告和忠告，几乎未能使它改变习惯；同样，在那情况下，即使有人施加再大的必要，也不能说服它违反这条法律。就像一个人绝不会选择服某种毒药，无论必要多么紧迫；同样，他也不会说出誓言！\par

23. 如果这项修正得以实现，它将鼓励并促使我们获得剩余部分的德行。因为那位一事无成的人会变得无精打采，迅速跌倒；但那位自觉至少履行了一条诫命的人，由此怀有良好的希望，将更敏捷地走向其余的诫命；这样，当他达到一条后，很快就会达到另一条；并且不会止步，直到获得一切冠冕。因为就财富而言，人得的越多，就越贪求；同样，在属灵的成就上更是如此。因此，我急忙并恳切希望这项工作得以开始，并且德行的根基能在你们灵魂中奠定。我们祈祷并恳求，你们不仅会在现在记住这些话，而且在家里、在市场以及你们度过的任何地方都会记住。哦！但愿我能与你们亲密交谈！那么我这次冗长的训话就不必要了。但现在既然不可能，你们要记住我的话，而不是记住我：当你们坐在餐桌旁时，设想我进来，站在你们旁边，将我现在在这个地方向你们说的话反复灌输给你们。无论你们中间有什么关于我的谈论，最重要的是记住这条诫命，并以此作为你们对我爱你们的回报。如果我看你们已履行了它，我就得到了完全的回报，并为我劳苦得到了足够的酬报。为了使你们既能使我们更加活跃，又使你们自己能享有良好的希望，并更容易地为完成剩余诫命做准备，请以极大的谨慎将这条诫命珍藏于你们灵魂中，你们就会明白这劝告的益处。正如一件金线刺绣的外衣是美丽而引人注目的，但当我们穿在自己身上时，它更显得如此；同样，天主的诫命在被赞美时是美丽的，但当被正确实践时，它们显得更加可爱。因为现在你们确实称赞那片刻所说的话，但如果你们付诸实践，你们将整天、终生称赞自己和你们自己。这不是重点，即我们将互相称赞；而是天主将接纳我们；不仅接纳我们，还要用那些伟大而不可言喻的恩赐赏报我们！愿我们众人被相称为配得这些恩赐，因我们的主耶稣基督的恩宠和慈爱，藉着祂、并与祂，偕同圣神，归于圣父，光荣现在及永远，世世无尽。阿们。\par

\SourceRecord{Translated by W.R.W. Stephens.；Nicene and Post-Nicene Fathers, First Series；Vol. 9.；Edited by Philip Schaff.；Buffalo, NY: Christian Literature Publishing Co.,；1889.；<http://www.newadvent.org/fathers/190105.htm>.}

% structure: p0001 source=h1 target=section reason=page-title

\section{第六篇讲道：论雕像}

\Emph{本篇讲道旨在表明对官吏的恐惧是有益的。其中也叙述了那些将叛乱消息传达给皇帝的人在旅途中所发生的事。又引约纳的事例作为说明。继续劝勉对死亡的恐惧；并指出，一个遭受不义之苦却因天主许可而感恩的人，就是为天主而受苦。再次从三个青年和巴比伦火窑的历史中举出例证。讲道以关于必须戒除誓言的劝言结束。}\par

1. 我们已花费多日向你们的爱德说安慰的话。然而，我们并不因此就将这主题搁置；只要沮丧的创伤仍在，我们就要不断敷上慰藉的药膏。因为论到身体的创伤，医生们若不等到疼痛消退，就不会停止敷药；对待灵魂，更应如此。沮丧是灵魂的创伤；因此我们必须不断用温柔的话语来敷它。因为温水软化肉体的硬瘤，其效力根本比不上安慰的话语平息灵魂激荡情感的能力。这里不需要医生的海绵，而是用舌头来代替。这里不需要火来烧热水；而是用圣神的恩宠来代替火。请容许我们今天这样做。因为如果我们不安慰你们，你们又能从哪里得到安慰呢？官吏恐吓，因此司铎必须安慰！统治者威胁，因此教会必须给予安慰！对待小孩子也是如此：教师恐吓他们，他们哭着跑到母亲那里；母亲将他们搂在怀里，抱着亲吻，擦干眼泪，抚慰忧伤的心灵，用言语说服他们，让他们明白害怕教师是有益的。既然统治者也使你们害怕，使你们焦虑，教会——我们众人的共同母亲——敞开怀抱，将我们抱在怀里，每日给予安慰，告诉我们：对统治者的恐惧是有益的，而来自这里的安慰也是有益的。因为前者的恐惧不容我们因懈怠而松弛，后者的安慰也不容我们在悲伤的重压下沉沦；天主通过这两者来保护我们的安全。祂亲自武装了官吏，使他们能够震慑放荡的人；祂也设立了司铎，使他们能安慰忧伤的人。\par

2. 圣经以及最近发生的事件的实际经验都教导我们这两件事。因为，当官吏和武装士兵存在时，少数人——一群乌合之众的疯狂——竟能在如此短暂的时间内在我们中间燃起大火，掀起风暴，使我们所有人都面临船毁人亡的恐惧；假如对官吏的恐惧完全消除，他们会疯狂到什么地步呢？他们岂不会把城市从根基上推翻，使一切天翻地覆，甚至夺走我们的生命？你若废除公共法庭，就是废除了我们生活中的一切秩序。正如你夺去船的舵手，船就会沉没；你从军队中撤去将军，士兵就会束手被敌人俘虏；同样，你若从城市中撤去统治者，我们的生活就会比野兽还不如，互相撕咬吞噬：富人吞噬穷人，强者吞噬弱者，大胆的吞噬温和的。但如今，因天主的恩宠，这些事一件也没有发生。因为生活在虔敬中的人不需要官吏的纠正；正如经上说：\Quote{法律不是为义人设立的。}\BibleRef{弟前1:9}有人这样说。但多数人倾向邪恶，如果没有恐惧悬在他们头上，他们就会给城市带来无数灾祸。保禄知道这一点，他说：\Quote{没有权柄不是从天主来的，现有的权柄都是由天主规定的。}\BibleRef{罗13:1}因为如同房屋中的横梁一样，统治者在城市中也是如此；你若抽去横梁，墙壁就会分离，自行倒塌；同样，你若从世上除去官吏以及对他们的恐惧，房屋、城市和民族立刻就会陷入无节制的混乱，互相倾覆，因为没有人能用惩罚的恐惧来遏制、阻止或说服他们保持和平！\par

3. 所以，亲爱的诸位，我们不要因对统治者的恐惧而悲伤，而要感谢天主，因为祂除去了我们的懈怠，使我们更加勤勉。请告诉我，这种担忧和焦虑带来了什么害处？我们不是变得更庄重、更温和、更勤勉、更热心了吗？我们看不到有人醉酒、唱淫荡的歌曲了吗？或者，不是有不断的祈求、祈祷和泪水吗？不合时宜的欢笑、污秽的言语和一切放荡不都被驱逐了吗？城市如今在各方面都像一位端庄贤淑的妇人的榜样？我问，你们会为这些事中的哪一件而悲伤呢？这些事，无疑应当欢喜，并感谢天主，因为几天的恐惧就终结了如此的愚昧！\par

\Quote{说得很好，}有人说，\Quote{如果我们的危险不超出恐惧，我们本已获得足够的益处；但现在我们担心祸患会发展得更远，我们所有人都将陷入极端的危险。}\par

然而，我説，不要害怕。\Person{保祿}安慰你説：\Quote{天主是忠信的，祂不許你們受試探超過你們所能受的，但在試探的同時，也要開一條出路，叫你們能忍受得住。}\BibleRef{格前 10:13} 祂親自説過：\Quote{我決不離開你，也決不棄舍你。} 因為，如果祂定意要我們在實際行動中受罰，就不會讓我們在這麼多天內一直處於恐懼之中。因為當祂不願處罰時，祂才令人驚恐；如果祂打算懲罰，恐懼便是多餘的，威脅也是多餘的。但如今，我們經歷了比無數次死亡更痛苦的生活：這麼多天恐懼戰兢，連自己的影子都懷疑，承受了\Person{加音}的刑罰；睡夢中也因心靈持續的痛苦而驚醒。因此，如果我們激起了天主的義怒，我們在承受這樣的懲罰時已平息了祂的怒氣。因為即使我們沒有完全賠補所犯的罪，但這已足夠滿足天主的仁慈了。\par

4. 但不僅如此，我們還應有許多其他信心的理由。因為天主已經給了我們不少有利希望的保證。首先，那些攜帶噩耗的人如飛般離開此地，以為早已抵達軍營，卻仍然在旅途中耽擱。如此多的障礙和阻滯出現了；他們丟了馬匹，現在乘車前行；因此他們的到達必定會推遲。因為天主在這裡激動了我們的司鐸和共同的父親，説服他出發並承擔這項使命，同時在半路上攔住了那些使者，免得他們先到而點燃怒火，使我們老師的努力在御耳發怒時化為無用。因為這旅途中的阻礙顯然不是沒有天主的干預。這些人一生熟悉這樣的旅途，又常騎馬，如今卻因騎馬疲勞而倒下；所以現在發生的事與\Person{約納}的情況相反。因為天主催促不願意的\Person{約納}去完成使命；但這些人願意去，祂卻阻止了。哦，奇妙非凡的事！祂不願宣講毀滅，卻強迫\Person{約納}違背自己的意願前往；這些人急忙出發去傳報毀滅的消息，祂又違背他們的意願阻止了他們！你認爲是什麼原因？因為在這個情況下，匆忙是有害的；而在那個情況下，匆忙帶來益處。因此，祂藉著鯨魚催促\Person{約納}前進，又藉著馬匹阻止了這些人。你看到天主的智慧嗎？每一方希望藉以達成目的的手段，反而成了他們的障礙。\Person{約納}指望藉船逃脫，船卻成了他的鎖鏈；這些使者指望藉馬更快見到皇帝，馬卻成了障礙；或者更確切地說，既不是馬，也不是船，而是天主的眷顧處處按自己的智慧引導一切！\par

5. 也請思量祂對我們的眷顧，以及祂如何既驚嚇我們又安慰我們。因為祂在發生所有暴行的當天就允許他們出發，好像他們要向皇帝報告一切；祂使我們所有人都因他們突然離去而驚惶。但當他們離開兩三天後，我們以爲我們司鐸的行程現在已無用，因為他到得太晚，祂卻通過在半路上留住他們（如我所觀察的），並通過從那條路上來的人向我們報告他們旅途中所遇的一切困難，使我們脫離了恐懼，得了安慰，從而我們可以稍微喘口氣，事實上也的確如此，我們的憂慮大大減輕。聽到這事，我們朝拜了成就此事的天主，祂現在比任何父親更溫柔地為我們安排一切，以某種看不見的力量延遲了那些惡信的使者，並似乎向他們説：\Quote{你們爲何著急？爲何急於前行，要去傾覆這樣一座大城市？難道你們要給皇帝帶去好消息嗎？等候在那裡，直到我預備好我的僕人，如同一位傑出的醫生，趕上你們，搶在你們前面。} 如果在這罪惡傷口初次爆發時就有如此多的眷顧，那麼在悔改、痛悔、如此多的恐懼、眼淚和祈禱之後，我們必將得到更大的平安。因為\Person{約納}確實被迫如此，以便他得以被強行帶到悔改；但你們已經表現出悔改和皈依的顯著標記。因此，你們應當得到安慰，而非威脅的信息。為此緣故，祂也派我們的共同父親從這裡出發，儘管有許多阻礙。如果祂不關心我們的安危，祂就不會説服他這樣做，反而會阻止他，即使他多麼願意踏上旅程。\par

6. 还有第三个理由，或许能说服你们心怀信赖；那就是当前的神圣时节，连几乎所有不信者也尊重它；但我们这位蒙天主眷顾的皇帝，对它表现出的敬畏与尊荣，竟超越了他以前所有秉持宗教热情统治的皇帝。证据是：他在这些日子，为庆祝这节日而发出一封诏书，几乎释放了所有被关押在监狱中的人；我们的司铎一到，就会向他宣读这封诏书，提醒他自己的法令，并对他说：\Quote{你要劝勉自己，记住你自己的作为！你在自己的家里就有一个仁爱的榜样！你选择了饶恕一场正当的杀戮，难道你竟会容忍施行一场不义的杀戮吗？你因尊重节日而释放了那些已被定罪判刑的人，请问，你现在还会定那无辜者、那没有行任何暴力的罪吗，何况神圣的时节当前？皇帝，千万不可！你通过这封书信对各城说：\InnerQuote{但愿我甚至能复活死人。}我们现在所需要的，正是这仁爱与这些话语。战胜敌人，并不能使君王如战胜愤怒与怒气般荣耀；因为前者之功归于武器与士兵，而这里的凯旋纯粹属于你自己，无人能与你分享这道德智慧的荣耀。你已战胜了野蛮人的战争，再战胜皇帝的怒气吧！愿所有不信者学习到，对基督的敬畏能抑制一切权柄。藉着宽恕同仆的过犯来光荣你的主，使祂也能更加光荣你；这样在审判之日，祂会以仁慈与平静的目光注视你，记念你的仁爱！\Quote} 他将说这些话，以及更多的话，并必能挽救我们脱离皇帝的愤怒。不仅这斋戒期在影响皇帝善待我们方面对我们有极大助益，而且对于坚忍地承受所遭遇之事，我们也从这个季节获得不小的安慰。因为我们每天聚集在一起，得以聆听神圣圣经，彼此见面，一同哭泣，一同祈祷，领受祝福，然后回家，这已卸去了我们大部分的痛苦。\par

7. 因此，我们不要沮丧，也不要因痛苦而放弃自己；让我们等待，期望一个有利的结局；并留意现在就要说的话。因为我打算今天再次与你们谈论轻看死亡。我昨天对你们说，我们怕死，不是因为它真正可怕，而是因为对天国的爱没有点燃我们，对地狱的恐惧没有抓住我们，而且除此之外我们没有良好的良心。你们愿意我谈谈这不合时宜之痛苦的第四个原因吗？它不比前几个小，而且更加真实：我们不过与基督徒相称的克修生活。相反，我们喜爱追随这种纵欲、放荡、懒惰的生活；因此我们依恋现世之物也是自然的；因为如果我们把此生用在禁食、守夜、粗茶淡饭上，斩断一切过分的欲望，约束我们的享乐，经历德行的辛劳，像保禄一样制服身体，使它服从，不为肉体的私欲作准备，并走窄门窄路，我们就会热切渴望未来之事，急切想摆脱目前的劳苦。为证明我所说的不假，请登上山顶，观察那里的隐修士：有的身着苦衣，有的身负锁链，有的在禁食，有的被关在黑暗中。你们会发现，所有这些人都急切地渴望死亡，称之为安息。因为正如拳击手急于离开竞技场，以摆脱伤口；摔跤手盼望赛事结束，以从劳累中解脱；同样，依靠德行过着克修与苦行生活的人，也急切地渴望死亡，为从目前的劳苦中得释放，并且能对储存在那里的冠冕充满确信，因为藉着抵达宁静的港湾，迁移到不再有失事忧虑的地方。因此，天主也为我们提供了天生劳苦烦扰的生活，为的是我们在此受磨难催促时，能对未来福乐产生渴望；因为如果现在，在这么多忧愁、危险、恐惧、焦虑从四面环绕我们时，我们尚且这样依恋现世生活，那么，如果我们的现世存在完全没有忧伤与痛苦，我们何时才会渴望来世生活呢？\par

8. 天主也这样对待犹太人。祂愿意激发他们回归（客纳罕）的渴望，并说服他们憎恨埃及，便容许他们在黏土和制砖的劳苦中受压迫，好让他们被这沉重的辛劳和苦难所压迫时，为回归而向天主呼号。因为如果他们在这些事发生后离开时，竟然又想起埃及和那残酷的奴役，急于回到那先前的暴政之下；那么假如他们从未受过这些野蛮人的如此对待，他们何时才会愿意离开那异乡呢？因此，为使我们不至过分依恋尘世，因贪恋眼前事物而变得可怜，并忘记将来，天主使我们今世的生活充满了劳苦。那么，让我们不要怀抱超过必要限度的对现世生活的爱。它对我们有何益处？或者，紧紧铆钉于对现世状态的渴望有何好处？你们愿意知道这生活有何益处吗？它是来世生活的基础和起点；是摔跤学校和日后赢得冠冕的竞技场！因此，如果它没有为我们提供这些，它就比千次死亡更糟。因为如果我们不愿生活以悦乐天主，那么死更好。有何增益？我们多得什么？我们不是每天看见同样的太阳和月亮，同样的冬天和夏天，同样的事物进程吗？\Quote{已有的事，后必再有；已行的事，后必再行。} \BibleRef{训 1:9} 那么，让我们不要立刻称活着的人为有福，并为死者哀哭；而要哀哭那些处于罪恶中的人，无论他们是死是活。另一方面，让我们称那些处于正义状态的人为有福，无论他们处于何种状况。你确实害怕并哀叹\Quote{一种}死亡；但保禄天天冒死，\BibleRef{格前 15:31} 他不但没有为此流泪，反而喜乐和欢欣！\par

9. 有人说：\Quote{但愿我为天主受苦，我便无忧无虑了！} 但即使在现在，也不要陷入沮丧；因为不仅为天主受苦的人蒙受悦纳，而且那无故受苦、并以高尚精神忍受、且感谢允许此事的天主的人，并不逊于为天主承受这些考验的人。蒙福的约伯就是证明，他因魔鬼无益、徒然和无缘无故的谋害而遭受了那么多难以忍受的创伤。然而，因为他勇敢地承受了它们，并感谢允许此事的天主，他便获得了完美的冠冕。因此，不要因死亡而悲伤；因为死亡是自然的；但要为罪而忧伤，因为那是意志的过错。但如果你为死者悲伤，也要为那些出生到世界上的人悲伤；因为正如一件事是出于自然，另一件事也是出于自然。因此，如果有人以死亡威胁你，要对他说：\Quote{基督教训我\InnerQuote{不要害怕那些杀害身体，却不能杀害灵魂的。}} \BibleRef{玛 10:28} 或者他威胁要没收你的财产，你要对他说：\Quote{我赤身出于母胎，也必赤身归回。我们没有带什么到世上来，也不能带什么去。} \Quote{纵然你不夺去我，死亡也会来夺去我；纵然你不杀我，自然律也会立刻介入并带来结局。} 因此，我们不应害怕那些因自然秩序而临到我们的事，而应害怕那些由我们自己的邪恶意志所产生的事；因为这些事会带来惩罚。但让我们不断思考这一点：就那些突然临到我们的事件而言，我们不会因悲伤而改善它们，因此我们将停止悲伤。\par

10. 此外，我们还应当这样想：如果我们今生不义地遭受任何灾祸，我们便清偿了许多罪恶。因此，在这里而非在那里承受我们罪的惩罚，实为莫大的裨益；因为那个富人在这里没有遭受任何灾祸，所以他在那里被火焰灼烧；而这就是他得不到任何安慰的原因，请听亚巴郎所说的话作为证明：\Quote{孩子，你已得到了你的福乐；因此你受痛苦。}\newline{}至于赐予拉匝禄的福乐，不仅他的德行，而且他在这里忍受了千般痛苦也作出了贡献，这你也可以从圣祖的话中得知。因为他向富人说完\Quote{你已得到了你的福乐，}紧接着就说：\Quote{而拉匝禄受了苦，因此他得了安慰。} \BibleRef{路 16:25}\newline{}因为正如那些生活圣洁却受磨难的人，从天主那里获得双倍报酬一样，那生活在邪恶中且享尽奢华的人，也将受双倍的惩罚。我再说，我说这话并非为了指责那些已经逃走的人，因为经上说：\Quote{不要使忧伤的心增添更多的烦恼；} \BibleRef{德 4:3}\newline{}我这样说也不是为了谴责（因为病人需要安慰），而是为了努力推动改善。我们不要将我们的安全托付给逃跑，而要逃离罪恶，离开我们的邪道。如果我们逃离了这些，即使身处千万士兵之中，也没有一个能打击我们；但如果不逃离这些，即使我们登上山巅，也会在那里遇到数不尽的敌人！让我们再次想起那三位青年，他们在火窑中却未受任何伤害，而将他们投入火窑的人，那些坐在周围的全都被烧死了。还有什么比这更奇妙的呢？那火焰释放了它所抓住的人，却猛烈地抓住它未曾抓住的人，为要教导你：带来安全或惩罚的不是居所，而是生活方式。在火窑内的人得以逃脱，而外面的人却被烧死。对两者来说，身体都一样，但心境却不同。因此，他们身上的效果也不相同；因为干草即使不在火焰中，也会很快被点燃；但金子即使留在火中，也会变得更加璀璨！\par

11. 现在，那些说\Quote{让皇帝拿走一切，只给我们自由的身体？}的人在哪里？让他们去学习什么是自由的身体。使身体自由的不是免于惩罚，而是坚持正义的生活。例如，这些青年的身体是自由的，尽管他们被投入火窑，因为他们已经摆脱了罪恶的奴役。因为这才是真正的自由，而不是免于惩罚或免于遭遇可怕之事。但是，听到火窑，你要想起在那可怕之日将要出现的\Quote{火河}。因为正如上面所发生的，火焰抓住了某些人，却敬畏了另一些人，那些火河也将如此。如果任何人那时有干草、木柴、碎秸，他就增加了火焰；但如果他有金子、银子，他就变得更加光耀。因此，让我们积聚这种材料，并勇敢地承受目前的处境；因为我们知道，如果我们懂得如何实践真正的智慧，这苦难不仅能带我们脱离那惩罚，而且能在此世使我们变得更好；并且不仅是我们，若我们警醒，常常连那给我们带来麻烦的人也会变好；这属神智慧的力量如此丰富，正如当时发生在暴君身上的事。因为当他得知他们未受任何伤害时，请听他如何改变了他的话语。\Quote{至高天主的仆人，出来，到这里来。} \BibleRef{达 3:26}\newline{}你刚才不是说过\Quote{谁能从那一位天主手中救出你们？} \BibleRef{达 3:15}\newline{}发生了什么事？这转变从何而来？你看见外面的人被毁灭，却呼唤里面的人？你怎么在这些事上变得明智了？你看君主发生了多么大的变化！当他尚未对他们行使权力时，他亵渎了天主；但一旦将他们投入火中，他就开始表现出道德智慧。因此，天主也容许一切发生，暴君想做什么就做什么，为要显明：凡被祂保守的人，无人能加害他们。祂对约伯所做的，在这里也成就了。因为在那一场合，祂也容许魔鬼施展其全部力量；直到他耗尽所有箭矢，再无计谋可施时，那战士才被带出战场，使胜利辉煌而无可置疑。所以在这里，祂也行了同样的事。他（暴君）愿意推翻他们的城，天主没有阻止他；他要把他们掳去，天主没有拦阻他；他要捆绑他们，天主容许；要把他们投入火窑，天主也允许；要把火焰烧得超过限度，天主也任凭；当暴君再没有什么可做的，耗尽了他所有的力量时，天主才彰显了自己的力量和青年的忍耐。你看见了吗？天主容许这些磨难直到最后，为要向攻击者显明被攻击者的属神智慧，以及祂自己的眷顾。那时，那人也看清了这两者，便喊道：\Quote{至高天主的仆人，出来，到这里来。}\par

12. 但请你与我一同思考这些青年的宽宏大量；因为他们既不在召唤之前跳出来，免得有人以为他们害怕火；也不在被召唤时仍留在里面，免得有人认为他们好胜好斗。他们说：\Quote{当你们知道我们是哪位仆人的时候，当你们承认我们的主的时候，我们就出来，向在场所有人成为天主能力的传报者。}不如说，不仅他们自己，连敌人也亲口，是的，既亲口又通过他的谕令，向所有人宣告了战士们的坚贞和主导这场竞赛的那位的伟力。正如传报者在剧场中央宣告胜利者的名字时，也会提到他们所属的城邦：\Quote{某人，某城的人！}同样，他代替他们的城邦，宣告了他们的主，说：\Quote{\Person{沙得辣客}、\Person{默沙客}、\Person{阿贝得乃哥}，至高天主的仆人，出来，到这里来。}发生了什么事，使你称他们为天主的仆人？他们不是你的仆人吗？他说：\Quote{是的，但他们推翻了我的主权；他们践踏了我的骄傲。他们用行动表明，他才是他们真正的主。如果他们是人的仆人，火就不会畏惧他们；火焰就不会给他们让路；因为受造物不懂得尊重或荣耀人的仆人。}因此他又说：\Quote{愿\Person{沙得辣客}、\Person{默沙客}、\Person{阿贝得乃哥}的天主受赞美。}\par

13. 请也与我一同思考，他首先如何宣告竞赛的仲裁者。\Quote{愿天主受赞美，祂派遣了自己的天使，拯救了祂的仆人。}\BibleRef{达 3:28} 这是关于天主的权能。他也谈到战士们的德行。\Quote{因为他们信赖祂，改变了国王的命令，并交出了自己的身体，以免敬拜任何神，只敬拜他们自己的天主。}有什么能与这样的德行相比？在此之前，当他们说：\Quote{我们不会侍奉你的神们}时，他比火炉本身更猛烈地燃烧起来；但现在，当他们的行为教导了他之后，他不但不愤怒，反而赞美和钦佩他们，因为他们没有服从他！德行是如此美好，以至于它甚至让敌人自己也来鼓掌和钦佩它！这些人都战斗并得胜了，但失败的一方却感恩，因为火焰的景象没有吓倒他们，而是对主的希望安慰了他们。他以这三个青年来命名普世的天主，丝毫没有限制祂的主权，因为这三个青年等同于整个世界。因此，他既为那些藐视他的人喝彩，又越过众多总督、国王和王子，那些服从他的人，却对这三个俘虏和奴隶，这些嘲笑他的暴政的人，表示赞叹！因为他们做这些事，不是为了争辩，而是出于对智慧的热爱；不是出于反抗，而是出于虔诚；不是出于骄傲自负，而是出于热忱。因为对天主的希望确实是极大的祝福；这一点连那蛮族人也学到了，并清楚表明正是从那里他们逃脱了迫在眉睫的危险，他大声喊道：\Quote{因为他们信赖了祂！}\BibleRef{达 3:28}\par

14. 但我现在说这一切，并选取所有包含考验和磨难、国王的愤怒及其恶谋的历史，为使我们除了得罪天主之外，一无所惧。因为那时也有火炉在燃烧；然而他们嘲笑了它，却惧怕犯罪。因为他们知道，如果他们在火中被烧尽，他们不会遭受任何可怕的事；但如果他们陷于不敬，他们将遭受极端的痛苦。最大的惩罚就是犯罪，即使我们可能不受罚；另一方面，最大的荣誉和安息就是过有德行的生活，即使我们可能受罚。因为罪恶将我们与天主隔绝；正如祂亲自所说：\Quote{你们的罪恶不是将你们与我分开了吗？}但惩罚将我们引领回天主。如一人所说：\Quote{赐予平安，因为祢为了所有事补偿了我们。}假设有人受了伤；最可怕的是什么：是坏疽，还是外科医生的刀？是钢铁，还是溃疡的吞噬性进展？罪恶是坏疽，惩罚是外科医生的刀。那么，正如一个患坏疽的人，即使没有被割开，也必须承受疾病，而且当他不被割开时，情况更糟；同样，罪人即使不受惩罚，也是最不幸的人；而当他没有惩罚、没有遭受痛苦时，他就尤其不幸。正如那些患有脾病或水肿的人，当他们享受丰盛的餐桌、冷饮和各种美味佳肴时，那时他们尤其处于最可怜的状态，因为他们的奢华加剧了疾病；但如果他们按照医学定律严格地约束自己忍受饥饿和干渴，他们可能有一些康复的希望；同样，那些生活在不义中的人，如果他们受到惩罚，可能会有有利的希望；但如果他们除了邪恶之外，还享有安全和奢华，他们比那些虽然水肿却填饱肚子的人更不幸；而且灵魂比身体更优越，所以程度更深。因此，如果你看到一些人在同样的罪中，有些人不断与饥饿和无数疾病搏斗，而另一些人则开怀畅饮、生活奢侈、暴饮暴食；那么你要认为那些忍受痛苦的人更好。因为不仅情欲的火焰被这些不幸切断，而且他们离开现世，前往未来的审判和那令人畏惧的法庭时，也带着不小的解脱；他们在今世通过所遭受的苦难偿还了大部分罪的惩罚。\par

15. 但安慰已足够。现在，我们终须转向关于避免起誓的劝诫，并揭穿那些起誓者看似合理却徒劳无益的托词。因为我们指控他们时，他们便举出同样行事的人为例，说：\Quote{某某人也起誓。} 我们当对他们说：尽管如此，那人不一定起誓；天主将凭善行来审判你们，因为罪人不会因同犯过犯而彼此得益；而行义的人却定罪罪人。因为许多人不曾给基督食物或饮料，但他们彼此无益。见：\BibleRef{玛 25:35} 五个童女的情况也类似，她们并未因同伴而获赦免，见：\BibleRef{玛 25:10} 反因与明智者的比较而被定罪，这两种人都同样受罚。\par

16. 因此，让我们抛弃这虚假自欺的辩论，不要注视那些跌倒的人，而要注视那些行为端正的人；并努力在封斋结束后仍念念不忘此次斋戒。正如我们常有的情形：当我们买了一件衣裳、一个奴隶或一件珍贵的器皿时，我们会回想购买的时刻，彼此说道：\Quote{那个奴隶是我在某节日买的；那件衣裳是我在某时买的；} 同样，如果我们现在将此法则付诸实践，我们便要说：\Quote{我在那次四旬期改掉了起誓的恶习；因为在那之前我常起誓；但仅仅听了劝诫后，我便戒绝了这罪过。}\par

但有人反驳说：\Quote{习惯难以改正。} 我知道是这样；因此我催促你们养成另一种良好的有益习惯。因为当你们说\Quote{我难以戒除习惯}时，正因如此，我说你们应当赶紧戒除，因为你们确知，一旦你们为自己养成了不起誓的习惯，之后便不费功夫。哪样更难：不起誓，还是整天不吃不喝，仅靠饮水粗食度日？显然后者更难；然而，习惯已使这变得如此可行和容易，以至于当斋期来临时，即使有人千次劝勉，或屡次强逼人喝酒或品尝斋戒所禁之物，一个人宁愿受苦也不愿碰那禁物；这并非由于不喜爱筵席，但我们凭良心的习惯都坚忍承受。起誓之事亦然；正如现在，即使有人施加极大压力，你们也会因固守习惯而坚定不移；同样，在那事上，即使有人万次催促，你们也不会背离习惯。\par

18. 因此，你们回家后，要将这一切与家中人谈论；正如许多人常做的：从草地归来时，摘一朵玫瑰、紫罗兰或此类花朵，在指间捻玩；又如有人离开花园回家时，带着带果的树枝；还有人从盛宴中为家人带回残羹；同样，你们也要从此处带走劝诫，带给妻子、儿女和全家。因为这劝诫比草地、花园或宴席更有益。这些玫瑰永不凋谢；这些果实永不坠落；这些美味永不腐朽。前者带来短暂的欢愉，后者却带来持久的益处，不仅在改过之后，甚至在改过之中。试想，若抛开一切公私事务，不断只谈论神圣律法\BibleRef{申 6:7}，在餐桌、广场及其他聚会中，这是何等的美事！若我们专注于此，便不会说出危险或有害的话，也不会无意中犯罪。若我们闲暇时谈论这些，便能将灵魂从笼罩我们的沮丧中抽离，而不像现在这样焦虑，彼此说道：\Quote{皇帝听说了所发生的事吗？他发怒了吗？他下了什么判决？有人向他请愿了吗？什么？他自己会忍心毁灭一座如此广大繁华的城市吗？} 让我们将这一切忧虑交托给天主，只关心他所命令的！这样我们便能摆脱所有愁苦；即使我们中只有十人成功，十人会很快变成二十；二十变为五十；五十变为一百；一百变为一千；一千充满全城。正如点燃十盏灯，便能轻易照亮全屋；同样，若只有十人行善，我们便能在全城点燃普遍的火焰，照耀并带给我们安全。因为火落在森林中，会依次点燃邻近树木，而美德的热忱一旦抓住几个心灵，其进展之强大，足以蔓延至整个团体。\par

19. 那么，请赐我因你们而欢欣的缘由，无论是在今生，还是在将来那个日子——那时，那些受托付了塔冷通的人将被召来！你们的好名声已足以酬报我的辛劳；若我见你们度虔敬的生活，我便心满意足。所以，请做我昨日所劝勉、今日将重复、且永不止息地言说的事：为那些发誓的人定一个惩罚——这惩罚是获益而非损失；并且从今以后，你们要预备自己，好能向我们证明你们的成功。因为当这次聚会解散时，我将设法与你们每人进行长谈，以便在持续的交谈中发现哪些人行得正直，哪些人没有。若我发现有人仍在发誓，我要向所有已改正的人指明他，以便通过责备、训斥和纠正，我们赶快使他脱离这恶习。因为他在这里受责备而改过，远胜于在那一日，当我们的罪过在众人眼前被揭露时，在全宇宙的集合面前蒙羞受罚！但天主不许在这美好集会中的任何人，在那里遭遇这样的事！而是借着圣教父们的祈祷，纠正我们所有的过犯，并结出丰硕的德行之果，使我们满怀信心地离开此地，借着我们的主耶稣基督的恩宠和慈爱，愿光荣归于圣父及圣神，于无穷世之世。阿们。\par

\SourceRecord{Translated by W.R.W. Stephens.；Nicene and Post-Nicene Fathers, First Series；Vol. 9.；Edited by Philip Schaff.；Buffalo, NY: Christian Literature Publishing Co.,；1889.；<http://www.newadvent.org/fathers/190106.htm>.}

% structure: p0001 source=h1 target=section reason=page-title

\section{第七篇讲道：论雕像}

\Emph{前番劝勉的综述。罪恶为世界带来死亡与悲伤，而此二者有助于医治罪恶。悲伤唯独为消灭罪恶而有用。对经文\BibleRef{创 1:1}的评论。\Quote{在起初，天主创造了天地。}论证天主在创造工程中为人所做的眷顾提供安慰；即使在惩罚中也显示慈悲，正如这句话：\Quote{亚当，你在哪里？}最后劝勉避免起誓。}\par

1. 昨天，我向你们的爱德讲了诸多语句，涉及许多主题；如果在这许多内容中你们无法全部记住，我尤其想提醒你们记住这一观察：天主在我们本性中植入悲伤这种情感，唯一的原因就是罪恶，祂已从实际经验中证明了这一点。因为当我们因失去财富而悲伤痛苦时，或因疾病、死亡及临于我们的其他灾祸而悲伤时，我们非但不能从忧愁中得到安慰，反而加剧了这些灾祸的力量。但如果我们因自己的罪而痛苦悲伤，我们就减轻了罪的重量；我们使大的变小，常常甚至将它完全涂去。我再说一遍，你们应当不断记住这一点，以便你们只为罪而哀伤，不为别的；还有这一点：罪虽然为我们的生命带来了死亡和悲伤，却又被这两者所消灭，这是我最近已经证明的。因此，让我们什么也不怕，只怕罪与过犯。让我们不要怕惩罚，这样我们就会免于惩罚。正如三位青年不怕火窑，因而逃离了火窑。天主的仆人正应当如此。因为那些在旧约时代长大的人，那时死亡尚未被消灭，他的\Quote{铜门尚未被破坏}，他的\Quote{铁闩尚未被砍断}，他们却如此英勇地面对死亡；我们既已蒙受了如此伟大的恩宠，却连与他们同等的德行也达不到，那我们将是多么没有辩护或借口呢？况且如今死亡只是个名字，毫无真实。因为死亡不过是睡眠、旅程、迁移、休息、宁静的港湾；是脱离烦恼，摆脱此生劳虑！\par

2. 但让我们在此搁置安慰的话题；我们已是第五天向你们的爱德讲安慰的话语，可能已显得烦人了。因为对于留心倾听的人，已经说的足够了；但对于胆小怯懦的人，即使我们再添加什么，也无济于事。现在是时候将我们的教导转向圣经的讲解了。因为，如果我们对当前的灾祸只字不提，人们可能会指责我们冷酷无情、缺乏人道；但如果我们总是谈论此事，我们也会被公正地指责为胆小怯懦。那么，将你们的心交托给天主，祂能在你们心中说话，驱散一切忧愁，现在就让我们按往常的方式开始教导；尤其是因为每一段圣经讲解都是安慰和缓解。所以，虽然我们似乎放弃了安慰的话题，但通过讲解圣经我们又会触及同一主题。因为全部圣经都为留意它的人提供安慰，我要从圣经本身的证据向你们证明这一点。我不打算在圣经叙事中搜罗某些安慰的论证；而是为了使我所承担的事的证明更清晰，我们要拿起今天读给我们听的那本书；如果你们愿意，我们提出它的引言和开端，这似乎尤其不带有安慰的痕迹，而是全然与安慰的话题无关，我将使我所肯定的事显明出来。\par

3. 那么，这引言是什么呢？\Quote{在起初，天主创造了天地，地是无形无状的，黑暗笼罩在深渊之上。}这些话在你们有些人看来不能给困境带来安慰吗？难道这不是一段历史叙述，关于创造的教导吗？\par

那你们愿意我指出隐藏在这句话里的安慰吗？那么就警醒自己，以热忱聆听即将讲述的事。因为当你们听说天主为了你们、为了你们的得救和尊荣而造了天、地、海、空气、水、繁星、两大光体、植物、四足动物、游鱼和飞鸟，以及你们所见的万物，无一例外时；当你们想到祂为你这样一个微小的人，创造出如此美丽、广阔而奇妙的宇宙，难道不立刻得安慰，并将此视为天主最强烈的爱的证明吗？因此，当你们听到\Quote{在起初天主创造了天地}时，不要草率地读过这句话；而要在心里遍历大地的宽广，默想祂如何为我们摆设如此丰盛精美的筵席，赐予我们如此丰盈的喜乐。而这确实是最奇妙的事：祂赐予我们这样的世界，并非因我们有所服事，也不是作为善工的酬报；相反，祂在造我们之时，就以这个王国尊荣了我们的种族。因为祂说：\Quote{让我们照我们的肖像，按我们的模样造人。}\BibleRef{创 1:26} 这话所指为何：\Quote{照我们的肖像，按我们的模样}？这里指的是统治的肖像；正如天上无人高于天主，地上也无人高于人。这是我们受尊荣的第一方面：使人依照祂的肖像；第二方面，祂赐予我们这治权，非因服事而偿付，而是完全出自祂对人的爱；第三方面，祂将此作为本性的事赐予我们。因为统治有些是自然的，有些是选择的——自然的统治如狮子对四足动物，或如鹰对飞鸟；选择的统治如皇帝对我们；他并非因自然权柄统治其同仆。因此他常会丧失主权。因为非自然固有的事易于变化和转移。但狮子并非如此：它依自然统治四足动物，正如鹰统治飞鸟。因此，统治的特性常赋予其种族，从没见过哪只狮子被剥夺此种统治。天主从起初就赐予我们这种统治，使我们统管万物。祂不仅在这方面尊荣我们的本性，还以我们所处之地的卓绝——选择乐园为我们的佳美居所——并赐予理性与不朽灵魂。\par

4. 但我却不愿谈论这些事；因为我说：天主的眷顾如此丰盛，使我们不仅从祂尊荣我们的方式，也从祂惩罚我们的方式，知道祂的良善与祂对人的爱。这一点我特别劝诫你们认真思考：天主既是良善的，不仅在祂以尊荣和恩惠待我们时，也在祂惩罚和责打我们时。无论我们是要与异教徒或异端者争辩关于天主的慈爱和良善，我们都将使祂的良善显明，不仅从祂施恩的情形，也从祂施罚的情形。因为如果祂只在尊荣我们时是良善的，而在惩罚我们时不是良善的，那么祂只是半个良善。但事实并非如此。天主禁止！在人中，这大概会发生，因为他们在忿怒和情欲中施罚；但天主既不受情欲影响，无论祂施恩还是施罚，祂都是良善的。地狱的威胁丝毫不比天国的应许更少显示祂的良善。\BibleRef{迦 3:24} 这是如何呢？我回答：如果祂没有威胁地狱，如果祂没有预备惩罚，可能没有多少人能达到天国。因为福乐的应许并不像恶事的威胁那样，借恐惧强制并唤醒他们照顾灵魂，从而有力地促使众人趋向美德。因此，虽然地狱与天国对立，但两者都朝向同一目的——人的救恩：一个吸引人，另一个驱赶人走向其反面，并借恐惧的作用纠正那些漫不经心的人。\par

5. 我并非无缘无故详述这主题；而是因为有许多人常常在饥荒、干旱、战争，或皇帝发怒临头，或其他类似意外事件发生时，欺骗较单纯的人，说这些事与天主的眷顾不相称。\par

因此，我不得不在此部分讲论中多加停留，以免我们被言语所迷惑，而能清楚明白：无论祂给我们带来饥荒、战争或任何灾祸，都是出于祂极其伟大的眷顾与慈爱。因为即使是那些特别疼爱子女的父亲，也会禁止他们上桌，施加鞭打，以羞辱惩罚他们，并用无数其他方式管教行为不端的儿女；然而，他们仍然是父亲，不仅在他们给予尊荣时，而且在如此行事时也是如此；是的，当他们如此行事时，他们尤其身为父亲。\BibleRef{希 12:9} 但如果人常常被愤怒的情感所驱使而越过分寸，却被认为惩罚他们所爱的人并非出于残忍和不人道，而是出于某种关怀与顾念；那么，对于天主，更应当怀有这样的心意——祂的良善极其丰盈，远超任何父爱的程度。为使你们不以为这只是我的揣测，让我们，我恳求你们，将我们的论述转向圣经本身。当人被邪恶的魔鬼欺骗和迷惑后，让我们观察天主在犯下如此大罪之后如何对待他。祂是否彻底毁灭了他？然而公义本身的要求本应如此：一个没有表现出任何善行，反而在享受如此大恩之后，从起初就放纵无度的人，应当被除掉并彻底毁灭；但天主并没有如此行事；祂也没有对那位如此忘恩负义于恩主的人感到厌恶和厌弃，反而如同医生来到病人身边一般，来到他面前。\par

6. 亲爱的，不要轻率地忽略刚才所说的话！但请思考：这并非派遣天使、总领天使或任何其他同仆，而是主亲自降下，到那偏离正道、跌倒的人那里，扶起那被抛下的人，并接近他，一对一，如同朋友来到不幸且深陷极大痛苦的朋友身边！因为祂如此行事是出于祂极大的慈爱，祂对他所说的话也清楚显明了祂难以言喻的深情。我为何要说\Emph{所有}话呢？第一句话立刻显明了祂的温柔。因为祂没有说出一个被如此轻慢的人可能会说的话：\Quote{啊，邪恶至极的人！你从我这里享受了如此大的恩宠，被授予如此大的权柄，因你自身的功德而超乎地上一切受造物；你实际领受了我的眷顾的保证，以及我天主眷顾\OriginalTerm{天主眷顾}{Providence}的真实显现，却认为一个邪恶、害人的魔鬼，你救恩的敌人，比你的主和恩主更可信？他为你显示了什么样的关怀，像我这样？我岂没有为你造了天、地、海、太阳、月亮和所有星辰？的确，没有一位天使需要这创造之工；但为你，并为你的欢愉，我造了如此伟大而美好的世界；而你却认为仅仅言语、虚假的承诺和充满欺骗的应许，比以行动显明的慈爱与眷顾更值得相信，以致你将自己交给了他，践踏了我的律法！}这些话以及更多此类话，一个被轻慢的人很可能会说。但天主并未如此行事；恰恰相反。因为祂的第一句话立刻将他从沮丧中扶起，使那恐惧战栗的人获得信心，因为祂自己首先呼唤他，或者更确切地说，不仅首先呼唤他，而且用他熟悉的称呼叫他，说道：\Quote{亚当，你在哪里？}如此祂显明了祂的温柔，以及对他极大的关怀。因为你们必须知道，这正是亲密友谊的标志。同样，那些呼唤死者的人也惯于如此，不断重复他们的名字。另一方面，那些对某人怀有仇恨和敌意的人，无法忍受提及那些冒犯他们的人的名字。例如，撒乌耳虽然从未受到达味的伤害，反而大大亏负了他，但因憎恶和仇恨他，便不能忍受提及他的本名；当众人聚坐一处，不见达味在场时，他说了什么？他没有说：\Quote{达味在哪里？}而是说：\Quote{叶瑟的儿子在哪里？}\BibleRef{列上 20:27}，用他父亲的名字称呼他。同样，犹太人对待基督也是如此：因为他们憎恶和仇恨祂，他们没有说：\Quote{基督在哪里？}而是说：\Quote{那个人在哪里？}\BibleRef{若 7:11}\par

7. 但天主愿意藉此显示，罪恶并未熄灭祂的慈爱，悖逆也未夺去祂对亚当的眷顾，祂仍施以眷顾并照顾那堕落者，便说：\Quote{亚当，你在哪里？}并非不知他在何处，而是因为犯罪之人的口被封闭；罪恶扭转了舌头，良心抓住了它，以致这样的人缄默不语，如同被某种锁链束缚在沉默中。天主因此愿意邀请他畅所欲言，给予他信心，引导他为自己的过犯辩护，好使他能得到些许宽恕，便亲自先呼唤他；用亲昵的称呼消除了亚当许多的苦恼，驱散了他的恐惧，藉此呼唤开启了那紧闭的口。因此祂也说：\Quote{亚当，你在哪里？}\Quote{我把你留在一种境况中，}祂说，\Quote{却发现你在另一种境况中。我留下你时满有信心和光荣；如今却发现你蒙羞受辱、沉默不语！}请注意天主在此事中的眷顾。祂没有呼唤厄娃；祂没有呼唤蛇；而是将犯罪最轻的那一个首先带到审判台前，为的是从那个尚能找到些许借口的人开始，祂能对那犯罪最重的人也作出更仁慈的判决。事实上，法官们不屑于亲自审问他们的同仆和共享同一本性的人，而是派一个侍从介入，让他将自己的问题传达给罪犯；并通过他来传达和听取他们想知道的一切。但天主与人打交道时无需中介；祂亲自同时审判和安慰他。这不仅奇妙，而且祂还纠正已犯下的罪行。因为一般的法官，当发现小偷和盗墓者时，他们不考虑如何使他们变好，而是如何让他们为所犯的罪行付出代价。但天主恰恰相反，当祂遇到一个罪人时，祂考虑的不是如何让他付出代价，而是如何纠正他，使他变得更好，并使他将来不可战胜。因此，天主同时是法官、医生和教师：作为法官，祂审问；作为医生，祂纠正；作为教师，祂教导那些犯罪的人，引导他们获得一切属灵的智慧。\par

8. 但是，如果一段简短的话语就如此证明了天主的眷顾，那么如果我们通读这整个审判，并展开其全部记录，又会如何呢？你看见全部圣经都是安慰和鼓励吗？但关于这些记录，我们将在适当的时候谈论；然而在此之前，有必要说明这本书是在什么时候赐下的；因为这些事情并不是在起初、也不是在亚当被造时立刻写下的，而是在许多世代之后；值得探究的是，为什么会有这样的延迟，以及为什么最终只赐给犹太人，而不是所有人；为什么用希伯来文写成；为什么在西乃旷野？因为宗徒并非仅仅随意提及这个地方；而是表明在那情形中，对我们也有一个伟大的默观对象，当时他对我们说：\Quote{因为有两个盟约，一个出自西乃山，生于为奴。}\BibleRef{迦 4:24}\par

9. 此外，还有其他事情值得我们探究。但我看到时间不允许我们将言论投放到如此广阔的海域；因此，明智地将这些保留到适当的时候，我们再次就戒除誓愿向你们讲话；我们恳求你们的爱德在这件事上要十分勤勉。因为，有什么比这更荒谬呢：连仆人都不敢直呼主人的名字，也不敢随意不拘地提及他，而你却随意而轻慢地四处滥用天使之主的名字！如果必须拿起福音书，你会用先洗净的手来接；并恐惧战兢，怀着极大的尊敬和虔诚；而你却随意地在舌头上滥用福音之主！你愿意学习天上的权能如何称呼那名字吗？带着怎样的敬畏、怎样的恐惧、怎样的惊奇？先知说：\Quote{我看见主坐在高高的宝座上，祂的衣裳垂下，遮满圣殿。在祂周围站着的色辣芬，彼此呼喊说：\InnerQuote{圣！圣！圣！万军的上主！祂的光荣充满大地！}}\BibleRef{依 6:3}你明白吗？他们称呼那名字时带着怎样的恐惧和敬畏，同时光荣赞美祂？但你在祈祷和恳求中，却十分漫不经心地呼求祂；而当你本当充满敬畏、警醒和清醒时！在誓愿中，完全不适合引入这奇妙名字的地方，你却连珠炮似地发出各种诅咒形式！那么，我们有什么宽恕或借口呢？无论我们怎样以这\Quote{习惯}为借口？据说，有一个外邦演说家，由于某种愚蠢的习惯，走路时不停地抖动右肩。他通过把锋利的刀片绑在肩膀两侧来克服这个习惯，这样害怕被割伤就控制住了那不适当的动作，因怕受伤而制止了它！那么，你也要这样对待你的舌头：以天主的惩戒代替刀片悬在它之上，你一定会得胜！因为看来不可能，完全不可能，那些对此关心、认真并真正致力于此事的人会被胜过。\par

10. 你们为现在所说的鼓掌，但当你们改正后，你们将更大地鼓掌——不仅为我们，也为你们自己；你们会更乐意听所讲的话；并以纯洁的良心呼求天主，祂对你们何等宽容，人啊！祂说：\Quote{你不可指着你的头起誓。} \BibleRef{玛 5:36} 但你竟如此藐视祂，甚至指着祂的光荣起誓。\Quote{但我该怎么办，}有人说，\Quote{那些强加于我的人呢？} 有什么必要呢，人啊？让所有人都明白：你宁可忍受任何事也不愿违犯天主的法律；他们就会停止强迫你。因为事实证明，使人可信的不是誓言，而是生活的见证、言行的正直和良好的声誉；许多人曾起誓到喉咙嘶哑，却仍不能说服任何人；而另一些人仅凭一句肯定的话，就比那些发誓无数的人更被相信。因此，知道这一切，并把为起誓者和发假誓者所预备的惩罚放在我们眼前，让我们戒除这邪恶的习俗，从而由此前进，改正其余的一切，使我们能享受来世的真福——愿天主恩赐我们众人堪当获得，因我们的主耶稣基督的恩宠和爱人，藉着祂、偕同祂、偕同圣神，归于圣父，光荣、权能、尊威，现在及永远，世世无尽。阿们。\par

\SourceRecord{Translated by W.R.W. Stephens.；Nicene and Post-Nicene Fathers, First Series；Vol. 9.；Edited by Philip Schaff.；Buffalo, NY: Christian Literature Publishing Co.,；1889.；<http://www.newadvent.org/fathers/190107.htm>.}

% structure: p0001 source=h1 target=section reason=page-title

\section{讲道集第八篇：论雕像}

\Emph{劝勉修德——特别是关于经文\BibleQuote{天主在乐园中，晚凉之际散步}（\BibleRef{创 3:8}）——以及再次论戒发誓的主题。}\par

你们最近听到，全部圣经虽然是一部历史叙述，却带来安慰和慰藉。例如，\BibleQuote{在起初，天主创造了天地。}（\BibleRef{创 1:1}）是一个历史性的宣告；但在我们的讲道中，已证明这句话充满了安慰；例如，天主为我们设了双重盛宴，同时铺开了海洋和陆地；并在高处点燃了双重光明，即太阳和月亮；又规定了它们运行的双重季节，即白昼和黑夜，一个用于劳作，另一个用于休息。因为黑夜带给我们的益处不亚于白昼。但正如我论树木时所说的，那些不结果的树木，在用途上堪与结果的树木媲美；因为我们不必为建造而触动那些果实可口的树木。野生的和未驯服的动物同样服务于我们的需要，其程度不亚于驯养的动物；它们因使我们畏惧而驱使我们聚集到城市，使我们更加谨慎，并使我们彼此联结；它们锻炼了一些人的力量，又使另一些人摆脱了疾病；因为医生从它们身上配制了许多药物；它们还使我们记起古代的罪。因为当我听到：\BibleQuote{地上一切的野兽，对你们都要感到恐惧和惊怕。}（\BibleRef{创 9:2}）随后又看到这尊荣后来被削减，我就记起罪，它驱散了我们对野兽的恐惧，削弱了我们的权威。这样，我变得更好、更明智，因为我学会了罪给我们带来的伤害。正如我所说，那些被提及的事物以及其他类似的事物，造物主天主知道这些，对我们现世的生活贡献不小；所以现在我也说，黑夜带来的益处不亚于白昼，它是劳作的休息和疾病的良药。医生们虽然多方努力，配制了无数药物，却往往不能治愈患病的人；但睡眠自动临到他身上，却彻底消除了疾病，为他们省去了无尽的麻烦。黑夜不仅是身体劳作的良药，也是心灵疾病的良药，它给痛苦的灵魂带来安息。常常有人失去了儿子；无数安慰者都无法使他脱离泪水和呻吟。但当黑夜来临，他被睡眠的绝对力量征服，闭上了眼睑，得到一些白天苦难的缓解。\par

2. 现在，我请求你们，让我们继续探讨引出这些观察的主题。因为我深知，你们都在热切期待这事；你们每个人都痛苦不堪，直到他明白为何这本书从一开始就没有给出。但即使现在，我看谈论此事的时机仍未成熟。为什么呢？因为一周几乎已经在我们中间结束，我害怕触及一个主题，它的阐述随后不久就不得不中断。因为这主题需要我们连续几天和持续的记忆努力；因此我们必须再次推迟。但请不要见怪！我们一定加倍偿还债务；因为这对你们和负责偿还的我们都有好处。与此同时，让我们现在谈论昨天遗漏的那个主题。我们昨天遗漏了什么？经文说：\BibleQuote{天主在乐园中，晚凉之际散步。}（\BibleRef{创 3:8}）请问这里是什么意思？\Quote{天主散步！}天主并没有散步；祂无所不在、充满万有，又怎能散步呢？但祂在亚当内引起了这种感知，为使他收敛心神，为使他不再粗心大意；为使他在逃跑和躲藏时，甚至在说话之前，就预先呈现一部分理由。因为就像那些即将被带到法庭，承受所犯之罪的指控的人，他们面带污秽、蒙尘、忧伤和谦卑的容貌出现在审判者面前，以便从外表上使他们倾向慈爱、怜悯和宽恕；同样的事也发生在亚当身上。因为他必须谦卑地被带到这法庭。因此天主预先取了他，并使他谦卑。但他察觉到有人在散步；但他从哪里想到是天主在那里散步呢？这是犯罪之人的习惯：他们怀疑一切，对影子战栗，被每个声音惊吓，以为每个人都以敌对的方式接近他们。因此，有罪的人常常看到别人忙别的事，便以为他们是来对付自己的；当别人就完全不同的话题彼此交谈时，那些自知有罪的人就以为他们在谈论自己。\par

3. 罪的本性正是如此：它出卖人，却无人指责；它定人罪，却无人控告；它使罪人成为胆怯者，闻声颤抖；义德则产生相反的效果。至少请听圣经如何描述前者的怯懦和后者的勇敢。\Quote{恶人虽无人追赶，也逃跑。}\BibleRef{箴 28:1} 他如何无人追赶却逃跑？他内在有一个驱赶他的东西——良心中的控告者；他到处携带它，正如他不能逃避自己，同样他也不能逃避内在的追逼者；无论他到哪里，都受鞭打，且带着不治的创伤。但义人却不是这样。他是什么样的呢？请听：\Quote{义人却勇敢如狮子！} 厄里亚就是这样的一个人。例如，他看见国王向他走来，当国王说\Quote{你为何使以色列陷于混乱？}时，他回答说：\Quote{我没有使以色列陷于混乱，而是你和你的父家。}\BibleRef{列上 18:17} 的确，义人勇敢如狮子；因为他对抗国王，如同狮子对抗卑鄙的野狗。虽然一个身披紫袍，另一个只披羊皮，但羊皮是两者中更可敬的；因为紫袍带来了严重的饥荒，而这羊皮却使人从那灾祸中解脱！它分开了约但河！它使厄里叟成了双倍的厄里亚！哦，圣人的德能何其伟大！不仅他们的言语，不仅他们的身体，就连他们的衣服也总是受整个受造物的尊敬。这人的羊皮分开了约但河！三圣童的凉鞋践踏了火焰！厄里叟的话改变了水，使铁浮在水面上！梅瑟的杖分开了红海，劈开了磐石！保禄的衣服驱除了疾病！伯多禄的影子赶走了死亡！圣殉道者的骨灰驱逐了魔鬼！为此，他们行事都带着权柄，正如厄里亚一样。因为他不是看王冠，也不是看国王的外在威仪，而是看那灵魂衣衫褴褛、肮脏不堪，比任何罪犯更悲惨；看见他被情欲俘虏和奴役，便蔑视他的权力。因为他仿佛看见一个舞台上的国王，而非真实的国王。外在的丰裕有何益处，内在的贫穷如此巨大？外在的贫穷又有何害处，内在的财富宝藏如此丰厚？有福的保禄也是这样一头狮子；因为他进入监狱，只一发声，便摇动了所有根基；他用话语，而非牙齿，咬碎了镣铐；因此，称这样的人为狮子还不恰当，应称他们为超过狮子；因为狮子常常落入网中被捕，而圣人们被捆绑时，反而更有力量；正如这有福之人在监狱中所做的，他释放了囚犯，摇动了墙壁，捆绑了狱警，并以虔诚的话语征服了他。狮子吼叫，使所有野兽逃散。圣人发声，驱散四方的魔鬼！狮子的武器是鬃毛、尖爪和利牙。义人的武器是属神的智慧、节制、忍耐、轻看现世一切。谁拥有这些武器，不仅能讥讽恶人，甚至能对抗那些敌对权势本身。\par

4. 人哪，要研习按照天主的生活，那么任何时候无人能征服你；即使你被视为最微不足道的人，你也要比所有人更有力量。另一方面，如果你对灵魂的德行漠不关心，即使你是最有力量的人，你也将轻易被所有攻击者击败。以上所举的例子证明了这一点。但如果你愿意，我也将试图用实际事例教导你义人的不可征服和罪人的脆弱状态。那么请听先知如何说明这两点。\Quote{恶人}，他说，\Quote{却不是这样，而是像糠秕，被风吹散。} 因为正如糠秕暴露在风中，易被卷起吹走，同样罪人被各种诱惑驱赶；当他与自己交战，内心带着战争时，他还有什么得救的希望？他在家被出卖，随身带着那常作敌人的良心。义人却不是这样。他是怎样的呢？请听同一先知说：\Quote{信赖上主的人如同熙雍山。} 那么，\Quote{如同熙雍山}是什么意思？\Quote{他永不动摇。}他说，\Quote{永远。} 因为无论你运来什么机械，投掷什么标枪，想要倾倒一座山，你永远无法成功；怎么可能呢？你将击碎所有机械，耗尽自己的力量。义人也是如此。无论他遭受什么打击，他都不受伤害，反而摧毁那些图谋反对他之人的力量，不仅是人，还有魔鬼。你常听见魔鬼对约伯使用了什么机械；但他不但没有推翻那座山，反而筋疲力尽地后退，他的标枪粉碎，他的机械在这次攻击中变得无用！\par

5. 知道这些事以后，让我们留心自己的生活；不要执著于会朽坏的财物，不要执著于会过去的荣耀，不要执著于会衰老的身体，不要执著于会凋谢的美貌，不要执著于易逝的享乐；但要把我们所有的关心都用在灵魂上，并以各种方式为它的福祉作准备。因为治愈有病的身体，对每个人来说都不是一件容易的事；但治愈有病的灵魂，对所有人来说都是容易的；身体的疾病需要药物和金钱才能治愈；但灵魂的医治却是容易获得且没有花费的。肉身的本性要摆脱那些麻烦的伤口，需要付出很大的辛劳；因为常常要动刀，要用苦药；但对于灵魂，却没有这类事。只需运用意志和愿望，一切便可成就。这是天主眷顾的工作。因为从身体的疾病中不会产生多大的伤害（即使我们没有生病，死亡无论如何也会来临，摧毁并分解身体）；但一切都取决于灵魂的健康；既然灵魂更加宝贵和必要，祂便使其医治变得容易，且没有花费或痛苦。因此，如果我们当身体生病时，要为其花费金钱、请医生、忍受许多痛苦，我们如此关心它（尽管从那种疾病中可能不会给我们带来多大的伤害），却忽视灵魂，我们将得到什么借口或宽恕呢？而这种情况，我们既不需要付钱，也不需要给别人添麻烦，也不需要承受任何痛苦；相反，不需要所有这些，仅凭选择和意愿，我们就能完成其全部的改正；并且确知，如果我们不这样做，我们将承受最终的判决、惩罚和刑罚，这些是无法挽回的！请告诉我，如果有人许诺在短时间内、不花金钱和劳力就教会你医术，你难道不认为他是恩人吗？你难道不会去做、去忍受他所命令的一切吗？请看，如今你可以不劳而获地找到一种药，不是医治身体的伤口，而是灵魂的伤口，并且毫不痛苦地恢复其健康！让我们不要漠不关心！请问，对得罪自己的人放下愤怒有什么痛苦呢？记住伤害而不和解，那才是痛苦！祈祷、向乐意赐予万善的天主祈求，有什么劳力呢？不诽谤任何人，有什么劳力呢？摆脱嫉妒和恶意有什么困难呢？爱自己的近人有什么麻烦呢？不说污言秽语、不咒骂、不侮辱别人，有什么痛苦呢？不发誓有什么疲劳呢？因为我又回到这同一劝诫上。发誓的劳苦确实极大。常常，在愤怒或暴怒之下，我们发誓说，绝不与伤害我们的人和好。但后来，当怒气平息、愤怒消减时，我们希望和好，却被这些誓言所约束，我们便遭受同样的痛苦，如同陷入网罗、被无法挣脱的锁链捆住一样。魔鬼知道这一点，并且清楚愤怒是火，它容易熄灭，一旦熄灭，和好与爱便会随之而来；为了不让这火熄灭，魔鬼常以誓言束缚我们；这样，即使愤怒停止，誓言的责任仍然存留，在我们内保持那火；结果两件事必有一件发生：要么我们和好而背誓，要么我们不和好而承受怀恨的惩罚。\par

6. 因此，知道这些事以后，让我们避免发誓；让我们的口常练习说\Quote{相信我}；这将成为我们一切虔诚行为的基础；因为舌头一旦被训练使用这一表达，就会感到羞愧，耻于说出可耻和丑恶的话；即使有时因习惯而被拖走，也会因有许多指责者而再次受到约束。因为当有人看到不发誓的人说出污言秽语时，他就会利用这机会嘲笑他，并讥讽地说：\Quote{你在所有事情上都说\InnerQuote{相信我}，从不敢发誓，却用这些可耻的话玷污你的舌头吗？}这样，即使我们不愿意，也会因身边人的强迫而重回虔诚的行为。\Quote{但是什么，}有人说，\Quote{如果必须发誓呢？}在违法的地方，就没有所谓必须。\Quote{那么可能吗，}有人回答，\Quote{完全不发誓？}你说什么？天主命令了，你还敢问祂的律法能否遵守吗？其实，祂的律法不可能不被遵守；我愿意从目前的情况说服你这一点：非但不是不可能不发誓，反而是不可能发誓。看，城中的居民奉命缴纳一笔黄金，这在许多人看来可能超出能力；但大部分款项已经收齐；你可以听到收税的人说：\Quote{你为什么拖延，人？为什么一天天推诿？不可能避免。这是皇帝的法律，不容拖延。}你说什么？皇帝命令你交出金钱，不可能不交！天主命令你避免发誓，你怎么能说不可能避免呢？\par

7. 我现在是第六天就这条诫命劝诫你们。从今以后，我渴望向你们告别，打算不再谈论这个话题，使你们警惕。你们将不再有任何借口或宽容；因为按理说，即使没有说过这事，你们自己也应当改正，因为这不是一件复杂的事，或需要大量准备的事。但既然你们已经受益于如此多的劝诫和劝言，当你们站在那可畏的审判台前受控告，被要求为这次过犯交账时，你们还有什么借口可以提出？不可能想出任何借口；但你们必须要么从今以后改正，要么，如果没有改正，就受惩罚，并承受最严厉的刑罚！因此，思想这一切，并怀着对此极大的忧虑离开这里，你们要彼此劝勉，使这么多天所讲论的事得以以一切警醒存记在你们心中，这样，当我们沉默时，你们彼此教导、建立、劝勉，可以表现出极大的长进；并且遵守了所有其他的诫命，得以享受永生的荣冠。愿天主使我们众人藉着我们的主耶稣基督的恩宠和慈爱，获得这一切。愿光荣藉着祂并偕同祂，归于圣父，偕同圣神，直到永远。阿们。\par

\SourceRecord{Translated by W.R.W. Stephens.；Nicene and Post-Nicene Fathers, First Series；Vol. 9.；Edited by Philip Schaff.；Buffalo, NY: Christian Literature Publishing Co.,；1889.；<http://www.newadvent.org/fathers/190108.htm>.}

% structure: p0001 source=h1 target=section reason=page-title

\section{论雕像讲道集 第九篇}

称赞那些已戒除起誓恶习的人。论证饭后在教会聆听天主圣言无需顾虑。回应为何圣经迟迟未赐下的问题。讲解\BibleQuote{诸天宣示天主的荣耀}以及受造界的描述。最后是反对起誓的劝诫。\par

1. 不久之前我对你们说的话，现在我又对你们重述！但愿我能常与你们同在——是的，我确实常与你们同在，尽管不是身体临在，而是藉着爱的力量！因为除了你们和你们的救恩，我没有别的生命。正如农夫只关心他的种子和收成，舵手只关心波涛和港口，同样，宣讲者也关心他的听众及其进步，如同我现在一样！因此，我不仅在这里，在家里也把你们所有人放在心上。因为即使人数众多，我心胸狭窄，但爱是宽广的；并且\BibleQuote{你们并没有在我们身上受限制}。我不会接着说出下面的话，因为我们也没有在你们身上受限制。何以见得呢？因为我遇到许多人说：\Quote{我们遵守了诫命，彼此订立规则，为起誓者定下处罚，并对违反此法者施加惩罚}。这一惩罚确实与你们相称，是最大爱德的标记。因为我不以在这些事上忙碌为耻，这种干预并非出于无聊的好奇，而是出于温柔的关怀。因为如果医生询问病人的情况并不算羞耻，那么我们不断询问你们的救恩也不算过错；因为由此了解到什么已完成、什么未完成，我们就能凭必要的知识施用进一步的补救。我们通过询问确认了这些事；我们感谢天主，因为我们没有把种子撒在石头上，也没有撒在荆棘中；我们既不需要很多时间，也不需要长期延迟就能收获庄稼。因此，我恒常把你们放在心上。因此，我不感到教导的辛劳，因为听众的获益减轻了负担。这报酬确实足以恢复我们的力量，给我们翅膀，提升我们，并说服我们为你们承受最大的劳苦。\par

2. 既然你们表现出如此慷慨的情怀，请允许我们偿还前几天承诺的另一笔债务；尽管我看到当时在场的人现在并不都在这里。请问，这是什么原因？是什么使他们离开了我们的餐桌？看来，享用过身体食粮的人认为，在接受物质食物之后，再来聆听天主的圣言是件有失体面的事。但他们这样想是不对的。因为如果这是不合适的，基督就不会在那次奥秘晚餐后发表他那长篇大论的讲道；如果这是不适宜的，他也不会在旷野中喂饱群众之后，在饭后对他们讲话。因为（如果必须在此说句令人惊讶的话），那时聆听天主圣言尤其有益。当你决定饭后还要前往聚会时，你一定会谨慎行事，尽管或许不情愿，也要保持节制；你绝不会任何时候醉酒或暴食。因为进入教堂的思想和期待，训练你以合宜的节制享用饮食；以免你进去与弟兄们相聚时，因酒气和粗俗的呃逆而在众人面前显得可笑。我现在说的不是你们这些在场的人，而是那些缺席的人；好让他们通过你们得知这些事。因为妨碍聆听的，不是吃了东西，而是懈怠。然而，你们认为不守斋是罪过，却又加上另一项更大更重的过错，即不分享这神圣的食粮；你们滋养了身体，却以饥饿消耗灵魂。对此，你们有什么可辩解的？因为关于守斋，你们或许可以借口身体软弱，但对于聆听，你们有什么可说的呢？身体软弱决不阻碍你们参与天主圣言！如果我说：\Quote{吃过早餐的人不要加入我们}，\Quote{吃过饭的人不要听讲}，你们或许有些借口；但现在，我们愿意拉拢、劝诱、恳求你们来，你们有什么理由拒绝我们？不合适的听者不是那些吃了喝了的人，而是那些不留意所讲的内容、打哈欠、懒散、身体在此而心思他往的人；这样的人，即使守斋，也是无益的听者。反之，认真、警醒、保持思想专注的人，即使吃了喝了，也是我们最合适的听者。确实，这条规矩恰当地适用于世俗的法庭和议会。因为他们不知道如何属神地明智，所以他们的饮食不是为滋养，而是为胀腹；他们常常过量饮酒。因此，他们使自己不适于处理事务，便在傍晚和中午关闭法院和议会厅。但这里绝无此事——天主不许！而是那吃过的人，在灵魂的节制方面，将与守斋者争胜；因为他饮食，不是为了鼓胀肚腹或蒙蔽理智，而是为了恢复已衰弱的身体的力量。\par

3. 不过，这劝诫已足够。现在是讨论我们主题的时候了；尽管我们的心因那些未到的人而退缩，不愿给予这训导。正如一位慈母在准备摆桌时，若她的孩子不全在场，便悲伤哀叹；我现在也正是如此受苦；当我想到弟兄们的缺席，我便不愿清偿我的债务。但你们有能力使我摆脱这迟延。如果你们答应我，将把我所说的一切准确报告给他们，我就欣然将全部付给你们；因为这样，你们以爱德提供的训导，将弥补他们的缺席；而你们会更留心听我，知道你们必须向他人交代这些事。为使我们的主题更加明了，让我们从开始重新提起。那天我们探究：为什么圣经在这么多年后才被交付？因为这部书既不是在\Person{亚当}时代、\Person{诺厄}时代、\Person{亚巴郎}时代，而是在\Person{梅瑟}时代交付的。我听到许多人说，如果这部书有益，它应该从一开始就交付；如果它无用，那么后来也不应该交付。但这是一个过时的论点；因为并非任何有益的东西都必须从一开始就交付，也并非任何从一开始交付的东西都必须后来继续。例如：奶有用，但并非总是给予；它只在我们是孩童时给予；固体食物有用，但没有人会在生命之初就给我们，而是当我们过了童年之后。同样，夏季有用，但它不常出现；冬季有益，但它也要让位给其他季节。那么怎样？他们说圣经没有用吗？我回答：它们最有用，最必要。那么如果如此有用，为什么他们说不从一开始就交付给我们？那是因为天主渴望以事物，而非文字，来教导人的本性。但\Quote{藉着事物}这个说法是什么意思呢？藉着受造物本身。\par

4. 请看，宗徒如何论及同一主题，并针对那些声称他们从一开始并非从圣经认识天主的希腊人，给出了他的回答。他曾说：\BibleQuote{天主的忿怒，从天上显明在一切不虔不义的人身上，就是那些以不义压制真理的人。}\BibleRef{罗 1:18} 当他看到自己遭遇反对，并且许多人仍会追问外邦人从哪里得知天主的真理时，他接着加上说：\Quote{因为天主的事情，在他们里面是明显的。} 但如何在他们里面明显？他们如何能认识天主？又是谁指教的？请说明这一点。他说：\Quote{天主已经给他们显明了。} 以什么方式？借着派遣了哪一位先知？哪一位福音作者？哪一位教师？既然圣经尚未赐下。他说：\BibleQuote{祂那看不见的，即祂永远的大能和神性，自从创造世界以来，从受造物上可以看得清清楚楚，被明白了解。}\BibleRef{罗 1:20} 他所说的就是：天主把自己的受造物放在众人眼前，使他们从祂的工程推测造物主；另一位作者也引用了这一点：\BibleQuote{从受造物的伟大和美，可以相应地看到造物主。}\BibleRef{智 13:5} 你看见伟大吗？惊叹那创造者的能力！你看见美吗？惊讶那装饰它的智慧！这正是先知所说的：\Quote{诸天宣示天主的荣耀。} 那么，请问，它们如何宣示？它们没有声音，没有口，没有舌头！它们如何宣示？借由景象本身。因为当你看见其美、其广阔、其高度、其位置、其形状、以及如此长久的稳定时，你仿佛听到一个声音，受到景象教导，你就朝拜那创造如此美丽而奇妙形体的那一位！诸天可能沉默，但看见它们就发出比号角声更响亮的声音；不是通过耳朵教导我们，而是通过眼睛；因为后者是比前者更可靠、更清晰的感官。\par

5. 因为如果天主通过书籍和文字来教导，那么认识字的人就会知道所写的内容；而不认识字的人则不会从这来源得到任何益处，除非有人引导他。富人能购买圣经，但穷人却无法得到它。再者，懂得书写所用语言的人可能知道其中的内容；但\OriginalTerm{斯基泰人}{Scythian}、蛮族、\OriginalTerm{印度人}{Indian}和\OriginalTerm{埃及人}{Egyptian}，以及所有不懂那种语言的人，都会得不到任何教导而去。然而，对于天穹却不能这样说；因为\OriginalTerm{斯基泰人}{Scythian}、蛮族、印度人、埃及人和每一个行走在大地上的人，都会听到这声音；因为不是通过耳朵，而是通过视觉，它到达我们的理智。对于所见的事物，有一种统一的感知；不像语言那样有差别。在这卷书上，无学问的人和有智慧的人都一样能看；穷人和富人也一样；无论一个人来到哪里，抬头仰望天穹，就会从它们的景象中得到足够的教训。先知本人也暗示并指出了这一事实：受造物发出这声音，使蛮族、希腊人和全人类都能理解，当他这样说时：\Quote{没有言辞，没有话语，也听不到它们的声音。}他的意思是这样：没有哪个民族或语言不能理解这语言；而是它们的宣告，全人类都能听到。这不仅关乎天穹，也关乎白昼和黑夜。但天穹藉其美丽和宏大以及其他一切让观者惊叹，并使他衷心赞美\OriginalTerm{造物主}{Creator}；至于白昼和黑夜，它们能向我们展示同样的事吗？当然不是同样的，而是其他不逊于它们的事；例如：它们如此精确遵守的和谐与秩序。因为当你思量它们如何分配全年，并如同用横梁和天平相互分割整个空间的长度，你会对那安排它们的一位感到惊讶！正如某些姊妹在彼此深情中分配父亲的遗产，丝毫不互相冒犯，同样，白昼和黑夜以如此均等的部分、以极大的精确度分配全年；并坚守自己的界限，从不互相推挤。白昼从未在冬季变长；同样，黑夜也从未在夏季变长，尽管如此多的世代已经过去；但在如此长的时间间隔中，一个未曾欺骗另一个，甚至半分半秒都没有！\par

6. 因此，\OriginalTerm{圣咏作者}{Psalmist}也对这分配的均等感到惊讶，喊道：\Quote{夜复一夜传递知识。}如果你知道如何明智地默想这些事，你就会钦佩那从起初就设定了这些不动界限的一位。愿贪婪的人和觊觎他人财富的人听到这些；愿他们效法白昼与黑夜的均等。愿傲慢自大的人，以及不愿把首位让给他人的人也听到！白昼给黑夜让路，并不侵犯他人的领地！但你，常享尊荣，却不能忍受与弟兄分享吗？请也与我一同思考\OriginalTerm{立法者}{Lawgiver}的智慧。在冬天，他安排黑夜变长；此时嫩芽娇嫩，需要更多凉爽，无法承受太阳更热的射线；但当它们稍微长大，白昼又随之增长，并在果实成熟时变得最长。这不仅对种子有益，也对我们的身体有益。因为在冬天，水手、舵手、旅行者、士兵和农夫大多因霜冻而困守家中，季节是闲暇的；天主就安排这段时间的大部分在黑夜中消耗，以免白昼在人们无所作为时显得多余。谁能描述季节的完美秩序？它们如何像一些围成圆圈的少女，以最愉快的和谐相互接替？那些中间的季节如何不停顿地、以逐渐无声的过渡走向对立的一侧？因此，我们既不会在冬天之后立刻遇到夏天，也不会在夏天之后立刻遇到冬天；而是中间插入了春天；这样，当我们温和而逐渐地接受一个接一个的季节时，我们的身体就能被锻炼得安然承受夏天的炎热。因为突然转向对立极端会造成最严重的伤害和疾病，天主便安排我们在冬天之后迎接春天，春天之后是夏天，夏天之后是秋天，然后这样把我们送回到冬天；因此，这些从对立季节而来的变化，通过中间季节的帮助，无害而逐渐地临到我们。那么，谁这样可怜可悲，看见天穹、看见海洋和陆地、看见季节的精确调整和白昼黑夜的不变秩序，却认为这些事是自动发生的，而不去崇拜以相应智慧安排这一切的祂呢？\par

7. 但在这件事上，我还有话要说。因为不仅受造界的宏伟与美丽，而且其本身的方式，都彰显了一位作为宇宙工匠的天主。既然我们在起初，当祂从事形成和创造万物工程时，我们并不在场；即使我们在场，我们也无法知道万物如何生成，因为安排万物的能力是看不见的；祂使这个受造界的方式成为我们最好的老师，以一种超越自然过程的方式组合万物。或许我所说的，不够清楚。因此，我必须以更清晰的方式再重复一遍。那么，所有人都必须承认，自然过程是水被地承载，而非地被水承载。因为地是一种密集、坚硬、不屈、坚固的实体，容易承载水的本性；而水是流动、稀薄、柔软、扩散，对一切所遇之物让路，不能承载任何固体，即使是最轻的。当小石子落在水上时，水往往会让步、让路，将其沉入水底。因此，当你看见不是小石子，而是整个大地承载在水上，却没有沉没时，当钦崇那位以超自然方式施行这些奇事者的能力！大地承载在水上，这从哪里显明呢？先知宣布这事，他说：\BibleQuote{祂将地建立在海上，预备在洪流之上。} 又说：\BibleQuote{祂将地建立在诸水之上。} 你说什么？水不能使小石子浮在水面，却承载着如此巨大的大地、山脉、丘陵、城市、植物、人和牲畜，它却没有沉没！我说什么？不是没有沉没吗？这是怎么发生的？既然水从下方与它紧密接触如此长久，它却没有溶解，整个变成泥浆？木材的质料，在水中浸泡片刻，便会腐烂溶解；我何谈及木头？有什么比铁更坚固？然而，它若长时间留在水中，也常变软；这也情有可原。因为它来源于大地。因此，许多逃亡的仆人，在逃跑时拖着脚镣和锁链，前往溪流，将戴镣的脚伸入水中，以此法使铁变软后，便用石头轻易击碎它。的确，铁变软，木腐烂，石头被水的本性磨损；然而如此巨大的大地，长久以来躺卧在水上，却既未沉没，也未溶解或毁灭！\par

8. 有谁不会对这些事感到惊奇与赞叹，并确信地说，这些不是本性之作，而是超越本性的天主眷顾之作？因此，有人这样说：\BibleQuote{谁将大地悬在虚无之上。}\BibleRef{约 26:7} 另一个人观察说：\BibleQuote{祂手中持有地极。} 又说：\BibleQuote{祂将地基立在海上。} 这些宣言虽看似相互矛盾，却完全一致。因为那说\BibleQuote{祂将地基立在海上}的人，与那宣称\BibleQuote{祂将它悬在虚无之上}的人，意思相同。那么，它悬在何处，置于何处？听同一个人说：\BibleQuote{祂手中持有地极。} 并非天主有手，而是为使你知道，是祂的能力照顾一切，维系并支撑着大地之体！但若你不信我此刻所说的，就相信你所见的，因为在另一种元素中也能找到这令人赞叹的工艺。因为火的本性趋向上升，总是向上攀登；即使你极力逼迫约束它，它也不能顺从向下。因为我们时常举着点燃的火炬，即使将头向下倾，也不能迫使火焰转向地面；它仍然向上转，从下往上升。但对于太阳，天主却使之完全相反。祂使太阳的光线朝向大地，使其光芒向下照射，仿佛藉着（穹苍的）形状对它说：\Quote{向下看！照射世人，因为你为他们而造！} 蜡烛的光固然不能服从于此；但这颗伟大神奇的星辰却向下弯曲，朝向大地，这与火的本性相反；这是由于命令它的天主的能力。你愿意我再说一件类似的事吗？水从各方拥抱可见之穹苍的背面，却既不流下，也不离开原位，尽管水的本性并非如此。因为水容易聚集在凹处；但当物体呈凸形时，它会从各方滑落，甚至连一小部分也不能停留在这样的形状上。但是，看哪！这奇事竟存在于穹苍中；此外，先知为暗示这一情形，说：\BibleQuote{天上的诸水，赞美主！} 此外，水没有熄灭太阳；太阳在下方运行如此之久，也没有晒干上面的水。\par

9. 你愿意我们再次带你下到地上，指出这奇事吗？难道你没有看见这大海波涛汹涌、狂风大作，却广阔浩瀚、狂暴异常，仅以微弱的沙粒为堤岸吗？也请留意天主的智慧：祂不允许大海静止或平静，免得你以为它的秩序仅是自然规律所致；但大海留在其界限内，发出声音，喧闹咆哮，掀起巨浪高耸入云。但一到岸边，看见沙粒，它便破裂，退回到自身之内；以此两件事教导你：它留在边界内并非本性的工，而是那位以权能约束它的主的工！因此，祂使堤岸脆弱；并未以林木、石头或群山环绕这些海岸，免得你将元素的秩序归于此类事物。因此，天主自己曾以此事斥责犹太人说：\BibleQuote{你们不惧怕我吗？我以沙粒为海的界限，使它不能越过。} \BibleRef{耶 5:22} 但奇事不仅在此，即祂创造了伟大而可赞叹的世界，并以超乎寻常自然秩序的方式将其组合；而且祂也以相反的事物构成它，如热与冷、干与湿、火与水、地与气，而这些相反的元素——整个宇宙由之组成——虽彼此不断争斗，却不互相毁灭。火没有蔓延烧尽一切；水没有泛滥淹没全地。然而，在我们的身体上，这些效应确实发生：胆汁增多时，热症产生，整个动物机体受损；痰液过多时，许多疾病产生，毁灭动物。但在宇宙中，此类事从不发生；每一事物仿佛被缰绳和束缚所持，按造物主的旨意守住自己的界限；它们的争斗成为整体的和平之源。这些事即使对于盲人岂不也是明显的吗？即使愚钝之人岂不也易于理解：它们是由某种眷顾所造、所持守的？因为谁是如此愚蠢无知，当看到如此庞大的物质、如此美丽、如此结合、如此巨大元素的不断争斗、它们的对立与持久，却不自问说：如果没有某种眷顾来持守这物质之体，不让宇宙分崩离析，它怎能存留？它怎能持久？季节的秩序如此完美，昼夜的和谐如此美妙，野兽、植物、种子、草本种类如此之多，它们保持自己的进程，而直到今日，没有一样曾陷入衰败或突然解体。\par

10. 我们不仅可以继续谈论这些事，还可以谈论许多其他更深奥的事；甚至可以对受造界本身作道德训诲；但将这些题目留待明日，让我们热切努力记住所说的话，并把它传达给其余的人。我确实知道，这些思辨的深奥对你们的耳朵似乎很陌生；但如果我们稍加警惕，并使自己习惯它们，我们就能轻易地教导他人。与此同时，还有必要向你们的\OriginalTerm{爱德}{Charity}说这些话。正如天主藉此伟大受造物赐给我们荣耀，我们也当以纯洁的品行荣耀祂！\Quote{诸天宣示天主的荣耀，}虽然只是被看见；因此我们不仅应在言语中，也应在沉默中，并以我们生活的光辉使所有人惊异，来宣示天主的荣耀。因为祂说：\Quote{你们的光也当这样照在人前，叫他们看见你们的好行为，便将荣耀归给你们在天上的父。}\BibleRef{玛 5:16}因为当一位不信者看见你——一位信徒——是温顺、谦逊、举止有节的，他会惊奇并说：\Quote{基督徒的天主真是伟大！祂造了什么样的人？祂从什么中造了他们，又把他们造成了什么？祂岂不是把他们从人变成了天使？若有人轻慢他们，他们不辱骂！若有人打他们，他们不愤怒！若有人伤害他们，他们为那使他们痛苦的人祈祷！他们没有仇敌！他们不知道怀恨！他们不犯虚浮的闲谈！他们没有学会说谎！他们不能忍受假誓，或者更确切地说，他们根本不发誓，宁愿割掉舌头，也不愿让誓言从口中发出！}这些就是我们应该使他们对我们说的话；我们应当根除发恶誓的恶习，至少像尊敬我们更贵重的衣服那样尊敬天主。因为这是多么荒谬：我们有一件比其余更好的衣服，就不允许自己一直穿着它；然而我们却随处漫不经心、毫无礼仪地拖着天主的名字！我恳切祈求并恳求你们，不要这样轻视自己的救恩；而是从起初我们为这条诫命所付出的关心，要一直坚持到底。因为我这样不断地劝戒你们关于发誓的事，并不是认为你们懈怠，而是因为我看见你们大部分已改正，所以我催促并敦促你们，使整个工作得以完成，达到圆满。正如观看竞技的观众，他们更激烈地鼓励那些接近奖品的人。所以，让我们绝不要疲倦；因为我们几乎已经接近这改正的完成；困难只是在起初。如今，恶习的主要部分已被切除，剩下要改正的已较少，不再需要劳苦，只需一段短时期内的适度警醒和勤奋，以便我们自己改正后，也能成为他人的导师；并能以极大的信心瞻仰神圣的逾越节，以极大的喜乐收获节日惯常喜乐的双倍或三倍。因为使我们喜悦的，与其说是从守斋的辛劳和疲劳中得解脱，不如说是带着光荣和应得的冠冕迎接那神圣的季节——这冠冕是永不衰残的！\par

11. 但为使改正更快发生，请做我所告诉你们的事。在你们房屋的墙上，并在你们心壁上，刻写那\Quote{\OriginalTerm{飞镰}{flying sickle}}；并思想它因咒诅而飞出，常常记住它。如果你们看见别人发誓，要制止、禁止，并为他小心，也要为你们自己的家人小心。因为如果我们留意这一点，不仅纠正自己，也带他人到同样的地步，我们自己就会很快到达目标；因为在我们承担教导他人时，若在我们自己身上似乎留下了未行我们所吩咐他们的事，我们就会羞愧脸红。不必多说；因为这些事已经讲了很多；现在说这些只是为提醒。但愿那比我们自己更爱护我们灵魂的天主，使我们在这事和一切善工上得以成全；好使我们完成所有义德之果，堪当承受天国，藉着我们的主耶稣基督的恩宠和慈爱，祂和祢及圣神，永生永王。阿们。\par

\SourceRecord{Translated by W.R.W. Stephens.；Nicene and Post-Nicene Fathers, First Series；Vol. 9.；Edited by Philip Schaff.；Buffalo, NY: Christian Literature Publishing Co.,；1889.；<http://www.newadvent.org/fathers/190109.htm>.}

% structure: p0001 source=h1 target=section reason=page-title

\section{关于雕像的讲道集 第十篇}

\Emph{称赞那些在用餐后前来听道的人；对自然界生理的观察；反对那些神化受造物的人；以及关于不可发誓的义务。}\par

1. 我欢喜，并同你们一起喜乐，因为你们确实实践了我们最近对那些缺席者所作的劝诫，因为他们没有守斋。因为我想，许多用了餐的人今天都来了，充实了这个美好的集会；这一点，我从周围所见更辉煌的景象和更多的听众可以推断出来。看来，我最近为他们说了那么多话并非徒然，向你们的爱德呼吁，要他们回到他们的母亲那里，并说服他们，即使在肉体食粮之后，也参与那属灵的是合法的。亲爱的，我问你们，你们在哪种情形下做得更好：是在上次集会时，用餐后便去睡觉；还是现在，用餐后前来聆听神圣的法律？最好的是当你们在广场上闲逛，参加毫无益处的聚会；还是现在，当你们与自己的弟兄们站在一起，聆听先知的神谕？亲爱的，用餐并非耻辱，但用餐后待在家里，从而被剥夺这神圣的宴席，才是耻辱。因为你们若待在家里，就会更加懒惰和消沉；但来到这里，你们将摆脱一切瞌睡和倦怠；不仅摆脱倦怠，也摆脱一切忧愁，你们会在一切可能发生的事上更加自在和振奋。\par

2. 那么还有什么更多可说的呢？只要站在那守斋的人旁边，你立即就能分享他的馨香；因为守斋是一种属灵的香料；通过眼睛、舌头和每一部分，它显示出灵魂的良好状态。我说这话，不是为了谴责那些用了餐的人，而是为了指出守斋的益处。然而，我并不把仅仅戒除肉类称为守斋；而是更早之前，戒除罪恶；因为那用餐后带着适当的节制来到这里的人，并不远远落后于那守斋的人；正如那继续守斋的人，若他不认真勤勉地留意所讲的话，就不会从他的斋戒中得到什么大益处。那吃饭却以适当的认真态度参与神圣集会的人，比那根本不吃且缺席的人好得多。这戒除绝不可能像参与属灵教导所赐予我们的益处和优点那样有益于我们。此外，你在哪里能听到你在这里所默想的事呢？若是你去审判台？那里是争吵和争辩！或是进入议事厅？那里是对政治事务的焦虑思虑！或是回到你的家？对你私事的忧虑从各方面折磨你！或是你到会场的会议和辩论中去？那里的一切都是属世的、可朽坏的！因为在那些聚集的人中所谈论的一切，都是关于商品、税务、丰盛的餐桌、土地买卖、其他合同、遗嘱、遗产，或其他类似的事情。即使你进入王宫的大厅，在那里你同样会听到所有人谈论财富、权力、或在此地受尊敬的荣耀，却没有一件是属灵的。但在这里，相反，一切都与天堂和天上的事有关；与我们的灵魂、我们的生命、我们为何出生、为何在世上度过规定的时间、我们如何从这迁移、以及在这些之后我们将进入什么状态、为何我们的身体是由泥土造成的、死亡的性质是什么、简言之，今生是什么、未来是什么。我们在这里所作的讲论丝毫不包含属世的东西，而全部是与属灵事物相关的。因此，我们将为我们的救恩做充分的准备，并怀着美好的希望离开这里。\par

3. 既然我们没有白白撒种，而是你们按我的劝告，把所有缺席的人都搜寻出来；现在请允许我们回报你们，提醒你们先前所说的一些事，并补充其余的部分。那么，先前所讨论的是什么呢？我们曾探究，在赐下圣经之前，天主是如何以怎样的方式安排他对我们的救恩计划；我们说过，祂藉着受造物来教导我们的种族，铺开诸天，在那里公开展开一部巨大的书卷，对愚笨者和智慧者、对穷人和富人、对西古提人和外邦人，以及所有居住在地上的人都有用；这部书卷远大于受它教导的人数。我们还详细谈论了黑夜、白昼及它们的秩序，以及它们所严格保持的和谐；又说了许多关于一年季节有规律的轮转及其均等的事。因为正如白昼在全年中不欺骗黑夜哪怕半小时，同样，这些季节也彼此均分所有的日子。但正如我先前所说，不仅受造物的伟大和美丽显明了神圣的工匠，而且它结合的方式以及运作方法也超越了自然的寻常过程。因为按照自然，水应当承载在地上；但现在我们却看到相反的情况：地由水支撑。按照自然，火应当向上；但现在我们却看到太阳的光线向下射向大地；水在天上，却不坠落；太阳在它们下面运行，却不被水熄灭，也不驱散它们的水汽。此外，我们还说，整个宇宙由四种元素构成，这些元素彼此敌对和冲突；但一个并不消耗另一个，尽管它们相互毁灭。由此可见，某种不可见的力量约束着它们，天主的旨意成了它们的束缚。\par

4. 今天，我想在这个主题上稍多停留。但请你们振奋精神，专心听我们讲！为使这奇事更清楚显明，我将从我们自己的身体来汲取关于这些事的教训。我们这身体，如此短小，由四种元素组成：即热的（血）、干的（黄胆汁）、湿的（痰）、冷的（黑胆汁）。但愿没有人认为这个主题与我们手头的事无关。\BibleQuote{因为属神的人能判断万有，但他自己却不受任何人判断。} \BibleRef{格前 2:15} 同样，保禄在给我们讲论复活时，也触及了农业的原理，他说：\BibleQuote{愚昧的人哪！你所播的种子，若不先死，决不会活过来。} \BibleRef{格前 15:38} 如果那位蒙福的人引进了农业的问题，那么若我们处理医学科学方面的事，也不应有人指责我们。因为我们现在的讲论是关于天主的创造；这个思想基础对我们的目的是必要的。所以，正如我先前所说，我们这个身体由四种元素组成；如果其中任何一个反叛整体，死亡就是这反叛的结果。例如，由于过热（胆汁过多）就产生热病；如果这超过一定限度，就导致迅速的瓦解。同样，当冷的元素过多时，就产生瘫痪、寒热、中风以及无数其他的疾病。每一种疾病都是这些元素过多的结果；当其中任何一个越过自己的界限，对其他元素实行暴政，破坏了整体的协调。那么，请质问那些说万物都是自发自生的人：如果这微小而短促的身体，有药物、医术、内在调节的灵魂、以及许多道德智慧和无数的其他帮助，却仍不能常保有序状态，而常常因内部某种扰乱而毁坏、消亡；那么像这样包含如此巨大实体、由同样元素复合而成的世界，如何能在如此长的时间内没有任何扰乱，除非它享有多种上智的安排？我们不可合理地认为，这个在内外都有照管的身体，几乎不足以自保；而像这样的世界，却没有这样的照管，竟能在如此多年中不遭受我们身体所遭受的那类事。因为，我问，为什么这些元素中没有哪一个越出自己的界限，吞没所有其他的？谁从起初把它们聚在一起？谁束缚了它们？谁约束了它们？谁在如此长的时期内使它们保持在一起？因为如果世界的身体是单纯均匀的，我所说的就不会如此不可能。但既然从起初就有元素之间的这种冲突，谁竟如此愚昧，以为这些事会自动聚集，并且在没有一位促成这结合的情况下，结合后还会继续保持？因为如果我们这些彼此怀有恶意的人（不是由于本性，而是由于意志）在彼此分歧、不友善时不能自发地达成一致，如果还需要另一个人使我们结合，并在结合之后进一步巩固我们，说服我们保持和好而不再次争吵；那么，那些既无感觉又无理性，且本性敌对、彼此为仇的元素，怎能聚集、同意并彼此保持，如果没有一种不可言喻的能力促成这结合，并在结合之后，一直用同样的纽带约束它们呢？\par

5. 难道你没有觉察到，灵魂离开后，这身体就衰败、枯槁、朽坏，而它的每个元素都归于其指定的地方吗？同样的事情，实际上也会发生在世界上，如果那常治理它的权能竟使其失去自己的上智安排。因为一艘船没有舵手就不能保持完整，而会很快沉没，那么世界若无人治理其航程，又怎能持续如此之久？为免我延长篇幅，假设世界是一条船：大地置于下面作为龙骨，诸天作为帆，人类作为乘客，下面的深渊作为海。那么，在如此漫长的时期，为何从未发生船难？现在，让一条船在没有舵手和船员的情况下航行一天，\BibleRef{宗 27:30}，你马上就会看到它沉没！但世界虽然已经存在了五千年，甚至更久，却从未遭受这样的事。但我为何谈论船？假设有人在葡萄园里搭了一个小棚，当果实采摘后，就任它空着；然而它连两三天都撑不住，很快碎裂倒塌！难道一个小棚没有照管就站不住？那么，这样美好奇妙的世界的工程——夜与昼的律则，季节的交替轮舞，整个大地、海洋、天空中的自然进程，处处五彩缤纷、变化万千，在植物、飞鸟、游鱼、行走、爬行的动物中，在比这些都尊贵得多的人类种族中——怎能没有某种上智安排而在如此漫长的时期中持续不坠呢？但除了已说的之外，请随我列举：草原、花园、各类花朵；各种药草及其用途；它们的香气、形态、排列，甚至它们的名称；结果实和不结果实的树；金属和动物的本性——在海中的或陆地上的；游水的和翱翔的；群山、森林、树丛；下面的草原和上面的草原；因为地上有草原，天上也有草原；星辰的各式花朵；下面的玫瑰和上面的彩虹！你还要我指出鸟类的草原吗？请看孔雀那斑斓的身体，胜过一切染料，以及紫色羽毛的鸟。请与我一起观想天空的美丽：它如何这么久没有被遮蔽，仍然像今日才造的那样明亮清澈；此外，大地的力量，它的子宫如何没有因这么长时间的生育而衰竭！请与我一起观想泉源：它们如何从被造以来就不断涌流，昼夜不息！请与我一起观想海洋：接受如此多的河流，却从未超越它的界限！但我们还要追逐不可企及的东西多久呢？的确，关于以上所说的每一件，我们都应该说：\BibleQuote{上主，你的化工何其浩繁，你以智慧造成一切。}\par

6. 当我们与不信者逐一细述这些细节时——受造界的伟大、美丽、丰饶、处处彰显的慷慨——他们所谓的智慧论证是什么呢？他们说，正是这一点是最坏的缺点，天主使世界如此美丽、如此广阔。因为，若他不曾使它美丽广阔，我们就不会把它神化；但如今，我们被它的宏伟所震撼，惊叹于它的美丽，就以为它是神祇了。但这样的论证毫无价值。因为使这世界既非其伟大也非其美丽成为这种不敬的原因，而是他们自己的缺乏理解力，这就是我们准备要证明的，用我们自身的情况来证明，我们从未受过这种影响。那么为什么\Quote{我们}没有把它神化？难道我们不与他们用同样的眼睛看？难道我们不从受造界享有同样的利益？难道我们没有同样的灵魂？难道我们没有同样的身体？难道我们不踏着同一块大地？为什么这美丽和伟大没有说服我们与他们想的一样？但这一点不仅从这证据，也从另一个证据显明。因为，作为证明他们并非因美丽而神化世界，而是出于自己的愚昧，为什么他们敬拜猿猴、鳄鱼、狗和最卑贱的动物？的确，\BibleQuote{他们思念致成虚妄，他们那愚昧的心就被蒙蔽了；他们自称有智慧，却成了愚拙。}\BibleRef{罗 1:21}\par

7. 尽管如此，我们不仅要以这些事来作答，还要进一步说明。因为天主从古时就预见了这些事，并以祂的智慧摧毁了他们的这个借口。为此，祂造的世界不仅奇妙宏大，而且也是可朽坏和易逝的，并在其中放置了许多它软弱的证据。祂对宗徒所做的，也对整个世界做了。祂对宗徒做了什么呢？因为他们常行许多伟大惊人的奇事和征兆，祂就容许他们不断被鞭打，被驱逐，住在监牢，遇到身体的疾病，处于持续的困苦中，免得他们奇迹的伟大使他们在人中被视为神。因此，当祂赐予他们如此大的恩宠时，祂容许他们的身体是可死的，并且在许多情况下易受疾病侵扰；祂没有除去他们的软弱，为要完全证明他们的本性。这并非我个人的断言，而是\Person{保禄}自己的话，他说：\BibleQuote{我即使愿意夸耀，也不算是狂妄；但我现在绝口不谈，免得有人把我估计得超过他所看见我在他身上的，或由我所听到的。} \BibleRef{格后 12:6} 又说：\BibleQuote{但我们是在瓦器中存有这宝贝。} \BibleRef{格后 4:7} 所谓\Quote{瓦器}是什么意思呢？他是指这身体，是可死和易逝的。正如瓦器由泥土和火形成，同样，这些圣人的身体是泥土，接受了属神之火的能量，就成了瓦器。但为什么如此，将这样大的宝贝和如此丰富的恩宠托付给可死和易朽的身体呢？\BibleQuote{为要使这超越的能力归于天主，而不是出于我们。} 因为当你看见宗徒们复活死人，而他们自己却生病，不能除去自己的软弱，你就能清楚明白，使死人复活并非出于复活者的能力，而是出于圣神的德能。为证明他们常生病，请听\Person{保禄}关于\Person{弟茂德}所说的话：\Quote{为了你的胃和你的屡次虚弱，应用少许酒。} 又说关于另一个人：\BibleQuote{\Person{特洛斐摩}我留在\Place{米肋托}病了。} \BibleRef{弟后 4:20} 在写给斐理伯人的书信中，他说：\BibleQuote{\Person{厄帕洛狄托}病得要死。} \BibleRef{斐 2:25} 因为，如果在这种情况下，他们尚且被当作神，并准备向他们献祭，说：\BibleQuote{神取了人形降到我们这里来了；} \BibleRef{宗 14:11} 如果这些软弱不存在，当人们看到他们的奇迹时，会堕入何等的不虔敬呢？所以，正如在这种情况下，由于这些异能之伟大，祂让他们本性停留在软弱状态，并容许那些反复的考验，以免他们被认为是神；同样，祂对受造界也这样做了，这是一件近乎平行的事。因为祂使受造界美丽而宏大，但另一方面又可朽坏。\par

8. 这两点圣经都教训我们，因为一位作者论及天空的美妙时这样说：\BibleQuote{诸天述说天主的荣耀。}\newline{} 又说：\BibleQuote{祂铺张穹苍如幔子，展开诸天如可住的帐棚。}\BibleRef{依 40:22}\newline{} 又说：\BibleQuote{祂坐在地球大圈之上。}\newline{} 但另一位作者指出世界虽大而美，却是可朽坏的，这样说：\BibleQuote{主啊，你起初立了地的根基；诸天是你手的工程。它们都要灭亡，你却长存；它们都要像衣服渐渐旧了，你要将天地卷起来，像一件外衣，天地就都改变了。}\newline{} 达味论及太阳又说：\BibleQuote{它如同新郎出洞房，又如勇士欢然奔路。}\newline{} 你看见他如何将这颗星辰的美妙和伟大摆在你面前吗？因为正如新郎从华丽的洞房出现，太阳从东方射出它的光线；它仿佛用藏红花色的面纱装饰天空，使云彩如玫瑰，终日无阻地奔驰，没有障碍阻挡它的行程。你看见它的美丽吗？你看见它的伟大吗？也看看它软弱的证明！因为一位智者为使这事显明，说：\BibleQuote{比太阳更光明的是什么？然而这光也会被遮蔽。}\BibleRef{德 17:31}\newline{} 不仅从这一情况可看出它的软弱，而且在云层汇聚时也是如此。至少，当一片云从它下面经过时，它虽射出光芒，努力穿透云层，却没有力量做到；因为云太密，不让它透过去。\Quote{它滋养种子，}有人回答说——是的——但它并非独自滋养它们，还需要大地、露水、雨水、风和季节的适当配合的帮助。除非这一切都协同，太阳的帮助便是多余的。但这似乎不像神性，因为它需要别人的帮助才能做它想做的事；因为天主的特有属性是无所需；祂自己至少没有用这种方式使种子从地里长出；祂只一吩咐，它们都发芽了。再者，为使你明白不是元素的本性，而是祂的命令成就一切；祂使这些先前不存在的元素产生；并且无需任何帮助，祂为犹太人降下玛纳。因为经上说：\BibleQuote{祂从天上赐给他们食粮。}\newline{} 但我为何说太阳为了果实的完美需要其他元素的帮助来维持它们；当它自己需要许多事物的帮助来维持自己，且它本身不足以自给呢？因为为了能前进，它需要天空如同铺在它下面的道路；为了能照耀，它需要空气的清澈和稀薄；因为如果空气变得异常浓密，它就不能显示它的光；另一方面，它需要凉爽和湿润，免得它的光线对一切变得难以忍受，烧毁一切。因此，当其他元素克制它、纠正它的软弱（克制，例如云、墙和某些其他物体拦截它的光；或纠正它的过度，如露水、泉水和凉爽的空气），这样的一个东西怎能是神呢？因为天主必须是独立的，不需要帮助，是万善的源头，不受任何阻碍；正如保禄和依撒意亚先知论及天主所说的；后者让祂亲自这样说：\BibleQuote{我充满天地，上主说。}\BibleRef{耶 23:24}\newline{} 又说：\BibleQuote{我岂是近处的天主，不也是远处的天主吗？}\BibleRef{耶 23:23}\newline{} 达味又说：\BibleQuote{我曾对上主说：你是我的主，你不需要我的善。}\newline{} 但保禄证明这种独立而不需帮助，并表明这两件事特别属于天主：一无所需，并亲自供给一切给万物，这样说：\BibleQuote{创造天、地、海的天主，本身不需要任何东西，将生命和一切赐给万物。}\par

9. 我们很容易考察其他元素——天、空气、地、海——并显示它们的软弱，以及每个元素如何需要邻舍的帮助，没有这帮助就会消亡。因为就地而言，如果泉水对它失效，从海和河流注入的湿气消失，它就很快因干涸而毁灭。其余元素也相互需要：空气需要太阳，太阳也需要空气。但为了不延长这篇讲道，在已经说过的内容中，我们为愿意领受的人提供了足够的理由作为出发点，我们就满足了。因为如果太阳，这整个创造中最惊人的部分，已被证明如此软弱和需要帮助，那么宇宙的其他部分更是如此。因此，我所提出的（将这些事提供给学生思考），我将再次从圣经中用讲论向你们展示，并证明不仅太阳，而且整个宇宙都是如此可朽坏的。因为既然元素相互毁灭，当大量寒冷介入时，它抑制太阳光线的力量；另一方面，热量占优势时，消耗寒冷；既然元素是彼此相反特质和倾向的原因和主体，那么很明显，这些事物提供了极大可朽坏性的证据；并且所有这些可见之物都是属于物质的实体。\par

10. 但因这主题对于我们的卑微过于高深，现在请允许我引领你们到圣经的甘泉，使我们的耳朵得以舒畅。因为我们不向你们分开谈论天和地，而是要展示宗徒以这些直白的言语向我们宣告关于整个受造物的这事：即整个受造物如今处于败坏的束缚中；以及它为何如此被束缚，以及何时将从中获得释放，并被转化为何种状态。因为在他说了：\BibleQuote{现时的苦楚，与将来在我们身上要显示的光荣，是不能相比的；}之后，他接着又加上：\BibleQuote{凡受造之物都热切地等待天主子女的显扬。因为受造之物被屈伏在败坏的状态之下，并不是出于自愿，而是由于使它屈伏的那一位的决意；但受造之物仍怀有希望。}\BibleRef{罗 8:21}。但他的意思是：\BibleQuote{受造物}，他说，\BibleQuote{成了可朽坏的}；因为这已隐含在\BibleQuote{被屈伏于虚空}这句话中。因为它是由天主的命令成为可朽坏的。但天主如此命令是为了我们人类；因为既然它要养育可朽坏的人，它本身也必须是同样的性质；因为当然可朽坏的身体不能居住在不可朽坏的受造物中。但他告诉我们，它不会始终如此。\BibleQuote{受造物本身也将从败坏的束缚中获得释放；}后来，为了显示这事件何时以及藉着谁发生，他补充说：\BibleQuote{进入天主子女的光荣自由。}因为当我们复活，按他的意思，并接受不可朽坏的身体时，那么这天的身体、地和整个受造物，也将成为不可朽坏、不可消灭的。因此，当你看到太阳升起时，要赞美造物主；当你看到它隐藏自己并消失时，要学习它本性的软弱，使你不致将它作为神祇来崇拜！因为天主不仅在元素的本性中植入了这软弱的证据，还命令祂的仆人——不过是人——去命令它们；因此，纵然你无法从它们的外表知道它们的奴役状态，你仍可从那些命令过它们的人得知，它们全都是你的同仆。因此，纳维之子\Person{若苏厄}说：\BibleQuote{太阳停在\Place{基贝红}，月亮停在\Place{阿雅隆谷}。}再者，先知\Person{依撒意亚}在\Person{希则克雅}王统治时，使太阳倒退；\Person{梅瑟}命令了空气、海、地和岩石；\Person{厄里叟}改变了水的性质；\Person{三少年}战胜了火。你们看，天主如何双方面地为我们预备：藉着元素的美引领我们认识祂的神性；又藉着它们的软弱，不让我们陷入对它们的崇拜。\par

11. 为此，让我们光荣祂，我们的守护者；不仅用言语，也用行为；让我们表现出优秀的生活方式，不仅普遍而言，尤其关于戒除发誓。因为并非每样罪都带来相同的惩罚；而是那些最容易改正的，带给我们最大的惩罚：这确实是\Person{撒罗满}所暗示的，当他说：\BibleQuote{若有人偷窃，并不足怪，因为他偷窃是为满足饥饿的灵魂；但奸淫者因缺乏理智，毁灭自己的灵魂。}但他的意思是：窃贼是严重的犯罪者，但不像奸淫者那样严重：因为前者，虽然其行为的理由可悲，但仍有贫困所致的必要性可以申辩；但后者，没有必要性强迫他，仅凭疯狂就坠入罪恶的深渊。这同样可以说是关于那些发誓的人。因为他们没有任何借口可提，唯有他们的轻视。\par

12. 我确实知道，我可能显得过于啰嗦和令人生厌；并且可能被认为因继续这劝诫而令人恼火。但我仍不放弃，为使你们甚至因我的厚颜而感到羞愧，从而戒除发誓的习惯。因为如果那不仁慈且残酷的判官，因尊重寡妇的纠缠而改变了习惯，你们更会这样做；尤其当那劝诫你们的人，不是为自己，而是为你们的得救。或更确切地说，我不能否认我这样做也是为自己；因为我把你们的益处视为我自己的成功。但我希望你们，正如我为你们的得救而劳苦费力一样，你们也当为自己的灵魂挂虑；那么这改良的工作必然完成。何需多言？因为若没有地狱，没有为悖逆者的惩罚，也没有为服从者的赏报；并且我来向你们，以请求恩惠的方式提出这事，你们难道不会同意吗？当我请求如此微小的恩惠时，你们难道不会应允我的请求？但当这恩惠是天主所请求，且是为你们这些施予者的益处，而非为祂自己领受者的益处；谁会是如此无礼，如此可怜和悲惨，而不肯应允天主的请求呢？尤其当那施予者将来要享受这恩惠的益处？考虑这些事，那么，当你们离去时，要对自己重复一切所说的话；并用各种方法纠正那些不留意的人；为使我们可以因他人和自身的善行得到报偿，藉着我们的主\Person{耶稣基督}的恩宠和慈爱，愿光荣归于祂和\Person{圣父}及\Person{圣神}，至于无穷之世。阿们。\par

\SourceRecord{Translated by W.R.W. Stephens.；Nicene and Post-Nicene Fathers, First Series；Vol. 9.；Edited by Philip Schaff.；Buffalo, NY: Christian Literature Publishing Co.,；1889.；<http://www.newadvent.org/fathers/190110.htm>.}

% structure: p0001 source=h1 target=section reason=page-title

\section{讲道集 第十一：论雕像}

\Emph{感谢天主救我们脱离因叛乱而预期的灾祸；并追忆当时所发生的事件。同时反驳那些挑剔人体结构的人，并一般地论及人的受造；最后，论及避免起誓的成功。}\par

1. 当我回想过去的暴风雨和现在的宁静时，我不断说：\Quote{愿天主受赞美，祂创造万有，改变万物；祂从黑暗中引出光明；祂引到地狱之门，又带回；祂责罚，但不杀死。} 我也希望你们时常重复这句话，永不停止。因为如果祂以行动恩待了我们，我们若不以言语还报，又怎能得到宽恕？因此，我劝勉我们永不停止感谢祂；因为如果我们对先前的恩惠感恩，显然我们也必会享受其他更大的恩惠。那么，让我们不断地说：愿天主受赞美，祂准许我们安全地在你们面前摆设惯常的筵席，同时也恩赐你们安心地聆听我们的话语！愿天主受赞美，我们不再因外在的危险而逃到这里，而只是出于听道的渴望；我们不再以痛苦、战栗和焦虑彼此相遇，而是怀着极大的信心，摆脱了所有恐惧。事实上，我们前些日子的处境，并不比那些在深海中颠簸、时刻期待船难的人更好。我们整天被无数的谣言惊吓，四面受扰，动荡不安；每天忙碌好奇地想知道：谁从宫廷来了？他带来了什么消息？以及所报告的是真是假？我们的夜晚也彻夜不眠，当我们看着这城时，我们为她哭泣，仿佛她即将毁灭。\par

2. 正因如此，你们先前也保持沉默，因为全城空虚，所有人都迁移到了旷野，而那些留下的人被沮丧的阴云笼罩。因为灵魂一旦充满沮丧，就不易听取任何话语。因此，当\Person{约伯}的朋友们到来，看见他家的惨剧，那义人坐在粪堆上，满身疮痍，他们就撕裂了外衣，叹息，并沉默地坐在他旁边；这表明对于受苦者，起初最适合的莫过于安静和沉默。因为灾难太大，无法安慰。因此，当犹太人受奴役，从事泥土和制砖工作时，他们看见\Person{梅瑟}来到他们中间，却因精神沮丧和痛苦，无法留意他的话。胆怯的人有这样的感受，不足为奇，因为我们发现门徒也陷入同样的软弱。因为在那个奥秘的晚餐之后，当基督把他们带到一边，与他们交谈时，门徒起初不止一次地问祂：\Quote{祢往哪里去？} 但当祂告诉他们不久后将遭遇的祸患——战争、迫害、普遍的敌意、鞭打、监禁、法庭、见官——那时，他们的灵魂如同被沉重的负担压迫，因祂所说之事的恐惧和这些临近事件的悲伤，从此陷入麻木状态。因此，基督看到他们的惊惶，就责备他们说：\Quote{我往父那里去，你们中没有人问我：\InnerQuote{祢往哪里去？} 只因我给你们说了这些，忧愁就充满了你们的心。} 因此，我们先前也沉默了一段时间，等候当前的时机。因为如果一个人要向某人求恩惠，即使请求是合理的，也会等待适当的时机提出，以便找着那乐意施恩的人处于温和良好的心境，并借助良机获得恩惠；那么，说话者岂不更应该寻找合适的时机，以便向一位心情良好、摆脱一切忧虑和沮丧的听众发言？我们正是这样做了。\par

3. 如今你们既已摆脱沮丧，我们愿意唤起你们对先前事物的回忆，好让我们的讲论对你们更为清晰。因为我们关于受造物所说的，即天主不仅使之美丽、奇妙、伟大，而且也使之脆弱、可朽；此外，祂还为此设立了各种证据，为我们的益处安排这两种情况：借其美丽引导我们钦佩造物主，借其脆弱引导我们远离对受造物的崇拜；这一点我们也可以从身体上看到。因为关于身体，真理的敌人以及我们自己队伍中的许多人，也提出疑问：为何它被造成可朽和脆弱的？许多外邦人和异端者甚至断言身体并非天主所造。他们宣称身体不配天主的创造技艺，并大肆渲染其不洁、汗水、眼泪、劳苦、痛苦，以及身体的其他所有遭遇。但对于我来说，当听到这类议论时，我首先要作如下答复：不要对我谈论那堕落、卑下、被定罪的人。但若你愿知道天主起初以怎样的身体塑造了我们，就让我们进入乐园，察看那起初被造的人吧。因为那身体并非如此可朽和会死；反而如同刚从炉中取出的金像，灿烂发光，那形体毫无朽坏。劳苦不烦扰它，汗水不玷污它。忧虑不攻击它，悲伤不围困它，也没有任何其他类似的情感折磨它。但当人不能以节制承受他的幸福，反而蔑视他的恩主，认为一个欺骗的魔鬼比照顾他、提升他尊位的天主更可信，并期望自己也成为神，怀着超越自身尊严的念头时——那时，确实那时，天主为以决定性的行动使他谦卑，便使他变得会死和可朽，并用如此多样的需要束缚他；这不是出于仇恨或厌恶，而是出于对他的关怀，并从一开始就抑制那邪恶而毁灭性的骄傲；而不让它再进一步，祂以实际经验告诫他：他是会死和可朽的；这样便说服他绝不能再像以前那样思想或梦想这类事。因为魔鬼的暗示是：\BibleQuote{你们将像天主一样。} \BibleRef{创 3:5} 于是，为彻底根除这个念头，天主使身体遭受许多痛苦和疾病；以其本性教导他绝不能再怀有这样的思想。这事真实无伪，从所发生的事可最清楚地看出：因为经过如此期待之后，他被判处了这刑罚。也请与我一同思考天主在此事上的智慧。祂没有让他第一个死去，却允许他的儿子经历这死亡；以便他在眼前看见身体腐烂朽坏时，从那景象中领受鲜明的智慧教训，明白所发生的事，并在离世之前受到适当的惩戒。\par

4. 确实，正如我所说，这一点从已发生的事上已显而易见；但从将要说明的事上也将同样清楚。因为若我们被如此多的身体需要束缚，若死亡、腐朽、当众腐烂、化为尘土是所有人的命运，以至于外邦哲学家为人类作出一个共同的定义（当被问及人是什么时，他们回答：人是理性的、会死的动物）；倘若，尽管如此，所有都承认这一点时，竟有人敢在众人面前将自己神化；而且尽管视觉本身见证了其死亡，他们仍野心勃勃地被称为神，并被尊为神；那么，若死亡没有不断教导所有人我们本性的死亡和腐朽，许多人会趋向何种不敬的程度呢？例如，请听先知论及一个外邦国王，当他染上这种疯狂时所说的话：\BibleQuote{我要高举我的宝座在星辰之上，我要与至高者相似。} \BibleRef{依 14:13} 之后，先知嘲笑他，谈及他的死亡说：\BibleQuote{朽坏在你下面，蛆虫是你的覆盖。} \BibleRef{依 14:11} 但他的意思是：\Quote{人哪，这样的结局在等待你，你竟敢怀有这些想象吗？} 又论及另一个人，即提洛王，当他怀有类似的目标并野心勃勃被视为神时，先知说：\BibleQuote{你不是神，而是人，刺透你的人必这样说。} \BibleRef{则 28:9} 因此，天主使我们的身体成为如今这样，从一开始就彻底消除了偶像崇拜的一切机会。\par

5. 然而，若是这事发生在身体上，又何必惊奇呢？因为连关于灵魂也是同样明显的事已经发生了。因为天主并没有使它成为可死的，却允许它成为不死的；然而祂却使它受制于遗忘、无知、忧伤和挂虑；这是为了怕它由于自身的尊贵出身，而对自己应有的尊严抱持过高的自负。因为即使是在这种情形下，有些人还敢妄称它属于\OriginalTerm{天主的本性}{Divine essence}；那么，倘若它没有这些缺陷，他们岂不会陷入何等疯狂的境地？然而，关于受造之物我所肯定的，我同样也肯定关于身体：这两者都同样激发我对天主的赞美；就是祂使它成为可朽的；而正是在它的可朽性中，祂彰显了祂自己的德能和智慧。因为祂本可用更好的材料来造成它，祂已从诸天和太阳的实质证明了这一点。因为那使那些东西成为现在这样的，若祂愿意，也能使这个像它们一样。但它的不完美原因是我前面已经提及的。这种情况绝不降低对造物主工艺的赞叹，反而增加它；因为物质的卑贱，彰显了祂技艺的资源和适应性；因为祂在泥土和灰烬中引入了如此和谐的部分，以及如此多样、繁多且能具有如此精神智慧的感官。\par

6. 因此，你越是批评物质的卑贱，就越该惊叹所展现艺术的伟大。为此，我并不那么钦佩那位用金子塑造美丽形象的雕刻家，而是钦佩那位藉艺术资源甚至能在易碎的泥土中展示出奇妙而无可仿效的美丽造型的人。在前一种情况下，材料给艺术家一些帮助；但在后一种情况下，却纯粹是他的艺术的展现。那么，你若要学习造物主何其伟大的智慧，请想一想合土造成的究竟是什么？除了砖和瓦，还有什么？然而，天主的至高艺术家，用只形成砖和瓦的同一材料，竟能造出如此美丽的眼睛，使所有看见它的人都惊讶，并赋予它如此力量，使它能同时遍览高空，并借助小小的瞳孔拥抱群山、森林、丘陵、海洋，甚至苍天，凭着如此渺小之物！那么，不要向我提泪液和分泌物，因为这些都是你罪的果实；但要考虑它的美丽和视觉能力；以及它如何在如此广阔的空间中漫游，却不感到疲倦或不适！脚走过一小段距离就会疲劳无力；但眼睛在如此高远宽阔的空间中穿梭，却感觉不到任何软弱。因为这是所有肢体中对我们最必要的，祂没有让它被疲劳压垮，为的是它服务我们能自由而不受束缚。\par

7. 然而我更要说，什么语言能完全描述这个肢体的一切卓越呢？我为何只谈论瞳孔和视觉能力？因为如果你审视所有肢体中看似最卑微的部分，我指的是睫毛，你也会在其中看到造物主天主的多种智慧！正如关于麦穗那样，麦芒像某种长矛一样挺拔，驱赶鸟类，不让它们落在果实上，折断那太过柔弱而无法承受它们的麦秆；眼睛的情形也是如此。眼睑的毛发排列在前方，起着类似麦芒和长矛的作用：将灰尘和轻小之物阻挡在眼睛之外，以及任何可能妨碍视线的东西；并且不让眼睑受到骚扰。另一个同样值得注意的智慧实例可以在眉毛中观察到。谁能不为它们的位置而惊叹呢？因为它们既不会过分突出以致遮蔽视线，也不会退后太多而不适当；而是像房屋的屋檐一样，从上方突出，承接从额头流下的汗水，不让它干扰眼睛。为此目的，它们上面还长有毛发，藉其粗糙起到阻挡上方来物的作用，并提供所需的精确保护，也为眼睛增添了不少美观。这还不是唯一值得惊奇之处！还有另一件事同样如此。我问，为什么头上的毛发会生长并被剪断，而眉毛却不是这样？因为这事也并非无意的或偶然发生的，而是为了不让它们因变得太长而过度遮蔽视线；那些到了极老年龄的人就会遭受这种不便。\par

8. 谁能完全穷尽藉着大脑所彰显的一切智慧呢！首先，祂使它柔软，因为它作为一切感官的源泉。其次，为避免它因其特殊性质而受损，祂用骨骼从四面八方将其加固。再者，为免它因骨骼的坚硬而摩擦受损，祂在其中间了一层薄膜；并且不只一层，而是两层：第一层紧贴在头骨下方，第二层包裹大脑的上层物质，而第一层比第二层更坚硬。这样做既为了上述原因，也为了使头部受到打击时，大脑不致首当其冲，而是由这些薄膜首先承受打击，使它免于一切伤害，保持完好。此外，覆盖大脑的骨头并非单一连续的一块，而是四面有许多骨缝，这大大有助于其安全。因为通过这些骨缝，周围的气体容易向外散发，以防止它窒息；而且如果某一点受到打击，损伤不会扩散到整个大脑。因为如果骨骼是一整块连续的，即使打击只落在一处，也会伤及全部；但如今由于它分成许多部分，这种事情绝不会发生。因为如果某一部分不幸受伤，只有靠近该部分的骨头受损，其余部分安然无恙；打击的连续性因骨骼的分隔而被截断，无法扩展到相邻部分。为此，天主用许多骨骼为大脑建造了覆盖物；正如人建造房屋时，在屋顶铺上瓦片，天主也将这些骨骼置于头顶上方，并使头发长出，充当一种帽子。\par

9. 对于心脏，祂也做了同样的事。因为心脏在我们身体的所有肢体中具有优越性，我们整个生命的最高权力都托付给它，而且只要受到轻微打击就会导致死亡；祂用坚硬粗糙的骨骼将其四面围护，前面以胸骨保护，后面以肩胛骨保护。至于祂对大脑的薄膜所做的，在此处也做了。为免心脏因搏动以及在愤怒和类似情绪下的快速跳动，与包围它的坚硬骨骼碰撞而摩擦疼痛，祂在中间安置了许多薄膜，并将肺部放在心脏旁边，像一张柔软的床垫承受这些跳动，使心脏能在其上释放其力量而不受损伤或痛苦。\par

但我为何谈论心脏和大脑呢？倘若有人甚至考察指甲，也会在这些指甲上看到天主多方面的智慧，无论是其形状、质地还是位置。我还可以提及为何我们的手指并不全都相等，以及其他许多细节；但对于愿意留意的人而言，创造我们的天主的智慧，从以上所述已足够清楚。为此，我将这部分留给那些渴慕这项任务的人去勤勉探究，而转向另一个异议。\par

10. 确实，除了已经提到的之外，有许多人还提出这样的异议：如果人是走兽的君王，为什么许多动物在力量、敏捷和速度上胜过人呢？因为马更快，牛更耐久，鹰更轻盈，狮子更强壮。那么，我们对此说法有何回应呢？答案是：正因如此，我们尤其能辨识天主的智慧和祂赐予我们的尊荣。诚然，马比人跑得快，但就赶路而言，人比马更合适。因为即使是最快最强的马，一天也几乎走不了二百斯塔迪亚；但人通过连续套用多匹马，能完成两千斯塔迪亚的路程。因此，敏捷给予马的优势，理智和技艺以更大的程度给予人。人的脚确实不如马脚强壮，但他有马脚为他服务，如同己有。因为没有一种走兽能够征服另一种为己所用；但人却掌握着它们全部，并藉着天主赐予他的多样技能，使每种动物服务于最适合它的用途。因为如果人的脚像马脚那样强壮，它们在其他用途上就无用了，比如崎岖地面、山顶、爬树；蹄子通常妨碍在这些地方行走。因此，虽然人的脚比它们的柔软，却适用于更多样的用途，并不因缺乏力量而更差，同时还有马的力量为他们服务，并且他们在行走的多样性上胜过马。再者，鹰有轻盈的翅膀；但我有理智和技艺，能降服并掌控所有有翼的动物。但如果你想看我的翅膀，我有一双比它轻盈得多的翅膀；它能翱翔，不仅十或二十斯塔迪亚，甚至高到天上，而且超越天本身，超越诸天之上的天；直至\Quote{基督坐在天主的右边！}\par

11. 再者，无理性的动物身上自有其武器：牛有角，野猪有獠牙，狮子有爪。但天主并未让我的身体本性配备武器，反使之与身体分离，为要表明人是温和的动物；并且我不必时刻使用武器，因为我能时而放下，时而拿起。为使我在这事上得以自由、不受束缚，不必始终携带武器，祂便使武器与我的本性分离。因为我们不仅因拥有理性本性而超越禽兽，在身体上也胜过它们。因为天主使身体与灵魂的尊贵相应，并适宜执行灵魂的命令。祂并非随意将身体造成如此，而是照其应有的样式，为服务于理性的灵魂；倘若身体不是如此，灵魂的运作必大受阻碍：这从疾病中便可看出。因为身体的精妙调节若稍有偏离正常状态，灵魂的许多能力便受阻碍，例如脑部过热或过冷。因此，从身体便容易看出天主眷顾的许多迹象：不仅因为祂起初造得比现在更好，也不仅因为如今祂为有益的目的而改变了它，更因为祂将使它复活，得享更大的荣耀。\par

12. 但若你愿意以另一种方式学习天主关于身体所显示的智慧，我要提及保禄似乎最常感到惊奇的一点。那是什么呢？就是祂使各肢体互相优越，却不在同一事上。祂命一些肢体在美丽上超越其余，一些在力量上。眼睛是美丽的，但脚更强壮。头是尊贵的，但不能对脚说：\BibleQuote{我不需要你。}\BibleRef{格前 12:21}这在无理性的动物身上也可看到；在生活的一切关系中也是如此。例如，国王需要臣民，臣民也需要国王；正如头需要脚。再者，关于禽兽：有些比其余更强壮，有些更美丽。有些使我们愉悦，有些滋养我们，有些给我们衣服。孔雀使人愉悦；家禽和猪滋养人；绵羊和山羊提供衣物；牛和驴分担我们的劳动。还有其他动物不提供这些，却唤起我们的能力。例如，野兽增强了猎人的力量；它们所激起的恐惧教诲我们种族，使我们更加谨慎；并且从它们身上为医药目的提供了不少贡献。因此，若有人对你说：\Quote{你既是禽兽的主子，为何还怕狮子？}你当回答他说：\Quote{起初并非如此安排，那时我蒙天主恩宠，住在乐园里。但当我得罪了我的主人，就落到了那些本是我仆役的权势之下！然而，甚至现在也不完全如此，因为我掌握了一门技艺，能战胜野兽。}大户人家也是如此：儿子们年幼时，怕许多仆役；但若犯了错，恐惧就大大增加。关于蛇、蝎子和蝰蛇，我们也可说：因着罪，它们使我们感到可畏。\par

13. 这种差异不仅见于我们的身体和生活的各种境遇，也不限于禽兽，还可在树木中看到；最卑微的树木也显示其优于高大者的长处，如此万物并非事事相同，为使我们都需要一切，并体认主的多重智慧。因此，不要因身体的可朽坏而归咎于天主，反而要为这事向祂致敬，并因祂的智慧和慈爱而钦佩祂：因祂的智慧，在于能在如此可朽坏的身体中展现如此和谐；因祂的慈爱，在于为灵魂的益处而使之朽坏，为要抑制她的虚荣，制服她的骄傲！有人问：为什么祂起初不这样造它呢？我回答：正是为了藉这些事工在你们面前证明自己，正如以结果本身说：我召叫你们获得更大的尊荣，但你们使自己不配这个恩赐，自绝于乐园！然而，即使现在我也不轻视你们，反要纠正你们的罪，带领你们回归天堂。因此，为了你们的益处，我容许你们长久地败坏朽坏，为在时期圆满时，谦卑的纪律得以确立，使你们永不再重拾昔日的高傲。\par

14. 为这一切，让我们感谢爱人的天主；并为祂对我们的慈爱，向祂献上那于我们自己也有益的回报；而对于我常向你们讲论的诫命，让我们竭尽全力！因为我不会停止劝勉，直到你们改过：因为我们所追求的，不是少向你们说话或多说话，而是继续讲论，直到说服你们。天主曾藉先知对犹太人说：\BibleQuote{你们若在争吵和辩论中禁食，何必为我禁食呢？}\BibleRef{依 58:4}而现在祂藉我们向你们说：你们若在发誓和伪誓中禁食，何必禁食呢？因为我们将如何瞻仰圣逾越节？我们将如何领受圣祭？我们将如何以同一舌头——我们曾用以践踏天主的法律、玷污灵魂的舌头——分享那些奇妙的奥迹？因为若无人敢以污秽的手接受皇袍，我们怎敢以被污染的舌头领受主的身体！因为发誓出于恶者，但圣祭属于主。\BibleQuote{光明与黑暗有什么相通？基督与贝里雅耳有什么和谐？}\BibleRef{格前 6:14}\par

15. 我知道你们确实渴望摆脱这种不虔；但由于每个人独自可能不易做到这一点，让我们就此组成团体和伙伴关系；如同穷人在他们的宴席中所为，当每个人单独无法备办一席完整的盛宴时，他们就聚集一堂，各自带来自己的那份贡献；我们也这样行。因为我们自身过于懈怠，让我们彼此结为伙伴，承诺贡献劝告、训诫、劝勉、责斥、提醒和警告；这样，我们每个人都能因着彼此的勤勉而得到矫正。因为我们观察邻人的事比观察自己的事更为敏锐，让我们留心他人的安全，把自己的监护权托付给他们；让我们投身于这种虔敬的竞争，以便如此胜过这种恶习，得以勇敢地参与这神圣的筵席；怀着良好的望德和纯洁的良心，分享这神圣的祭献；靠着我们的主\Person{耶稣基督}的恩宠和慈爱，愿光荣借着祂并偕同祂，归于圣父及圣神，于无穷世之世。阿们。\par

\SourceRecord{Translated by W.R.W. Stephens.；Nicene and Post-Nicene Fathers, First Series；Vol. 9.；Edited by Philip Schaff.；Buffalo, NY: Christian Literature Publishing Co.,；1889.；<http://www.newadvent.org/fathers/190111.htm>.}

% structure: p0001 source=h1 target=section reason=page-title

\section{讲道集第十二篇：论雕像}

\Emph{为得罪皇帝者所得赦免向天主谢恩。关于创造的有形讲论。证明天主在造人时将自然律铭刻在他心中。应极力避免起誓。}\par

1. 昨日我说\Quote{愿天主受赞美！}今日我仍说同样的话。因为虽然我们所惧怕的祸患已经过去，我们却不应让对它们的记忆消失；不是为悲伤，而是为感恩。因为若这些恐怖的记忆常在我们心中，我们便永不会实际遭遇这样的恐怖。我们何须经历，当记忆充当了警戒的角色？既然天主没有允许我们在临近时被那祸患的洪流吞没，我们就不应允许自己在这些事过后变得漫不经心。那时，当我们忧伤时，祂安慰了我们，现在让我们在喜乐中向祂谢恩。在痛苦中祂抚慰了我们，没有抛弃我们；因此让我们不要在顺境中因懈怠而背叛自己。\Quote{不要忘记，}有人说道，\Quote{在丰年时记住饥荒的日子。}\BibleRef{德 18:25}因此让我们在得救援之日记住受试探之时；对待我们的罪也该如此。你若犯了罪，天主赦免了你的罪，接受宽恕并感恩；但不要忘记罪；不是要你为此苦恼，而是为教训你的灵魂，不要放纵而再次陷入同样的罗网。\par

2. 保禄也是如此；他说了\Quote{他以为我忠信，委派我服役}之后，又加上\Quote{我从前是亵渎者、迫害者和施暴者。}\BibleRef{弟前 1:12}\Quote{让仆人的生活，}他说，\Quote{公开显露，以便显明主人的慈爱。因为我虽已获得罪的赦免，却不拒绝纪念那些罪。}这不仅显明了主的慈爱，也使这人更加光荣。因为当你得知他从前是怎样的人时，你会更加惊叹；当你看见他从什么成为怎样时，你会更加称赞他；若你犯了重罪，一旦改变，你也会从这个例子得到好希望。因为除了所说的，这样的榜样安慰了那些绝望的人，使他们重新站立。同样的事也将发生在我们城市：所发生的一切显示了你们的德行，你们藉着悔改得以避免这样的义怒，同时也显明了天主的慈爱，祂因着一点行为的改变就撤去了那样威胁的阴云，并因此使所有沉于绝望的人重新振作，当他们从我们的案例中得知，那仰望天主帮助的人，即使有无数波浪四面环绕，也不会被淹没。\par

3. 有谁见过、听说过像我们所受的苦难？我们每天预期城市连同居民将被连根拔起。但当魔鬼希望使船沉没时，天主却带来了完全的平静。让我们不要忘记这些恐惧的伟大，以便我们记得从天主所受的恩惠之大。那不知疾病性质的人，不会懂得医生的技艺。让我们也将这些事告诉我们的子女，并传给最遥远的世代，使所有人都知道魔鬼如何曾企图摧毁城市的根基，而天主又如何能够在其倒塌和俯伏时显然地重新扶起它，不允许它遭受丝毫损害，反而除去了恐惧，迅速驱散了它所处的危险。因为在过去的一周里，我们都曾预期财产将被没收，士兵会向我们冲来，我们幻想了上千种恐怖。但看！这一切都已过去，像云彩或飞逝的影子；我们只在可怕的预期中受了惩罚；或者更确切地说，我们没有被惩罚，而是受了管教，变得更好；天主软化了皇帝的心。让我们时常并每天说\Quote{愿天主受赞美！}并以更大的热忱留心于我们的聚会，让我们赶往教会，从这里我们获得了这份恩惠。因为你们知道起初你们逃往何处，你们聚集何处，我们的安全来自何处。让我们抓住这神圣的锚；就像在危险时刻它没有出卖我们，如今在缓解时刻也不要离开它；但让我们精确地等候定期的聚会和祈祷，并每天聆听天主圣言。我们曾忙于奔波打听那些从宫中来的人，为所威胁的祸患而焦虑，所花费的闲暇时间，让我们完全用于聆听天主法律，代替不合时宜和无意义的消遣，免得我们再次将自己置于需要那种忙碌的境地。\par

四、我们已在之前的三天里探求了一种认识天主的方法，并已得出结论；解释了\Quote{诸天述说天主的荣耀}；以及保禄所说的意思，即\BibleQuote{祂那看不见的事，自从创造天地以来，借着所造之物，就可以明明看见、被人领悟}。\BibleRef{罗 1:20}我们也展示了如何从创造世界，以及如何借天、地、海，使造物主受荣耀。但今天，我们在同样主题上稍作沉思之后，将转向另一个话题。因为祂不仅创造了它，而且当它被创造之后，还规定它继续运作；不许它完全静止，也不命令它完全处于运动状态。例如，诸天保持不动，正如先知所说：\BibleQuote{祂铺张诸天如穹苍，展开它如帐幕覆盖大地}。\BibleRef{依 40:42}而另一方面，太阳连同其余星辰，每日运行其轨道。再者，大地是固定的，而水不断流动；不单是水，还有云彩及频繁而至的时雨，按季节回归。云彩的本性是一，但从它们产生的东西却不同。因为雨水在葡萄中变成酒，在橄榄中变成油，在其他植物中化为汁液；大地的母胎是一，却结出不同的果实。阳光的热力也是一，却以不同方式使万物成熟；有的缓慢，有的快速。那么，有谁不会对这些事感到惊异和赞叹呢？\par

五、不仅如此，这还不是唯一的奇事，即祂以如此大的多样性和差异创造了它；而且，祂还把它普遍地呈现在所有人面前：富人和穷人，罪人和义人。正如基督也曾宣告：\BibleQuote{祂使太阳上升，光照恶人，也光照善人；降雨给义人，也给不义的人}。\BibleRef{玛 5:45}此外，当祂用各种动物充满世界，并在受造物中植入各种性情时，祂命令我们效法其中一些，而避免另一些。例如：蚂蚁勤劳，从事辛苦的工作。因此，你若留心，便会从这动物得到最强烈的告诫，不要沉溺于懒惰，也不要逃避劳作和辛苦。为此，圣经也打发懒汉到蚂蚁那里，说：\BibleQuote{懒汉，你去看看蚂蚁，效法它的道路，变得比它更聪明}。\BibleRef{箴 6:6}他的意思是，你不愿意从圣经学习劳动是好的，以及谁若不愿意工作，就不应当吃饭吗？\BibleRef{得后 3:10}那就从无理性的动物学习吧！我们在自己家中也是这样做的：那些年长且被视为优越的人做错了事，我们便叫他们注意有思想的孩童。我们说：\Quote{注意那个比你小的人，他是何等诚恳和警醒。}那么，你也照样从这动物接受关于勤勉的最好的劝诫；并惊叹你的主，不仅因为祂创造了天和太阳，也因为祂创造了蚂蚁。因为这动物虽小，却提供了大量证据，证明天主智慧的伟大。试想，蚂蚁是何等精明，天主在如此小的身体里植入了何等不懈的工作欲望！然而，当你从这动物学习勤勉时，还要从蜜蜂同时学习整洁、勤勉和社会的和谐！因为蜜蜂每日劳作，不仅为自己，更为我们；这确实特别适合基督徒：不求自己的事，只求别人的事。因此，正如它遍历所有草地，为他人预备筵席，你，人啊，也要照样做；如果你积累了财富，就花在别人身上；如果你有教导的才能，不要埋没那塔冷通，而要为那些需要的人公开运用！或者，如果你有任何其他优势，就对那些需要你劳动益处的人成为有用的人！你难道没看到，正因为这个原因，蜜蜂比其他动物更受尊敬吗？不是因为它劳作，而是因为它为他人劳作？因为蜘蛛也劳作、辛苦，在墙上铺开它的细网，胜过最巧的妇女；但这类动物却不受重视，因为它的工作对我们毫无益处；那些劳苦却只为自己的人也是如此！还要效法鸽子的纯朴！效法驴子对主人的爱，以及牛！效法鸟儿无忧无虑！因为，从无理性的动物那里可以获得多么大的益处，来修正品行啊！\par

6. 基督也由这些动物教导我们，祂说：\Quote{你们要像蛇一样机警，像鸽子一样纯洁。}\BibleRef{玛 10:16} 又说：\Quote{你们看天空的飞鸟，它们不播种，也不收割，也不收积在仓里，你们的圣父尚且养活它们。}\BibleRef{玛 6:26} 先知也为了羞辱忘恩的犹太人这样说：\Quote{牛认识自己的主人，驴认识自己主人的槽；但以色列却不认识我。}\BibleRef{依 1:3} 又说：\Quote{斑鸠、燕子和鹤都遵守来去的时令，但我的百姓却不知道主他天主的法令。}\BibleRef{耶 8:7} 从这些动物和类似者，学习修德，并受教藉相反者避免邪恶。因为正如蜜蜂跟随善，毒蛇则具有毁灭性。所以你要躲避邪恶，免得听见这话：\Quote{毒蛇的毒液在他们嘴唇之下。}再次，狗不知羞耻。因此，你要憎恨这类的邪恶。狐狸也狡猾诡诈。不要效法这种恶习；但要像蜜蜂飞过草地时，不选择每一种花，而是选取有用的，留下其余的；你也要如此；当你审视整个无理性的动物族类时，若能从它们中提取任何有益的东西，就接受；它们天性所有的长处，你要凭自己的自由抉择努力实践。因为在这方面你也受到天主的尊荣：祂准许你凭自己的自由抉择获得它们天性所有的长处，为使你也能得到赏报。因为它们中的善行不是出于自由意志和理性，而只出于本性。换句话说，蜜蜂酿蜜，不是因为它通过理性和反思学会，而是受本性教导。因为如果这工作不是天生的，不是这个族类所分得的，那么它们中一定会有不擅长此艺的；然而从世界初创至今，没有人见过蜜蜂停止劳作而不酿蜜。因为这类天性特征是全族共有的。但那些依赖我们自由抉择的事却不是共有的，因为需要努力才能完成。\par

7. 你要取用一切最好的，并用它们装备自己；因为你实为无理性之物的君王；但君王们，若其臣民拥有任何优异之物，无论是金、银、宝石或华丽的服装，他们通常也拥有更多。也要从受造界学习赞美你的主！若你看见的任何事物超出你的理解，你无法找到其理由，仍要为此光荣造物主，因为这些工程的智慧超越你的悟性。不要说：为什么是这样？或：有何目的？因为一切都是有益的，即使我们不知其理由。因此，如同你进入一间手术室，看见许多器械摆在你面前，你因器具的多样性而惊奇，即使不知其用途；对受造界也应如此。虽然你看见许多动物、草木、植物及其他事物，你不知其用途，但要欣赏它们的多样性；并为此惊叹天主的完美工程：祂既没有使所有事物都显明给你，也没有允许所有事物都被你知晓。因为祂没有允许所有事物都被你知晓，免得你说：现有之物并非出于天主的眷顾。祂没有允许所有事物都被你知道，免得你知识的伟大激起你的骄傲。至少那恶者就是这样诱惑了第一个人，并以更大知识的希望剥夺了他已拥有的。因此，有位智者劝诫说：\Quote{不要探究过于艰难的事，也不要追寻过于深奥的事。但所命令你的，你要虔敬地思考；因为祂的大部分作为是隐秘的。}又说：\Quote{所启示给你的，比人所能理解的更多。}但他说这话，是为了安慰那因不知一切而忧伤烦恼的人；因为你所观察到的、被允许知道的事，也远超过你的理解；因为你不能靠自己发现它们，而是蒙天主教导。因此，要满足于所赐你的财富，不要寻求更多；但为你已领受的，要感恩；不要因你未领受的而生气。为你所知的，要光荣；不要因你所不知的而绊倒。因为天主使两者同样有益：祂启示了一些事，却隐藏了另一些，为你的安全着想。\par

8. 认识天主的一种方式，便是我所说通过受造界的认识，这可能需要很多日子。因为为了精确地检查人的形成（我所说的是我们所能达到的精确，并非真正的精确；因为我们已为受造工程提出许多理由，但还有更多不可言传的理由，只有创造它们的天主知道——我们当然不完全知道），为了精确地审视整个人的塑造；并发现每个肢体中的巧妙；检查筋、静脉和动脉的分布与位置，以及每一其他部分的构造，即使一整年也不足以完成这样的探讨。\par

9. 为此，我们在此搁置这一主题；藉着上述之言，已为勤勉好学之人提供了机会，去同样考察其他受造之物；现在我们将把论述转向另一点，这本身也证明了天主的眷顾。那么，这第二点是什么呢？就是：当天主造人时，他从起初就将自然法植入人心。那么，这自然法又是什么呢？他在我们内赋予了良心；并使人对美善之事及其对立之事的认识成为自发的。因为我们无需学习便知奸淫是恶事，贞洁是善事，而是从一开始就知道这一点。为使你明白我们是从一开始就知道这一点，立法者后来颁赐法律时，说：\Quote{不可杀人}，\BibleRef{出 20:13} 并没有加上\Quote{因为杀人是恶事}，而是简单地说：\Quote{不可杀人}；因为他仅仅禁止了罪恶，并未加以教导。那么，当他说\Quote{不可杀人}时，为何没有加上\Quote{因为杀人是邪恶之事}呢？原因是，良心早已教导了这一点；他如此说话，是如同对知晓并理解此事的人说的。因此，当他向我们讲述另一条并非由良心直觉所知的诫命时，他不仅禁止，还附加了理由。例如，当他命令关于安息日的事时：\Quote{第七日你不可工作}；他又为这休息附加了理由。这理由是什么？\Quote{因为第七天天主歇了祂一切创造的工，就安息了。} \BibleRef{出 20:10} 又说：\Quote{因为你曾在埃及地为奴。} \BibleRef{申 21:18} 那么我问，为何他对于安息日附加了理由，而对于杀人却未如此行？因为这条诫命并非首要的诫命之一，并非我们良心中精确界定的诫命之一，而是某种局部的、暂时的诫命；为此，它后来被废除了。但那些必要且维系我们生命的诫命是：\Quote{不可杀人；不可奸淫；不可偷盗。} 因此，他在此情况下未加理由，也未对这事进行任何教导，而仅以简单的禁令为满足。\par

10. 不仅从那里，我还将藉着另一番考量，向你表明人如何在德行认识上是自发的。亚当犯了第一罪；犯罪后他立即藏匿起来；但若他不知道自己在做错事，他为何要藏起来呢？因为那时既无文字，也无法律，更无梅瑟。那么，他从何处认识到罪的？为何要藏起来？他不仅藏起来，而且在受质问时，还试图将过错推给别人，说：\Quote{祢所赐给我的女人，她把树上的果子给了我，我就吃了。} 那女人又将指控转嫁给另一者，即那蛇。也请观察天主的智慧；因为当亚当说：\Quote{我听见祢的声音，就害怕，因为我赤身露体，我便藏起来了，} 天主并未立即定罪他所做的事，也没有说：\Quote{你为何吃了那树上的果子？} 而是如何说的？\Quote{谁告诉你，} 他问道，\Quote{你是赤身露体的，除非你吃了那棵我唯独吩咐你不可吃的树上的果子？} 他既没有沉默，也没有公开定罪他。他没有沉默，为要召唤他前来承认自己的罪行。他没有公开定罪，免得一切来自他自身，而人因此被剥夺了因告解而来的宽恕。因此，他没有公开宣示这种知识从何而来，而是以询问的形式继续对话，让那人自己去承认罪行。\par

11. 再者，在加音和亚伯尔的事上，也可观察到同样的过程。因为首先，他们将各自劳作的果实奉献给天主。我们不仅要从他的罪，也要从他的德行表明，人有能力认识这两者。因此，亚当显明了人知道罪是恶事；亚伯尔又表明了人知道德行是善事。因为他并未从任何人学习，也未听到任何关于初熟之果的法律被颁布，而是由内在、由其良心所教导，献上了那祭品。为此，我不将论证推至较晚的时期，而是聚焦于这些早期人物的时代，那时尚无文字，尚无法律，尚无先知和民长；只有亚当和他的子女存在；好叫你明白，善恶的知识早已植入他们的本性。因为亚伯尔从哪里学到献祭是善事？学习到尊崇天主、凡事感恩是善事？有人回答说：\Quote{那么，加音没有奉献他的供物吗？} 此人也献了祭，但方式不同。而从这里，良心的知识再度显现。因为当他嫉妒那受尊荣的人，并蓄谋杀人时，他隐藏了自己奸诈的企图。他说了什么？\Quote{我们到田间去吧。} 外表是一回事，是爱的伪装；心思是另一回事，是杀兄弟的意图。但他若不知道这图谋是邪恶的，为何要隐藏呢？再者，杀人之后，受天主质问：\Quote{你的兄弟亚伯尔在哪里？} 他回答说：\Quote{我不知道；我岂是看守我兄弟的吗？} 他为何否认罪行？岂不是因为他极其自责吗？因为正如他的父亲藏匿起来，此人也否认自己的罪，在定罪之后又说道：\Quote{我的罪过太大，无法得赦免。}\par

12. 但或许有人会反对说，\OriginalTerm{外邦人}{Gentile}全然不承认这一点。那么，让我们来讨论这个问题；正如我们藉着圣经和理性的帮助，在创造之事上向这些反对者作战一样，如今我们也当在良心之事上如此行。因为\Person{保禄}与这类人争论时，也曾讨论过这个论点。那么，他们所提出的是什么？他们说，在我们\OriginalTerm{良心}{conscience}中并不存在不言自明的法律；天主并没有将这一法律植入我们的本性。但果若如此，我问：立法者为何要制定那些关于婚姻、杀人、遗嘱、信托、彼此不可侵犯以及成千上万其他事项的法律？因为现今活着的人或许是向他们的长辈学习的；而长辈又是向在他们之前的人学习的，这些人又是向更早的人学习的。但那些最初的立法者和最初制定法律的人，他们是从谁那里学来的呢？难道不是从良心学来的吗？因为他们不能说，他们曾与\Person{梅瑟}交往，或曾听过先知的话。他们既是外邦人，又怎能如此呢？显然，正是因为天主在起初造人时放在人里面的那同一法律，才有了法律的制定、技艺的发现以及一切其他事物。因为技艺也是这样建立起来的，它们的创始人是以自学的方式认识它们的。\par

13. 同样，法庭的出现、刑罚的界定，正如\Person{保禄}所观察的。因为许多外邦人准备反驳这一点，并说：\Quote{天主将如何审判在\Person{梅瑟}之前生活的人类？祂没有派遣立法者；祂没有颁赐法律；祂没有派遣先知、宗徒或传福音者；那么祂怎能追究这些人的责任呢？} 既然\Person{保禄}想证明他们拥有一个自学的法律，并且他们清楚知道该做什么，请听他如何说：\BibleQuote{因为没有法律的外邦人，如果顺着本性去行法律所规定的事，那么他们虽然没有法律，但对自己就是法律；如此他们表明了法律的作用刻在他们心中。} \BibleRef{罗 2:14} 但没有文字，如何呢？\BibleQuote{良心也作证，并且他们的思想互相控告或加以辩护；在天主藉耶稣基督审判人隐秘事的那一天，按照我的福音，也是如此。} \BibleRef{罗 2:16} 又说：\BibleQuote{凡没有法律而犯了罪的，也将没有法律而丧亡；凡在法律之下犯了罪的，将依据法律受审判。} \BibleRef{罗 2:12} 所谓\Quote{没有法律而丧亡}是什么意思？法律不控告他们，而是他们的思想和良心控告他们；因为如果他们没有良心的法律，就不必因行恶而丧亡。如果他们是在没有法律的情况下犯罪，又怎能如此呢？但当他说\Quote{没有法律}时，他并不是说他们没有法律，而是说他们没有成文的法律，虽然他们有自然法律。又说：\BibleQuote{然而光荣、尊贵与平安，加给一切行善的人，首先是犹太人，然后也是外邦人。} \BibleRef{罗 2:10}\par

14. 但这些话是他指着基督来临之前的早期时代说的；他这里所称的外邦人，不是拜偶像的人，而是事奉唯一天主的人；不受遵守犹太规条（即安息日、割损和种种取洁礼）的必要约束，却表现出各种智慧和虔诚。再者，论及这样的敬拜者，他说道：\BibleQuote{患难和困苦、忿怒和恼恨，加给一切作恶的人，首先是犹太人，然后也是外邦人。} \BibleRef{罗 2:9} 他在这里又以\OriginalTerm{外邦人}{Greek}一名来称呼那不受犹太习俗约束的人。那么，如果他既未听过法律，也未与犹太人交往，怎能因行恶而有忿怒、恼恨和困苦加于他呢？原因是，他拥有一个内在的良心，劝诫他、教导他、训诲他一切事物。这从何处显明呢？从他处罚他人作恶的方式，从他制定法律的方式，从他设立公义法庭的方式。为使这一点更加清楚，\Person{保禄}提到了那些生活在邪恶中的人。\BibleQuote{他们虽然知道天主的法令，行这样事的人是当死的，然而他们不但自己去行，而且还赞同作恶的人。} \BibleRef{罗 1:32} \Quote{但那是从何而来的？}有人说，\Quote{他们从何处知道，是天主的旨意，那些生活在邪恶中的人应被处死？} 从何处？难道不是从他们审判犯罪之人的方式吗？因为如果你不认为杀人是邪恶的，那么当你将一个杀人犯带到你的法庭上，你就不该惩罚他。同样，如果你不认为犯奸淫是邪恶的，那么当犯奸淫者落入你手中时，就释放他，不予惩罚！但如果你制定法律，规定刑罚，对他人之罪严厉审判，那么当你自己行恶时，你能以不知何事当做为借口为自己辩护吗？假设你和另一个人同样犯了奸淫。你有什么理由惩罚他，而认为自己应得宽恕？因为如果你不知奸淫为邪恶，就不该惩罚他人。但如果你惩罚他，却想自己逃脱惩罚，同样罪行不受同样处罚，这怎能合理呢？\par

15. 这正是保禄所斥责的，当他说：\BibleQuote{人哪，你判断行这样事的人，自己却行同样的事，你以为你能逃脱天主的审判吗？}\BibleRef{罗 2:3} 这是不可能的，绝对不可能；因为，他意思是说，从你加于别人的判决，天主也将依此判决审判你。难道你是公义的，而天主不公义吗？如果你不宽容别人受的冤屈，天主怎能宽容你呢？如果你纠正别人的罪，天主怎不纠正你呢？即使祂不立刻惩罚你，也不要因此自信，反而要更加惧怕。正如保禄所吩咐的：\BibleQuote{你藐视祂丰富的恩慈、宽容与忍耐，不晓得天主的恩慈是领你悔改吗？}\BibleRef{罗 2:4} 因此，他说，祂容忍你，不是要你变得更坏，而是要你悔改。但你若不悔改，这忍耐便成为你更重惩罚的原因，因为你继续硬着心不悔改。这正是他所说的：\BibleQuote{你竟硬着心，不肯悔改，为自己积蓄忿怒，直到忿怒的日子，就是天主公义审判显明的时候。祂要照各人的行为报应各人。}\BibleRef{罗 2:5} 既然祂照各人的行为报应各人，祂便将自然律放在我们心里，后来又赐予成文法，为要我们为罪恶交账，并奖赏行善的人。因此，我们当极其谨慎地生活，如同即将面对那可畏的审判台；知道我们若在自然律和成文法之外，又受了那么多教诲和不断的劝诫，却忽略自己的救恩，就必不得宽赦。\par

16. 我愿再对你们谈论关于起誓的事，但我感到羞愧。因为对你们日日夜夜重复同样的话，我并不厌倦。但我担心，若我连续多日跟踪你们，恐怕会判定你们过于懈怠，竟在如此容易的事上需要不断的劝诫。我不但羞愧，也为你们担忧！因为对留心听的人而言，频繁的教导是有益且有益的；但对懈怠的人，却是有害且极其危险的；因为人听得越多，若不行所听的，就给自己招来更大的惩罚。天主曾以此责斥犹太人，说：\BibleQuote{我差遣我的先知，清早起来差遣他们，你们仍不听。}\BibleRef{耶 29:9} 我们这样做，是出于对你们的大关怀。但我们担心，在那可畏的日子，这劝诫和忠告会起来攻击你们。因为当任务容易达成，而负有不断劝诫之责的人又不懈怠时，我们还有什么可辩护的？什么理由能救我们免受惩罚？请告诉我，若有人欠你一笔钱，你每次遇见债主，岂不提醒他那笔借款吗？你也当这样做；让每个人都以为邻人欠他钱，即遵守这条诫命，遇见他时，就提醒他还债，知道若我们忽略弟兄，就有不小的危险临到我们。因此，我也不停地提起这些事。因为我怕在那日听见这话：\BibleQuote{又恶又懒的仆人，你应当把我的银子放给兑换银钱的人。}\BibleRef{玛 25:26} 看，我不止一次、两次，而是多次讲明了。剩下的就是你们去偿还利息。听道的利息，就是借行为表现出来；因为存款是主的。让我们不要疏忽地接受所托付的，而要殷勤地保存，使我们在那日能连本带利地归还。除非你带别人也行同样的善工，否则你将听见那埋藏银子的人所听见的声音。但愿这声音不临到你！但愿你们听见基督对那赚了利润的人所说的不同的声音：\BibleQuote{好，善良忠信的仆人，你在小事上忠信，我要派你管理许多大事。}\BibleRef{玛 25:21}\par

17. 我们若表现出与他同样的热忱，就必听见这声音。我们若照我所说的去做，就必表现出这热忱。你们离去时，趁所听的还在心里火热，彼此劝勉！正如你们分别时互相问候，愿每个人从这里带着一句劝诫离去，对他的邻人说：\Quote{留心记住要遵守诫命。}这样我们必能得胜。因为当朋友也以这样的劝诫送别时，回家后妻子又提醒同一件事，我们独处时这话也抓住我们，我们就必很快摆脱这恶习。我知道你们惊讶我为何对此诫命如此热衷。但请你们履行所吩咐的，然后我就告诉你们。同时，我要说：这条诫命是神圣的法律，违背它并不安全。但若我见你们正确遵行，我将讲另一个理由，其重要性不亚于此，使你们知道，我如此重视这条法律是公义的。现在是时候以祈祷来结束这篇讲道了。因此，让我们共同说：\Quote{天主，祢不愿罪人死去，而愿他悔改得生；求祢赐给我们，在遵守了这条及所有其他诫命之后，得以堪当站立在祢的基督的审判台前，好使我们因着极大的胆量，获得祢荣耀的国度。因为光荣归于祢，及祢的独生子，及圣神，从今世直到永远。}阿们。\par

\SourceRecord{Translated by W.R.W. Stephens.；Nicene and Post-Nicene Fathers, First Series；Vol. 9.；Edited by Philip Schaff.；Buffalo, NY: Christian Literature Publishing Co.,；1889.；<http://www.newadvent.org/fathers/190112.htm>.}

% structure: p0001 source=h1 target=section reason=page-title

\section{\WorkTitle{论雕像讲道集 第十三篇}}

\Emph{对天主使晚期悲惨局势转变的进一步感恩。对被拖走并因叛乱受罚之人的回忆。关于人被创造并领受自然律的论述。关于完全戒绝发誓的完成。}\par

1. 以我昨天和前天开始的同一绪言和前奏，我今天也要开始。现在我又要说：\Quote{愿天主受赞美！}上周三我们见到了何等的日子！如今又是何等景象！那天阴霾何等沉重！今日的平静何等明亮！那天，那可怕的法庭设在城中，震撼了所有人的心，使白日无异于黑夜；这不是因为太阳的光芒熄灭了，而是因为沮丧和恐惧蒙蔽了你们的眼睛。因此，为使我们获得更多的喜乐，我愿讲述当时发生的几件事；因为我觉得叙述这些事对你们和后来的人都有益处。此外，对于脱离船难的人，当他们进入港口后，回忆波涛、风暴和狂风是甘甜的。对于患病的人，当疾病痊愈后，与他人谈论差点将他们带入坟墓的高烧也是一件愉快的事。当恐惧过去后，讲述那些恐惧有乐趣；灵魂不再害怕它们，反而从中获得更多的欢愉。对过去灾祸的回忆总能使眼前的幸福显得更加突出。\par

2. 当时，城中大部分人因那场恐惧和危险而躲入隐秘处、荒野和洞穴中；恐怖从四面八方包围他们；房屋空无妇女，广场空无男子，难得见两三人一起走过广场，即便有，也像行尸走肉般游荡。在这时，我前往法庭，为要看这些事情的结局；在那里，看见城中残存的人群聚集，我最惊奇的是：虽然门口有众多人，却有最深沉的寂静，仿佛没有人在那里，所有人都面面相觑；无人敢询问邻人，也不敢从他那里听到什么；因为每个人都以怀疑的目光看待邻人；因为许多人已经出乎意料地从广场中央被拖走，现在被关在里面。于是，我们所有人都同样仰望天空，静默地伸出双手，期待从上面来的帮助，恳求天主扶持那些受审的人，软化法官的心，使他们的判决仁慈。正如一些人在陆地上看见他人遭遇船难，却无法靠近他们，伸手援救，被波浪阻隔；然而在岸上，他们伸手流泪，哀求天主救助溺水者；同样，在那里，所有人都静默无声地以心神呼求天主，为法庭上的那些人，如同为被波浪包围的人代祷，求他伸出援手，不让船只覆没，也不让受审者的判决最终彻底沉沦。门外的情形如此；但当我进入法庭内，我看到了更可怕的景象：士兵佩带刀剑和棍棒，严格为里面的法官维持秩序。因为所有受审者的亲属，无论是妻子、母亲、女儿还是父亲，都站在审判座的门前；为防万一有人被带去处决时，有人看见惨状而激动，引起骚乱或动乱，士兵们将他们全部驱赶到远处；以此先以恐惧占据他们的心。\par

3. 有一个景象比一切都更可怜：某位正在法庭内受审之人的母亲和姐妹，坐在法庭门廊前，在石板地上打滚，成为所有旁观者的公共景象；她们蒙着脸，毫不顾惜羞耻，只出于灾祸的急迫才如此。没有婢女、邻居、女友或其他亲属陪伴她们。她们被一群士兵围住，孤身一人，衣衫褴褛，匍匐在地，就在门口处，比那些在里面受审的人更加可怜。她们听到行刑者的声音、鞭子的抽打声、受鞭打者的哀嚎、法官的可怕威胁，每次鞭打，她们自己所受的痛苦比那些被鞭打的人更尖锐。因为，当听到有人被鞭打以招出有罪者，并发出哭喊时，她们仰望上天，恳求天主赐予受苦者一些忍耐的力量，以免自己亲属的安全因他人的软弱而被出卖，因为那些人无法承受鞭打的剧痛。同样的情况也发生在与风暴搏斗的人身上。正如他们远远看见一个巨浪抬头，逐渐增大，准备吞没船只，在浪头靠近船之前，他们就几乎吓得半死；这些妇女也是如此。每当她们听到声音和哭喊声传来，眼前就看见千百种死亡，恐惧不已，担心那些被逼作证的人，因忍受不住折磨，会说出某个自己亲属的名字。因此，人们看到内外都有痛苦。行刑者在里面折磨人；这些妇女则被自然的力量和情感的同情所折磨。里面有哀号，外面也有！里面是被判有罪之人，外面是他们的亲属。不仅如此，连法官们自己也在内心哀叹，比任何人都更痛苦，因为他们被迫参与如此悲惨的场面。\par

4. 至于我，当我坐着看到这一切：主妇和贞女们，惯于深居简出，如今成了众人围观的景象；那些习惯躺在软榻上的人，如今以地板为床；那些享受婢女、阉仆及各种外在尊荣的人，如今被剥夺了一切；她们匍匐在每个人的脚下，恳求他尽己所能帮助受审的人，希望从所有人那里得到普遍的慈悲援助；我引用\Person{撒罗满}的话喊道：\BibleQuote{虚而又虚，万事皆虚。}\BibleRef{训 1:2}因为我看到这预言和另一预言在每件事上都应验了，它说：\BibleQuote{人的一切光荣有如草上的花；草必枯萎，花必凋谢。}\BibleRef{依 42:6}因为那时，财富、高贵、名望、朋友的庇护、亲属关系以及一切世俗事物，都显得毫无价值；所发生的罪和违法使这一切援助都逃跑了。正如幼鸟被带走后，母鸟回来发现巢空，无法救回被囚的幼雏，只能在捕鸟者手边盘旋，以此显示它的悲伤；同样，这些妇女也是如此，当她们的子女从家中被夺走，关在里面，如同在网或陷阱中。她们不能进去解救囚徒，只能通过在门口附近的地上打滚、哀号和呻吟，以及尽力靠近那些逮住他们的人，来表现她们的痛苦。看到这一切，我心中想起那可怕的审判台，心里说：如果现在，当人作审判官时，母亲、姐妹、父亲或其他任何人，即使对所犯的罪行毫不知情，也不能拯救罪犯；那么，当我们站在基督可畏的审判台前受审时，谁能为我们辩护？谁敢出声？谁能拯救那些将被带往不可忍受的惩罚中的人呢？尽管他们是当时受审的城中首要人物，贵族中的首领，但他们情愿失去一切财产，甚至如果需要，连自由也失去，只要能继续享受这现世的生命。\par

5. 但继续讲。白日将近，傍晚来临，法庭的最终判决即将宣布，所有人都更加痛苦，祈求天主赐予一些拖延和喘息，并感动法官们将已调查的事实呈报皇帝裁决，因为这样或许能产生一些益处。此外，民众向慈悲的天主发出普遍的恳求，祈求祂拯救这城的剩余部分，不让它完全从根基上被摧毁。没有一个人不是流着泪加入这呼声。然而，这些事一点也没有打动里面的法官，虽然他们听到了。他们只考虑一件事，就是要对所犯的罪行进行严格的追查。\par

6. 最后，他们给犯人戴上锁链，用铁链捆绑，穿过广场中央，将他们送往监狱。那些曾拥有马群、担任过竞技会主席、能列举上千种显赫职位的人，财产被没收，可以看见在所有门上贴了封条。他们的妻子也被逐出娘家，每个人都实实在在地扮演了\Person{约伯}妻子的角色。因为他们\Quote{从这家流落到那家，从这处流落到那处，寻找栖身之所}。要找到这样的地方并不容易，因为人人都战战兢兢，不敢收留或以任何方式帮助那些被控告者的亲属。尽管如此，虽然发生了这样的事，受苦者在这一切中仍保持忍耐，因为他们没有被剥夺现世生命。无论是财富的损失、羞辱、如此公开的暴露，还是任何其他类似的事，都没有使他们烦恼。因为灾祸的巨大，以及他们遭受这些时曾预料更坏的情况，这一切已使灵魂准备好操练智慧的刚毅。如今他们学到，德行对我们来说是多么简单的事，实行起来多么容易和迅速，只是由于我们的疏忽，它才显得艰难。那些在此之前无法以温顺之心忍受一点金钱损失的人，如今虽已失去全部家产，却因更大的恐惧而感觉如同找到了宝藏，因为他们没有失去生命。因此，如果我们对未来的地狱有所意识，并想到那些难以忍受的惩罚，即使为了天主的法律我们须舍弃家产、身体甚至生命，我们也不会悲伤，因为我们知道将要获得更大的东西，即从未来的恐怖中获得解脱。\par

7. 或许我所讲述的这一切悲剧，已大大软化了你们的心。但请不要见怪。因为我即将尝试一些更微妙的思想，需要你们有更敏锐的心境，我故意这样做，为的是经由描述的恐怖，使你们的心灵摆脱一切懈怠，从一切世俗忧虑中退出来，并以更大的准备，将即将言说之事的效力，传达至你们灵魂的深处。\par

确实，我们近来的论述已向你们充分证明，一个关于善恶的自然法就安放在我们内。但为使我们的证明更加清晰，我们今天将再次竭力专注于同一论题。因为天主从起初造人时，就使人有能力辨别这两者，这是所有人都显明的。因此，当我们犯罪时，在属下面前都会羞愧；常常一个主人去娼妓家的路上，如果这时看到自己一个更受尊敬的仆人，就会羞红着脸从这条不义之路折回。同样，当别人指责我们，把特定恶习的名称加给我们时，我们称之为侮辱；如果我们受了委屈，就把那些做错事的人拖到公共法庭。这样我们就能明白什么是恶，什么是善。为此，基督为了宣告这一点，并表明祂不是引入一条陌生的法律，或一条超越我们本性的法律，而是祂从起初就预先放在我们良心里的法律，在宣布了那许多真福之后，这样说：\Quote{凡你们愿意人给你们做的，你们也要照样给人做。}\BibleRef{玛 7:12}\Quote{许多话}，祂说，是不必要的，也无需冗长的法律，或多样化的教训。让你的意愿成为法律。你希望得到仁慈吗？就对别人仁慈。你希望得到怜悯吗？就怜悯你的近人。你希望受到称赞吗？就称赞别人。你希望被爱吗？就实践爱德。你希望享有首位吗？就先把这个位置让给别人。你自己成为你自己生活的审判者，你自己成为立法者。又说：\Quote{你不愿别人对你做的，也不要对别人做。}\BibleRef{多 4:16}藉着后一条诫命，祂要引人远离邪恶；藉着前一条，引人实践德行。\Quote{你不愿别人对你做的，}祂说，\Quote{也不要对别人做。}你憎恨受侮辱吗？就不要侮辱别人。你憎恨被嫉妒吗？就不要嫉妒别人。你憎恨被欺骗吗？就不要欺骗别人。总之，在所有事情上，如果我们坚守这两条诫命，就不需要任何别的教训。因为对德行的知识，祂已植入我们的本性；但德行的实践与修正，祂则托付给了我们的道德选择。\par

8. 或许所说的这些有些隐晦；因此我将再次努力使之更加清晰。为了知道实践节制是一件好事，我们不需要言语，也不需要教导；因为我们自己就在本性中拥有这种知识，无需劳力或疲乏地四处探询节制是否美好和有益；但我们所有人都一致承认这一点，没有人对这项德行存有疑问。同样，我们也认为奸淫是一件恶事，这里也不需要任何麻烦或学习去认识这罪的邪恶；但我们全都在这类判断上是自学的；我们称赞德行，尽管并不追随；另一方面，我们憎恨恶行，尽管我们实践它。这是天主极其美好的作为：祂使我们的良心和选择能力，在行动之前就已经与德行认同，并与邪恶为敌。\par

9. 正如我当时所说，这些事的知识都存在于所有人的良心中，我们无需老师来教导我们这些事；但行为的规范则交由我们的选择、热忱和努力。为何如此？因为如果祂使一切都出于本性，我们就会空手而去，没有冠冕，也没有奖赏；如同禽兽一样，它们因天性拥有的优势既得不到奖赏也得不到赞美，我们也不会享受这些事；因为天赋的优势不是拥有者的赞美和称誉，而是赐予者的赞美。为此，祂没有把一切都交给本性；另一方面，祂也没有让我们的意志独自承担知识和正确规范的全部重担，免得它在德行的劳作中绝望。良心提示它该做的事，而它则贡献自己的努力来完成。节制是好的，这一点我们都毫不困难地理解，因为知识是本性的；但若不克制私欲，不付出很多努力，我们便无法实践节制的规条；因为这一点不像知识那样由本性而来，还需要一个甘愿的心志和热忱。不但如此，祂还在另一方面减轻了我们的负担，就是让某些善的倾向也自然存在于我们内。我们所有人都自然地对那些受轻慢的人感到愤慨（因此我们成为傲慢者的敌人，即使我们自己没有遭受任何委屈），并且同情那些得到帮助和保护者的喜乐；我们被别人的灾难所打动，也被相互的温情所触动。因为尽管灾难事件可能引起某种怯懦，我们仍然对彼此怀有普遍的喜爱。一位智者意味深长地说：\Quote{一切动物都爱自己的同类，人也爱自己的近人。}\BibleRef{德 13:19}\par

10. 但是，天主除了良心之外，还为我们预备了许多其他的导师：即父亲教导儿女，主人教导仆人，丈夫教导妻子，教师教导学生，立法者和法官教导被治理者，朋友教导朋友。而且，我们常常从敌人那里获得的益处并不少于从朋友那里获得的；因为当敌人指责我们的过失时，即使我们不愿意，也被激起去改正它们。祂为我们设立了如此多的导师，目的是使发现有益之事和规范我们的行为变得容易，那些敦促我们向善的众多事物不容许我们偏离有益于我们的事。因为即使我们藐视父母，但当我们畏惧长官时，我们无论如何也会比不畏惧时更顺从。即使我们在犯罪时可以无视他们，却永远逃不过良心的责备；如果我们羞辱并拒绝良心，那么因着畏惧众人的意见，我们也会因此获益。即使我们对这一点毫无羞耻，法律的恐惧也会压迫我们，即使不情愿，也会约束我们。\par

11. 同样，父亲和教师管教年轻人，使他们步入正轨；立法者和法官管教成年人。仆人们更容易懈怠，除了上述之外，还有主人约束他们节制；妻子有丈夫管教。许多墙壁从四面八方环绕我们的人类，免得它轻易滑落，陷入邪恶。除此之外，疾病和灾难也在教导我们。因为贫穷约束我们，损失使我们清醒，危险使我们谦卑，还有许多类似的事。难道父亲、教师、王子、立法者、法官都不使你畏惧吗？难道没有朋友使你羞愧，没有敌人刺痛你吗？难道没有主人责罚你吗？难道没有丈夫教导你吗？难道没有良心纠正你吗？然而，当身体患病时，常常会纠正一切；一场损失使狂妄的人变得温和。更重要的是，严重的灾祸，不仅临到我们自己，也临到他人，时常对我们大有裨益；而我们自己虽未遭受任何痛苦，但看到别人受罚，也和他们一样变得清醒。\par

12. 至于善行，任何人都可以看到同样的事发生；因为正如坏人受罚时其他人变得更好，同样，每当善人做了一些正确的事，许多人就被激励起同样的热忱：这在戒除发誓方面也发生过。因为许多人看到别人已经戒除了发誓的恶习，就以他们的勤勉为榜样，战胜了这罪；因此我们更愿意再次提起这劝诫的主题。因为不要有人对我说\Quote{许多人}已经做到了；我所希望的并非如此，而是\Quote{所有人}都这样做；直到我看见这事，我才得以喘息。那位牧者有一百只羊，当其中一只走失时，他不顾那九十九只的安全，直到找到那只迷失的，并把它带回羊群。\BibleRef{玛 18:12} 你们难道没有看见，在身体上也是如此吗？如果我们撞到障碍物，只弯回了一根指甲，整个身体都会与那肢体一同感受。不要说只有少数人失败了；但要考虑这一点：这些少数人若不改过，将会败坏许多其他人。虽然在格林多人中只有一人犯了奸淫，但保禄却哀叹，好像全城都丧失了。这很合理，因为他知道，如果那肢体不受责罚，疾病逐渐蔓延，最终会攻击所有其他肢体。不久前，我在法庭上看见那些显贵被捆绑着押解过广场；有人对这种极不寻常的屈辱感到惊奇，另有人说这没什么可惊奇的；因为，当有叛国罪时，地位便无足轻重。那么，当有不敬之罪时，地位岂不更是无足轻重吗？\par

13. 默想这些事，让我们警醒自己；因为若你不付出自己的努力，我们这边的一切都是徒劳。为什么如此？因为教导的职务不像其他技艺。银匠制造了某种器皿后，将它搁置一旁，次日会发现它与他离开时一样。铜匠、石匠和所有其他工匠，都会再次接手自己留下的作品，无论是什么，都保持原样。但我们并非如此，全然相反；因为我们锻造的不是无生命的器皿，而是有理智的灵魂。因此，我们发现你们并非我们离开时的样子；当我们接纳你们，以多种劳苦塑造、改造你们，并在你们离开此地时增加了你们的热情，事务的紧迫从四面围困你们，再次使你们偏离，给我们带来更大的困难。所以，我恳求和祈求你们也亲自参与这项工作；当你们离开这里时，要对自己得救表现出同样热切的关心，正如我在这里为你们的改正所表现的那样。\par

14. 但愿我能替你们行善，而你们能领受这些善工的赏报！那么我就不会这样麻烦你们了。但我怎能做到呢？这是不可能的事；因为祂要按各人的行为报答各人。因此，正如一位母亲看见儿子发烧，目睹他因窒息和发炎而痛苦，常常为他哀哭，对他说：\Quote{我儿，但愿我能承担你的热病，把它的火焰引到我身上！} 现在我也说：但愿我能替你们劳苦，为你们所有人行善！但这不可能。每个人都必须为自己的行为交账，一个人不能看见另一个人代他受罚。为此，我痛苦哀伤，因为在那一日，当你们被传唤受审时，我将无法帮助你们；因为说实话，我在天主面前没有这样讲话的自信。但即使我有许多自信，我也不比梅瑟更圣洁，或比撒慕尔更正义；经上论到他们说，虽然他们达到了如此伟大的德性，却丝毫不能帮助犹太人，因为那百姓过于疏忽自己。\BibleRef{耶 15:1} 既然我们将因自己的行为受罚或得救；因此，我恳求你们，连同其他所有诫命，努力实践这一条；这样，我们最终怀着美好的希望离世，就能获得所应许的美善事物，藉着我们的主耶稣基督的恩宠和慈爱，祂和圣父及圣神，永生永王，世世无尽。阿们。\par

\SourceRecord{Translated by W.R.W. Stephens.；Nicene and Post-Nicene Fathers, First Series；Vol. 9.；Edited by Philip Schaff.；Buffalo, NY: Christian Literature Publishing Co.,；1889.；<http://www.newadvent.org/fathers/190113.htm>.}

% structure: p0001 source=h1 target=section reason=page-title

\section{讲道集第十四篇 论雕像}

\Emph{全体民众脱离一切忧苦，并得安全之后，又有人捏造虚假报告，再次扰乱城市，且被定罪。因此，本篇讲道涉及此事，并论及关于誓言的告诫；故此也引用了\Person{约纳堂}和\Person{撒乌耳}的历史，以及\Person{耶弗他}的历史；并显示从一次誓言产生了多少伪证。}\par

1. 昨日魔鬼大大扰乱了我们的城市；但天主也大大安慰了我们；以致我们每人适切地引用那先知的话说：\Quote{我心里的愁苦众多，祢的安慰使我的灵魂舒畅。}祂不仅安慰我们，甚至允许我们受苦，也显示了祂对我们的慈爱。因为今天我将重复我从未停止说过的话：我们脱离邪恶和邪恶的容许，都出于天主的仁爱。因为祂看见我们陷入懈怠，远离与祂的共融，不关心属灵之事，就暂时舍弃我们；好使我们清醒过来，更热切地归向祂。何况我们如此懈怠，祂这样待我们，有什么可惊奇的呢？因为连保禄也声明，他自己和他的门徒的试炼也是如此之因由。他写给格林多人第二封书信时说：\Quote{弟兄们，我们不愿你们不知道我们在亚细亚所受的苦难，我们被压得太重，力不能胜，甚至连活命的指望都绝了；但我们自己认为必死无疑。\BibleRef{格后 1:8}}仿佛他说：\Quote{如此巨大的危险压在我们身上，我们自认失落，不再指望任何转机，而完全期待死亡。}因为\Quote{我们自己认为必死无疑}这句话就是此意。然而，在如此绝望之后，天主驱散了风暴，拨开乌云，从死亡之门抢救了我们。后来，为表明允许他陷入此危险也是出于对他的极大关怀，他提到试炼带来的益处，即他得以不断仰望祂，既不傲慢，也不自信。因此说了\Quote{我们自己认为必死无疑。\BibleRef{格后 1:9}}他又加了原因：\Quote{叫我们不依靠自己，只依靠使死人复活的天主。}因为试炼的本性就是当我们打盹或跌倒时唤醒我们，激励我们，使我们更虔诚。因此，亲爱的！当你看到试炼时而熄灭，时而复燃，不要沮丧！不要绝望，而要持守有利的希望，这样对自己说：天主把我们交在敌人手中，不是因祂恨我们或抛弃我们，而是因为祂渴望使我们更热切、更与祂亲密。\par

2. 因此，我们不要沮丧；不要对好转绝望；而要盼望很快会有平静；总之，将困扰我们的一切事端的结局交托给天主，让我们再次处理惯常的要点；再次提出我们通常的教导主题。因为我渴望进一步与你们谈论同一主题，以便彻底从你们的灵魂中根除发誓的恶习。因此，我必须再次采用我之前的同样恳求。因为我最近恳求你们，每人拿起若翰刚被砍下的头，还滴着温热的鲜血，就这样回家，想象它就在你们眼前，而它发出声音说：\Quote{憎恨我的凶手——誓言！}责斥未能成就的事，誓言却成就了；暴君的愤怒所不能的事，遵守誓言的必要性却带来了！当暴君在众人面前受到公开责斥时，他高尚地承受了谴责；但当他陷入由誓言导致的致命必要性时，他便砍下了那有福的头。因此，我恳求同一件事，并不停恳求：无论我们到哪里，都要携带这颗头；并把它展示给所有人，如同它大声呼喊，谴责誓言。因为即使我们从未如此懈怠和疏忽，但只要看见那颗头的眼睛可怕地怒视我们，在我们发誓时威胁我们，我们就会被这恐惧比任何缰绳更有效地约束；并能轻易地约束舌头，使其避免发誓的倾向。\par

3. 誓言不仅有此大恶——它惩罚那些犯此罪的人，无论他们违背誓言还是遵守誓言（我们未见其他罪有这种现象）；而且还有另一个同样大的恶伴随着它。那是什么呢？即常常即使那些渴望并决心遵守誓言的人，也根本无法遵守。因为，首先，不断发誓的人，无论自愿或不自愿，有意或无意，玩笑或认真，常常被愤怒和其他许多事情所驱使，必然成为背誓者。没有人能反驳这一点；事实如此明显且普遍认可：经常发誓的人必定也是背誓者。其次，我断言，即使他不被情感驱使，并非不情愿和不知不觉地成为背誓的牺牲品，单凭事情本身的特质，他也必然有意识且自愿地背誓。例如，我们常在家吃饭时，若有一个仆人做错了事，妻子发誓要鞭打他，而丈夫则发誓相反，反对并禁止。在这种情况下，无论他们怎么做，背誓总是结果；因为他们愿意并努力遵守誓言，却不再可能；但无论发生什么，其中一人（或两人）必陷入背誓的网罗。\par

4. 我要解释这是怎样的矛盾。那曾发誓要鞭打男仆或女仆的人，事后又被禁止这样做，他便发了虚誓，因为他没有行他所发誓要行的事；并且，他也使那禁止并阻止誓言履行的人陷于虚誓的罪中。因为不仅那些发假誓的人，而且那些将这种必要性强加于他人的人，都同样难逃指控。不但在家庭中，而且在市场上，我们也可以看到这种情况；尤其是在打斗时，那些互相拳击的人发誓要做相反的事。一个发誓要打，另一个发誓不被打。一个发誓要夺走斗篷，另一个发誓不容此事发生。一个发誓要追索金钱，另一个发誓不支付。还有许多其他互相矛盾的事情，那些好争辩的人发誓去做。同样在店铺里、在学堂中，通常也可观察到同样的事发生。于是工匠常发誓，除非学徒完成所有指定工作，否则不许他吃喝；保姆也常对少年如此行；主妇也对女仆如此行。而当夜幕降临，工作仍未完成时，要么那些未完成任务的人必然挨饿，要么那些发过誓的人必然完全背弃誓言。因为那恶毒的魔鬼，时常潜伏暗算我们的福乐，他临在并听见誓言的约束，便使那些负有责任的人漠不关心，或制造其他困难，以致工作未完成，殴打、辱骂、虚誓和成千上万的其它祸患便可能发生。正如当孩子们竭尽全力朝相反方向拉扯一根长而腐烂的绳子时，如果绳子从中间断裂，他们全都仰面跌倒，有的撞到头，有的撞到身体其他部位；同样，那些各自发誓要做相反之事的人，当誓言因事势所迫而被打破时，双方都落入同一虚誓的深渊：一方是实际发了虚誓，另一方则是为他人提供了发虚誓的机会。\par

5. 为使这事不仅从每日在私宅和公共场所发生的事，而且从\WorkTitle{圣经}本身得以显明，我要向你们讲述一段古代史，与刚才所说的相关。有一次，犹太人遭到敌人入侵，\Person{约纳堂}（他是\Person{撒乌耳}的儿子）已杀死了一些人，并将余众击溃；他的父亲\Person{撒乌耳}渴望更有效地鼓舞军队去对付残敌，并且为了使他们不停止追击，直到他征服所有敌人，竟做了一件与他所愿完全相反的事，即发誓：任何人不得在晚上以前、在向敌人复仇以前吃任何食物。请问，还有比这更愚昧的吗？因为正当他该使疲惫困乏的士兵休整，并以重新振作的力量派去对付敌人时，他却以誓言的约束使他们忍受更甚于敌人的折磨——将他们交给过度饥饿。的确，任何人就与自己有关的事发誓都是危险的，因为环境紧迫会迫使我们做许多事；但以自己的誓言约束他人的决定则更危险，尤其是当一个人发誓不是关乎一个、两个、三个，而是无数的大众时——撒乌耳当时就轻率地这样做了，没有想到在如此众多的人中，至少可能有人会违背誓言；或者，兵士，尤其是出征的兵士，远未具备道德智慧，根本不懂克制肚腹；更何况他们的劳苦极大。然而，他忽略了这一切，仿佛只是对一名容易管束的仆人发誓一般，同样对待他的全军。结果是，他为魔鬼打开了一道门，以致短时间内，从这个誓言中编织出不是两三个、四个，而是多得多的虚誓。因为正如当我们完全不起誓时，便关闭了魔鬼的全部入口；若我们只发出一个誓言，便给了他极大的自由，去构造无尽的虚誓。正如那些绕线团的人，若有人握着线头，便能精巧地绕完整根线；但若无人这样做，他们甚至无法开始；同样，魔鬼要绕我们罪愆的线团时，若不能从我们的舌头得到开端，便无法着手这工作；但我们只要一开端，如同用手握住誓言在舌头上，他便完全自由地在后续工作中施展其恶毒伎俩，从一个誓言构建并编织出成千上万的虚誓。\par

6. 这正是他现在在\Person{撒乌耳}身上所做的。然而，请注意，这誓言立刻被设下了何等样的陷阱：\Quote{军队穿过一座树林，林中有一个蜂巢，蜂巢就在众人面前，众人来到蜂巢前，边走边谈论。}你看出此处有何陷阱吗？一张摆好的餐桌，使得易于取得、食物的甘甜以及隐藏的希望，可能引诱他们违反誓言。因为饥饿、疲劳和时辰（据说，\Quote{全地}人\Quote{正在用餐}）立即催促他们违反誓言。此外，蜂房的景象从外面吸引他们放松决心。因为甘甜、摆好的餐桌以及偷吃难以被察觉，足以诱捕他们至高的智慧。若是需要煮或烤的肉，他们的心就不会如此着迷；因为在烹饪这些食物并准备食用时耽搁之际，他们可能预计会被发现。但如今没有这类事；只有蜂蜜，无需如此劳动，只需指尖蘸一下便足以享用，而且可以秘密进行。尽管如此，这些人克制了食欲，没有在心里说：\Quote{这关我们何事？我们中有人发过此誓吗？他为自己轻率的誓言受罚吧，因为他何必发誓？}他们没有这样想；而是虔敬地继续前行；虽有如此多的诱惑，他们仍明智行事。\Quote{众人边走边谈论。}这个字\Quote{谈论？}是什么意思？原来，他们彼此交谈，目的是用言语缓解痛苦。\par

7. 那么，当所有的人都如此明智行事时，就没有更多的事发生吗？那誓言果然被遵守了吗？甚至也没有被遵守。相反，它被违反了！如何，以何种方式？你立刻就会听到，为使你也彻底辨明魔鬼的全部伎俩。因为\Person{约纳堂}没有听见他父亲发这誓，\BibleQuote{他伸出手中的杖头，蘸在蜂房里，他的眼睛就明亮了。}见：\BibleRef{撒上 14:27}请注意，他怂恿谁违反誓言？不是士兵中的任何一个，而是那发誓者的亲生儿子。因为他不仅渴望导致背誓，而且还在策划杀害儿子，并事先为此作准备；他急于使天性反对自身。他在\Person{耶弗他}身上从前所做过的，他现在希望再次实现。因为他也同样，当他许诺在战斗胜利后，最先迎接他的东西他将献祭，就落入了杀子的陷阱；因为他的女儿首先迎接他，他就献了她，而天主没有禁止。事实上，我知道，许多不信者因这祭献指责我们残忍和无人性；但我不得不说，这祭献的容忍是天主眷顾和仁慈的一个显著例子；而且是为眷顾我们人类，祂没有阻止那祭献。因为如果在那个誓愿和许诺之后祂禁止了那祭献，许多\Person{耶弗他}之后的人，期待天主不会接受他们的誓愿，就会增加这类誓愿的数目，并继续前行而陷入杀子之罪。但如今，通过让这个誓愿实际实现，祂停止了将来所有这类情况。为了证明这是真的，在\Person{耶弗他}的女儿被杀之后，为使那灾难总是被记念，并使她的命运不被遗忘，在犹太人之中成了一条法律：童女们在同一季节聚集，哀悼四十天那已经发生的祭献；为的是通过哀哭更新其记忆，使所有人将来更加明智；并使他们知道这事不是按照天主的旨意行的，因为那样祂就不会允许童女们哀悼痛哭她了。我所说的不是揣测，事件证明了这一点；因为这祭献之后，没有人再向天主发过这样的誓愿。因此，祂确实没有禁止这事；但祂在\Person{依撒格}的情况下明确吩咐的事，祂直接禁止了；见：\BibleRef{创 22:12}清楚地显示出通过两个案例，祂并不喜悦这样的祭献。\par

8. 但那恶毒的魔鬼此时又竭力制造这样的悲剧。因此它怂恿约纳堂犯下罪过。因为若只是士兵中有人违犯了法律，在他看来自不算什么大恶；但他既对人的灾祸贪得无厌，永不能从我们的苦难中得到满足，便认为若只成就一件普通的凶杀，算不上什么大事。若不能同时用杀死他儿子的方式玷污君王右手，他便觉得一无所成。我何必提杀子之事？因为那恶者认为，借此他能完成比那更受诅咒的杀戮。因为若他明知故犯而被祭杀，那只是杀子之事；但如今他无知而犯（因为他尚未听闻那誓言），若他因此被杀，会使他父亲的痛苦加倍；因他既要牺牲一个儿子，又是一个无过错的儿子。但现在继续叙述历史：经上说：\Quote{他吃了以后，} \Quote{他的眼就明亮了。} 这里便指责君王大愚，表明饥饿几乎使全军失明，使他们的眼睛布满黑暗。随后有一士兵见到此事，便说：\Quote{你父亲曾向全民起誓说：今日谁若吃任何食物，必受诅咒。百姓都疲惫不堪。约纳堂说：我父亲把这片土地糟蹋了。} 这话中\Quote{糟蹋}何意？即他毁灭或败坏了他们。因此，当誓言被违犯时，众人默不作声，无人敢指出罪犯；这后来成了不小的指责，因为不仅那些背誓者，连那些知情而隐瞒的人也分沾了罪过。\par

9. 但我们来看下文：\BibleQuote{撒乌耳说：我们下去追赶外邦人，掳掠他们。司祭说：我们当在这里接近天主。}\BibleRef{撒上 14:36} 因为古时天主率领百姓出战；没有祂的许可，无人敢投入战斗，战争对他们而言是宗教事务。因为他们战败，不是由于身体软弱，而是由于自己的罪过；他们得胜，也不是凭勇力，而是靠上天的眷顾。胜利与失败对他们而言也是一种训练，是修德的学校。不仅对他们，对敌人也是如此；因为这使敌人也清楚看到，与犹太人交战的结局不是由兵器的性质决定，而是由战士的生活与善行决定。米德杨人至少看清了这一点，知道那民族是不可战胜的，用兵器与战车攻击他们是徒劳的，惟有藉罪恶才能战胜他们，于是他们装扮美貌的少女，列阵向前，激发士兵的邪淫，企图以淫乱剥夺天主对他们的助佑；结果果然如此。因为他们一犯罪，便轻易成为众人的猎物；而那些兵器、战马、士兵、及许多器械都不能俘获的人，罪恶的本性却将他们捆绑交予敌人。盾牌、长枪、标枪都同样无用；而容貌之美与灵魂之荡却胜过了这些勇士。\par

10. 因此，有人这样告诫：\BibleQuote{不要注视淫妇的美貌，不要接近好色的妇人。\BibleRef{训 9:8} 因为淫妇的嘴唇滴下蜂蜜，当时似乎甘甜入喉，但日后你必发现它比苦胆更苦，比双刃的剑更锋利。\BibleRef{箴 5:3}} 因为淫妇不知真爱，只知诱捕；她的吻含毒，她的口有致命之药。如果这毒不立即显现，就更当远离她，因她遮掩了那毁灭，隐藏了那死亡，不使它一开始就暴露。所以，若有人追求快乐和充满喜乐的生活，就当避免与淫妇交往，因为她们以言语和行为，在情人们心中激起无数争斗与骚动，挑起不断的争吵。正如那些最凶恶的仇敌，她们行为与计谋的目标就是使情人们陷入耻辱、贫穷与最坏的绝境。又如猎人布下网罗，驱赶野兽入内以杀死它们，这些女人也是如此。她们以眼神、装饰和言语，在四面撒下淫荡之网，然后驱赶情人们进来，捆绑他们，直至吸尽他们的血，最后还羞辱他们，嘲笑他们的愚昧，对他们倾泻嘲弄。这样的人实在不再值得怜悯，反该受讥笑，因为他比一个女人、一个妓女更无理性。因此，智者再次劝勉说：\BibleQuote{你当喝你自己池中的水，饮你自己井里的活水。\BibleRef{箴 5:15}} 又说：\Quote{愿你所爱的牝鹿，你宠爱的母鹿与你相伴。} 他这话是指依婚姻律与丈夫结合的妻子。为何你离开那帮助你的，反去投奔那谋害你的？为何你远离生活伴侣，去追求那颠覆你生命的人？她是你的肢体，你的身体；另一个却是利剑。所以，亲爱的诸位，要远避邪淫，既为现世的祸害，也为将来的惩罚。\par

11. 我们或许显得偏离了主题；但说到这个程度，并未脱离主题。因为我们并非只是为了阅读历史而向你们讲述，而是为使你们能纠正困扰自己的每一种激情；因此我们才这样频繁地劝勉，用各种文体风格为你们准备我们的训诲；因为在如此庞大的集会中，很可能有各种不同的病症；我们的职责不仅是治愈一种，而是许多不同的创伤；因此教育之药必须多样化。然而，让我们回到我们由此离题之处：\BibleQuote{司铎说：\InnerQuote{我们到天主面前去。}撒乌耳求问天主说：\InnerQuote{我该下去追赶外邦人吗？你将他们交在我手中吗？}但那一天主没有回答他。}请注意天主爱人的仁慈与温和。因为祂没有发出雷霆，也没有震动大地；而是像朋友们受到轻慢时所做的那样，主对仆人做了同样的事。祂只是沉默地接纳了他，以沉默说话，藉沉默表达了祂全部的愤怒。撒乌耳明白了这一点，便如经上所记说：\BibleQuote{把百姓各支派都叫到这里来，要知道和看看今天这罪是在谁身上。因为主，拯救以色列的那位是永生的，即使答案是指向我的儿子约纳堂，他也定要死。}\BibleRef{撒上 14:38}你看到他的轻率了吗？他意识到第一个誓言已被违背，却仍未学会自制，反而又加了第二个。也请考虑魔鬼的恶意。因为它知道，当儿子被发现并公开受审时，仅凭一瞥就能使父亲心软，并可能缓和王的愤怒，因此它藉第二个誓言的约束抢先宣判；用它作为一种双重锁链抓住他，不容他成为自己决定的主人，反而迫使他四面走向那不义的谋杀。甚至当犯事者尚未被带出来时，他就已下了判决；在不知罪犯是谁时，他就宣判了。父亲成了行刑者；在调查之前就宣布了定罪！还有什么比这更不合理的事呢？\par

12. 撒乌耳这样宣告后，百姓比以前更加害怕，所有人都陷入极大的战栗和恐惧中。但魔鬼却欢喜，因它使所有人都这样焦虑。据记载，百姓中没有人回答。\BibleQuote{撒乌耳说：\InnerQuote{你们将成为奴仆，我和我的儿子约纳堂也将成为奴仆。}}但他真正的意思是：\Quote{你们别无企图，不过是要把自己交给敌人，从自由人变成奴隶；同时你们惹天主发怒，因为不交出有罪的人。}也请注意誓言产生的另一个矛盾。如果他想要找出这罪的制造者，本应不作这样的威胁，也不应以誓言把自己绑在复仇上；这样，他们就不会那么害怕，反而更愿意把犯事者揭发出来。但在愤怒、极大的疯狂和先前的不理智影响下，他又做了与愿望直接相反的事。何须多说？他将此事交给抽签决定；签落在撒乌耳和约纳堂身上；\BibleQuote{撒乌耳说：\InnerQuote{你们在我和约纳堂之间抽签。}他们就抽签，约纳堂被抽中了。撒乌耳对约纳堂说：\InnerQuote{告诉我，你做了什么？}约纳堂告诉他说：\InnerQuote{我仅用手中的杖头尝了一点蜂蜜，看，我就要死了。}}有谁听了这些话会不被感动而心生怜悯呢？想想撒乌耳当时经历了怎样的风暴：他的五脏六腑被痛苦撕裂，两边都出现了最深的悬崖！然而他并没有学会自制，因为他说了什么？\BibleQuote{天主对我这样做，且加倍地做，因为你今天必死无疑。}\BibleRef{撒上 14:44}看，这是第三个誓言，而且不仅是第三个，还是一个时间极其紧迫的誓言；因为他不仅说\InnerQuote{你要死}，而且说\InnerQuote{今天}。因为魔鬼在催促他、催逼他、强迫他走向这亵渎的谋杀。因此它不许他为判决指定任何将来的日子，以免因拖延而纠正这恶行。百姓对撒乌耳说：\BibleQuote{天主对我们这样做，且加倍地做，如果执行死刑，这个人行了这样大的拯救在以色列中。主是永生的，他头上的一根头发也不可落地，因为他今天从天主行了仁慈的事。}\BibleRef{撒上 14:45}看，其次，百姓也起了誓，而且起了与王相反的誓。\par

13. 现在，我请求你们回想一下，孩子们拉扯的绳子，它断裂了，将拉扯的人摔倒在地上。\Person{撒乌耳}发誓不止一次两次，而是多次。百姓发相反的誓，朝着相反的方向拉。于是必然的结果是，无论如何，\OriginalTerm{誓愿}{oath}必须被打破。因为不可能所有人都遵守他们的誓愿。现在不要告诉我这件事的结果；但要思考由此产生了多少罪恶；以及魔鬼如何由此预备了\Person{阿贝沙隆}的悲剧和篡位。因为如果国王选择抵抗并执行他的誓愿，百姓就会列阵反对他；一场严重的叛乱就会发动。再者，如果儿子为了自身安全而选择投靠军队，他立刻就会成为\OriginalTerm{弑亲者}{parricide}。你们难道没有看见，叛乱、杀子、弑亲、战斗、内战、屠杀、流血、无数尸体，都是同一个誓愿的后果吗？因为如果战争偶然爆发，\Person{撒乌耳}可能被杀，\Person{约纳堂}也可能被杀，许多士兵会被砍成碎块；而最终誓言的遵守也无法推进。所以，你们不要考虑这些事件没有发生，而要注意这一点：事情的本性必然导致这类事情发生。然而，百姓占了上风。那么，让我们数一数随之而来的\OriginalTerm{伪誓}{perjury}。\Person{撒乌耳}的誓愿首先被他儿子违背；第二次和第三次，关于他儿子的死，被\Person{撒乌耳}自己违背。百姓似乎遵守了誓愿。但如果有人仔细审查，他们也都成了伪誓的罪人。因为他们迫使\Person{约纳堂}的父亲发伪誓，没有把儿子交给父亲。你们看到没有，一个誓愿使多少人自愿或不自愿地犯了伪誓？它造成了多少罪恶，带来了多少死亡？\par

14. 现在，在这篇讲道的开头，我曾许诺要证明，相反的誓愿必然导致伪誓；但历史的进程确实证明了比我所确立的更多。它展示的不仅是一个、两个或三个人，而是全体百姓；不仅是一个、两个或三个誓愿被违犯，而是更多。我也可以举另一个例子，从中表明一个誓愿如何造成了更大更严重的灾难。因为一个誓愿使全体犹太人遭致城邑被夺、妻儿被掳、火吞噬、蛮族入侵、圣物被亵渎，以及成千上万其他更令人痛苦的祸患。但我看出讲道已经太长了。因此，在此略过这段历史叙述，我恳求你们，连同\Person{若翰}被斩首之事，彼此也讲述\Person{约纳堂}的被杀以及全体百姓的普遍毁灭（虽然这并未实际发生，但已包含在誓愿的义务之中）；无论是在家中、在公共场合、与妻儿、与朋友、与邻居、与所有人，都要认真对待此事，不要认为我们以习惯作借口就足够了。\par

15. 因为这种借口只是托辞，过失并非来自习惯，而是来自懈怠，我将努力从已经发生的事来说服你们。皇帝关闭了城中的浴室，下令不得有人沐浴；没有人敢违犯法律，也没有人指责所发生的事，或者援引习惯。但即使身体虚弱，男女、儿童、老人，还有许多刚摆脱分娩痛苦的妇女，尽管所有人都将此视为必要的药方，却情愿或不情愿地忍受了禁令；既不求告身体的软弱，也不抱怨习惯的专横，更不申辩他们受罚而他人是肇事者，或其他诸如此类的事，而是心甘情愿地忍受这种惩罚，因为他们预期有更大的祸患；并且每日祈祷，愿皇帝的怒气不再继续。你们看见吗？哪里有恐惧，习惯的束缚就容易放松，尽管它由来已久，且十分必要。被剥夺使用浴室的权利无疑是一件痛苦的事。因为即使我们多么富于哲学精神，身体的本性也无法从灵魂的哲学中获得任何健康益处。但是关于戒除\OriginalTerm{起誓}{swearing}，这是极其容易的，且毫无伤害；对身体无伤，对心灵无损；相反，带来巨大收益、许多安全和丰裕的财富。那么，当皇帝命令时，我们忍受最大的艰难，而当天主命令的并不是什么沉重或困难的事，而是非常可容忍和容易的事时，我们却轻视或嘲笑它，并以习惯为借口，这难道不是荒谬吗？我恳求你们，不要如此轻视我们自己的得救，但要敬畏天主如同敬畏人一样。我知道你们听了这话会战栗，但真正值得战栗的是你们没有给予天主同样的尊重；你们殷勤遵守皇帝的谕令，却践踏来自天上的神圣诫命，并将对这些诫命的殷勤视为次要之事。如果我们经过如此多的劝告仍然坚持同样的行为，还能留给我们什么借口和宽恕呢？因为我在这场灾难一开始就开始了这劝告，这场灾难抓住了这座城市，而它现在即将结束；但我们甚至连一条\OriginalTerm{诫命}{precept}还没有完全付诸实践。那么，当我们未能遵守一条诫命时，我们又如何能祈求解除仍然困扰我们的灾祸呢？我们如何能期望情况好转？我们该如何祈祷？用哪条舌头呼求天主？因为如果我们遵守法律，当皇帝与城市和解时，我们将享受许多快乐。但如果我们仍处在\OriginalTerm{违犯}{transgression}中，羞耻和责备将从各方临到我们，因为当天主使我们脱离危险之后，我们却仍留在同样的懈怠中。\par

16. 啊！但愿我能脱去那些常发誓者的灵魂，将他们每日从誓言中所受的伤痕和瘀伤显露出来！那样我们便不再需要劝诫或劝导；因为这些伤痕的目睹比一切言语更有力，甚至能使那些最沉溺于此恶习的人远离他们的邪恶。然而，若不能将灵魂的羞耻状态展现在眼前，或许可以将其暴露在思想中，并在其腐朽和败坏中显示出来。正如经上所说：\Quote{不断受打的仆人不会没有瘀伤，同样，不断发誓和呼号天主之名的人不会洗净他的罪。}\BibleRef{德 23:10}不可能，完全不可能，那惯于发誓的口不常常犯虚誓。因此，我恳求你们所有人，抛弃这可怖的邪恶习惯，赢得另一顶冠冕。既然处处都在歌颂我们的城市，说她比世上所有城市更先在自己额上佩戴了基督徒之名，愿所有人都这样说：\Place{安提约基雅}是世上所有城市中唯一从自己边界驱逐了一切发誓的城市。不仅如此，如果这样做了，她不仅自己得冠冕，还会带动他人达到同样的热忱。正如基督徒之名源于此地，仿佛从某种泉源涌流至全世界，同样，这善功从此扎根和发端，将使地上所有居民成为你们的弟子；这样，你们将获得双倍甚至三倍的赏报，既因你们自己的善功，也因给予他人的教导。这对你们将是最灿烂的冠冕！这将使你们的城市成为一座母亲之城，不是在地上，而是在天上！这将在那一日为我们辩护，带给我们正义的冠冕；愿天主恩赐我们都能获得，藉着我们的主耶稣基督的恩宠和慈爱，祂与圣父及圣神，同享荣耀，自今而始，及于永远，世世无尽。阿们。\par

\SourceRecord{Translated by W.R.W. Stephens.；Nicene and Post-Nicene Fathers, First Series；Vol. 9.；Edited by Philip Schaff.；Buffalo, NY: Christian Literature Publishing Co.,；1889.；<http://www.newadvent.org/fathers/190114.htm>.}

% structure: p0001 source=h1 target=section reason=page-title

\section{论雕像 第十五讲}

\Emph{ \Emph{再论\Place{安提约基雅}城的灾难。恐惧在各方面都有益。忧伤比欢笑更有用。关于那句训言，\Quote{要记住你是在陷阱中行走。}} \Emph{而且，强迫人起誓比杀人更恶劣。} }\par

1. 今天和上一个安息日，我们本应开始谈论守斋；但不要有人认为我所说的话不合时宜。因为在守斋的日子里，关于这方面的劝诫和训导确实并不必要；这些日子本身就足以激发那些最懈怠的人努力守斋。然而，许多人，当即将进入守斋时，仿佛肚腹即将被交付给一种长期的围困，便预先囤积暴饮暴食的酒肉；而一旦获得自由，如同从长期的饥荒和痛苦的监牢中出来，又带着不合宜的贪婪冲向餐桌，好像他们试图通过暴饮暴食来抵消守斋所获得的益处；那时和现在，我们都有必要谈论节制。尽管如此，我们最近没有说过这类话，现在也不准备说。因为对迫在眉睫的灾难的恐惧，足以代替最有力的劝诫和训导，使每个人的灵魂清醒。因为在这样的风暴中，有谁如此可怜和堕落，竟至醉酒？有谁如此麻木，当城市如此动荡，这样的海难迫在眉睫时，竟不变得节制和警醒，并因这种痛苦而比任何其他训诫和劝告更彻底地改过自新？因为言语的效果不如恐惧。而现在，从所发生的事件中，这一点也可见一斑。在此之前，我们曾花费多少言语劝诫许多散漫的人，劝他们远离剧院和这些场所的污秽！但他们仍然没有远离；而总是在这一天，他们聚集在一起观看舞者的非法表演；他们举行魔鬼的集会，对抗天主教会充满的会众；以至于从那个地方传来的他们高声呼喊，在空中响起，与我们在这里咏唱的圣咏相对抗。但是看，现在我们保持沉默，对此事一言不发，他们却自行关闭了他们的乐队席；赛马场已变得无人问津！在此之前，我们中的许多人常赶去他们那里；但现在他们都从那里逃到教会，所有人都一样加入赞美我们的天主！\par

2. 你看到了恐惧带来了什么益处？如果恐惧不是一件好事，父亲们就不会为子女设置导师，立法者也不会为城市设立长官。还有什么比地狱更令人痛苦？然而，没有什么比对其的恐惧更有益；因为对地狱的恐惧将为我们带来天国的冠冕。哪里有恐惧，哪里就没有嫉妒；哪里有恐惧，贪爱钱财就不会扰乱；哪里有恐惧，愤怒被熄灭，邪恶的欲望被抑制，每一种不合理的激情都被消灭。如同在一所房子里，如果总是有一个武装士兵在场，就没有强盗、窃贼或任何此类恶徒敢于出现；同样，当恐惧占据我们的心灵时，任何邪恶的情感都不容易攻击我们，反而都被恐惧的专制力量驱散和放逐。不仅如此，我们从恐惧中还得到另一个更大的益处。因为恐惧不仅驱逐我们的邪恶情感，而且极其容易地引入各种美德。哪里有恐惧，哪里就有施舍的热忱、祈祷的热切、温暖而频繁的眼泪、充满痛悔的叹息。因为没有什么比持久的状态更能吞噬罪、使美德增长和兴旺的。因此，不活在恐惧中的人不可能行为正直；反之，活在恐惧中的人不可能走错路。\par

3. 因此，亲爱的，我们不要因目前的苦难而悲伤或沮丧，而要赞叹天主智慧的美妙计划。因为魔鬼希望藉以颠覆我们城市的手段，天主却藉以恢复和纠正了我们。魔鬼煽动一些不法之徒轻蔑地对待皇帝的雕像，意欲将城市夷为平地。但天主却利用同一个环境来纠正我们，用对预期震怒的恐惧驱散一切懈怠；而且，事情的结果与魔鬼所愿望的恰好相反，藉着他自己所准备的手段。因为我们的城市每天都在被净化；小巷、十字路口和公共场所，摆脱了淫荡和放纵的歌曲；无论我们转向何处，都有祈祷、感恩和眼泪，取代了粗鲁的笑声；有明智的言辞取代了污言秽语，我们的整座城市成了教会，作坊关闭，所有人整天都在参与这些公共祈祷，并用一个声音极其恳切地呼求天主。有什么讲道、劝诫、训导、多长的时间曾成就这些事呢？\par

4. 为此，我们应当感恩，不要急躁或不满；因为我们所说的，已显明恐惧是好的。但请听撒罗满关于恐惧所发出的智慧教训，他是在各种奢华中被养育，享受极大安全的撒罗满。那么他说了什么？\Quote{往哀悼之家，胜于往欢笑之家}。你说什么？请问，往那有哭泣、哀叹、呻吟和悲痛之处，胜于往那有舞蹈、铙钹、欢笑、奢华、饱食与醉酒之处吗？是，的确如此，他回答。请告诉我，为什么这样，理由何在？因为在前者，滋生傲慢；在后者，生发清醒。当一个人去一个更富有者的筵席时，他再不会以同样喜悦看待自己的家，而是带着不满回到妻子身边；不满地享用自家餐桌；对自己的仆人、儿女和家中每个人都变得暴躁；因他人的财富而更强烈地感受到自己的贫穷。不仅此恶，他还常常嫉妒那邀请他赴宴的人，回家后毫无益处。但论到哀悼之家，绝无此类事。相反，在那里可获得许多属灵的智慧，以及清醒。因为一旦一个人跨过停放尸体的房屋门槛，看见那已逝者无言躺卧，妻子撕扯头发、抓伤面颊、创伤手臂，他便被折服；面容变得忧伤；同坐的每个人都能对邻人说的只是：\Quote{我们算不了什么，我们的邪恶无法言表！}还有比这些话更充满智慧的吗？当我们承认本性的微不足道，指责自己的邪恶，视眼前事物为虚无时，我们实际上以不同的言辞表达了撒罗满那令人惊叹、充满天主智慧的心志：\BibleQuote{虚而又虚，万事皆虚}\BibleRef{训 1:2}。进入哀悼之家的人，立刻为逝者哭泣，即使他是仇敌。你看到那家比另一家好多少吗？因为在那边，即使是朋友，他也嫉妒；但在这里，即使是仇敌，他也哭泣。这是天主要求我们超过一切的事：不要凌辱那些曾使我们悲伤的人。不仅如此，我们还能收集到其他不亚于这些的好处。因为每个人也被提醒自己的罪过、可畏的审判台、大清算和审判；尽管他可能遭受他人加给的千般苦难，或因家中的理由而悲伤，他仍会得到并带回去这一切的药方。因为他沉思自己以及所有骄傲自大的人，不久也将遭受同样的事；一切现世之事，无论愉快还是痛苦，都是短暂的；于是他回家时，卸下了所有悲伤和嫉妒，心情轻松愉快；从此他将对所有人更加温和、善良、仁慈，也更加智慧；因为对未来的恐惧已进入他的灵魂，烧尽了所有荆棘。\par

6. 撒罗满觉察到这一切，当他说：\BibleQuote{往哀悼之家，胜于往饮酒之家}\BibleRef{训 7:3}。从前者生出懈怠，从后者生出恳切的焦虑。从前者生出轻视，从后者生出恐惧；这恐惧引导我们实践各种美德。如果恐惧不是好事，基督就不会在讲论未来的惩罚与报应上耗费如此长期而频繁的言论。恐惧无异于一道墙、一个防御、一座不可攻破的堡垒。因为我们确实需要许多防御，因为四面八方有许多埋伏。正如同一位撒罗满又劝诫说：\BibleQuote{你当知道，你行走在网罗之中，你行走在城垛之上}\BibleRef{德 9:13}。哦，这句话包含着多少美善！不亚于前者！让我们各人把这写在我们心版上，永远铭记在心，我们就不会轻易犯罪。让我们写在那里，首先极其准确地学会它。因为他不是说\Quote{注意}你行走在网罗之中，而是\Quote{知道}！他为什么说\Quote{辨别}？他告诉我们网罗是隐藏的；这确实是网罗，当毁灭不公开显现，伤害不明显，四面埋伏时。因此他说：\Quote{知道！}你需要许多反省和细心的审视。因为如同男孩用土覆盖陷阱，魔鬼也用今生的享乐遮盖我们的罪过。\par

7. 但要\Quote{辨别}；仔细审视；若遇到任何利益，不要只看利益，而要仔细检查，唯恐死亡与罪潜伏在利益之中；若察觉此事，就逃离它。同样，当某种娱乐或快乐偶然呈现时，不要只看快乐；唯恐在快乐的深处有某种邪恶隐藏着，要仔细搜索，若发现它，就赶快离开！若有人劝告、奉承、哄骗、承诺荣誉或任何类似之事，让我们作最仔细的调查；从各方面看这件事，唯恐有害的、危险的事，会因这劝告、荣誉或关注而偶然临于我们，我们便匆忙地、不知不觉地撞上去。因为若只有一两个陷阱，防范就容易了。但现在，请听\Person{撒罗满}如何说，当他希望陈述这些陷阱的众多时：\Quote{你要知道，你走在网罗之中}；他并没有说，你\Quote{经过}网罗，而是\Quote{在}网罗\Quote{中间}。两边是坑穴；两边是欺骗。一人去广场；看见仇敌；仅仅看见他就被激怒！一人看见朋友受尊敬；他就嫉妒！一人看见穷人；他就轻视，不加理会！一人看见富人；他就嫉妒！一人看见某人被侮辱；他就厌恶退缩！一人看见某人行侮辱之事；他就愤慨！一人看见美貌妇女，就被捉住！诸位，你们看见有多少陷阱吗？因此经上说：\Quote{记住，你走在网罗之中。}家中有网罗，餐桌上有网罗，交际中也有网罗。一个人往往无意中，出于友谊的信任，说出一些不应再重复的私事，于是招致极大的危险，致使全家因此毁灭！\par

8. 因此，让我们从各方面仔细搜查这些事。妻子、儿女、朋友、邻居，常常成为疏忽之人的网罗！有人会问：为什么有这么多网罗？为叫我们不低飞，而寻求上面的事。正如鸟儿，只要翱翔在高空，就不易被捕捉；照样，你只要仰望上面的事，就不易被网罗或任何其他诡计所捕获。魔鬼是捕鸟者。那么，飞得高过他的箭。那已升到高处的人，将不再仰慕今生任何事。但就如同我们登上山顶时，城墙与城市在我们看来很小，人们在地上行走如蚂蚁一般；同样，当你升至属灵智慧的高处，地上没有任何事物能迷住你；一切事物，甚至财富、光荣、尊荣，以及诸如此类的一切，当你注视天上之事时，都将显得微不足道。依\Person{保禄}看来，今生的一切荣耀都显得微不足道，比死物更为无用。因此他喊道：\BibleQuote{世界于我已被钉在十字架上。}\BibleRef{迦 6:14}因此他也劝诫说：\BibleQuote{你们要思念上面的事。}\BibleRef{哥 3:2}上面？请问你指的是什么事？是太阳所在之处，月亮所在之处？不，他说。那是在哪里？天使所在之处？总领天使？革鲁宾？色辣芬？不，他说。那是在哪里？\Quote{那里有基督坐在天主的右边。}\par

9. 那么，让我们听从，并时常思想这事：正如落入网罗的鸟，翅膀毫无用处，只是徒然拍打；同样，当你被邪恶的情欲牢牢捕获时，你的推理也无益，无论你怎样挣扎，你都被抓住了！为此，鸟被赐予翅膀，为叫它们能躲避网罗。为此，人拥有思考的能力，为叫他能避免罪。那么，当我们变得比禽兽更无知时，我们将有何等宽恕或借口？因为鸟一旦被网罗捕获，后来又逃脱，鹿落入网中又挣破后，都难以再次被同样的东西捕获；因为经验成了每个人谨慎的教师。但我们虽常被同一网罗所缠，仍坠入同样的陷阱；虽蒙恩宠拥有理性，却不效法无理性动物的先见之明和谨慎！因此，我们多少次因看见一个女人而遭受千万痛苦；回家后，怀着过度的欲望，受多日的煎熬；然而，我们仍不觉悟；当一创伤刚刚痊愈，我们又落入同样的恶行，被同样的手段捕获；为了一瞥短暂的快乐，我们忍受一种延长而持续的折磨。但若我们学会不断对自己重复这句话，就能避免这一切严重的邪恶。\par

10. 女子的美貌是最大的陷阱。更确切地说，不是女子的美貌，而是放荡的目光！因为我们不应指责对象，而应指责我们自己，以及我们自己的疏忽。我们不应说：不要有女子，而应说：不要有奸淫。我们不应说：不要有美貌，而应说：不要有淫乱。我们不应说：不要有肚腹，而应说：不要有贪饕；因为肚腹并不制造贪饕，而是我们的疏忽。我们不应说：吃喝是这一切罪恶的原因；因为这不是原因，而是我们的疏忽和贪得无厌。魔鬼既不吃也不喝，却堕落了！保禄吃了也喝了，却升上了天堂！我听见多少人这样说：不要有贫穷！因此，我们要堵住那些抱怨此事之人的口。因为发出这种怨言就是亵渎。对此，我们要对他们说：不要有狭隘的心胸！因为贫穷为我们的生活状态带来无数益处；没有贫穷，财富就毫无益处。因此，我们不应指责其中任何一个；因为贫穷和财富同样都是武器，若我们愿意，都会有助于德行。正如勇敢的士兵，无论拿起哪种武器，都展示自己的勇敢；而懦弱胆怯的人，无论哪种武器都会成为他的累赘。为使你明白这是真实的，请记住约伯的例子；他既富有，又贫穷，同样善用这两种武器，并在两者中都取得了胜利。当他富有时，他说：\BibleQuote{我的门向每位来者敞开。}\BibleRef{约 31:32} 但当他变得贫穷时，\Quote{主赐予，主收回。主看为好的，就这样成就了。} 他富有时，显出极大的好客；他贫穷时，显出极大的忍耐。那么，你呢——你富有吗？显出慷慨大度！你变贫穷了吗？显出坚忍和忍耐！因为财富本身不是恶，贫穷本身也不是；这些事之所以成为恶，取决于使用者的自由选择。因此，我们要训练自己，不对这些事情怀有那样的看法；也不要指责天主的作为，而要指责人的邪恶选择。财富不能有益于心胸狭隘的人；贫穷也绝不能损害宽宏大量的人。\par

11. 那么，让我们辨识这些陷阱，远远地离开它们！让我们辨识这些悬崖，甚至不要靠近它们！这将是我们最大安全的基础：不仅要避免有罪之事，也要避免那些似乎无关紧要、却容易使我们绊倒而陷入罪恶的事。例如，欢笑、戏谑的言谈，似乎不是明显的罪，但却引向明显的罪。欢笑常常生出污言秽语，污言秽语又生出更污秽的行为。言语和欢笑常常产生辱骂和侮辱；从辱骂和侮辱产生殴打和伤害；从殴打和伤害产生杀戮和凶杀。因此，你若为自己善加考虑，不仅要避免污言秽语、污秽行为、殴打、伤害和凶杀，还要避免不合时宜的欢笑本身和戏谑的言谈；因为这些事已成为后续罪恶的根源。因此，保禄说：\Quote{不可有愚妄的言谈和戏谑从你们口中出来。}因为虽然这本身似乎是一件小事，却为我们造成许多祸害。再者，过奢侈的生活似乎不是明显和公认的罪行；但它在我们的内产生巨大的恶——醉酒、暴力、勒索和抢夺。因为挥霍无度、生活奢华的人，过分地服侍肚腹，常常被迫偷窃、夺取他人的财物、使用勒索和暴力。因此，你若避免奢侈的生活，就除去了勒索、抢夺、醉酒和一千种其他邪恶的根源；从根部切断了不义。因此保禄说：\BibleQuote{那生活在奢侈中的寡妇，虽然活着，却已经死了。}\BibleRef{弟前 5:6} 再者，去戏院、观看赛马、掷骰子，对多数人来说似乎不是公认的罪行；但它们给我们的生活带来无穷的苦难。因为花时间在戏院里产生淫乱、无节制和各种不洁。观看赛马也带来争斗、辱骂、殴打、侮辱和长久的仇恨。对掷骰子的激情常常造成亵渎、伤害、愤怒、责骂和一千种更可怕的事。\par

12. 因此，我们不仅要避免罪恶，也要避免那些看似无所谓、却逐步引导我们犯下这些错事的事。的确，走在悬崖边上的人，即使没有掉下去，也会颤抖；而这颤抖常常使他倾覆，坠入深渊。同样，不远离罪恶，反而靠近它的人，将生活在恐惧中，并常常陷入其中。此外，贪婪地注视陌生美貌的人，即使没有犯奸淫，也已经怀有欲望；并且按照基督的宣告，他已经成为奸淫者；\BibleRef{玛 5:28} 而常常正是这欲望将他带入实际的罪行。因此，让我们远远地离开罪恶。你希望过节制的生活吗？不仅要避免奸淫，也要避免放荡的目光！你希望远离污言秽语吗？不仅要避免污言秽语，也要避免过分的欢笑和各样的情欲。你希望远离凶杀吗？也要避免辱骂。你希望远离醉酒吗？避免奢侈和丰盛的宴席，把恶习连根拔除。\par

13. 舌头的放任是一个大陷阱，需要强力的缰绳。因此，有人也说：\BibleQuote{人的口唇是自己有力的陷阱，他必因自己口中的话而受困。}\BibleRef{箴 6:2} 那么，在所有的肢体之上，让我们控制这个，让我们勒住它；让我们从口中驱逐辱骂、侮辱、污秽和毁谤的言语，以及起誓的恶习。因为我们的讲论再次将我们带到了同样的劝诫。但是昨天我已与你们的爱德约定，不再对此诫命多说，因为前些日子已经说得够多了。但我能怎么办呢？我不能停止这劝告，直到我看见你们付诸实践；因为\Person{保禄}对迦拉达人说话时，也说：\BibleQuote{从今以后，再无人能烦扰我。}\BibleRef{迦 6:17} 他似乎又与他们相遇并训诫他们。这就是慈父的心肠；虽说他们将要离开，却仍不离开，直到看见他们的儿子受到管教。你们今天听到了先知关于起誓对我们说什么？他说：\BibleQuote{我举目观看，看见，看，一把飞行的镰刀，长二十肘，宽十肘；他对我说：你看见什么？我说：我看见一把飞行的镰刀，长二十肘，宽十肘。它也要进入凡指着我的名起誓者的房屋，住在中间，拆毁石头和木头。} 这到底是什么意思？何以惩罚要以\Quote{镰刀}的形式，而且是\Quote{飞行的镰刀}来追赶起誓者？为叫你们看见审判是不可避免的，惩罚是无法逃避的。因为或许有人能从飞行的剑下逃脱，但从勒住脖子的镰刀，如同绳子一样，没有人能逃脱。再加上翅膀，还有什么得救的希望呢？但为什么它要拆毁起誓者房屋的石头和木头？为叫那废墟成为对众人的警戒。因为死者必须被泥土掩埋；他那已成废墟的房屋，变成了土堆，将警告所有路过看见的人，不要冒险做同样的事，免得遭受同样的惩罚；它将是对逝者之罪的持久见证。刀剑不如誓言的本质锋利！军刀不如誓言的打击那样具有毁灭性！起誓者虽然看似活着，其实已经死亡，受到了致命的打击。正如那个被套上绳索的人，在他出城走到坑边、看见刽子手站在他面前之前，从他走出法庭大门那一刻起就已死了；起誓者也是如此。\par

14. 让我们思考这一切，不要使我们的弟兄起誓。人啊，你做什么？你在祭台前逼人起誓，而就在基督被宰杀的地方，你杀害了自己的弟兄。强盗确实在大路上谋杀；而你在母亲面前杀死儿子：犯下比\Person{加音}更为可咒的谋杀；因为他在僻静处杀了他的兄弟，只带来现世的死亡；而你却在教堂中杀死你的兄弟，还带来未来不死的死亡！你以为教堂是为了这个目的而建立的，好让我们起誓吗？不，它建立是为了让我们祈祷！祭台摆在那里，是为了让我们宣誓吗？它摆在那里，是为了解除罪过，不是捆绑罪过。而你，若你不顾其他，至少该敬畏那本书，就是你在起誓时伸出的书；打开福音书，那是你命令人起誓时拿在手中的；当你听到基督在那里关于起誓所说的话时，应当战栗并停止！那么他在那里关于起誓说了什么？\BibleQuote{我却对你们说：总不可发誓。}\BibleRef{玛 5:34} 而你竟将禁止起誓的法律转变为誓言。哦，何等轻蔑！哦，何等暴行！因为你所做的正如同有人命令那位禁止杀人的立法者自己成为杀人的同谋。当我听说有人在路上被杀时，我并非如此哀叹哭泣，犹如当我看见任何人走近这祭台，将手放在上面，触摸福音书而发誓时，我呻吟、流泪、惊骇一样！我问你，你是在为钱财疑惑，却要杀害一个灵魂？你得到什么，能与你对自己灵魂和你邻人所造成的伤害相称？你若相信那人是诚实的，就不要施加起誓的义务；你若知道他是说谎者，就不要强迫他发假誓。有人说：\Quote{但我要求充分的保证。} 确实，当你没有让他起誓时，你才会获得良好而充分的保证。\par

15. 因为当你回家后，你会不断受到良心的折磨，内心这样思量：\Quote{那么，我让他发誓是毫无目的的吗？他岂不是真的发了虚誓？我岂不是成了这罪的原因？} 但如果你不让他发誓，你回家后就会得到极大的安慰，感谢天主，并说：\Quote{愿天主受赞美，因我克制了自己，没有强迫他无谓地、毫无目的地发誓。金子去罢！金钱消亡罢！} 因为特别给我们带来确证的是，我们没有违反法律，也没有强迫他人去违反。想一想，你是为了谁的缘故没有让人发誓；这就足以使你得到安慰和慰藉。的确，常常在争斗发生时，我们勇敢地忍受侮辱，并对侮辱者说：\Quote{我该拿你怎么办？某人阻止了我，他是你的保护人；他拦住了我的手。} 这足以安慰我们。同样，当你准备让人发誓时，要约束自己，停下来，并对那个即将发誓的人说：\Quote{我该拿你怎么办？天主禁止我让人发誓。他现在拦住了我。} 这既足以尊崇立法者，又足以保障你的安全，并使那准备发誓的人有所畏惧。因为当他看到我们如此害怕让人发誓时，他自己就更加害怕轻率地发誓了。你若这样说，你回到家时就会带着极大的确证。所以，要在祂的诫命中听从天主，好使祂自己在你的祈祷中垂听你！这句话要写在天上，在审判之日站在你身边，并赦免许多罪过。\par

16. 不仅关于发誓，而且关于每一件事，我们也应当这样思想。当我们为了天主的缘故要行一件善事，而它似乎会带来损失时，我们不要只看事情带来的损失，而要看我们因着为天主而行所将获得的收益。也就是说：有人侮辱了你吗？要高尚地忍受！你若不仅想到侮辱，而且想到那命令你忍受的天主的高贵，并温柔地忍受，你就能做到。你施舍了财物吗？不要只想到支出，而要想到从支出中产生的果实。你被人夺去了钱财吗？要感谢，不要只看到损失带来的痛苦，而要看到感恩带来的收益。如果我们这样约束自己，那么任何可能降临的沉重事件都不会使我们痛苦；相反，从那些看似痛苦的事中，我们甚至会受益，损失将比财富更甜美、更令人渴求，痛苦将比快乐更甜美，嘲笑和侮辱将比尊荣更甜美。这样，一切逆境都将转为我们的收益。在此世我们将享受极大的平安，在那世我们将获得天国；愿天主恩准我们众人配得获得，借着我们的主耶稣基督的恩宠和慈爱，祂和祂一起，偕同圣神，归于圣父，享受光荣、统治和尊威，从如今直到永远，永世无穷。阿们。\par

\SourceRecord{Translated by W.R.W. Stephens.；Nicene and Post-Nicene Fathers, First Series；Vol. 9.；Edited by Philip Schaff.；Buffalo, NY: Christian Literature Publishing Co.,；1889.；<http://www.newadvent.org/fathers/190115.htm>.}

% structure: p0001 source=h1 target=section reason=page-title

\section{论雕像第十六篇}

\Emph{\Emph{这篇讲道是在总督进入教堂时发表的，目的是安抚民心，因为有人向他报告了即将发生劫掠的谣言}\Emph{，当时所有人都考虑逃跑。它还涉及避免发誓的主题，以及宗徒的话：\Quote{保禄，耶稣基督的囚犯。}}}\par

1. 我赞赏总督的体恤，他看到城市动荡，人人都打算逃跑，便来到这里安慰了你们，并使你们怀有希望。但我却为你们脸红羞愧，经过这些长期频繁的讲道之后，你们竟然需要外来的安慰。我曾希望大地裂开将我吞没，当我听他（总督）与你们交谈，时而安慰，时而责备这种不合时宜且愚蠢的胆怯时。因为你们本不该由他教训；你们自己应该成为所有不信者的老师。保禄甚至不允许在不信者面前诉讼；\BibleRef{格前 6:1} 但你们，在接受了我们教父们那么多劝诫之后，却需要外来的老师；某些流浪汉和恶棍再次搅乱了这座大城，使之逃散。今后我们将以怎样的目光看待那些不信者，我们这些如此胆小懦弱的人？我们将以怎样的口舌对他们说话，并劝他们勇敢面对迫近的灾祸，当我们通过这次警报变得比任何兔子都胆小时候？\Quote{但我们能做什么呢？}有人说，\Quote{我们不过是人！}这正是我们不该害怕的原因，因为我们是人，不是畜生。因为畜生会被各种声音和噪音惊吓；因为它们没有推理能力来驱除恐惧。但你们，既然被赐予言语和理性的恩赐，为何沉沦到它们那种卑贱的状态？难道有人进了城，宣布士兵逼近？不要害怕，而是离开他，屈膝跪下：呼求你的主：痛苦呻吟，祂会阻止那可怕的事件。\par

2. 你们的确听到了关于行军的虚假报告，并面临与现世生命分离的危险。但那位有福的约伯，当使者接踵而来，听到他们宣布可怕的消息，并加上子女无法忍受的毁灭时，他既不哭泣也不呻吟，而是转向祈祷，感谢主。你也要效法他；当有人来报告士兵已包围城市，即将掠夺财富时，逃向你的主，并说：\BibleQuote{主赐予，主收回；按主所喜悦的，就如此成就。愿主的名永远受赞美。}实际事件的经历并没有吓到他；但仅仅一个报告就吓到了你们。我们这些被命令勇敢面对死亡本身的人，却因一个虚假的谣言而如此惊恐，又该如何看待呢？迷惑之人制造不真实的恐惧和看不见的麻烦；但灵魂安定宁静的人，甚至能粉碎真实的恐惧。你们没有看到舵手吗？当大海翻腾，乌云汇集，雷电爆发，船上所有的人都慌乱时，他们安坐舵旁，没有骚动，专注于自己的技艺，考虑如何抵挡即将来临的风暴。以他们为榜样；抓住神圣的锚，即对天主的希望，坚定不移。\BibleQuote{凡听了我这些话而不实行的，就好像一个愚昧人，把房屋盖在沙土上；雨淋，水冲，风吹，袭击那座房屋，它就坍塌了，且坍塌得很惨。}\BibleRef{玛 7:26}你们看见，愚昧的性情就是一头栽倒、被推翻吗？或者，我们不仅陷入了那愚昧人的状态，而且我们的跌倒更为悲惨。因为那人的房屋是在河水雨水降下、风吹打之后才倒塌的；但我们，在没有任何风吹、水冲、袭击的情况下，在灾难临到之前，就因一个谣言而倾覆，立刻放弃了我们所有思考的哲学。\par

3. 你们以为我现在在想什么？我该如何隐藏——是的，埋葬自己？我该如何羞愧脸红？若不是教父们极力催促，我不会站起来，不会发言，因为我的心思因你们的怯懦而充满悲伤。但即使现在我也无法恢复，因为愤怒和悲伤围攻了我的灵魂。谁会不感到愤慨和愤怒：经过这么多教导，你们竟需要外邦人的指导，以便得到安慰并勇敢地忍受当前的警报？所以你们要祈祷，使我们在开口时能获得发言的权能；使我们能摆脱这悲伤，稍加振作；因为这因你们怯懦而生的羞愧确实大大挫败了我们的精神。\par

4. 近日，我向你们的爱德讲述了许多关于我们四围陷阱的事；也论及恐惧与悲伤、忧愁与喜乐；以及那飞降在发誓者房屋上的镰刀。如今，在这一切众多事中，我特别希望你们记住我论及\Quote{带翅膀的镰刀}及其落在发誓者房屋中，拆毁石头与木头，并吞灭整个建筑所说的话。此外，也要留意这一点：以拿着福音书发誓，将禁止发誓的律法本身转为誓言，这是极端的愚昧；而宁可遭受财产损失，也不愿向邻人强加誓言，这才是对天主应做的极大尊荣。因为当你对天主说：\Quote{因祢的缘故，我没有让那抢夺并伤害我的人起誓}，天主必因这尊荣，在此世与来世赐予你丰厚的赏报。将这些事告诉别人，自己也当遵守。我知道在此处我们会更加敬畏，并抛弃一切恶习。但所当追求的是，不仅在此处热爱智慧，更要在离去时，将这敬畏带出去，尤其是当我们特别需要它的时候。因为那些打水的人，并非只在泉水旁装满器皿，到家时却倒空；相反，他们在家时会格外小心地放好，以免打翻，使辛劳白费。让我们都效法这些人：回家后，严格持守所讲的道理；因为倘若你在此处装满，回家时却空空如也，使理解的器皿缺乏所听的内容，那么在此处的充满便无益处。不要只向我展示竞技场中的摔跤手，而要展示实际竞赛中的；不要只在听道时显虔诚，而要在实行时显虔诚。\par

5. 你对现在所说的话表示赞许。但当需要你发誓时，便要记住这一切。若你速速成就这条律法，我们便会将教导推进到其他更重大的事。看！这是我向你们的爱德讲道的第二年，而我尚不能解释一百行圣经经文。原因是，你们需要从我们这里学习那些你们在家中自己便能付诸实践的事；因此，我们劝勉的大部分内容都消耗在伦理训导上。但这本不应当如此；品德的规范你们本应在家中自己学习，而圣经的意义及其思辨，则可以托付给我们。然而，倘若你们必须从我们这里听到这等事，那也不需要超过一天：因为所要说的并非多样或困难，也不需任何精细阐述。当天主宣告祂的判决时，微妙的争论便不合时宜。天主曾说：\BibleQuote{你不可发誓。}那么，不要向我追问这诫命的原因。这是君王的律法。设立它者，知道律法的原因。若它无益，祂便不会禁止。君王制定法律，但或许并非全部有益；因为他们是人，不能像天主一样有能力发现何谓有益。然而，我们仍服从他们。无论我们结婚、立遗嘱，或要购买仆人、房屋、田地，或做任何其他事，我们都按他们所订的法律而行，而非凭自己的心意；我们并非完全能按己意处置关乎自己的事；在许多事上，我们服从他们的意愿；若我们做了任何违反他们判断的事，那事便无效且无用。那么，请告诉我，我们既如此尊重人的法律，却践踏天主的法律，这行为配得什么辩护或宽恕呢？祂曾说：\BibleQuote{你不可发誓。}为使你一切言行安稳，你不可在实际生活中制定一条相反的法律。\par

6. 但这些事已经说得够了。现在我们来看看今天所读经文中的一句话，并以此结束这次讲道。他说：\Quote{\Person{保禄}，\Person{耶稣基督}的囚犯；以及兄弟\Person{弟茂德}}。保禄的称号是何等伟大：没有权柄与能力的头衔，却提到捆绑与锁链！真是伟大啊！虽然许多其他事使他显赫——他被提到第三层天，被带到乐园，听见不可言说的话语——但他没有列出这些，反以锁链代替一切，因为这比那些事更使他显赫和荣耀。为什么呢？因为前者是主慈爱的无偿恩赐，而后者是仆人恒心与忍耐的标记。凡有爱德的人，常以他们为所爱者所受的苦难而夸耀，胜过接受他们的恩惠。一个君王不会像保禄因自己的锁链而夸耀那样以王冠为荣。这很合理。因为王冠只是给戴冠之头增添装饰，但锁链既是更大的装饰，又是保障。王冠常出卖它所环绕的头，引来无数叛徒，怂恿他们贪图权势。在战斗中，这种装饰如此危险，以致必须隐藏和放下。因此君王在战斗中改变外表，混入战士群中；王冠带来如此多的背叛。但锁链不会给佩戴者带来任何这类事，反而完全相反：因为若有战争，与魔鬼交战，敌对势力，那如此被环绕的人，只要举起他的锁链，就能击退他们的攻击。而且许多世俗官员不仅在掌权时持有官职头衔，甚至在卸任后也是如此。这样的人被称为\OriginalTerm{前执政官}{ex-consul}，这样的人被称为\OriginalTerm{前裁判官}{ex-praetor}。但他代替所有这些称号说：\Quote{囚犯保禄}。这很恰当。因为那些官职并非心灵德行的完全证据；它们可以凭金钱买到，凭朋友的请求得到；但这由捆绑得到的殊荣是灵魂热爱智慧的证明，是渴慕基督的最强标记。前者很快消失，但这殊荣无与伦比。请看，从那时至今，多么长的时间过去了，这位囚犯的名字却日益显赫。至于所有以前的执政官，无论他们是谁，都已归于沉默；甚至他们的名字也不为一般人所知。但这囚犯、可敬的保禄的名字，在此处仍是伟大的，在异邦人之地是伟大的，在斯奇提亚人和印度人中也是伟大的；即使你到有人居住的世界的尽头，你也会听到这个称呼，无论到哪里，都会发现保禄的名字挂在所有人的口中。这有何惊奇呢，若在陆地与海上如此，甚至在诸天中保禄的名字也是伟大的；在天使、总领天使和天上的掌权者中，并在这一切的君王，即天主面前！但有人会问：这些锁链是什么样的？难道不是铁制的吗？它们确实是铁制的，但它们包含着圣神的恩宠，在其中丰盛地兴盛，因为他为基督的缘故佩戴它们。哦，奇事！仆人被捆绑，主人被钉十字架，然而福音的宣讲日日增长！并且借着那些被认为会熄灭它的方法，恰恰通过这些方法它被点燃了；十字架与锁链，曾被视为可憎之物，如今却成了救恩的象征；那铁链为我们比全金更宝贵，不是因为它本身的本质，而是因为这个原因与根基！\par

7. 但在此我看到一个问题从这一点产生；如果你们注意听，我将准确陈述问题，并加上解答。那么探讨的主题是什么呢？这位保禄曾有一次来到\Person{斐斯托}面前，向他讲话，为自己辩护，反驳犹太人对他提出的指控，讲述他如何看见了\Person{耶稣}，如何听到了那有福的声音；如何暂时失明又恢复视力，跌倒又起来；如何不带锁链地被俘进入大马士革；同样也讲了律法和先知，并表明他们早已预言了这一切。他赢得了审判官，几乎说服他归向自己。因为圣人们的灵魂正是如此：当他们陷入危险时，他们不考虑如何脱离危险，而是竭力想方设法俘获迫害他们的人。当时正是如此。他进来为自己辩护，离开时却带走了审判官！审判官为此作证说：\BibleQuote{你几乎说服我成为基督徒。}\BibleRef{宗 26:28} 这件事本应在今天发生；这位总督来到你们中间，本应赞赏你们的慷慨大度、你们的坚毅、你们的完全平静；并带着对你们良好秩序的教训离开，赞赏你们的集会，赞扬你们的聚会，并从事实中学习到外邦人与基督徒之间有多大的差别！\par

8. 但正如我所说的：当保禄捉住他，他说：\Quote{你几乎说服我成为基督徒}，保禄这样回答说：\BibleQuote{我祈求天主，不但你，而且今天所有听我的人，都像我所是的，几乎并完全如此，除了这些锁链之外。}\newline{}\BibleRef{宗 26:29} 保禄，你说什么？当你致书厄弗所人时，你说：\BibleQuote{所以我这为主被囚的，劝你们：行事为人要与你们所蒙的召叫相称。}\newline{}\BibleRef{弗 4:1} 当你对弟茂德说话时：\BibleQuote{我在这事上为福音受苦，甚至被锁链束缚，如同作恶的人一样。}\newline{}\BibleRef{弟后 2:9} 再次，当对费肋孟时，这样：\BibleQuote{保禄，为基督耶稣被囚的。}\newline{}\BibleRef{费 1:1} 再次，当与犹太人辩论时，你说：\BibleQuote{我为了以色列的希望，被这条锁链捆绑。}\newline{}\BibleRef{宗 28:20} 当致书斐理伯人时，你说：\BibleQuote{许多弟兄因我受的锁链，在主内更加勇敢地宣讲圣言而无惧。}\newline{}\BibleRef{斐 1:14} 你到处带着锁链，到处展示你的锁链，并以此为夸耀。但当你来到法庭，你却出卖了你的哲学，本应最勇敢说话的地方，你对判官说：\Quote{我祈求天主，你成为基督徒时\InnerQuote{没有}这些锁链！}然而，如果锁链是好的，且如此之好，以致能他人在真宗教事业上更加勇敢（因为你自己先前已经宣布了这事，当你说：\Quote{许多弟兄因我的锁链得了鼓励，竟更勇敢地宣讲圣言而无惧}），那么，你为什么不在判官面前以此夸耀，反而做相反的事呢？\par

9. 我所说的难道不像是问题吗？然而我立刻提出它的解答。因为保禄如此行，不是出于痛苦或恐惧，而是出于丰富的智慧和属灵的领悟。我接着解释这是如何的。他正在对一个外邦人，一个不信者说话，那人对我们的事一无所知。因此他不愿用令人不愉快的事情引他入门，而是如他所说：\BibleQuote{对那没有法律的人，我就如同没有法律的人}；\newline{}\BibleRef{格前 9:21} 他在当前的情况也是如此行。他的意思是：如果外邦人听到锁链和患难，他会立刻逃跑；因为他不知道锁链的能力。首先，让他成为一个信徒；让他品尝所宣讲的圣言，然后他甚至会自己急切地奔向这些锁链。我曾听见主说：\Quote{没有人把新布补在旧衣服上，因为所补上的新布会扯破那衣服，破绽就更大了。也没有人把新酒装在旧皮囊里，否则，皮囊会破裂。}这个人的灵魂是旧衣服，旧皮囊。它没有被信德更新，也没有被圣神的恩宠革新。它仍是软弱的、属世的。它贪恋今生的事物。它急切地追逐世俗的炫耀。它喜爱现世的荣耀。如果他立刻，甚至从一开始就听说，若成为基督徒，就会立即成为囚犯，并被锁链环绕；他会感到羞愧和愤慨，从而从所宣讲的圣言前退缩。因此，他说：\BibleQuote{除了这些锁链}\newline{}\BibleRef{宗 26:28} 这不是贬低锁链本身（决不允许！），而是迁就对方的软弱；因为他自己热爱并欢迎他的锁链，正如一个爱装扮的女人喜爱她的金饰一样。哪里显明这一点？他说：\BibleQuote{我在为你们受苦中喜乐，并在我肉身上补充基督苦难所欠缺的}\newline{}\BibleRef{哥 1:24} 再次：\BibleQuote{因为为了基督的缘故，赐给你们的恩赐，不仅相信他，而且也要为他受苦}\newline{}\BibleRef{斐 1:29} 再次：\BibleQuote{不但如此，我们在苦难中也喜乐}\newline{}\BibleRef{罗 5:3} 因此，如果他为此喜乐和夸耀，并称之为恩宠的恩赐，那么显然当他向判官说话时，他之所以那样说，是出于刚才所说的理由。此外，在另一处，当他碰巧有夸耀的需要时，他表明完全相同的事，说道：\BibleQuote{所以我甘心情愿夸耀我的软弱……在凌辱、艰难、迫害、困苦中，因为基督的能力要停留在我的身上}\newline{}\BibleRef{格后 12:9} 再次：\BibleQuote{若必须夸耀，我就要夸耀那关乎我软弱的事}\newline{}\BibleRef{格后 11:30} 在其他地方，当他将自己与他人比较，并在比较中向我们展示他的优越时，他这样说：\BibleQuote{他们是基督的仆役吗？（我如愚昧人说，）我更是}\newline{}\BibleRef{格后 11:23} 为了显示这优越，他没有说他使死人复活，没有说他驱逐了魔鬼，没有说他洁净了癞病人，也没有说他做了任何其他类似的事，而是说他忍受了那无数的艰难。因此当他说\Quote{我更是}时，他立即列举他的许多试炼：\BibleQuote{受鞭打是过量的，冒死是屡次，坐监是多次……被犹太人鞭打五次，每次四十减去一下，被石头砸过一次，遭遇船难三次，在深海中度过了一夜一日}以及所有其余。因此保禄到处在患难中夸耀；并为此事极为自豪。这十分恰当。因为正是此事特别显示出基督的能力，即宗徒们藉着这些方式得胜：藉着锁链、患难、鞭打和最坏的恶事。\par

10. 因为基督曾预告过这两件事：苦难与赦免，辛劳与冠冕，劳苦与奖赏，令人愉快与令人悲伤的事。然而，祂将悲伤的事归于现世生活，却将愉快的事存留于来世；既显示祂无意欺骗人，又愿以此安排减轻人类苦难的重担。因为欺骗者首先提供愉快的事，随后才带来令人不愉快的事。例如：绑架者想要偷走并拐带幼童时，并不向他们许诺拳打脚踢或任何此类之事，而是提供糕点、糖果以及诸如此类，这些通常是儿童年龄所喜好的；为的是他们被这些东西引诱后，出卖自己的自由，陷入极大的危险。此外，捕鸟者和渔夫也这样引诱他们所追捕的猎物，先提供它们的惯常食物及合其口味的东西，以此掩盖陷阱。所以，这尤其是欺骗者所为：先提供令人愉快的事，然后才带来令人不愉快的事。但对于那些真正关怀并远虑他人的人，情况完全相反。父亲们至少以与绑架者相反的方式行事。当他们送孩子上学时，给孩子设立师长，以鞭挞相威胁，并用恐惧四面环绕他们。但当他们这样度过人生的最初部分，习惯养成后，就使他们获得尊荣、权力、奢华以及属于他们的一切财富。\par

11. 天主正是这样行事。祂效法远虑的父亲，不像绑架者，先使我们经受困苦之事；将现世的苦难交给我们，如同交给师长和教师一般；为使我们借此受到管教和清醒，在显示一切忍耐并学习一切正确纪律之后，当我们养成合适的习惯时，就能继承天国。祂先预备我们，使我们适合管理祂将赐予的财富，然后才使我们实际拥有财富。因为若非如此，赐予财富便不是恩赐，而是惩罚和报复。正如一个愚昧挥霍的儿子，继承父业后，恰恰因这财富而自取灭亡，因为他毫无管理财富所需的实践智慧；但如果他聪明、温良、节制、适度，适当地管理父业，便会因此更加显赫出众：同样，在我们身上也必须如此。当我们获得属灵的理解，当我们都达到了\Quote{完美的成年男子,}并达到完全身量的尺度时，祂才使我们拥有祂所许诺的一切；但现在，祂像对待幼童一样管教我们，同时给予安慰和抚慰。预先经历苦难不仅有此益处，还有另一项同样不小的益处。因为那先享奢华，然后必须在奢华之后等待惩罚的人，甚至因对即将来临的苦难的预期而无法享受眼前的奢华；但那些先处在忧伤状态、并期待日后享受美善的人，因着对将来美善的希望，会忽视眼前的困难。因此，不仅为了我们的安全，也为了我们的愉悦和安慰，祂安排了先经受困苦的事；为使我们因着对未来的希望而轻松，对眼前的苦难变得麻木。保禄宗徒正想表明并阐明这一点，当他说：\BibleQuote{我们这至暂至轻的苦难，要为我们成就极重无比、永远的荣耀。我们不是顾念所见的，乃是顾念所不见的。}见：\BibleRef{格后 4:17}他称苦难为轻，并非因为困苦之事本身的性质，而是因为对未来美善的期待。正如商人因着对货物的希望而对航海的辛劳漠不关心；又如拳击手勇敢地承受击打头部，期待未来的冠冕；同样，我们确实也应当如此：热切仰望天堂及天上的美善，无论什么恶事临到我们，都靠对未来美善的希望而坚强承受，以勇气支撑一切。\par

12. 因此，让我们带着这句话回家；因为它虽然简单简短，却包含着许多灵性智慧的道理。处于忧伤和困苦中的人，有了足够的安慰；生活在奢侈和富足中的人，也有了能使他们大大清醒的东西。因为当你坐在餐桌旁时，一想起这句话，你就会迅速远离酗酒和贪食；通过这句话，你学会我们是多么需要努力奋斗；你会对自己说：\Quote{保禄生活在锁链和地牢中，而我却在酗酒和丰盛的筵席中！那么我将获得什么宽恕呢？}这句话也适合妇女；那些喜欢装饰、昂贵衣着、用黄金把自己从头到脚缠绕起来的妇女，当她们想起这条锁链时，我相信她们会憎恨和厌弃那样的打扮；并会渴望这样的锁链。因为那些装饰常常是许多邪恶的原因，给家庭带来一千次争吵，并滋生嫉妒、猜忌和仇恨。但这些锁链解开了全世界的罪恶，吓退了魔鬼，赶走了恶魔。带着这些锁链，他在狱中说服了狱卒；带着这些锁链，他吸引了阿格黎帕本人；带着这些锁链，他赢得了许多门徒。因此他说：\BibleQuote{我在这事上受苦，甚至被锁链捆绑，像作恶的，然而天主的言语却不被捆绑。}\BibleRef{弟后 2:9}正如不可能捆绑一束阳光，或将其关在屋里一样，言语的宣讲也是如此；而更甚的是，教师被捆绑，言语却四处飞翔；他住在监狱里，但他的教义却迅速插翅飞遍全世界！\par

知道这些事，让我们在遇到逆境时不要沮丧，反而要更坚强、更有力量；\BibleQuote{因为患难生忍耐。}\BibleRef{罗 5:3}让我们不为临到我们的灾祸而悲伤，反而要在一切事上感谢天主！\par

13. 我们已经完成了四旬期的第二周，但这不应成为我们的考量；因为度过四旬期并不在于仅仅度过时间，而在于以品行的改善来度过。让我们思考这一点：我们是否变得更勤勉？我们是否改正了任何缺点？我们是否洗去了我们的罪过？在四旬期内，每个人都会问自己守斋了几周；有些人可能会听说他们守了两周，有些人三周，还有些人守了整周。但如果我们的四旬期缺乏善工，又有什么益处呢？如果有人说：\Quote{我守了整四旬期，}你要说：\Quote{我有一个仇敌，但我与他和好了；我有说坏话的习惯，但我停止了；我有发誓的习惯，但我打破了这种邪恶的实践。}对于商人来说，航行了大片海洋却空手而归，毫无益处；必须负载货物和大量商品才有益。如果我们只是按部就班地度过四旬期，毫无结果，四旬期对我们毫无益处。如果我们只实行禁食肉食，四十天过去了，四旬期也就结束了。但如果我们戒绝罪恶，这即使在四旬期结束后仍然存在，并且从此成为我们持续的利益；在我们到达天国之前，它就会在此世给我们不小的赏报。因为正如生活在不义中的人，即使在地狱之前，也因良心责备而受惩罚；同样，富有善工的人，即使在天国之前，也将因被有福的希望滋养而享有极大的喜乐。\par

14. 因此基督说：\BibleQuote{我要再见到你们，你们的心要喜乐，并且你们的喜乐没有人能夺去。}\BibleRef{若 16:22}一句简短的话，却包含着许多安慰。那么这是什么意思呢？\Quote{你们的喜乐没有人能夺去？}如果你有钱，许多人能夺去你财富带来的喜乐；例如，小偷挖墙而入，仆人盗取所托付的财物，皇帝没收财产，嫉妒者以侮辱夺走。如果你拥有权力，许多人能剥夺你因权力而有的喜乐。因为当职务的条件结束时，快乐的条件也随之结束。即使在行使职务本身的过程中，也有很多意外发生，带来困难和忧虑，会铲除你的满足感。如果你有身体的力量，疾病的侵袭会阻止从那里来的喜乐。如果你有美貌和青春，衰老的临近会使其凋谢，并带走那喜乐。或者如果你享受丰盛的餐桌，夜幕降临宴会的喜乐就结束了；因为此生的一切都容易受损，无法给我们持久的快乐；但虔诚和灵魂的德行则完全相反。如果你做了施舍，没有人能夺去这个善工。即使军队、国王、无数的诽谤者和阴谋者从四面八方包围你，他们也不能夺去已存入天堂的财富；其喜乐持续常在；因为经上说：\BibleQuote{他施舍钱财，赒济贫穷，他的仁义存到永远。}这是非常公义的；因为在天堂的仓库里，没有贼挖窟窿，强盗抢夺，蛀虫咬坏。如果你献上持续热切的祈祷，没有人能夺去他们的果实；因为这果实也扎根于天堂；它不受任何伤害，不属凡人所能及。如果你在受虐待时行了善事；如果你耐心忍受别人的恶言；如果你以祝福回报辱骂；这些善工会持续常在，没有人能夺去它们的喜乐；而且每当你想起这些，你就欢喜快乐，收获巨大的喜悦果实。同样，如果我们成功避免发誓，并说服我们的舌头戒除这种有害的实践，善工将在短时间内完成，但由此产生的喜乐将是持续不断的。\par

15. 现在正是你们成为他人的导师和向导的时候；朋友应负责教导和引领邻人，仆人对同仆，青年对同龄人。若有人应许你，每改造一个人就给你一块金币，你岂不是要竭尽全力，整日坐在他们身旁，劝导和鼓励他们吗？\newline{}然而，现今天主应许你的不是一块金币，也不是十块、二十块、一百块、一千块，甚至不是整个世界，作为你劳苦的报酬，而是赐给你那比全世界更伟大的，就是天国；不仅如此，还另有赏报。那赏报是什么呢？祂说：\BibleQuote{你若从卑贱中取出宝贵的，必如我的口。}\BibleRef{耶 15:19}\newline{}\newline{}就尊荣和安全而言，有什么能与这相比呢？那些在如此重大应许之后仍忽视邻人得救的人，还有什么借口或宽恕可言呢？你若看见一个瞎子快要掉进坑里，你会伸手救他，并认为若坐视其死亡是件可耻的事吗？但每天看见你的弟兄们陷入发虚誓的恶习，你竟连一句话也不敢说吗？\newline{}你或许已经说过一次，他却未听从；因此，要说两次、三次，甚至多次，直到你说服他为止。天主每天都对我们说话，我们却不听；然而祂并没有停止说话。所以，你也要效法这种对邻人的温柔关怀。正因为如此，我们才彼此相聚；我们才住在城市里，在圣堂中聚会，为的是要彼此分担重担，纠正彼此的罪过。\newline{}如同同一店铺里的商人们，各自经营，却将所得全部投入公共基金；我们也该如此行事。每个人能给予邻人的任何益处，都不要吝啬或退缩，而要实行这种属灵的商贸和互惠；\newline{}这样，我们将一切存入公共库藏，获得巨大的财富，积攒丰厚的宝藏，一同分享天国。\newline{}这是藉着我们的主耶稣基督的恩宠和慈爱，愿光荣归于祂，偕同祂，在圣神内，归于圣父，从现在直到永远，万世无疆。阿们。\par

\SourceRecord{Translated by W.R.W. Stephens.；Nicene and Post-Nicene Fathers, First Series；Vol. 9.；Edited by Philip Schaff.；Buffalo, NY: Christian Literature Publishing Co.,；1889.；<http://www.newadvent.org/fathers/190116.htm>.}

% structure: p0001 source=h1 target=section reason=page-title

\section{讲道集十七：论雕像}

\Emph{关于皇帝狄奥多西因雕像被推翻事件为调查罪犯所派遣的专员（军队指挥官\Person{赫勒比库斯}和宫廷总管\Person{凯撒里乌斯}）。}\par

1. 今天我们最适时地众口同声地高唱：\Quote{上主以色列的天主当受赞美，唯独祂行奇事。}因为所发生的事何其奇妙，远超一切期望！整个城市，如此众多的人口，即将被淹没，沉入波涛，彻底而即刻地毁灭，祂却在片刻间将其完全从船难中拯救出来！我们不仅要感恩，因为天主平息了风暴，也因祂容许了风暴发生；不仅因祂将我们从船难中救出，也因祂允许我们陷入那样的困境，并让如此极端的危险悬在我们头上。同样，保禄也命令我们\BibleQuote{事事感谢}\BibleRef{得前 5:18}。但他说\BibleQuote{事事感谢}时，不仅指从恶事中得到解救，也包括遭受这些恶事的时候。因为\BibleQuote{天主使一切协助那些爱祂的人获得益处}\BibleRef{罗 8:28}。让我们为这从试炼中获得的解救而感恩，并永远不忘这些事。让我们专心祈祷，不断恳求，并加深虔诚。\par

2. 当这场灾难的悲伤之火初次燃起时，我曾说过，这不是讲道的季节，而是祈祷的季节。现在，当火已熄灭，我再重复同样的话：此刻尤其、甚至更需祈祷；现在是泪水与痛悔、焦虑的灵魂、极大的勤勉与极大谨慎的季节。因为那时，忧患的性质本身，无论我们是否愿意，都约束了我们，使我们倾向清醒，并引导我们更虔诚；但如今，约束的缰绳已卸下，阴云已过，我们恐怕会重新陷入懈怠，或因这喘息而松弛；恐怕有人也会这样论我们说：\Quote{祂击杀他们时，他们才寻求祂，回转过来，切切寻求天主。}因此，梅瑟也劝诫犹太人：\BibleQuote{当你吃喝饱足时，要记念你的天主上主。}\BibleRef{申 6:11}你心志的良善现在将得以显明，如果你继续践行同样的虔诚。因为那时，许多人将你的热忱归于恐惧和灾难的临近；但现在，如果你仍坚持保持这热忱，这将纯粹是你自己的成就。正如一个由他所畏惧的师傅引导的男孩，当他清醒温和地生活时，无可称赞，因为人们将少年的清醒归于他对师傅的恐惧。但当从那一方的约束消除后，他仍保持同样得体的行为，那么所有人也会将他在早年表现出的清醒归功于他。同样，让我们也这样做：让我们继续处于同样的敬畏状态，以便我们也能因先前的勤勉而从天主那里获得许多赞美。\par

3. 我们曾预料无数的灾祸：财产被掠夺，房屋连同居民一起被烧毁，城市从世间连根拔起，其碎片彻底毁灭，其土地置于犁下！但看！这一切只存在于预料中，并未实际发生。不仅如此，天主不仅除去了如此巨大的危险，还大大降福了我们，装饰了我们的城市，并通过这一试炼和灾难使我们更加受到赞许！但这些如何发生，我将说明。当皇帝所派的人设立那可怕的审判庭，调查所发生的事件，并传唤每个人为他们的所作所为交账，各种死亡的预想弥漫于众人心中时，那些住在山顶上的隐修士展现了他们真正的智慧。尽管他们在自己的斗室中隐居了多年，却未经任何人请求，无人劝告，当他们看见如此乌云密布在城市上空时，便离开洞穴和茅屋，从四面八方聚集而来，仿佛从天而降的许多天使。那时，人们可以看到城市如同天堂，这些圣人无处不在；仅凭他们的面容就安慰了哀悼者，并带领他们完全无视这场灾难。因为看见他们的人，谁会嘲笑死亡，谁会轻视生命？不仅如此，令人惊叹的是，当他们走近审判官本人时，他们为被告仗义执言，全都准备流血，牺牲性命，以便从他们所期待的可怖事件中夺回被捕者。他们还声明，除非审判官宽恕城市的居民，否则他们不会离开，或者将他们与被告一同送到皇帝面前。他们说：\Quote{统治我们这方世界的人是一位虔诚的人，一位信徒，一位活在虔诚实践中的人。因此，我们一定会说服他。我们不会允许你，也不会准许你动刀或砍头。但如果你不停止，我们也决心与他们同死。我们承认所犯的罪非常严重；但这些行为的邪恶并不超越皇帝的人性。}据说其中一人还说了另一句充满智慧的话，大意如下：\Quote{被推倒的雕像又重新立起，恢复了它们本来的样子，损害迅速得到了纠正；但如果你杀死天主的肖像，你如何能再次撤消此事？你如何能使那些被剥夺生命的人复活，并让他们的灵魂回归身体？}他们还向他们说了许多关于审判的事。\par

4. 谁能不感到惊讶？谁能不钦佩这些人的道德智慧？当其中一位被告的母亲，披散着头发，露出灰白的头发，抓住法官的马缰绳，跟着他跑过广场，就这样和他一起进入了审判的地方，我们都惊讶不已，我们都钦佩她那无比的温柔和宽宏大量。那么，我们难道不更应该对这些人的行为感到惊叹吗？因为即使她为自己的儿子而死，也没有什么奇怪的，因为自然的专制是强大的，而源于分娩之痛的义务是不可抗拒的！但这些人是以如此的爱去爱那些他们未曾生育、未曾抚养的人，是的，甚至他们从未见过、从未听说过、从未相遇过、仅从他们的灾难中认识的人，以至于如果他们有一千条命，他们也会选择全部献出以拯救他们。不要告诉我他们没有遭到屠杀，没有流血，而是他们对法官表现出的勇气，是其他任何人都不可能有的，除非是那些已经放弃了自己生命的人；而且带着这种心情，他们从山上跑到法庭。事实上，如果他们事先没有准备好面对各种屠杀，他们就不可能如此自由地向法官说话，或表现出如此的宽宏大量。因为他们整天坐在审判所门口，准备从刽子手手中夺走那些即将被带去受罚的人！\par

5. 现在那些穿着破旧斗篷、留着长胡子、右手拄着拐杖的人在哪里呢？那些世俗的哲学家，他们的性情比桌子底下的狗还要卑贱，所做的一切都是为了肚腹？所有这些人当时都抛弃了城市，都匆忙逃走，躲进了洞穴！但只有这些人，他们真正用行为证明了热爱智慧的人，毫无畏惧地出现在广场上，仿佛没有灾难降临城市一样。城里的居民逃到了山上和荒野，但荒野的居民急忙进了城；他们用行动证明了我在前几天不断说的话，那就是：即使火炉也不能伤害过着德性生活的人。这就是灵魂的哲学，它超越一切，超越所有顺境或逆境；因为前者不能削弱它，后者也不能击倒或贬低它，而是在整个过程中保持同一水平，展示其自身本来的力量和权威！的确，谁没有被当前的危机所暴露出的软弱所定罪？那些在我们城中担任最高职务的人，那些掌权的人，那些拥有巨大财富的人，那些深得皇帝宠信的人，完全抛弃了他们的房屋，各自寻求安全，所有的友谊和亲属关系都变得毫无价值，那些他们以前认识的人，在这灾难的时刻，他们不想认识，并且祈祷不被他们所认识！但隐修士们，虽然贫穷，除了一件破旧的衣服外一无所有，生活最为粗陋，以前似乎是无名小卒，习惯于山岭和森林，却像许多狮子一样，怀着伟大而崇高的灵魂，当所有人都恐惧战栗时，他们挺身而出，消除了危险，而且不是在多日之内，而是在短暂的一瞬间！就像杰出的战士，不与敌人近距离交战，只是出现在队伍中，呐喊一声，就使敌人溃败一样，这些人也在一天之内下山，说了他们该说的话，消除了灾难，然后回到了自己的帐篷。这就是基督带给人类的道德智慧有多么伟大。\par

6. 我何必提到富人和那些掌权者呢？当那些被授权审判罪犯的人，那些享有最高权威的人，被这些隐修士恳求准许赦免时，他们说他们对结果无能为力；因为不仅侮辱皇帝是危险和不安全的，而且释放那些已被抓获的侮辱者而不加惩罚也是不安全的。但这些人强大到无人能抗拒；他们以宽宏大量和坚持不懈围困他们，通过他们的恳求，使这些官员行使了他们未曾从皇帝那里得到的权力；甚至成功地说服了法官，当人们已被明显判明有罪时，不要宣判有罪，而是将最终结果推迟到皇帝的决定；他们保证一定能说服皇帝赦免那些冒犯他的人；他们即将启程前往他那里。但法官们敬畏这些人的道德智慧，被他们崇高的精神所打动，没有让他们进行这次长途旅行，而是承诺只要他们写下书面的话，他们就会亲自出发，成功地恳求皇帝平息一切愤怒（我们现在正期待他这样做）。因为当判决本应作出时，他们被允许进入法庭，说出了最高智慧的话语，并通过书信恳求皇帝施以仁慈；他们提醒他最后的审判，并说如果他的仁慈不被恩赐，他们将献出自己的头颅。法官们将这些话记录下来，然后离开了。这比最灿烂的冠冕更会装饰我们的城市。这里所发生的事，皇帝现在将会听到；是的，伟大的城市将会听到，全世界都会听到，居住在\Place{安提约基雅}城的隐修士们，是展现了宗徒般勇气的人；现在，当他们的信件在朝廷上被宣读时，所有人都会敬佩他们的宽宏大量；所有人都会称我们的城市有福；我们将摆脱我们的坏名声；并且将在各处皆知，所发生的事不是本市居民所为，而是陌生人和心地败坏的人做的；隐修士们的这个见证，将是城市品格的充分证据。\par

7. 所以，亲爱的诸位，我们不要忧愁，而要怀抱美好的希望；因为如果他们对人的勇敢能够避免这样的危险，那么他们对天主的勇敢将不会产生什么效果呢？当希腊人胆敢与我们争论他们的哲学家时，我们也应当把这些事告诉他们！由此可见，他们古代的传说都是假的，而我们中间所传的古代之事——即关于若望、保禄、伯多禄以及其余一切人的事——才是真的。因为这些隐修士继承了那些人的虔诚，所以他们也彰显了同样的勇敢。既然他们受教于同一法律，也就自然效法了他们的德行。因此，我们无需借助著作来显示德行，因为事实本身就在大声宣告，师傅由门徒彰显出来。我们无需辩论来揭露希腊人的虚妄和他们的哲学家的狭隘，因为他们的行为现在正像从前一样大声宣告：他们的一切都是寓言、戏剧、表演。\par

8. 同样的宽宏大量也显示在司铎和隐修士身上，他们共同分担了我们安全的责任。其中一位司铎前往宫廷，把一切视为次于对你们的爱；他本人已准备好，如果不能说服皇帝，就舍命以赴。而那些留在此地的人，也展现了与隐修士们同样的德行；他们亲手拉住法官们，不肯让他们进入法庭，直到法官们对审判结果作出承诺。当他们看到法官们做出拒绝的表示时，又再次以极大的勇气努力争取；而一旦看到他们同意，便抱住他们的脚和膝，亲吻他们的手，从而对两种德行——自由与温良——给出了极大的证明。因为他们那不是妄自尊大的勇敢，他们亲吻法官的膝盖、抱住他们的脚，便清楚地表明了这一点。同时，这并不是谄媚，也不是一种谄媚的奴性，更不是奴隶精神的结果，他们先前的大胆行为已经证明了这一点。我们从这个考验中不仅收获了这些益处，还收获了大量的清醒和温良；我们的城市瞬间变成了一座隐修院。即使有人在广场上竖起金像，也无法像现在这样装饰和荣耀它，因为它产生了这些美丽的德行之像，并展示了它真正的财富！\par

9. 或许皇帝所颁布的命令令人痛苦。不！这些命令其实一点也不沉重，反而带来了许多益处。因为请问，其中哪一项是难以忍受的呢？皇帝关闭了乐池，禁止了赛马场，封闭并堵塞了这些邪恶的源泉。愿它们永不再开启！从那里滋生出危害城市的恶行之根；从那里产生了败坏城市名声的人——那些向舞者出卖自己声音的人，为了三个小银币而出卖自己的得救，把一切搞得天翻地覆！亲爱的，你为这些事忧愁吗？其实你应当为此欢喜快乐，并向皇帝表示感谢，因为他的责罚成了纠正，他的惩罚成了训练，他的愤怒成了教训！至于澡堂关闭了？这也不是不可忍受的艰难，因为那些过着安逸、柔弱、放纵生活的人，虽然不情愿，也被带回到对真正智慧的热爱。\par

10. 有人抱怨说，皇帝剥夺了这座城市的尊严，不再允许它被称为大都市吗？但他又能做什么呢？他难道能赞美所发生的事，并承认这是恩惠吗？那样谁会不责备他，连愤怒的外表都不表露出来呢？你们没有看见父亲们对子女行许多类似的事吗？他们离开他们，禁止他们入席。皇帝也做了同样的事，施加了这些没有伤害性、却带有许多矫正的惩罚。想想我们所期望的和所发生的事，然后我们便会特别看出天主的恩宠！你们为城市的尊严被剥夺而悲伤吗？要明白城市的尊严是什么；然后你们就会清楚知道，如果居民不背叛它，就没有人能夺走一座城市的尊严！不是因为它是一个大都市；也不是因为它包含巨大而美丽的建筑；也不是因为它有许多圆柱、宽敞的柱廊和步道，也不是因为它在公告中先于其他城市被提及，而是居民的德行和虔诚；这才是城市的尊严、装饰和防卫；因为如果这些不在其中，那么它便是世上最微不足道的，即使它享有皇帝无上的尊荣！你们想了解你们城市的尊严吗？你们想知道它的祖先吗？我将准确讲述；不仅使你们知道，也使你们效法。那么，我们这个城市的尊严究竟是什么呢？\Quote{这事发生，门徒们在安提约基雅最先被称为基督徒。}\BibleRef{宗 11:26} 这个尊严，普世的城市没有一座拥有，甚至\Person{罗慕路斯}本人的城市也没有！因为这一点，它可以面对整个世界；由于那种对基督的爱，那种勇敢和德行。你们还想听属于这座城市的另一种尊严和称誉吗？一次严重的饥荒即将来临，\Place{安提约基雅}的居民决定，按每人能力，向住在\Place{耶路撒冷}的圣人们提供救济。\BibleRef{宗 11:28} 看，第二个尊严，饥荒时期的爱德！时令没有使他们吝啬，对灾难的预期也没有使他们退缩；但当所有人都倾向于搜刮不属于自己的东西时，他们却分施自己的，不仅给近处的人，也给远处居住的人！你们在这里看到对天主的信德和对邻人的爱德了吗？你们想了解这座城市的另一个尊严吗？有些人从\Place{犹太}下来到\Place{安提约基雅}，败坏所宣讲的道理，并引入犹太人的惯例。\Place{安提约基雅}的人没有沉默地忍受这种新奇。他们没有保持缄默，而是聚集起来，开了会议，派遣\Person{保禄}和\Person{巴尔纳伯}上\Place{耶路撒冷}，使宗徒们规定，纯洁的道理，清除了所有犹太人的缺陷，可以分送到世界各地！这就是城市的尊严！这就是它的优先地位！这使它成为大都市，不是在地上，而是在天上；因为其他所有尊荣都是可朽坏的、短暂的，与现世生命一同消亡，并且常常在生命结束之前就终结，就像现在所发生的那样！在我看来，一个没有虔诚居民的城市比任何一个村庄都更卑贱，比任何一个洞穴都更可鄙。\par

11. 我为什么提到城市呢？为了你们能确切明白，只有德行才是居民的装饰，我不再对你们谈城市，而是尽力通过提出比任何城市更可敬的事物来证明这一点——即在\Place{耶路撒冷}的天主圣殿。因为这座圣殿里有祭献、祈祷和礼仪；那里有至圣所、革鲁宾、盟约和金罐；是天主对那民族眷顾的伟大象征；那里不断接收到来自天上的神谕，先知在那里受默感，建造不是人工的成果，而是出自天主的智慧，墙壁四周闪耀着大量黄金，在那里，卓越超凡，材料的昂贵和艺术的完美结合在一起，证明地上没有像这样的另一座圣殿！是的，不仅如此，不仅是艺术的完美，而且天主的智慧也协助了那建造。因为\Person{撒罗满}学到这一切，不是凭直觉和靠自己，而是来自天主；从天上领受了设计后，他便规划并建造了它。然而，这座如此美丽、奇妙和神圣的圣殿，当使用它的人败坏时，竟被如此羞辱、轻视和亵渎，甚至在充军之前就被称为\Quote{强盗的巢穴，鬣狗的洞穴；}\BibleRef{耶 7:11} 此后，它被交到野蛮、污秽和亵渎的手中！\par

12. 关于城市，你愿意学习这同样的事实吗？有哪座城市比索多玛更辉煌呢？房屋和建筑都很壮观，城墙也是如此；土地肥沃富饶，\BibleQuote{像天主的乐园。} \BibleRef{创 13:10} 但\Person{亚巴郎}的帐幕简陋狭小，没有防御工事。然而当外敌战争发生时，那些异邦人攻破了设防的城市，带着居民俘虏离去。然而，\Person{亚巴郎}，这个旷野的居民，当他们攻击他时，他们无法抵挡！事实本该如此。因为他拥有真正的虔诚：一种比数量和城墙的防御更强大的力量。如果你是基督徒，没有一座地上的城市属于你。对我们的城市而言，\BibleQuote{建设者和创造者是天主。} \BibleRef{希 11:10} 即使我们获得了整个世界，我们在其中仍然是客旅和寄居者！我们已被登记在天上：我们的公民身份在那里！让我们不要像小孩子一样，轻视伟大的事物，而羡慕渺小的事物！不是我们城市的伟大，而是灵魂的德行，才是我们的装饰和防御。如果你以为尊严属于一座城市，想想有多少人必须分享这种尊严，他们却是淫乱者、柔弱者、堕落者、充满千万邪恶之事的人，最终会轻视这样的荣誉！但上面的那座城市不是这样的；因为凡是没有表现出一切德行的人，是不可能分享它的。\par

13. 因此，我们不要愚昧；但当有人剥夺我们灵魂的尊严时，当我们犯罪时，当我们得罪了共同的万物之主时，让我们忧伤；因为关于现时所发生在我们身上的事，它们不但没有伤害这座城市，反而如果我们警醒，将大大有益于我们。因为就连现在，我们的城市似乎像一位端庄、高贵、稳重的夫人。恐惧使她变得温和而更显庄重，并且把她从那些参与最近胆大妄为行为的恶徒手中解救出来。因此，我们不要像妇人那样哀叹。因为我听到许多人在广场上说：\Quote{哀哉！你，安提约基雅！你遭遇了什么！你如何被羞辱！} 当我听到时，我嘲笑那能说出这些话的幼稚心智！这样的话现在并不合适；但当你们看到人们跳舞、醉酒、歌唱、亵渎、诅咒、发誓、作假见证、说谎时，那时应用这样的话：\Quote{哀哉！你，城市啊，你遭遇了什么！} 但如果你们看到广场上有一些温和、谦逊、节制的人，那么就称这城为\Quote{有福的！} 因为稀少绝不会在任何方面伤害它，如果同时拥有德行；另一方面，纵然有众多人口，只要有邪恶，就毫无益处。\Quote{如果}先知说，\Quote{以色列子孙的数目如同海沙，剩下的部分将得救；}这就是说，\Quote{众多人口绝不会在我面前得胜。} 基督也是如此说的。他称城市为可悲的；不是因为它们的渺小，也不是因为它们在城市等级上不是首府。而耶路撒冷本身，他又因同样的理由称它为可悲的，这样说：\BibleQuote{耶路撒冷，耶路撒冷，你杀死先知，又用石头砸死那些被派到你这里来的人！} \BibleRef{玛 23:37} 我问，如果他们的生活方式是邪恶的，众多人口有什么好处呢？不，相反，甚至带来损害。还有什么造成了近期涌现的邪恶？岂不是居民的懒惰、鲁莽和堕落吗？城市的尊严、建筑的宏伟，以及作为首府的事实，对它有任何帮助吗？如果对于地上的君王，当它做了如此错事时，没有什么能保护它，反而所有这些特权都被剥夺了；那么对于天使的主，它的尊严怎能保护它呢？因为在那一日，我们曾居住在一个有许多宽敞柱廊和其他此类尊严的首府，这对我们毫无益处！我为什么说那一日呢？因为就现世生命而言，你的城市是首府对你有什么益处呢？请问，有人借此恢复了一个困苦的家庭吗？或者从这种尊严中获得任何收入吗？消除了忧伤吗？摆脱了任何身体的病痛吗？除去了灵魂的恶习吗？亲爱的！让我们不要轻浮，也不要看重群众的意见，而应理解什么是城市的真正尊严；是什么使一座城市真正成为首府。\par

14. 我说这一切，尽管我期望这城市将再次恢复这外在的荣誉，并出现在它自己应有的优先位置上。因为皇帝既是仁爱的，又是虔诚的。但我希望，如果它被恢复，你们不要对此过分看重，也不要因此自夸，更不要将我们城市的荣耀归于那上面。当你们想赞美这座城市时，不要对我提起达佛涅的郊区，也不要提起柏树的高度和数量，不要提起它的泉水，不要提起居住在这城中的众多人口，不要提起它的市场甚至到半夜都频繁出入的巨大自由，也不要提起货物的丰盛！这一切都是外在感官的事物，只存留于现世。但如果你们能提及德行、温良、施舍、夜间警醒、祈祷、节制、灵魂的真智慧，那就为这座城市称赞这些事吧！对于住在旷野的人来说，这些事物的临在使旷野比任何城市都更辉煌；反之，如果这些事物不在市民中显明，那么最可鄙的地方也是如此。让我们不仅在城市的事上，也在人的事上作此判断。如果你看见一个高大的人，被养得很好，身材高大，在四肢长度上超过别人，不要羡慕他，直到你确知这人的灵魂是怎样的。不该从外表的俊美，而应从属于灵魂的美善，才称任何人为有福的！达味身材矮小，体形短小；然而这样一个短小的人，赤裸着，没有任何武器，一击就击倒了那么庞大的一支军队，以及那座血肉之塔；而且这不是用投掷长矛、放箭或拔剑，而是用一块小石子完成这一切！因此，有一个人劝告说：\Quote{不要因人的美貌而称赞他，也不要因人的外貌而厌恶他。蜜蜂在飞禽中是最小的，但它的果实是甘甜之首。}见：\BibleRef{德 11:2}\par

15. 因此，让我们也这样谈论一座城市和人，互相说出这样的智慧，不断为现在和过去的恩惠感谢天主；并共同竭尽全力呼求祂，使那些现在被囚禁的人得释放，使那些即将被流放的人得以返回。他们也是我们的肢体。他们曾与我们一同搏击风浪，与我们一同抵挡风暴！让我们恳求仁慈的天主，使他们与我们一同享受平静！谁也不要说：\Quote{这与我还有什么关系？我已经脱离了危险；某人可以灭亡；某人可以被毁灭！}不要因这种冷漠而激怒天主；而要哀悼，好像我们自己处在同样的危险中。让我们以极大的热忱恳求天主，履行保禄的这句话：\Quote{你们要记念被囚禁的人，好像与他们一同被囚禁；记念受苦难的人，好像你们也亲身在肉身内。}\BibleRef{希 13:3} \Quote{与哭泣的人一同哭泣，屈就于卑微的人。}见：\BibleRef{罗 12:15} 这对我们自己也有最大的益处；因为没有什么比我们非常愿意为自己的肢体哀痛更常使天主喜悦的了。因此，让我们共同为现在和将来的事恳求祂，使祂能拯救我们脱离日后的惩罚。因为现在的事，无论怎样，都是可忍受的，且有终结；但那里的刑罚是不死的、无止境的！我们既得了安慰，也要亲自努力不再陷入这样的罪，知道今后我们不会再得赦免！让我们全体一同俯伏在天主面前；无论在这里或在家里，让我们说：\Quote{上主，在祢所做的一切事上，祢是正义的；因为祢以公义的审判，将祢所带来的一切都加给了我们。}\BibleRef{厄下 9:33} \Quote{如果我们的罪恶起来攻击我们，求祢为祢圣名的缘故为我们作主；}\BibleRef{耶 14:7} 不要再容许我们经历如此痛苦的事。不要让我们陷入诱惑，但救我们脱离凶恶，因为国度、权能和荣耀，都属于祢，直到永远。阿们。\par

\SourceRecord{Translated by W.R.W. Stephens.；Nicene and Post-Nicene Fathers, First Series；Vol. 9.；Edited by Philip Schaff.；Buffalo, NY: Christian Literature Publishing Co.,；1889.；<http://www.newadvent.org/fathers/190117.htm>.}

% structure: p0001 source=h1 target=section reason=page-title

\section{论雕像讲道集 第十八篇}

\Emph{继续先前关于叛乱的论述；也论及禁食；并关于宗徒的训导：\Quote{在主内常常喜乐。}}\par

1. 我注意到许多人在欢喜，彼此说：\Quote{我们胜利了；我们得胜了；封斋期已过了一半。}但我劝告这样的人不要因为封斋期已过一半而欢喜，倒要省察自己的罪过了半没有；若是如此，然后才可欢跃。因为这是值得满足的题目。这正是所追求的，也是所有事情为之而做的，为使我们能改正我们的缺点；并且使我们进入封斋期时是什么样的人，离开时不再是同样的人，而是处于洁净的状态；并且抛掉一切属于恶习的东西后，我们就能如此遵守神圣的庆节。因为如果情况相反，我们非但得不到任何益处，封斋期的完成反而会成为我们最大的损害。所以，我们不要因为度过了封斋期的长度而欢喜，这不算什么大事；但我们要欢喜，如果我们以新的收获度过了它，以致当这结束之时，它的果实能显现出来。因为冬季的收益更特别地在其季节过去之后才彰显出来。那时，茂盛的谷物和缀满叶子与果实的树木，以其景象宣告了它们从冬季所获得的益处！同样的事也要在我们身上发生。因为在冬季里，我们享受了各种各样频繁的雨水，在封斋期间领受了持续不断的教导，并接受了属灵的种子，砍掉了奢侈的荆棘。\par

2. 因此，让我们坚持到底，以最大的勤谨保留我们所听到的，以便封斋期结束时，封斋的果实能丰盛，并且藉着我们从封斋期收集的美好事物，我们能记念这封斋期本身。如果我们这样塑造自己，当封斋期再来时，我们将再次乐意欢迎它。因为我看到许多人心智如此软弱，以致在当前的季节，他们为下一个四旬期忧虑；并且我听到许多人说，在从封斋期解脱之后，由于对来年的忧虑，他们对这缓解感觉不到任何快乐。有什么比这更软弱的呢？我问；原因是什么？这是因为当封斋期来到时，我们并不努力使灵魂的事务得到良好的安排，而将封斋期仅限于戒除食物。因为如果我们能在行为的改造中充分受益，我们就会希望封斋期每天都来临，真正体验它的好处；而且我们永远不会抛弃对它的渴望，也不会在期待它时沮丧和焦虑。\par

3. 因为没有任何事物能使一个心灵有序、关心自己灵魂的人受苦；他反而会享受纯洁而持续的快乐。这是真实的，你们今天从保禄听到了，他劝勉我们说：\Quote{在主内常常喜乐，我再说，喜乐。}\BibleRef{斐 4:4}我的确知道，对许多人来说，这句话似乎不可能。\Quote{因为这是如何可能的呢？}有人说，\Quote{他只不过是一个人，怎能常常喜乐呢？}因为有许多悲伤的原因，从四面环绕我们。一个人可能失去了儿子、妻子、或心爱的朋友，比所有亲戚更亲近他；或者他遭受了财富的损失；或者陷入了疾病；或者要承受其他命运的改变；或者为不应得的轻蔑待遇而悲伤；或者饥荒、瘟疫、某些难以忍受的勒索、家庭中的情况困扰他——不，无法说清有多少公共或私人的情况习惯于使我们忧愁。那么，他怎能常常喜乐呢？是的，人啊！这是可能的；若不然，保禄就不会给予这劝勉；一个拥有属灵智慧的人也不会提供这样的忠告；正因为此，我常对你们说，并且将不停地说，你们不能从别处学到的智慧，你们可以在此默想。因为人类普遍渴望快乐和喜乐；为此，他们做一切事，说一切话，承担一切事。所以商人航海，为了积累财富；他积累财富，为了在他所储存的东西上喜乐。士兵也因此作战，农夫也因此耕作；为此每个人从事其技艺。那些喜爱统治的人，也为此目的而喜爱统治，为了获得荣耀；他们渴望获得荣耀，为了喜乐；任何人都可以察觉到我们的每项事业都指向这一点，并且每个人看着这一点，通过各种方式匆忙走向它。\par

4. 正如我所说，所有人都喜爱喜乐，但并非所有人都能获得它，因为他们不知道通往喜乐的道路。许多人以为喜乐源于富有。但若真是如此，就没有一个富人会是忧伤的。然而事实上，许多富人认为生命不值得活，一旦遭遇任何困苦，便宁愿死去；在所有的人中，他们最容易陷入过度的忧伤。你们不应只看到他们的筵席、奉承者和寄生虫，而应看到由此而来的烦恼——侮辱、毁谤、危险和痛苦，更糟的是，他们毫无准备地面对这些逆境，不知道如何以哲学家的态度承受，或如何勇敢地承担所发生的事。因此，灾祸在他们眼中并不以其本性显现，甚至连真正轻微的事情也变得难以忍受。反之，在穷人身上则相反：不可挽救的事似乎容易忍受，因为他们对此已习以为常。因为使降临在我们身上的邪恶显得或大或小的，与其说是事件的本性，不如说是受苦者的心态。为了不远离实例，我要对你们讲讲最近发生在我们身上的事。请看，所有的穷人都逃过了，民众脱离了危险，享受完全的自由！但那些管理城邦事务的人，那些养马并主持公共竞技的人，以及那些担任过其他公职的人，现在却成了监狱的囚徒，恐惧着最坏的结果；他们独自为所有人所犯的罪行付出代价，时刻处于恐怖之中；他们现在是世上最悲惨的人，并非因为危险之大，而是因为他们迄今所过的奢侈生活！至少，许多人在被我们劝勉和鼓励要勇敢地承受这些逆境时，曾这样说：\Quote{我们从未练习过这种事，也不知道如何实践这样的哲学；这就是为什么我们需要如此多的安慰。}\par

5. 另一些人则认为，拥有健康是快乐的源泉。但事实并非如此。因为许多健康的人曾千百次想死，无法忍受他人加给自己的侮辱。还有人断言，享有荣耀、获得权力、担任最高官职、受众人奉承，能带来持续的喜乐。但事实也并非如此。我何必提其他权力职务呢？因为即使我们思想上升至王权本身，到达那居于王位者，我们也会发现王权被各种烦恼所包围，且因其事务更重，忧伤的必要原因也更多。何需谈论战争、战事和蛮族的叛乱呢？他常常有理由惧怕身边的人。因为许多君主虽从敌人手中逃脱，却未能逃脱自己侍卫的阴谋。君王们必然拥有像海洋的波涛一样多的忧伤原因。但如果君主制不能使生活脱离悲伤，那么还有什么能做到呢？的确，今世之物无一能做到；唯独\Person{保禄}的这句话，简短而朴素，却能亲自为我们开启这宝藏。\par

6. 因为不需要许多言语，也不需要长篇大论的论证，只要我们思考他的话，就会发现通往喜乐的道路。他并没有简单地说：\Quote{要常常喜乐}，而是加上了持续喜乐的原因，说：\Quote{你们要常在主内喜乐}。那\Quote{在主内}喜乐的人，任何事都不能剥夺他的喜乐。因为其他一切让我们喜乐的事物都是可变的、易变的、会变化的。不仅如此，这些事物还存在一个可悲的情况：即使它们存留，也不能给我们足够的喜乐来抵抗和掩盖从别处而来的悲伤。但对天主的敬畏同时具备这两个条件。它是稳固不动摇的，并且倾注如此多的喜乐，以至于我们感觉不到其他邪恶。因为那按应当的方式敬畏天主并信赖祂的人，从喜乐的根本处汲取，并拥有整个喜乐的源泉。就如一颗火花落入广阔的海洋中，很快消失；同样，任何发生在敬畏天主者身上的事件，落在巨大的喜乐海洋中，都会被熄灭和摧毁！最令人惊叹的正是这一点：当引发悲伤的事物存在时，这人仍能保持喜乐。因为如果没有任何产生悲伤的事，那么他能常常喜乐并不算什么大事。但在许多事情的压力下催迫他悲伤时，他却超越这一切，在悲伤中仍欢欣，这才真正令人惊奇！正如无人会惊讶三位青年没有被烧死，如果他们远离巴比伦的火窑——因为使众人惊讶的是，他们如此长时间地与火紧密接触，却比那些未接触火的人更毫发无损；同样，我们也能论及圣人们：如果没有诱惑附着他们，我们就不会对他们持续不断的喜乐感到惊奇。但值得赞叹且超越人性的要点是：他们被无数波浪四面环绕，他们的处境却比那些享受完全平静的人更加安逸！\par

7. 由以上所述可见，在教会之外的人中，不可能找到任何一种生活方式是始终被外在事物的持续快乐所环绕的。但是，我现在要证明，信徒绝不会被剥夺持续喜乐的享受，为的是你们不仅能学习，还能效法这种无痛苦的生活状态。因为假设一个人没有什么可自责的，却保有良好的良心，渴望未来的状态和那些美好盼望的实现；我问，有什么能使这样的人忧伤呢？死亡不是最难以承受的事吗？然而对死亡的期待不但不会使他悲伤，反而使他更加喜乐；因为他知道死亡的到来是从劳苦中的释放，是朝着为那些在虔敬和德行竞赛中奋战的人所存留的冠冕和赏报飞奔。但他的子女夭折呢？不，他也会高尚地承受，并引用约伯的话：\BibleQuote{主给了，主又收回；照主看来是好的，就这样成就了。愿主的名永受赞美。}\BibleRef{约 1:21} 但如果死亡和丧子不能使人忧伤，那么金钱的损失、羞辱、责骂、诬告，就更不能在任何时候影响如此伟大而高贵的灵魂；不，连身体的痛苦也不能，因为宗徒们曾受过鞭打，却没有忧伤。这确实是一桩大事；但更了不起的是，他们非但没有忧伤，反而将鞭打本身视为额外喜乐的理由。\BibleQuote{他们离开公议会，心里欢喜，因被算是配为基督的名受辱。}\BibleRef{宗 5:41} 有人侮辱和谩骂这样的人吗？他受基督教导，要因这些辱骂而喜乐。祂说：\BibleQuote{欢喜快乐吧，当人们为了我的缘故，以各种恶言诬蔑你们时，并且要极其欢喜，因为你们在天上的赏报是大的。}\BibleRef{玛 5:11} 假定有人患病呢？他听过他人的劝诫说：\BibleQuote{在疾病和贫穷中，你要信赖祂；因为如同黄金在火中受试炼，蒙悦纳的人也在卑微的炉中受锻炼。}\BibleRef{德 2:4} 因此，既然死亡、金钱损失、身体疾病、羞辱、责骂，以及任何其他类似的事物都不能使他忧伤，反而使他更加喜乐，那么他任何时候又有何忧伤的根据呢？\par

8. 有人说：\Quote{那么怎样？}圣人岂不是常怀忧愁吗？你们没有听见保禄说：\BibleQuote{我大有忧愁，心里时常伤痛。}\BibleRef{罗 9:2} 这确实是奇妙之事：忧愁带来了利益，以及由利益产生的喜乐；因为正如鞭打没有给他们带来痛苦，反而带来了喜乐；同样，忧愁也给他们带来了那些伟大的冠冕。这就是矛盾之处：不仅世俗的忧愁，连世俗的欢乐也包含着巨大的损失；但在属灵的事上，恰恰相反；不仅喜乐，连忧愁也包含着丰富的美善宝藏！现在我来解释这是如何发生的。在世界上，人常常因见到仇敌遭难而喜乐；因这喜乐，他为自己招致了重大的惩罚。另一个人因见到弟兄跌倒而哀伤；因这哀伤，他将为自己获得天主极大的恩宠。你们看见合乎天主心意的忧愁如何比世俗的喜乐更美好、更有益吗？同样，保禄也为罪人和不信天主的人忧愁；这忧愁成了为他积蓄大赏报的方法。但为使我说得更清楚，并使你们知道虽然我所说的非常奇怪，却是真实的，即：悲伤常常能更新愁苦的灵魂，使沉重的良心轻省：试想，妇女们在失去最心爱的儿女时，如果被禁止哀悼和流泪，往往会心碎而亡。但如果她们做出所有悲伤之人惯常做的事，就会得到缓解和安慰。妇女尚且如此，当你们看见一位先知也以类似方式受影响时，又有何奇怪呢？因此他不断地说：\BibleQuote{容我痛哭吧！不要劳神安慰我，因我的人民女儿遭了蹂躏。}\BibleRef{依 22:4} 所以，悲伤常常是安慰的载体；如果对于今世尚且如此，对于属灵的事更是如此。因此他说：\BibleQuote{合乎天主心意的忧愁能产生令得救的悔改，不再后悔。}\BibleRef{格后 7:10} 这似乎令人费解；但他所说的是：\Quote{若你为财富而忧愁，一无所益；若为疾病而忧愁，也一无所获，反而增加了你的痛苦。}\par

9. 我听见许多人在经历这些事之后责怪自己说：\Quote{我悲伤有何益处呢？我既没有追回钱财，又伤害了自己。}但若你因罪而悲伤，你便涂抹了它，收获了最大的喜乐。若你为跌倒的弟兄而悲伤，你既鼓励安慰了自己，又复兴了他们；即使你未能使他们受益，你也有丰厚的赏报。为使你懂得：为跌倒者悲伤，即使我们丝毫不能帮助他们，仍带给我们巨大的报酬，请听\Person{厄则克耳}所说的——更确切地说，是天主亲自藉他所说的。因为当祂派遣使者去倾覆那城，用刀剑和火毁灭所有房舍及其居民时，祂这样吩咐其中一位说：\BibleQuote{要在那些叹息悲哀的人额上画一个记号。} 又吩咐其他人说：\BibleQuote{从我的圣者开始吧。} 接着祂补充说：\BibleQuote{凡有这记号的，不可挨近他们。} \BibleRef{则 9:4} 为什么呢？告诉我。因为他们虽然无济于事，仍为所行之事哀叹悲伤。再者，祂指责另一些人说：他们耽于安逸、贪食、享尽安稳，看见犹太人被掳去时，既不悲伤，也不分担他们的忧愁。请听祂责备时所说的话：\BibleQuote{他们毫不关心\Person{若瑟}的苦难。} \BibleRef{亚 6:6} 这里的\Person{若瑟}是指全体百姓。又说：\BibleQuote{\Place{厄南}的居民没有出来哀悼邻舍的家属。} 因为他们虽受公义的惩罚，天主仍愿我们同情他们，不可幸灾乐祸或侮辱。\BibleQuote{我虽惩罚，}祂说，\BibleQuote{不是喜乐地做，也不以他们的惩罚为乐；}因为\BibleQuote{我决不愿意罪人丧亡；} \BibleRef{则 18:32} 你理当效法你的主，正为此事而哀恸：罪人提供了公义惩罚的理由与契机。因此，若有人怀有合乎虔敬的忧愁，他必将从中获得极大的益处。\par

10. 既然受鞭打的人比鞭打者更有福，我们中信德之内的受难者比信德之外安逸的人更有福，忧愁的人比享乐的人更有福，那么，我们还有什么患难之源呢？因此，我们不应称任何人有福，唯独那按天主生活的人有福。圣经只称这些人为有福。\BibleQuote{有福的人，}经上说，\BibleQuote{是不随从恶人的计谋的人。有福的人，是祢所惩戒、教以祢律法的人。行为无瑕的人是有福的。凡信赖祂的人是有福的。以天主为天主的人民是有福的。良心不谴责自己的人是有福的。敬畏上主的人是有福的。} 再者，基督又说：\BibleQuote{哀恸的人是有福的，谦卑的人是有福的，温良的人是有福的，缔造和平的人是有福的，为义而受迫害的人是有福的。} \BibleRef{玛 5:3} 你看见了吗？神圣的律法处处宣告有福的，不是富人、出身高贵者或拥有荣耀的人，而是那握有美德的人。因为我们所要求的，是在一切所行或所受的事上，以敬畏天主为基础；你若以此作为根基，那么，安宁、尊荣、荣耀、关注，不仅会为你结出喜乐之果，而且连敌意、诽谤、轻蔑、羞辱、折磨以及一切事物，无一例外。正如树根本身是苦的，却结出最甘甜的果实，同样，合乎虔敬的忧愁必要带给我们丰厚的喜乐。那些曾以痛苦祈祷、流泪的人，知道他们收获了何等的喜乐，如何洁净了良心，如何怀着美好的希望起来！因为我常说：使忧愁或喜乐的，不是事物的本性，而是我们的心境。如果我们能使心境合宜，我们便拥有一切喜乐的保证。正如身体，损害或助益它的，不是空气的性质或外在遭遇，而是它自身的内在状况；灵魂也是如此，而且更加如此。因为一个受自然必然性支配，另一个则完全取决于意志的选择。因此，\Person{保禄}在忍受了无数恶事——海难、战争、迫害、阴谋、强盗袭击、数不胜数的事，且天天面对死亡——之后，非但没有忧愁或不满，反而夸耀、喜乐，且说：\BibleQuote{现今我在苦难中喜乐，并在我的肉身上补满\Person{基督}苦难所欠缺的。} \BibleRef{哥 1:24} 又说：\BibleQuote{不但如此，我们在患难中也夸耀。} \BibleRef{罗 5:3} 夸耀意味着喜乐的扩展。\par

11. 你若渴望喜乐，就不要追求财富、身体健康、荣耀、权力、奢华、丰盛的宴席、丝绸衣物、昂贵的土地、辉煌显赫的房舍，或任何此类事物；而要追求那合乎天主的属神智慧，握住美德；于是，现在或将来的事物，没有一样能使你忧愁。我说忧愁？诚然，那些使他人忧愁的事物，对你将成为喜乐的添加。因为鞭打、死亡、损失、诽谤、受虐待，以及诸如此类的事，当它们为天主之故而临到我们，且从这根基发出，必给我们的灵魂带来极大的喜乐。因为若我们不使自己成为不幸的，就没有人能使我们不幸；同样，若不使自己成为有福的——跟随天主的恩宠——也没有人能让我们有福。\par

12. 为使你明白惟独敬畏主的人是有福的，我现在就来向你证明这一点，不是藉着过去所发生的事，而是藉着发生在我们自己身上的事。我们的城市曾面临被彻底抹去的危险；富人、显贵、名流中，没有一个人敢出现在公众面前，反而都逃跑了，匆忙躲避。然而，那些敬畏天主的人，那些在隐修院里度过时光的人，却带着极大的胆量急忙下山，将众人从这恐惧中解救出来；所发生的可怕事件以及预期将付诸实施的威胁，非但没有使他们恐惧或焦虑，反而使他们虽远离灾难、未曾涉足其中，却甘愿投身火中，拯救了所有人；至于死亡——这普遍令人畏惧可怖的事——他们以最大的准备迎接它，比他人追求权位和荣誉更乐意地奔向它。这是因为他们知道：这才是最大的权位和荣誉。他们以实际行动表明：惟独把握从上而来的智慧的人是有福的，他不受变迁，不遇逆境，反而享有持续的安宁，并嘲笑一切看似悲伤的事。至少现在，那些一度掌权的人被深重的忧伤所压迫，囚禁在狱中，身负锁链，每日等待处决。但这些人却享有最纯粹的喜乐；即使他们注定要遭受可怕之事，这——以及他人看来可怕的事——对他们来说也是受欢迎的，因为他们深知自己奔向何方，以及离开这个世界后等待他们的命运。然而，尽管他们生活得如此严谨，对死亡报以微笑，他们仍为他人悲伤，并由此反过来获得最大的益处。那么，让我们认真关心我们的灵魂，这样任何意想不到的事都不能使我们忧伤。为那些在狱中的人，让我们恳求天主，求祂拯救他们脱离目前的灾难。因为天主原本可以立刻使我们脱离这可怕的灾祸，连一丁点痕迹都不留下；但为了使我们不再回到先前的懈怠，祂安排这恶事的洪流缓缓地、一点一点地消退，从而使我们坚守同样的虔诚决心。\newline{}\par

13. 这是真实的，而且如果立刻从整个困境中解脱，许多人会回到先前的懈怠，从以下情况可以显明：尽管灾难的余波仍在，皇帝的判决仍不确定，那些管理城市事务的人都在狱中，我们许多同城的居民却因过度渴望沐浴而跑到河边，在那里无尽地嬉戏、行为放荡、跳跃、跳舞，并拉着妇女跟在后面。这样的人配得什么宽恕？他们能提出什么藉口？或者，他们不该受怎样的惩罚和报应？城市的首领在公共监狱中；我们的成员在流放中；关于他们的判决仍不确定；而你，我问你，却跳舞、嬉戏、大笑？\Quote{我们无法忍受，}有人说，\Quote{不沐浴的生活？}无耻的心态，卑劣而败坏！我问你，多少个月，多少年过去了？你被禁止沐浴还不到二十天；你就如此痛苦不满，仿佛你整整一年没有洗澡！告诉我，当你预期军队的袭击时，当你每日等待被处死时，当你逃往旷野、奔向山顶时，你是否也这样？那时若有人提议你\Quote{一年}不沐浴，以求从迫在眉睫的苦难中获救，你难道不会欣然接受并服从吗？因此，当你本该感谢天主——祂使你毫无损失地脱离这一切——之时，你却再次变得狂妄自大；当恐惧过去之后，你回转到一个更糟的懈怠状态吗？这些可怕的事真的触及了你，然而你却如此渴望沐浴？如果沐浴被允许，那么那些仍在困禁中的人的灾难，难道不足以劝诫那些不在同样痛苦境况中的人忘却一切奢侈吗？生命本身处于危险之中，你还记得沐浴，并渴望奢侈？你因现已逃脱而轻视危险吗？当心，免得你将自己卷入更大惩罚的必然中，召唤那已移除的愤怒以更大的量返回，并经历基督论及魔鬼时所宣告的事。因为祂说：\BibleQuote{邪魔出去以后，又见房屋空着、打扫干净，就带了另外七个比自己更凶恶的魔鬼来，进入灵魂，那人最后的处境比先前更坏了。}因此，我们也当害怕，免得我们从先前的邪恶中解脱出来，后来却因懈怠招致更大的灾祸！我知道你们自己并不犯这种愚行；但你们应当约束、惩罚并使那些行为不检的人清醒，好使你们常能如\Person{保禄}所命令的那样喜乐，并因自己的善行和对别人的关怀，在此世和来生都享有丰厚的赏报；藉着我们的主耶稣基督的恩宠和慈爱，愿光荣、尊威和朝拜归于祂及圣父偕同圣神，自今至永远，世世无穷。亚孟。\newline{}\par

\SourceRecord{Translated by W.R.W. Stephens.；Nicene and Post-Nicene Fathers, First Series；Vol. 9.；Edited by Philip Schaff.；Buffalo, NY: Christian Literature Publishing Co.,；1889.；<http://www.newadvent.org/fathers/190118.htm>.}

% structure: p0001 source=h1 target=section reason=page-title

\section{论雕像讲道集 第十九篇}

\Emph{\Emph{论称为\Quote{集粮主日}的主日，}\Emph{给那些从乡间来到安提约基雅的人——兼论避免起誓。}}\par

1. 你们在过去的几天里，在圣殉道者身上欢庆了！你们饱享了这属灵的盛宴！你们全都以正直的喜乐而喜乐！你们看见了他们的肋骨暴露，腰背撕裂，鲜血四处流淌，以及成千上万种酷刑！你们看见了那超越自然的人性，以及用鲜血编织的华冠！你们在全城跳了一场美好的舞蹈，你们这位高贵的领袖引导着你们；然而，疾病迫使我留在家中，尽管我心有不甘。但我虽未参与这庆节，却分享了其中的喜乐。我虽未能享受你们公开集会的乐趣，却分享了你们的欢愉。因为爱的力量如此之大，它使那些并未实际享受的人与那些正在享受的人一同喜乐，使他们相信邻人的美善也是自己的。因此，即使我坐在家中，我也与你们一同喜乐；如今我尚未完全摆脱疾病，却已起身跑来迎接你们，为要看见你们那令人渴望的面容，并参与这眼前的庆节。\par

2. 因为我认为今日是一个极其盛大的庆节，为了我们的弟兄们，他们因临在而美化了我们的城市，装饰了教会；这些人在语言上与我们不同，但在信仰上与我们和谐一致；他们生活在安宁中，过着正直而清醒的生活。因为在这些人中，没有邪恶的表演——没有赛马，没有娼妓，也没有任何城市中的放荡行为，而是各种放荡都被驱逐，处处盛行着极大的清醒。原因是他们的生活是劳苦的，他们在耕种土地中拥有德行和清醒的学校，并且遵循天主在我们生命之初就引入的那门技艺。因为在亚当犯罪之前，当他享有极大的自由时，某种耕作土地的任务就交付给了他；这工作并非劳苦或麻烦的，而是给了他许多良好的训练，因为据说他受命\Quote{耕种乐园，看守乐园}。你们可以看到这些人，有时正在套上劳作的牛，引导犁头，切割深深的犁沟；有时又登上神圣的讲道台，用圣言耕耘那些属于他们权下之人的灵魂；有时用镰刀割去土壤中的荆棘，有时又用圣言清除灵魂的罪恶。因为他们不像我们城中的居民那样以劳动为耻，而是以懒惰为耻，知道懒惰从起初就教导了各种邪恶，并且对于喜爱它的人来说，它已成为邪恶的导师。\par

3. 这些人是我们的哲学家，他们拥有最佳的哲学，不是以外在显现其德行，而是以心灵。外教的哲学家在品格上绝不比那些从事舞台和演员表演的人更好；他们除了破旧的外衣、胡须和长袍之外，别无所有！然而这些人恰恰相反，他们告别了手杖、胡须和其他装束，他们的灵魂以真正哲学的教义为装饰，不仅以教义，而且以实际的实践。若你们询问这些在锄头和犁耙旁过着乡村生活的任何一位，关于那些外教哲学家曾经谈论无数、耗费大量言辞却未能说出任何正确见解的教义，其中一人会从他的智慧宝库中给出准确的答复。这不仅是值得惊叹的，而且他们还以行动证实了这些教义的可信性。因为他们坚信我们有不朽的灵魂，我们将来要为在此所做的事交账，并站在可畏的审判台前；他们以这样的希望作为整个生活的准则，并超越了一切世俗的炫耀，因为圣经教导他们：\Quote{虚而又虚，万事皆虚}\BibleRef{训 1:2}，他们并不贪婪地渴求那些看似如此光辉灿烂的事物中的任何一样。\par

4. 这些隐修士也懂得依照天主的定论来论及天主；如果你从中带出一位，并请一位外邦哲学家来与之辩论——其实如今你已找不到了！——但若你带出其中一位，然后翻开古代哲学家的著作，浏览一遍，并以平行对照的方式查究他们如今所回答的与当年那些哲学家所推论的，你就会看出前者的智慧何等深厚，后者的愚昧何等严重。因为其中一些人会断言现存的事物没有眷顾，受造物并非源自天主；美德不足以自足，反而需要财富、门第、外在的荣华以及其他更可笑的东西；而另一方面，这些人虽然完全没有世俗学识，却能明智地谈论有关眷顾、未来的审判台、天主从无中生出万物的创造之能，以及其他一切要点；有谁不能由此学到基督的能力呢？这能力证明了这些无学问的单纯之人比那些自夸智慧的人更加明智，正如有辨识力的人与幼童相比一样。因为当他们思想充满大智慧时，学识上的单纯能对他们造成什么伤害呢？而那些哲学家从学识中得了什么益处呢？当他们理智缺乏正确思想时。好比一个人有一把刀，刀鞘是银的，但刀刃却比最劣的铅还弱。事实上，这些哲学家的舌头装饰着辞藻和名词，但他们的理智却充满软弱和无用。但这些隐修士则不如此，而是完全相反：他们的理智充满属神的智慧，他们的生活方式是其教义的翻版。在他们中间，没有奢侈的妇女；没有服饰、颜色或脂粉的装饰；所有这类败坏的风气都受到谴责。因此，他们所管辖的民众更容易被训练成节制的人；而保禄所定的律法，即指示人有食物和遮蔽就当知足，不再追求更多，他们极其严格地遵守。\BibleRef{弟前 6:8} 在他们中间，没有迷醉感官的芳香油膏；大地生出香草，为他们预备了超越一切调香师手艺的百花馨香。因此，他们的身体和灵魂都享有健康的状况，因为他们摒弃了饮食的一切奢侈，驱散了醉酒的一切恶流；他们只吃足以维持生存的份量。我们不要因他们的外表而鄙视他们，而要欣赏他们的心灵。因为外在的衣著有什么益处呢？当灵魂比任何乞丐穿得更褴褛时！人应当受到称赞和钦佩，不是因衣著，更不是因形体，而是因灵魂。你若揭露这些人的灵魂，就会看见它的美丽和它所拥有的财富，表现在他们的言语、教义和整个行为体系中！\par

5. 愿外邦人羞愧，愿他们低头退避，因他们的哲学家和他们的智慧——其可怜超过一切愚昧！因为在他们有生之年，那些哲学家几乎无法将他们的教义传授给很少数、易于计数的人；而当任何微小的危险临到时，他们甚至失去了这些人。但基督的门徒——渔夫、税吏和制帐篷的人——在几年内使全世界归向真理；并且从那时起，虽有成千上万的危险不断兴起，福音的宣讲不但没有被压制，反而仍然繁盛增长；他们教导单纯的人、耕种土地的人和牧养牲畜的人成为爱智慧者。就是这些人，此外还在所有方面深深扎根了那作为万善之源的爱德，\BibleRef{弗 3:17} 他们急忙来到我们这里，跋涉如此漫长的旅程，只为前来拥抱他们的同肢。\par

6. 那么，为报答这些恩惠（我是指他们的爱德与善意），让我们给他们备办一些行装，送他们回家；并且让我们再次提出关于起誓的问题，好从所有人的心灵中根除这恶习。但首先，我今天愿意稍微提醒你们我们最近所讲的事。\par

当犹太人从波斯被释放、摆脱那暴政返回本国时，有一位说：\BibleQuote{我看见一把飞行的镰刀，长二十肘，宽十肘。}\BibleRef{匝 5:1} 他们也听见先知这样教训：\BibleQuote{这是诅咒，它行在全地的面上，进入发假誓者的房屋，并停留在其中，毁坏那房屋的木材和石头。}当我们读了这段经文后，当时也追问为何它不仅要毁灭起假誓者本人，也要毁灭他的房屋；我们解释了这个原因：天主让最严重之罪的刑罚持续可见，好使以后的人都能学习谨慎。既然有必要使发假誓者死后被埋葬、归还大地怀抱，为了不让他的邪恶与他一同埋葬，他的房屋被堆成废墟，使所有路过的人看见，得知倒塌的原因后，避免效法那罪。\par

7. 这事也发生在索多玛。当他们彼此欲火中烧时，当时土地本身也被从上而来的火点燃焚烧。因为天主安排，让这罪的惩罚永远存留。\par

且看天主的仁慈！祂并未使犯罪者一直焚烧至今，而是当他们一度被火焰吞噬后，便将他们埋葬；祂烧灼了地面，将其显现在所有后来渴望目睹这些事的人眼前；如今，这片土地在历代以来所呈现的景象，发出了一种超越一切言语的劝诫，仿佛在呐喊说：\Quote{切勿做索多玛的恶行，免得遭受索多玛的命运！}因为诫命通常不如一个恐怖的景象那样深深印在心灵上，这景象承载着历久不衰的灾难痕迹。曾到访那些地方的人可以作证：他们时常在听到圣经讲述这些事时，并不十分惊惧；但当他们亲身前往站在那遗址上，看到整个地表面目全非，目睹烈火的遗迹，不见土壤，尽是尘土和灰烬时，便因所见而惊愕不已，并带走了一堂深刻的贞洁功课。诚然，惩罚的性质正是罪恶性质的写照！正如他们发明了一种荒淫的交合，不以生育为目的，同样，天主也降下这样的惩罚，使那地之子宫永远荒芜，不结果实！为此，祂也威胁要摧毁起誓者的居所，好使他人因他们的惩罚而更加自制。\par

8. 然而今天，我准备展示的并非起誓导致一所、两所或三所房屋的毁灭，而是一座整城和一个为天主所爱的民族的毁灭；是一个向来深受天主眷顾的邦国，是一个曾逃脱许多危险的种族。因为耶路撒冷，那座天主的城，拥有圣约柜和一切神圣礼仪——那里曾有先知、圣神的恩宠和约柜，以及盟约的法版和金罐——那里常有天使降临——我说，这座城市，当众多战争发生、许多外邦民族进攻它时，它仿佛被金刚石城墙环绕，始终嘲笑他们；当整个地区被摧毁时，它却毫发无损！不仅如此，令人惊叹的是，它屡次驱逐敌人，给予他们重创，并深受天主眷顾的看顾，以致天主曾亲口说：\BibleQuote{在旷野里，我找到了以色列，如葡萄树初熟的果实；我看见了你们的祖先，如无花果树的最早果实。}\BibleRef{欧 9:10}关于这座城，又说：\BibleQuote{像高枝梢头的橄榄果，人们要说：不可伤害它们。}然而，这座为天主所爱的城，曾逃脱许多危险，在众多罪恶中得到宽恕，是唯一能在其余各族被掳（不止一次或两次，而是经常）时避免被俘的城，却仅仅因一个誓言而毁灭了。我接下来就说明是如何毁灭的。\par

9. 他们的一位君王是漆德克雅。这位漆德克雅曾向蛮族的君王拿步高宣誓，要与他保持同盟。之后他背叛了，投奔埃及王，藐视他誓约的义务，遭受了你即将听到的事。但首先有必要提及先知以谜语描述这一切的比喻：他说：\BibleQuote{上主的话传给我说：人子，你要出一个谜语，讲一个比喻，说：吾主上主这样说：有一只大鹰，有大翅膀，长阔，满有爪甲。}\BibleRef{则 17:2} 这里他称巴比伦王为鹰，说他\Quote{大而长翅}；他又称他为\Quote{长阔}、\Quote{满有爪甲}，因为他的军队众多，权力强大，且入侵迅速。因为正如鹰的翅膀和爪甲是它的武器，马匹和士兵对君王也是如此。这鹰，他接着说，\Quote{入了黎巴嫩之顶。} 这\Quote{顶}是什么意思？计谋——计划。犹太被称为黎巴嫩，因为靠近那座山。随后，为谈论起誓和盟约，他说：\BibleQuote{又从本地拿了种子，栽种在肥美的田里，更在多水之处，使它扎根。他使其得以看见，就生长起来，成了一棵低矮的蔓藤，向那鹰伸出枝条，它的根也在那鹰之下。}\BibleRef{则 17:5} 这里他称耶路撒冷城为葡萄树；但说它向鹰伸出枝条，根在鹰下，指的是与鹰订立的条约和盟约；并且它依附于他。接着，为宣告此事的不义，他说：\BibleQuote{又有另一只大鹰，}（指埃及王）\BibleQuote{有大翅膀，有许多爪甲；那葡萄树就向它弯垂，它的卷须也向它伸出，发出枝条，为要得水灌溉。因此我说：吾主上主这样说：它岂能兴盛呢？}\BibleRef{则 7:7} 也就是说，\Quote{在背弃了誓言和盟约之后，它岂能存留，或得安全，或免于跌倒呢？} 随后，为表明此事不会发生，而是因誓言必被毁灭，他谈论其惩罚，并指出原因。\BibleQuote{因为它的嫩根和果实必腐败，凡由此而生的都要枯干。}\BibleRef{则 17:9} 并且为表明它不是被人力摧毁，而是因这些誓言使天主成为它的敌人，他接着说：\BibleQuote{不在于强大的臂膀，也不在于众多的人民，要把它连根拔起。} 这比喻诚然如此，但先知再次解释时，他说：\BibleQuote{看，巴比伦王来到耶路撒冷。}\BibleRef{则 17:12} 然后，在说了一些中间的话之后，他提及誓言和盟约。他说：\BibleQuote{因为他要与他立约；}\BibleRef{则 17:14} 随后，谈到背弃盟约，他接着说：\BibleQuote{他要离开他，派人到埃及去，要他们给他马匹和众多人民。} 接着他继续表明，正是因着誓言，这一切毁灭才要发生。\BibleQuote{的确，在他立他为王的那君王所居住的地方，就是那轻视我的诅咒、违犯我的盟约的人，他要在巴比伦中间死亡；不在于大能，也不在于人多，因为他轻视了誓言，违犯了我的盟约；我必将这遭受他侮辱的誓言和破坏的盟约，报复在他头上；我要在他身上张设我的网罗。}\BibleRef{则 17:16} 你看见了吗？他说，不是一次两次，而是再三地说，因着誓言他必须受这一切。因为当天主受到藐视时，祂是不宽容的。不仅从誓言带给城市的惩罚，也从延迟和缓期中可以看出，天主是多么关心誓言的不可侵犯性。\BibleQuote{这事发生，}我们被告知，\BibleQuote{在漆德克雅王第九年十月十日，巴比伦王拿步高率领全军前来攻打耶路撒冷，安营围城，在其周围筑垒攻城，城被围困直到漆德克雅王第十一年正月初九，城里没有粮食给人民吃，城就攻破了。}\BibleRef{列下 25:1} 祂原本可以立刻从第一天就把他们交出去，把他们交在敌人手中；但祂准许他们先被荒废三年之久，经历最困苦的围城；目的是在这期间，他们或被外面的军队惊吓，或被城内的饥荒压迫而谦卑下来，迫使国王即使不情愿也要向蛮族投降；这样，所犯的罪也许能得到一些减轻。为证明这是真实的，并非我自己的臆测，请听祂借先知向他说的话：\BibleQuote{你若出去向巴比伦王的将帅投降，你必存活，这城也必不被火烧；你和你的全家必都存活。但你若不出去向巴比伦王的将帅投降，这城必被交在加色丁人手中；他们要放火烧城，而你也逃不出他们的手。君王说：我害怕那些已倒向加色丁人的犹太人，恐怕他们把我交在他们手中，他们讥笑我。但耶肋米亚说：他们不会把你交出。请你听从我向你所说上主的话，这事必为你带来好处，你的性命必得保全。但你若拒绝出去，这是上主指示我的话：凡留在犹大王宫中的妇女，都要被带出去交给巴比伦王的将帅；她们要说：你信任的人骗了你，而且胜过了你；当你的脚滑跌时，他们胜过了你；他们离开了你，并将你的所有妻子儿女带出去交给加色丁人；你也逃不出他们的手，因为你将被巴比伦王擒获，这城也必被火烧。}\BibleRef{耶 38:17}\par

10. 但是，当祂藉着这番话没有说服他，他仍停留在罪恶和过犯中时，三年之后，天主将城交了出来，同时显示了自己的仁慈和那国王的忘恩负义。他们极其轻易地进入，\BibleQuote{焚烧了上主的殿、王宫、耶路撒冷的房屋，以及所有高大的房屋，卫队长焚烧了，并拆毁了耶路撒冷的城墙；}到处是蛮族的火，誓言成为大火的引导，将火焰带往各方。\BibleQuote{卫队长将城中剩下的余民和投降巴比伦王的逃难者掳去。}\BibleRef{耶 39:9}\BibleQuote{加色丁人打碎了上主殿中的铜柱、底座和铜海，又将盆、肉钩、碗、香炉，以及所有行礼仪用的铜器都带走了。火盆和所有金碗银碗也被带走。此外，卫队长乃步匝辣当带走了所罗门在上主殿中所造的两根柱子、底座和铜海。他们带走了大司祭色辣雅、副司祭则法尼雅和三个守门的；从城中带走了管理军队的一个太监、王面前的五个人、军长沙番、书记长和六十个人。他将这些人带到巴比伦王那里，巴比伦王击杀了他们。}\BibleRef{列下 25:13}\par

11. 因此，我恳求你们，请记住那\Quote{飞镰}，它\Quote{停在起誓者的家中}，并且\Quote{摧毁墙壁、木料和石头。}请记住，这誓言如何进入城中，推翻了房屋、圣殿、墙壁和华丽的建筑，使城成为一堆瓦砾；无论是至圣所、圣器，还是任何其他东西，都无法抵挡那惩罚和报应，因为誓言被违背了！城确实就这样悲惨地被毁灭了。但国王所遭受的更加悲惨和可悲。\BibleRef{列下 25:4} 正如飞镰推倒了建筑，它也在他逃亡时将他砍倒。因为经上说：\BibleQuote{国王在夜间从城门出发，加色丁人包围了城，加色丁军队追赶国王，追上了他，他们抓住国王，带他到巴比伦王那里，巴比伦王对漆德克雅宣判，当着他的面杀了他的儿子，剜了漆德克雅的眼睛，用锁链捆住他，带他到了巴比伦。}\Quote{宣判}一词是什么意思？他要求他解释自己的行为，指控他；并先杀了他的两个儿子，让他亲眼目睹自己家庭的灾难，看那悲惨的悲剧；然后剜了他的眼睛。我再问，为何发生这事？为使他成为蛮族和住在他们当中的犹太人的导师，使那些有眼睛的人，借着这个失明的人，能看出誓言是多么大的罪恶！不仅如此，所有住在路边的人，看见这被锁链捆绑、失明的人，都能从他的灾难学到他的罪有多大。因此，有一位先知宣告：\BibleQuote{他不得看见巴比伦。}\BibleRef{则 12:13} 另一位说：\BibleQuote{他必被掳到巴比伦。}\BibleRef{耶 32:5} 这预言看似矛盾，其实不然；因为两者都是真的。他虽被掳到巴比伦，却未看见巴比伦。他是如何没有看见巴比伦呢？因为他在犹太地被剜了眼睛；在誓言被蔑视的地方，誓言也得到了维护，他自己受了惩罚。他是如何被掳到巴比伦呢？是在被掳的状态下。因为惩罚是双重的：失明和被掳，先知们分别取了其中之一。一位说\BibleQuote{他不得看见巴比伦，}指的是失去眼睛；另一位说\BibleQuote{他必被掳到巴比伦，}指的是他被掳。\par

12. 弟兄们，既然知道这些事，并集合起刚才所讲的和以前所说的，让我们最终戒除这邪恶的习惯，是的，我恳求和乞求你们大家！因为如果在旧约时代，犹太人没有被要求最严格的道德智慧，而是对他们有许多宽容，一个誓言竟引起了这样的愤怒，这样的被掳和囚禁；那么，如今有了明确禁止此事的法律，并增加了如此多的诫命，那些起誓的人现在可能会遭受怎样的惩罚呢？难道我们来到集会中，听所讲的话，就够了吗？实际上，我们不断听讲，却不遵行所吩咐的，这正是更大的定罪和更难免罚的理由！如果我们从幼年到老年都聚集在这里，享受如此多的教导，却仍像他们一样，不努力改正一个缺点，我们还有什么借口或宽恕呢？从此，谁也不要以习惯为借口。因为我正是为此感到愤慨和恼怒：我们不能胜过习惯。请问，如果我们不能胜过习惯，又怎能胜过那根植于我们本性原则中的私欲偏情呢？感到欲望是自然的，但邪恶地欲望则是出于选择。而起誓的习惯甚至不是源于本性，而只是出于疏忽。\par

13. 要使你知道，这罪之所以发展到如此地步，并非由于事情本身的困难，而是因我们的疏忽，让我们回想，人所完成的比这些更难的事有多少；而且还不指望任何报酬。让我们想想魔鬼所施加的劳役是何等艰辛、何等麻烦；然而困难并没有成为这些劳役的障碍。请问，还有什么比以下的事更困难呢？当一个年轻人把自己交给那些负责使他的四肢柔软灵活的人，用尽全身力气把整个身体弯成轮子形状，在铺道上翻转；同时通过眼睛、手的运动以及其他扭曲来模拟女性形象。然而这些技艺的困难以及由此而来的卑劣却无人理会。再者，那些被拖上舞蹈舞台、把身体肢体当作翅膀来使用的人，谁看见他们能不感到惊奇？同样，那些在空中轮流抛掷刀子、且都能抓住刀柄的人，他们使那些为德性不愿付出任何辛劳的人何等羞愧！人能对那些用额头平衡一根杆子、使其像扎根地里的树一样稳定的人说什么呢？这还不是唯一神奇之处，他们还让小孩子在杆子顶端相互摔跤；不用手或身体其他任何部位辅助，仅凭额头就稳稳地支撑着杆子，比任何固定方式都更稳。再如，另一个人在极细的绳索上行走，与人们在平地上奔跑一样无所畏惧。然而这些在思想中甚至显得不可行的事，却通过技艺成为可能。请问，关于发誓，我们有什么类似的事可以声称？有什么困难？什么辛劳？什么技艺？什么危险？我们这边只需一点热忱，整个任务就能迅速完成。\par

14. 不要对我说：\Quote{我已经完成了大部分；}但如果你没有完成全部，就当自己还未做任何事；因为这微小部分，若被忽视，就会毁掉其余的一切。确实，人建造房屋并盖上屋顶后，常因不在意一块松脱的瓦片而毁掉整个结构。在衣服上也是如此；因为那里若有一个小洞不修补，就会造成大裂口。洪水也常是这样；因为只要找到一个小入口，就会让整个洪流涌入。同样，你即使已四面加固，只留下一小部分未被巩固，也要把这部分对魔鬼堵住，使你在一切方面都坚固！你见过镰刀！你见过若翰的头！你听过关于撒乌耳的历史！你听过犹太人的被掳情况！此外，你听过基督的宣告，说不仅发虚誓，而且以任何方式起誓，都是魔鬼的事，完全是恶者的伎俩。你听过，虚誓到处都伴随着誓言。那么，把这一切都铭记在你的理智中。你没有看见妇女和小孩子把福音挂在颈上作为强有力的护符，到处携带吗？同样，你也要把福音的命令和法令写在你的心中。这里不需要金子、财产或买书；只需要意志和灵魂的觉醒，福音就会成为你更可靠的守护者，因为你携带它，不是在外面，而是珍藏在内里，即在灵魂的密室里。当你从床上起来，走出家门时，重复这条法令：\BibleQuote{我告诉你们：什么誓都不可起。}\BibleRef{玛 5:34}这句话将成为你的纪律；因为不需要很多劳力，只需中等程度的注意。这一点是真实的，可以这样证明：叫你的儿子，吓唬他，威胁他如果不好好遵守这条法令，就鞭打他几下；你就会看到他如何立刻戒除这个习惯。因此，我们使小孩子因恐惧而遵守这条诫命，而我们却不像儿子畏惧我们那样畏惧天主，这岂不是很荒唐吗？\par

15. 因此我先前所说过的话，现在再次重复。让我们在这件事上为自己订立一条法则：在履行这条法则之前，不干预公共或私人事务；然后，在这义务的压力下，我们必然轻易取胜，并立即美化自己，装饰我们的城市。试想，如果全世界到处都在传说：\Quote{在安提约基雅建立了合于基督徒的习俗，你听不到任何人发出誓言，即使有极大的必要加在他身上！}这是何等的一件大事！邻近的城市必定会听到；不，不仅是邻近的城市，这消息甚至会传到地极。因为与你们交往的商人以及从这地方来的其它人很可能都会报告这一切。因此，当许多人以赞美之词提及别的城市的港口、市场或丰富的货物时，使从这地方来的人能够说：在安提约基雅有别的城市所看不到的事物；因为那里的居民宁愿割掉舌头，也不容许誓言从他们口中发出！这将是你们的装饰和保障，不仅如此，它还会带来丰厚的赏报。因为其他人也必定会效法并模仿你们。但如果一个人赢得了哪怕一两个人，\BibleRef{雅 5:20} 他就要从天主那里得到如此大的赏报，那么当你们作为全世界的导师时，你们将得不到何等的赏报呢？因此你们有责任振奋自己，警醒谨慎，知道我们不仅要从自己的善工，也要从我们在他人身上所行的善工中，获得最好的赏报，并享受在天主面前的极大恩宠。愿祂使我们众人不断享受这恩宠，并在将来获得天国，在耶稣基督我们的主内。愿光荣与权能归于祂，以及圣父和圣神，从现在直到永远，世世代代。阿们。\par

\SourceRecord{Translated by W.R.W. Stephens.；Nicene and Post-Nicene Fathers, First Series；Vol. 9.；Edited by Philip Schaff.；Buffalo, NY: Christian Literature Publishing Co.,；1889.；<http://www.newadvent.org/fathers/190119.htm>.}

% structure: p0001 source=h1 target=section reason=page-title

\section{讲道第二十篇：论雕像}

\Emph{四旬期的斋戒不足以使我们有资格领受圣体，而圣德是首要条件。如何可能不怀怨恨，以及天主如何高度重视这条法律；怀怨恨的人在抵达受苦之地前便已惩罚了自己。——也论起誓的戒绝，以及那些未能戒除发誓的人。}\par

1. 终于，季节临近斋戒的尾声，因此我们应当更加热切地致力于圣德。因为正如赛跑的人，若错失了奖品，则所有的赛程都无济于事；同样，如果我们不能以良心得享神圣的祭台，那么这些关于斋戒的种种劳作与辛苦也将毫无益处。为此设立了斋戒与四旬期，以及如此多的庄严集会、听道、祈祷和教导，为使我们通过这些热忱，以一种可能的方式，从全年所犯的罪过中得到净化，从而以属灵的胆量虔诚地参与这无血之祭；倘若结果并非如此，那么我们所承受的一切辛劳便完全徒劳，毫无裨益。因此，让每个人反省自己：他改正了何种缺陷？获得了何种善工？抛弃了何种罪恶？净化了何种污点？在何处变得更好？如果他发现自己在这次美好的交易中因斋戒而获益，并且内心意识到对自己的伤口已多加照料，就让他上前！但如果他一直疏忽大意，除了仅仅斋戒之外别无所有，且未做什么正确的事，就让他留在外面；等他清除所有这些过犯后再进入。愿无人只依靠斋戒，而同时继续留在恶习中不改。因为不斋戒的人或许能得到宽恕，可借口身体虚弱；但未改正过错的人不可能有借口。你也许因身体虚弱而未斋戒。告诉我，你为何不与仇敌和好？你在这里当真能以身体虚弱为理由吗？再者，你若保持嫉妒与仇恨，请问你又有什么借口？对于这类过犯，无人能借身体虚弱之由逃避。这一切都是基督爱人工程的体现，即首要的诫命和那些维持我们生命的诫命，不应因身体的软弱而有丝毫受损。\par

2. 既然我们需要同样遵守所有神圣的法律，尤其是那条命令我们不视任何人为仇敌、不长怀怨恨而立刻和好的法律，请允许我们今天和你们谈谈这条诫命。因为正如行淫者和亵渎者不能参与圣体祭台一样，同样，怀有仇敌和怨恨的人也不可能享受圣体。这也是有充分理由的。因为一个人若犯了奸淫或通奸，在他满足了情欲的同时，也完成了罪过；如果他愿意通过警醒的生活从那跌倒中恢复，日后藉着显明的重大补赎，或可获得一些缓解。但怀恨的人每天都施行同样的不义，永不停息。在前者情形中，行为已过去，罪过已完成；但在这里，罪过每天都在持续发生。请问，我们还有什么借口自愿将自己交付给如此邪恶的怪物呢？当你对同仆如此苛刻和不肯宽恕时，你怎能祈求你的主对你温和慈悲呢？\par

3. 但是，你的同仆也许轻慢了你？是的！而你也多次轻慢了天主。然而，同仆与主之间有何可比呢？对于前者，当他或许在某方面受伤害时，他侮辱了你，你就被激怒了。但你却侮辱了主，而主既未对你不义，也未曾怀恶意，反而日复一日地赐福于你。那么，请想想：如果天主选择严厉追究你冒犯祂的行为，我们连一天也活不了。因为先知说：\BibleQuote{主，你若细察罪过，主，谁能站立得住？}再者，且不说其他一切事——每个罪人的良心都知道，而且无人见证，惟独天主知道——就算我们只被追究那些公开且公认的罪过，我们还能指望对这些罪得到什么宽恕呢？如果祂要审视我们在祈祷中的懈怠和疏忽：我们在天主面前站立并恳求祂时，竟连仆人对待主人、士兵对待长官、朋友对待朋友的那份敬畏与尊重都没有表现出来，那又会怎样？当你和朋友交谈时，你留意自己在做什么；但当你为你的罪过侍奉天主，求祂赦免如此多的过犯，并以为会得到宽恕时，你却常常心不在焉；你的双膝跪在地上，你的心却四处游荡，在市场、在家里，口中虚妄地喃喃自语，毫无意义！这并非一次两次，而是屡屡发生！如果天主只审视这一点，你认为我们能得到宽恕，或能找到任何借口吗？我真心认为不能！\par

4. 但是，如果我们每天恶毒地彼此口出恶言，以及我们无故论断邻人的轻率判断——只因我们爱挑剔、好指责——这些被带到我们面前，我们还能有什么可辩解的呢？再者，如果他审视我们游荡的目光和我们心中携带的邪恶欲望——我们屡次因眼睛无拘无束地游荡而接纳可耻和不洁的思想——我们将承受怎样的惩罚呢？如果他要求我们为辱骂（他说：\BibleQuote{谁若向自己的弟兄说\InnerQuote{傻子}，就要受地狱的罚。}）交账，我们岂能开口或动嘴唇，或说些什么大小话来回应呢？此外，关于我们在祈祷、禁食、施舍时所怀有的虚荣心——如果我们审视它们——我不是说天主，而是说我们这些罪人自己这样做——我们还能举目望天吗？然后，我们彼此间所设的诡计——当面称赞弟兄，交谈如朋友，背后却诋毁他——我们能承受这一切的惩罚吗？还有誓言呢？还有谎言呢？还有假誓呢？还有不义的愤怒，以及我们常对受尊荣者——不仅是仇敌，还有朋友——所怀的嫉妒呢？再者，我们见到别人遭祸便高兴，视他人的不幸为我们自己痛苦的安慰，这又如何呢？\par

5. 但假若惩罚是针对我们在大聚会中的懈怠，我们的情况又会如何呢？因为你们不能不知道，当天主亲自通过先知对我们全体说话时，我们却经常与邻座的人进行频繁的长谈，谈论与我们毫不相干的事。那么，撇开其他一切，如果他选择只追究这一罪的惩罚，哪里还有救恩的希望呢？因为不要以为这过失是小事；如果你要认识其严重性，就看看这件事在人间是如何看待的，你就会明白这罪的严重。你只要敢，当某位长官对你说话时，或者更确切地说，某位地位稍高的朋友对你说话时，你转身与他谈话，而与你的仆人交谈；你就会明白，你如此对待天主所冒的风险！因为如果他是较显贵阶层中的一员，他会因为这样的侮辱而要求你赔礼道歉。然而天主，每天都受到同样甚至更大的轻视，而且不是一个人、两个人或三个人，而是我们几乎所有人，仍然忍耐宽容，不仅对此，而且对更为严重的事情也是如此。因为这些都是必须承认、众目昭彰、几乎所有人都胆大妄为的事。但还有另一些事，只有当事人自己的良心知晓。确实，如果我们想到这一切；如果我们自我反省，即使我们是最残忍、最冷酷的人，但一看到自己罪过的众多，我们就会因恐惧和痛苦而无法记起别人对我们的伤害。记住火河、毒虫、可怕的审判，那时一切都将赤裸敞开！要思想，现今隐藏的事，那时都要显露出来！但如果你宽恕邻人，所有这些等待揭露的罪过，在此世就被消除；当你离开此生时，你就不会拖带着那条罪过的锁链；因此你所收到的比你给予的更大。因为这样的过犯，我们常常犯下，没有其他人知道；当我们想到，在那一天，我们的这些罪过将在宇宙的公共剧场中暴露在众人眼前，我们内心就感到痛苦超过任何惩罚，被良心窒息扼杀。然而，这羞耻虽大；这些罪，这些惩罚虽大；却有可能通过对邻人的宽恕来涤净。\par

6. 事实上，没有什么能与这种德行相比。你愿意学习这德行的力量吗？\Quote{纵有\Person{梅瑟}和\Person{撒慕尔}站在我面前，}天主说，\Quote{我的心仍不转向这人民。}\BibleRef{耶 15:1} 然而，\Person{梅瑟}和\Person{撒慕尔}不能从天主的义怒中夺回的人，这一诫命如果遵守，却能夺回。因此，他不断劝诫那些听了这些话的人说：\Quote{你们中间不可有人心怀 \BibleRef{耶 15:1} 恶念对付自己的兄弟，}又说\Quote{你们不可思念邻人的恶意。}这不只是说，不要发怒；而是说，不要将怒气存留在心中；不要思想它；放下一切怨恨；除去那创伤。因为你以为你是在报复对方，但首先你是折磨自己，在你内心各处设立你的愤怒作为刽子手，撕裂你自己的脏腑。还有什么比一个常常发怒的人更不幸的呢？正如癫狂者永不得安宁，心怀怨恨、记仇的人也永远享受不到平安；他不停地发怒，日复一日加重心中的风暴，思索对方的话语和行动，憎恨那得罪他之人的名字。你只要一提他的仇人，他立刻暴怒，内心承受极大的痛苦；若偶然仅瞥见对方，他就恐惧战栗，如同遭遇最大的灾祸。是的，若他看见对方的任何亲属，甚至他的衣服、他的住处、他的街道，他都会因所见而痛苦。因为正如对于所爱的人，他们的面容、衣服、鞋子、房屋或街道，一看见就振奋我们；同样，若我们看见我们所恨并视为仇敌之人的仆役、朋友、房屋、街道或任何属于他们的事物，我们都会被这一切所刺痛；因看见其中每一件而承受的打击是频繁而不断的。\par

7. 那么，何必忍受这样的围困、这样的折磨、这样的惩罚呢？即使地狱并不威胁那心怀怨恨者，但单是为了事情本身的痛苦，我们也该宽恕那得罪我们的人所犯的过错。但当永罚尚在后面等候时，还有什么比这样的人更愚昧呢？他在这里和来世都给自己带来惩罚，却自以为是在报复仇敌！假设我们看见他仍然兴旺，我们就会因烦恼而快要死；但若他处境不利，我们又担心会有什么转机发生。但这两者都为我们积蓄了不可避免的惩罚。因为经上说：\Quote{你的仇敌跌倒时，你不要高兴。}\BibleRef{箴 24:17} 不要对我说所受的伤害有多大；因为使你怀恨的不是这个，而是你忘记了自己的过犯，你眼中没有地狱或敬畏天主！为使你相信这是真的，我要从本城所发生的事来证明。当那些被控告犯有滔天罪行的人被拖到审判台前——当火在里面点燃，刽子手站在周围，撕裂他们的肋骨时，若有人站在旁边向他们宣告说：\Quote{如果你们有任何仇敌，就放弃怨恨，我们就能救你们脱离这刑罚；}——他们岂不会亲吻那人的脚吗？我为何说脚呢？若有人吩咐他们视仇敌为主人，他们也不会拒绝。如果人间的、有界限的刑罚能战胜一切愤怒，那么来世的刑罚，若常占据我们的思想，岂不更能驱散灵魂中不仅怨恨，而且一切邪恶的意念吗？请问，除掉对伤害者的怨恨，有什么更容易的呢？需要长途跋涉吗？需要花钱吗？需要求助于他人吗？只要下定决心，这善行便立刻达到目标。那么，我们该受怎样的惩罚呢？如果我们为了世俗事务而屈就于奴役般的劳作，表现出与己不相称的卑贱，并且花钱，与门房交谈以讨好不敬神的人，做尽一切事情以期达到目的；然而为了天主的诫命，却不能容忍去恳求那得罪我们的兄弟，反而认为先迈出一步是可耻的。告诉我，当你去率先获益时，你难道不感到羞愧？相反，你应当为坚持这种激情而感到羞愧，并且等着那犯罪的人来与你修和；因为这才是耻辱、责备和最大的损失。\par

8. 谁先来，谁就获得全部果实；当你因别人的恳求而放下怒气时，这善工要算在他的账上，因为你是作为施恩于他、而非服从天主而履行了律法。但是，若无人恳求，甚至那伤害了你的人也不来接近或恳求你，你自己却抛开一切羞耻与迟延，自由地跑向伤害者，完全平息怒气，那么这善行就完全属于你，你将领受全部赏报。我说：\Quote{操练禁食}，你也许推说身体软弱。我说：\Quote{施舍穷人}，你推说贫穷和抚养子女。我说：\Quote{为教会的聚会腾出时间}，你推说世俗的牵挂。我说：\Quote{留心听所讲的，并思考所教导的道理的能力}，你推说缺乏学问。我说：\Quote{规劝别人}，你说：\Quote{给他劝告时，他不理会，因为我已多次说话，反被轻蔑。}这些托词虽然苍白无力，但你终究有些托词可说。但假如我说：\Quote{放下你的怒气}，你那时又能以其中哪一条作答？因为无论是身体的软弱、贫穷、缺乏教养、没有闲暇，还是其他类似的事情，你都无可推诿；这罪比其他一切都更不可原谅。那时你怎能向天伸手？怎能开口祈求宽恕？即使天主愿意宽恕你的罪，你自己却不允许祂，因你仍怀恨你的同仆！岂不知那人是残忍、凶猛、野蛮、贪图报复与回击吗？正因如此，你更应当宽恕。你受了很大的冤屈、被掠夺、被诽谤、在最重要的事情上受了伤害，你希望看见仇敌受罚？然而即使为此，宽恕他对你也有益。因为如果你亲自报仇，以言语、行为或诅咒攻击敌人，那么天主此后就不会再追究，因为你已报了仇；祂不但不会为你追究此事，还要向你追讨罪罚，因你藐视了祂。因为在人间也是如此：若我们打了别人的仆人，主人会恼怒，称之为侮辱（即使我们受到奴隶或自由人的欺负，也应当等待官长或主人依法裁决）；那么若在人间，自行报复尚且不安全，更何况天主是复仇者呢！\par

9. 你的邻人亏负了你、令你伤心，使你遭受了千般祸害？就算如此，你也不要亲自报复，免得你竟敢对主怀恨！将此事交给天主，祂定会比你所愿的更好地处理。祂给你的命令只是为伤害者祈祷；但如何处理他，祂命你交由祂自己。你绝不可能像这样报复自己，正如祂预备为你报复一样——只要你让位给祂独自行事，不对得罪你的人发出诅咒，而是让天主独自作为审判的仲裁者。因为即使我们宽恕了得罪我们的人，即使我们与他和好，即使我们为他们祈祷，天主却不会宽恕，除非他们自己悔改并变得更好。祂不宽恕，是为了他们的益处。祂称赞你、认可你的神性智慧，却管教他，免得他因你的智慧而变得更糟。因此，关于此事的俗语并不恰当。因为有许多人，当我谴责他们不愿听从劝告与仇敌和好时，他们竟提出这样的辩解，这不过是他们罪恶的遮羞布。有人说：\Quote{我不愿意和好，免得使那人变得更坏、更暴躁、日后更轻视我。}除此之外，他们还提出另一理由：\Quote{许多人觉得我先去和好并恳求仇敌是软弱的表现。}这一切都是愚昧的；因为那不打盹的眼已看见了你的善意；所以，当你赢得了那位将要审判你的案件的法官的认可时，你就不该在乎同仆的意见。\par

10. 但如果你担心的是，你的宽仁会使你的敌人变得更坏——你要知道：并非如此；相反，如果你不和解，他反而会更坏。因为即使他是最卑鄙的人，即使他既不会承认也不会公开宣扬，但他会在心里默默赞许你的基督徒智慧，并在自己的良心中尊重你的温良。然而，如果他继续行同样的不义，而你在努力软化并劝解他时，他必将承受天主最严厉的惩罚。为了让你知道，即使我们应为敌人和伤害我们的人祈祷，但如果他们因我们的容忍而变得更坏，天主也不会宽恕，我要向你讲述一段古代史。\Person{米黎盎}曾诽谤\Person{梅瑟}。天主打发了什么？祂使\Person{米黎盎}患上了癞病，使她成为不洁，尽管她在其他方面一向温顺谦和。后来，当受损害的\Person{梅瑟}本人恳求息怒时，天主并未同意；反而说：\BibleQuote{如果她的父亲吐唾沫在她脸上，她岂不应当羞愧？让她在营外}，祂说，\BibleQuote{隔离七天。} \BibleRef{户 12:14} 但祂的意思是：\Quote{如果她有一位父亲，父亲将她逐出自己面前，她岂不承受责备？我确实赞许你的兄弟之爱，你的温良和宽仁；但我知道何时是减轻她惩罚的适当时候。} 因此，你当向你的兄弟展现全部的人性；不要为了对他施以更大的惩罚而宽恕他的过犯，而应出于你的温柔和善意；然而，你当非常清楚地明白，他越轻视你，而你在努力和解时，他将为自己招致越大的惩罚。\par

11. 你说什么？告诉我，你的关注使他变坏了吗？这固然是他的过失，却是你的称赞。你的称赞在于，虽然看到他的行为如此，你却没有停止遵行天主的旨意去与他和解。但对他而言，这是过失，因为他没有因你的宽仁而变得更好。然而，\Quote{让别人因我们而被责备，远比我们因他们而被责备更为可取。} 不要对我说这种冷淡的答复：\Quote{我害怕别人会认为我是出于恐惧才向他示好；因此他会更加轻视我。} 这样的答复显示出幼稚愚蠢的心灵，为人的称赞而焦虑。让他认为你是出于恐惧才主动接近他好了；你的赏报将更大，因为你在事先知道这一点的情况下，仍为敬畏天主而甘心忍受一切。因为追求人的称赞并为这个目的寻求和解的人，会削减赏报的酬报；但那些确知许多人会诽谤嘲笑他，却仍不停止寻求和解的人，将得到双倍，甚至三倍的花冠。这实实在在是为了天主而行这样的人。不要告诉我，那人在这一方面或那一方面亏待了你；因为即使他对你表现出了人类所有的邪恶，即便如此，天主也命令你宽恕他的一切！\par

12. 看哪！我预先警告、作证，并以众人能听见的声音宣告：\Quote{凡有仇敌的人，不可走近神圣的祭台，或领受主的身体！凡走近的人不可有仇敌！你有仇敌吗？不要走近！你要走近吗？先和好，然后走近，触摸圣物！} 这并非我自己的宣告，而是为我们被钉在十字架上的主亲自说的。祂为了使你与圣父和好，不惜被牺牲，倾流自己的血！而你却不愿说一句话，或主动迈出一步，好与你的同仆和好吗？请听主论及怀着这种心态的人所说的话：\BibleQuote{若你献礼物在祭台上，在那里想起你的弟兄有什么怨你的事}——祂没有说\Quote{等他来就你}，也没有说\Quote{托别人作中间人}，更没有说\Quote{求别人}，而是说\Quote{你自己主动去就他。} 因为劝勉是：\BibleQuote{先去，与你的弟兄和好。} \BibleRef{玛 5:23} 哦，不可思议的奇事！祂自己并不以为不献礼物是耻辱，而你却认为主动和好是丢脸吗？请问，这样的人还能被算为值得赦免吗？如果你看见自己身上有一肢体被砍断，你岂不竭尽全力使它重新接在身体上？对你的弟兄也要这样做；当你看见他们因友谊断绝而被割离时，要赶紧挽回他们！不要等他们先迈出一步，而是奋勇向前，以便率先获得奖赏。\par

13. 我们奉命只有一个敌人，就是魔鬼。你永远不要与他和好！但对弟兄，心里决不可怀敌意。若有什么心胸狭隘，只应是短暂之事，决不可超过一天。因为，\BibleQuote{不要让太阳，} 他说，\BibleQuote{落在你们的愤怒上。} \BibleRef{弗 4:26} 因为如果你在黄昏前和好，你将从天主那里获得一些宽恕。但如果你继续怀有敌意更久，那敌意就不再是突然被愤怒和怨恨冲昏头脑的结果，而是出于邪恶、污秽的精神，而且成为惯于行恶的习性！而且这不仅是可怕之事，即你剥夺了自己的宽恕，而且正确的道路变得更加困难。因为过了一天，羞耻感更大；到了第二天，它又增加；若到了第三天、第四天，就会加上第五天。如此五天变成十天，十天变成二十天，二十天变成一百天；从此伤口将变得不可医治，因为随着时间推移，裂痕越来越宽。但是，人啊，你绝不可屈服于这些非理性的情感；不要羞愧，不要脸红，也不要在心里说：\Quote{不久以前我们还互相辱骂，说了许多该说或不该说的话；难道现在我要立刻赶去和好吗？那样谁会不责备我的过分轻易呢？} 我回答：有理智的人不会责备你的轻易；但当你坚持不和解时，那么所有人都会嘲笑你。那时你便给魔鬼占了这大裂痕的便宜。因为那时敌意更难消除，不仅由于时间的流逝，也由于在此期间发生的情况。因为正如\BibleQuote{爱德遮盖许多罪过，} \BibleRef{伯前 4:8} 同样，敌意使不存在的罪过显得真实，从此所有成为控告者的人都值得信任，他们幸灾乐祸，并公开宣扬他人行为中可耻的事。\par

14. 既然知道这一切，就要主动向弟兄迈出第一步；在他完全远离你之前抓住他；必要时，当天跑遍全城；必要时，出城或长途跋涉；放下手头一切其他事务，只专注于和好你弟兄这一件事。因为如果这工作辛苦，就反思你是为天主而承受这一切，你必得到充分的安慰。当你的灵魂畏缩、退却、害羞、惭愧时，也要激励它，不断重复这个主题说：你为什么拖延？你为什么退缩、犹豫？我们所关心的不是金钱，也不是其他任何短暂之物，而是我们的救恩。天主命令我们做这一切，所有事情都应次于祂的命令。这件事是一种属灵的买卖。我们不可忽视它，不可懒惰。也让我们的敌人明白，我们费了许多心力，是为了做悦乐天主的事。即使他再次侮辱、打击我们，或做其他更恶劣的事，让我们勇敢地承受一切，因为我们这样做与其说是造福他，不如说是造福我们自己。在所有的善工中，这一件在那一日将最特别地帮助我们。我们在众多大事上犯罪冒犯，激怒了我们的主。藉着祂的慈爱，祂给了我们这修和之道。那么，我们不可辜负这美好的宝藏。难道祂没有能力简单地命令我们和好，而不赋予任何赏报吗？祂还需要向谁质询或纠正祂的安排吗？然而，因着祂的极大慈爱，祂应许了我们一个伟大而不可言喻的赏报，是我们必须特别渴望获得的，即我们罪过的赦免；这样也使我们的服从更容易实行。\par

15. 那么，我们还能得到什么宽恕呢？如果我们可能领受如此大的赏报却仍不服从立法者，反而坚持我们的轻蔑；因为这一点是轻蔑，从此处显而易见。如果皇帝制定法律，所有互为敌人者必须彼此和好，否则斩首，难道我们每个人不都会急忙与邻人和好吗？是的！真的，我想是的！那么，我们有什么藉口，不将那应给予我们同仆的同样尊敬归给主呢？为此，我们奉命说：\BibleQuote{宽免我们的罪债，如同我们也宽免得罪我们的人。} \BibleRef{玛 6:12} 还有什么比这条诫命更温和、更仁慈的呢？祂使你成为自己过犯宽恕的法官！如果你宽恕的少，祂也宽恕你少！如果你宽恕的多，祂也宽恕你多！如果你从心里真诚地宽恕，天主也同样宽恕你！如果你在宽恕他之外，还把他当作朋友，天主也会这样对待你；因此，他罪的越多，我们就越需要赶紧和好，因为这成为我们更大过犯被宽恕的原因。你愿意知道，如果我们记恨伤害，就没有宽恕给我们，也没有人能拯救我们吗？我举一个例子来清楚说明我的断言。假设一个邻人伤害了你，夺取了你的财物；没收或侵吞了它们；不仅如此，我还要加上更多更糟的事，随你想象；他渴望消灭你；他使你面临千种危险；他向你表现出各种恶意；没有留下任何人性的邪恶所能做的事？为免逐一列举，假设他伤害你的程度空前——那么，即使在这种情况下，如果你心怀怨恨，你也不配得到宽恕。我来解释为何如此。\par

16. 假如你的一个仆人欠你一百块金币，而另一个人又欠这个仆人几块银币；假如欠仆人债的那个人来恳求你，求你施恩，你就叫来自己的仆人，嘱咐他说：\Quote{宽免了这个人的债，你欠我的数目中，我扣除这笔债。}之后，那仆人竟邪恶无耻到抓住他的欠债人，那时还有谁能把他从你手中救出来？你不是肯定会狠狠地打他一千鞭，因他把你侮辱到了极点吗？这也很公正。天主也要这样做：祂在那一天要对你说：邪恶而卑鄙的仆人，难道你是用自己的东西赦免了他吗？你从我这里欠了债，我命令你向他算清。因为祂说：\Quote{宽免，我就宽免你！}虽然说实话，即使我没有加这个条件，你也应因你主的缘故而宽免。但在此事上，我不是作为主人命令你，而是作为朋友请求你；我用自己的财产请求你，并且应许给你更大的回报；即使如此，你也没有变得更好。此外，人这样行时，他们就把自己仆人的账记下，如同债务的数目。例如，仆人欠主人一百块金币，而欠仆人的债主欠十块，如果后者宽免了债，主人并不宽免他一百块，只宽免这十块，其余的他仍然要讨还。但天主并非如此；如果你宽免了同仆的几件小事，祂就宽免了你所有的债。\par

17. 这从哪里显明呢？从祈祷本身。祂说：\Quote{因为如果你们宽免人的过犯，你们的天父也必宽免你们的过犯。}\BibleRef{玛 6:14} 正如\Quote{一百个银币}和\Quote{一万个塔冷通}之间的差别，这边和那边的债务差别也是如此之大！\par

那么，一个人本要接受一万塔冷通，却只为一百个银币，竟连这区区小数也不肯宽免，反而献上反对自己的祈祷，该受什么惩罚呢？因为当你念\Quote{求你宽免我们，如同我们宽免别人}，过后却不宽免时，你向天主祈求的无非是祂彻底剥夺你一切推诿或宽免。有人说：\Quote{我不敢说\InnerQuote{求你宽免我，如同我宽免别人}，只说\InnerQuote{求你宽免我}。}但这有什么关系？因为即使你自己不说，天主仍照样行事：正如你宽免，祂就宽免。祂从下文已清楚地表明这一点，因为那里说：\Quote{如果你们不宽免别人，你们的天父也必不宽免你们。}所以，不要以为不念全句就是敬虔的谨慎；也不要半截儿地献上祈祷，而要照祂吩咐的那样祈祷，好使那话语的义务日日使你畏惧，迫使你对邻人施行宽免。\par

18. 不要对我说：\Quote{我已多次恳求他，我祈求，我哀求，但未能和好。}你切勿停止，直到你与他修好。因为祂没有说：\Quote{留下你的礼物，你走吧。}去求你兄弟。而是：\Quote{你去。先与他和好。}\BibleRef{玛 5:24} 所以，即使你曾多次恳求，若未说服他，你仍不可停止。天主每天恳求我们，我们却不听，但祂仍不停止恳求。你竟不屑去恳求你的同仆吗？那么你怎能得救呢？假设你已多次请求而遭拒绝，但为此你将得到更大的赏报。因为对方越争辩，你越坚持恳求，你的赏报就越增多。善工越艰难成就，和好所需劳力越大，对他人的审判也越大，对你的忍耐所得的胜利冠冕也越辉煌。我们不仅应称赞这一切，也应以行动效法它；绝不放弃这项工作，直到我们恢复起初的友谊。因为仅仅避免激怒仇敌、不伤害他、心里不怀恶意并不够，还必须使他对自己心存友善。因为我听见许多人说：\Quote{我没有敌意；我不恼怒；我也不与他有任何关系。}但这并非天主的命令——你与他无关，而是你与他大有关系。正因如此，他是你的\Quote{兄弟}。正因如此，祂没有说：你要宽免你的兄弟你对他所怀的怨恨。而是什么？\Quote{你去。先与他修好；}即使他\Quote{有什么对不住你}，你也切勿停止，直到你使那肢体在友好和睦中重新合一。然而，你为了获得一个有用的仆人，会数出金子，与许多商人交谈，常作长途旅行；告诉我，你岂不尽力而为，好把仇敌变成朋友吗？那么，你若如此忽略祂的法律，又如何能呼求天主呢？确实，拥有一个仆人不会给我们带来多大益处；但把仇敌变成朋友，却会使天主对我们仁慈眷顾，轻易宽免我们的罪，赢得人的称赞，并使我们一生得到极大的平安；因为没有什么比一个哪怕只有一个仇敌的人更不安全了。我们的世间声誉受损，因为那人向每个人说了我们无数坏话；我们的思想也沸腾，良心不安；我们暴露在不断的焦虑风暴之中。\par

19. 既然我们意识到这一切的真实性，就让我们使自己摆脱惩罚和报复；让我们藉着履行所有之前所说的，来显示我们对当前庆节的敬意；并且让我们自己也使他人享受到那些我们以为能因这庆节从皇帝那里获得的恩惠。因为我确实听见许多人说，皇帝出于对圣逾越节的尊敬，将与这城和好，并赦免它的一切过犯。那么，当我们不得不依靠他人来获得安全时，我们就提出庆节及其要求；但当我们被命令彼此和好时，我们却藐视这同一个庆节，毫不把它放在心上，这是多么荒谬啊！确实，没有人像那个人那样玷污这神圣的庆节，他在守节的同时却心怀恶意。或者更确切地说，我要说，这样的人不可能守节，即使他连续十天不进食。因为哪里有仇恨和争斗，哪里就没有斋戒或节庆。你不会敢用未洗净的手去触摸圣祭，无论需要多么紧迫。那么，就不要用未洗净的灵魂去接近！因为这比前一种情况更糟，并带来更重的惩罚。因为没有什么比持续存留在内的愤怒更能使灵魂充满污秽。温良的精神不会居住在存在恼怒或激情的地方；当一个人失去圣神时，他还有什么得救的希望？他又怎能正当地行走？所以，亲爱的，当你渴望向你的仇敌报复时，不要将自己头朝下摔下去，也不要使自己失去天主的护佑而孤独无依！因为事实上，即使这责任是艰难的，但由这悖逆行为所带来的惩罚之重大，也足以唤醒最懒惰和最麻木的人，并说服他们去忍受各种劳作。但现在我们的论证已表明这责任是极其容易的，只要我们愿意。\par

20. 那么，让我们不要疏忽我们的生命，而是要认真起来；做一切事，好使我们没有仇敌，并这样出现在圣桌前。因为天主所命令的一切，没有一样是困难的，只要我们留心——我再说一遍，没有一样。这一点从那些已经改过自新的人身上可以看出来。有多少人曾受发誓习惯的欺骗，以为这种习气的改正是极其困难的？然而，藉着天主的恩宠，当你稍加努力时，你在很大程度上就洗净了自己这一恶习。因此我恳求你们也放下其余的一切，并成为他人的导师。至于那些尚未做到的人，却向我们诉说自己以前发咒发誓的时间很长，并说多年根深蒂固的东西不可能在短期内拔除，我要这样回答：在天主所命令的任何一条规诫需要切实实行的地方，不需要长时间，不需要许多天，也不需要隔几年；只需要敬畏和灵魂的虔敬；那么我们就必定能在短时间内完成它。但恐怕你们以为我是在随意说这些，那就找一个你认为沉溺于起誓的人，一个发誓次数比说话次数还多的人；把这个人交给我十天，如果我在这几天内没有使他戒除全部习惯，就请对我下最严厉的判决吧。\par

21. 这些话并非虚夸，从已经发生的事将向你们显明。有什么比尼尼微人更愚蠢？有什么比他们更缺乏悟性？然而，这些野蛮、愚昧的人，从未听过任何人教导他们智慧，从未从他人领受过这样的规诫，当他们听见先知说：\BibleQuote{还有三天，尼尼微就要倾覆了}\BibleRef{纳 2:4}，就在三天之内，抛弃了他们全部的恶习。奸淫者成了贞洁的；蛮横者成了温良的；贪婪勒索者成了节制和仁慈的；懒惰者成了勤勉的。他们确实没有用补救的方式改正一、二、三、四种恶习，而是改造了全部的不义。但有人要说，从哪里看得出来呢？从先知的话中；因为那位曾是他们的控告者，并说过\BibleQuote{他们的邪恶的呼声上达于天}\BibleRef{纳 1:5}的人，他自己又作了相反的见证，说：\BibleQuote{天主看见人人都离开了自己的邪道。}他是怎样离开的呢？就如天主所知道的，而不是人判断的。既然如此，我们岂不羞愧，岂不脸红，如果结果是在三天之内野蛮人就抛弃了全部邪恶，而我们这些在这么多天里被敦促和教导的人，却连一个坏习惯也没有克服？此外，这些人以前曾走到了邪恶的极端；因为当你听见说：\BibleQuote{他们的邪恶的呼声上达于我面前}时，你只能理解为他们邪恶的过度。然而，在三天之内，他们竟能够转变到完全德行的状态。因为哪里有对天主的敬畏，哪里就不需要天数或时间间隔；同样，反之，缺少这种敬畏时，天数也无济于事。正如锈蚀的工具，用只用水擦洗的人，即使花费很长时间，也不能除去全部污垢；但将工具放入火炉的人，会立刻使它比新制造的还要光亮：同样，灵魂被罪恶的锈迹污染，如果它只是轻微地、漫不经心地洁净自己，并且每天悔改，也不会获得更多益处；但如果它仿佛把自己投入对天主敬畏的熔炉中，就会在极短的时间內净化一切。\par

22. 那么我们不要拖延到明天。因为我们\Quote{不知道明日将生出何事}；\BibleRef{箴 27:1} 也不要说：\Quote{我们将逐渐克服这习惯}；因为这个\Quote{逐渐}永远没有尽头。所以，我们要抛弃这个借口，说：\Quote{如果我们今天不改正发誓的习惯，我们就不会停止，直到我们改正为止，即使有一万件事逼迫我们；即使必须死、受罚或失去所有；我们不给魔鬼懈怠的机会，也不给拖延的借口。}如果天主看到你的灵魂火热，你的勤勉加快，那么祂自己也会帮助你改正！是的，我祈求并恳求你们，让我们认真起来，免得我们也听到对我们说的话：\Quote{尼尼微人要起来，审判这一代}；\BibleRef{路 11:32} 因为他们一听见就悔改了；而我们多次听闻却不悔改。他们在德行的每一部分都是精进的，而我们一无所成。他们听见自己的城将要倾覆就恐惧；而我们听见了地狱却不恐惧：他们是没有领受先知教导的人；我们却享有永恒教导和许多恩宠的益处。\par

23. 我现在对你们说这些话，并非为责备你们自己的罪，而是为了别人的缘故；因为我深知，你们（如我已指出的）已经成就了这条关于发誓的法律。但这并不足以保证我们的安全，除非我们通过教导改正别人，因为那领了一个塔冷通的人，虽然归还了托付给他的全部，却因没有使所托付的增值而受罚。所以，我们不要只看自己已从这罪中解脱；直到我们也使别人解脱，不要停止；让每个人向天主献上十个他所纠正的朋友：无论你有仆人、学徒，或没有仆人、学徒，你有朋友；你要改正他们。此外，不要这样回答我：\Quote{我们大部分时间已摒弃了誓言，很少陷入那罗网}；但连这稀少的过失也要除去。如果你丢失了一枚金币，难道你不四处寻找、查问，以便找到它吗？对待誓言也要这样。如果你发现自己被骗去一个誓言，就哭泣、哀叹，好像失去了全部家产。我再说我之前说过的话。把自己关在家里；与妻子、儿女和仆役一起练习和操练。首先对自己说：\Quote{在纠正我的灵魂之前，我绝不能插手私事或公事。}如果你这样教导你的儿子，他们也会依次教导他们的孩子，这样，这纪律一直延续到基督的圆满来临，将为那些从根本上着手的人带来全部大赏报。如果你的儿子学会了说\Quote{请相信我}，他就不敢去戏院、进酒馆或掷骰子消磨时间；因为那句话像缰绳一样放在他嘴上，会迫使他无论多么不情愿也感到羞耻和脸红。如果他出现在这些地方，会很快迫使他退却。假设有人嘲笑。你另一边为他们的过犯哭泣！许多人曾在诺厄建造方舟时嘲笑他；但洪水来时，他嘲笑他们；或者更确切地说，义人根本不曾嘲笑他们，而是哭泣哀悼！因此，当你看到有人嘲笑时，要想到那些现在咧着嘴笑的牙齿，有一天将要承受那最可怕的哀号和切齿，并且在那一天，当他们咬牙切齿时，会想起这同一的笑声！那时你也会想起这笑声！富人怎样嘲笑拉匝禄！但后来，当他看见他在亚巴郎怀中时，他除了哀叹自己别无他法！\par

24. 因此，要记住这一切，催促所有人迅速完成这条诫命。不要对我说你将逐渐去做；也不要拖延到明天，因为这个明天永远没有尽头。四十天已经过去了。如果神圣的复活节过去，我将不再宽恕任何人，也不再使用进一步的劝诫，而是使用命令的权威和不可轻蔑的严厉。因为这种来自习惯的借口毫无力量。为什么盗贼不能以习惯为由逃脱惩罚？为什么凶手和奸夫不能？所以我声明并警告所有人，如果当我私下与你们会面并加以检验时（我肯定会检验），我发现有人没有改正这恶习，我将对他们施加惩罚，命令他们被排除在神圣奥迹之外；并非要他们永远被排除，而是让他们改过自新，然后才进入，并以纯洁的良心享用圣桌；因为这才是分享共融！愿天主藉着领导我们的司祷者和所有圣人的祈祷，使我们改正所有这些和其他缺欠，获得天国，藉着我们的主耶稣基督的恩宠和慈爱，愿光荣、尊威和朝拜归于祂及圣父、圣神，从今世直到永远。阿们。\par

\SourceRecord{Translated by W.R.W. Stephens.；Nicene and Post-Nicene Fathers, First Series；Vol. 9.；Edited by Philip Schaff.；Buffalo, NY: Christian Literature Publishing Co.,；1889.；<http://www.newadvent.org/fathers/190120.htm>.}

% structure: p0001 source=h1 target=section reason=page-title

\section{雕像讲道 第二十一讲}

\Emph{论\Person{弗拉维安}主教之返回，以及皇帝与城市及与推翻雕像者之修好。}\par

1. 今日，我将以那一向在危险时节对你们开始讲话的习惯用语开头，与你们同声说：\Quote{愿天主受赞美，}祂恩赐我们今日得以极其喜乐欢欣地庆祝这神圣的庆节；并将首归还于身体，牧者归还于羊群，师傅归还于门徒，将军归还于士兵，大司祭归还于司祭！愿天主受赞美，\Quote{祂能成就一切，远远超过我们所求所想的！}\BibleRef{弗 3:20} 因为对我们而言，只要能摆脱迄今逼近的灾祸，就似乎已经够了；为此我们已献上所有恳求。但那位热爱人类的天主，祂的赐予总是以丰厚的恩惠超越我们的祈祷，祂已更早将我们的父带回来，远超我们所能预期。谁曾想到，在如此少的日子里，他就能前往，得到皇帝召见，使我们脱离灾难，又如此迅速地回到我们中间，以至能赶在圣逾越节之前，并与我们一起庆祝它？然而，这件如此出乎意料的事，业已实现！我们已迎回我们的父；我们因而享受更大的喜乐，因为我们如今出乎意料地迎回了他。为这一切，让我们感谢仁慈的天主，惊叹那为城市所彰显的能力、慈爱、智慧与温柔关怀。魔鬼曾企图通过所犯的狂妄罪行完全颠覆这座城市；但天主却借这同一场灾难，装饰了城市、司铎和皇帝，并使他们都更加显赫。\par

2. 城市赢得了声誉，因为当如此大的危险击打她时，她立刻越过所有当权者、那些拥有极大财富者、那些在皇帝面前有巨大影响力的人，而逃往教会和天主的司铎作避难所，并怀着极大的信德，完全依靠那从上而来的希望！确实，在共同的父离开后，许多人预备恐吓那些被囚禁的人，说道：\Quote{皇帝并没有放下他的怒气，反而更加愤怒，正在考虑城市的彻底毁灭。}但当他们低声说这一切及更多时，那些被捆绑者却毫不畏惧；当我们说：\Quote{这些都是虚假的，是魔鬼的诡计，牠想使你们惊慌失措；}他们回答我们说：\Quote{我们不需要对我们的劝慰；因为我们知道我们从一开始就逃往哪里，以及我们将希望寄托于什么。我们已将安全系于神圣的锚上！我们没有将此托付于人，而是托付于全能的天主；因此我们十分确信，结果必会顺利；因为这希望永远不可能蒙羞，确实不可能！}这为我们的城市相当于多少冠冕，多少赞颂？这也将在我们其他事务中为我们赢得多少天主的垂顾！因为那警戒诱惑的攻击、仰望天主、藐视一切属人的事物、渴慕神圣援助的，确实不是平庸灵魂所能做之事。\par

3. 城市就这样赢得了声誉；而司铎也同样不亚于城市，因为他为众人舍命；虽然有许多障碍，如冬天、他的年龄、庆节，以及同样不小的，他的姊妹正处在最后一息，但他却超越所有这些障碍，并没有对自己说：\Quote{这是什么事？我们唯一的姊妹，她曾与我同负基督之轭，并作我的家室伴侣这样久，如今正在最后一息；我们岂能撇下她，离开此地，不看她咽气，听她临终的话语？但她确实曾每日祈祷，要我合上她的眼睛，闭合并整理她的口，并照顾其他安葬诸事；但现在，如同被遗弃、失去保护者的人，她从她兄弟那里得不到任何这些照料；而正是她特别希望从他那里得到这些；当她离世时，她将看不到她最爱、超过所有人、希望与之同在的那个人？这难道不比死多次更沉重吗？是的，即使我远在他方，难道不该急速前来，做一切事、受一切苦，只为向她表示这份爱心吗？如今我就在近旁，难道我要离开她，启程抛弃她？她将如何度过剩余的日子？}\par

4. 然而，他非但没有说出这些事中的任何一件，甚至没有想过它们；但他将敬畏天主置于一切亲情之上，认识到这样一个事实：正如风暴显出舵手，危险显出将军，同样，考验的时刻也使司铎显明出来。\Quote{所有的人，}他说，\Quote{都急切地望着我们；犹太人和希腊人一样；不要辜负他们对我们的期望；不要忽视如此巨大的沉船；但将我们自身的一切托付给天主，我们也当冒生命的危险！}再者，请思考司铎的宽宏大量和天主的慈爱！所有他忽视的事物，他都享有了；以致他既能收到他随时准备好的赏报，又能因出乎意料地享受它们而获得更大的愉悦！他宁愿为了城市的安全而在异乡、远离自己的人民庆祝逾越节。但天主在逾越节之前就把他归还给了我们，好使我们共同参与逾越节的进行；使他既收获他选择的赏报，又享受更大的喜乐！他不怕一年的季节；在他旅行的整个期间都是夏天。他不考虑自己的年龄；他轻松地完成了这长途旅行，好像他年轻而活泼！他不考虑他妹妹的去世，也没有因此而软弱；当他返回时，他发现她仍然活着，所有他忽视的事物，他都得到了！\par

5. 这样，司铎确实赢得了天主和人的声誉！这事也以超越王冠的光彩装饰了皇帝！首先，因为当时显明，他会赐给司铎们他不会赐给任何其他人的东西；其次，他毫不延迟地赐予了恩惠，平息了他的怨恨。但为使你们更清楚地理解皇帝的宽宏大量和司铎的智慧，以及比这两者更重要的是天主的慈爱；请允许我向你们讲述所进行的会谈的一些细节。但我现在要讲述的，是从宫内的一人那里听来的；因为司铎关于这事告诉我们不多也不少；但他总是效法保禄的宽宏大量，隐藏自己的善行；对那些从四面八方问他向皇帝说了什么、如何说服了他、如何完全平息了他的愤怒的人，他回答说：\Quote{我们在这事上毫无贡献，但皇帝本人（天主软化了其心），甚至在我们开口之前，就消除了他的愤怒，平息了他的怨恨；并谈论所发生的事件，好像受了侮辱的是另一个人，他这样不愤怒地回顾了所发生的一切。}但那些他因谦卑而隐藏的事，天主使之显明了。\par

6. 这些是什么呢？让我从故事再往前一点讲述给你们听。当他离开城市，使众人处于如此巨大的沮丧中时，他忍受了比我们这些身处灾难中的人自己所受的更为痛苦的事。因为，首先，他在旅途中遇到了皇帝派来调查所发生事件的使者；从他们那里得知了他们的使命；并想到了城市将要面临的可怕事件——骚乱、混乱、逃亡、恐惧、痛苦、危险——他泪流如涌，内心因怜悯而撕裂；因为父亲们通常在他们无法陪伴受苦的儿女时更为悲伤；这正是这位最心肠柔软的人现在所忍受的；他不仅哀叹将要降临到我们身上的灾难，而且哀叹我们在忍受这些灾难时他远离我们。然而，这却是为了我们的安全。因为当他从他们那里得知这些事情后，他的泪泉涌流得更热烈，他以更炽热的祈祷转向天主；他整夜不眠，恳求天主支援正在忍受这些的城市，并使皇帝的心更温和。而当他来到那座大城，进入皇宫后，他远远地站在皇帝面前——无言——哭泣——垂目——掩面，好像他自己就是所有祸事的作俑者；他这样做，首先希望以他的姿态、表情和泪水打动皇帝，使他怜悯；然后开始为我们辩护；因为对犯罪者来说，唯一的宽恕希望就是保持沉默，不为所行之事作任何辩护。因为他希望一种情绪被消除，另一种情绪取而代之；愤怒被驱逐，悲伤被引入，以便他这样为他的辩护之言铺平道路；这确实发生了。正如梅瑟在百姓犯罪后上到山上，自己无言地站着，直到天主叫他，说：\Quote{你且由我，让我消灭他们；} \BibleRef{出 32:10} 他也这样做了：因此，皇帝当他看到他流泪、弯腰到地上时，自己走近；他从他对他说的话中显明了他看到司铎眼泪时的真实感受；因为那些话不是一个被激怒或被煽动的人的话，而是一个悲伤的人的话；不是愤怒的人的话，而是相当沮丧、被极度痛苦所迫的人的话。\par

7. 你若听见他说了什么话，便会明白此事属实。他并没有说：\Quote{这是什么意思？你们是来为那些不敬虔、可憎、甚至不配活着的人，为那些叛徒、革命者、罪该万死的人充当使节吗？}然而，他摒弃了此类言语，为自己作了一番充满敬意与尊严的辩护；他列举了在其整个统治期间对这座城市所赐的恩惠，并在每一件之后说：\Quote{难道我应得这样的回报吗？我做了什么错事，要遭到这样的报复？他们有什么大小抱怨，非但侮辱我，还要侮辱已故之人？难道向活人泄愤还不够吗？他们却以为，若不侮辱那些已入坟墓的人，便不算做了大事。就算我真的亏待了他们——如他们所想的那样——也理应对死者宽容，他们并未伤害过人；因为死者不能与他们有同样的怨言。难道我不是一直将此城置于一切之上，视其比我的故乡更为宝贵吗？难道我不是不断祈祷要访问此城，并向众人起誓吗？}\par

8. 司铎闻此，悲泣不已，热泪更为涌流，不再缄默：因他见皇帝的辩护使我们的罪过显得更为重大；但他从心底发出深沉悲痛的叹息，说道：\Quote{我们必须承认，皇帝啊，您对我们祖国所显示的这份爱！我们不能否认！正因如此，我们尤其哀伤，一座如此被爱的城市竟被\OriginalTerm{魔鬼}{demons}所迷惑；我们竟显得对恩人忘恩负义，激怒了她热心的爱人。即使您要推翻、焚烧、处死，或做任何其他事，也绝不能向我们施加我们应得的报复。我们自己已预先加给自己比千死还要痛苦的事！还有什么比被人发现我们无端激怒了我们的恩人，一位如此爱我们的人，而全世界都知道，并谴责我们犯下最可怕的忘恩负义，更为苦涩的呢？假如有蛮族入侵我们的城市，夷平城墙，焚烧房屋，将我们掳走为奴，祸患尚且较轻。为何如此？因为您尚在，并对我们保持如此慷慨的仁慈，我们仍有希望恢复原状，重获更辉煌的自由。但如今，失去了您的恩宠，熄灭了您的爱情——这爱情比任何城墙都更可靠——我们还能投靠谁？我们还能仰望何处？我们激怒了如此仁慈的主上，如此宽容的父亲！因此，他们似乎犯了最不可容忍的罪过，同时却遭受了最可怕的灾祸：不敢正视任何人，不能以自由的目光仰望太阳；羞耻压着他们的眼睑，迫使他们低头。失去了自信，他们比任何俘虏更为悲惨，蒙受最大的羞辱；一想到自己罪孽之重，狂傲之极，他们几乎无法呼吸；因为他们从全世界居民那里招来的斥责，比从那位被侮辱者那里招来的更为严厉。}\par

9. \Quote{但是，皇帝啊，如果您愿意，这伤口仍有药可医，这些灾祸——虽如此巨大——仍有补救之方！的确，在私人之间，常有这样一种情形：重大而不可容忍的冒犯反而成为深厚感情的基础。我们人类的情况也是如此。因为当天主造了人，将他安置在\OriginalTerm{乐园}{Paradise}中，并赐予他极大的尊荣时；\OriginalTerm{魔鬼}{devil}无法忍受他如此大的福乐，便嫉妒他，将他从所赐的尊严中驱逐出去。但天主远没有抛弃他，反而为我们开启了天堂来代替乐园；如此，既显示了祂自己的慈爱，也更严厉地惩罚了魔鬼。现在，您也当如此行事！那些\OriginalTerm{魔鬼}{demons}最近尽了一切努力，要从您的恩宠中撕裂那座您最心爱的城市。既知如此，您要什么惩罚都可以，但请不要让我们成为您以前爱情的弃儿！不仅如此，虽然说来奇怪，但我必须说：请向我们显示比以往更大的仁慈；再次将此城的名字写在你爱情的首位——如果您真的渴望报复那些煽动这些罪行的\OriginalTerm{魔鬼}{demons}！因为如果您拆毁、颠覆、夷平此城，您所做的正是他们长久以来所渴望的。但如果您息怒，并重新宣告您爱她如同从前，您就给了他们致命的一击。您对他们施予了最完美的报复，不仅表明他们的恶谋一无所成，而且结果完全与他们的愿望相反。您这样做是公正的，也是对一个城市施予怜悯，这城市因您的爱情而被魔鬼嫉妒；因为如果您不是如此极端地爱她，他们不会嫉妒她到这种程度！因此，即使我所说的似乎离奇，却仍然是事实：此城所遭受的一切，都是由于您和您的爱情！有什么焚烧，有什么毁灭，比您在辩护中所说的那些话更为苦涩呢？}\par

10. 你现在说，你受到了侮辱，遭受了连任何一位皇帝都未曾承受过的冤屈。但是，最仁慈、最智慧、最虔诚的君主，只要你愿意，这轻视将为你赢得一顶比你戴的皇冠更尊贵、更辉煌的冠冕！因为这皇冠是你君主德行的展现，但它也是赐予者慷慨的标记；而这由你的仁爱编织的冠冕将完全是你自己的善工，出于你自己的爱智之心；众人将因这些宝石而少钦佩你，却会因你胜过这愤怒而更多赞扬你。你的雕像被推倒了吗？你有能力重新竖立起更辉煌的雕像。因为如果你赦免那些伤害你的人的过犯，不报复他们，他们将为你竖立一座雕像，不是广场上的铜像、金像或镶嵌宝石的雕像，而是披着比任何材质更珍贵的袍子——即仁爱与慈悲的袍子！每个人都会在自己的灵魂中为你竖立雕像；你将拥有与现在及将来居住在全世界的所有人同样多的雕像！因为不仅我们，而且所有后来的人及他们的继承者，都将听到这些事，并会爱慕你，如同他们亲身经历了这恩惠一般！\par

11. 为了证明我这样说并非奉承，而是必定如此，我将引用一段古代历史，让你明白，使君主如此辉煌的，既不是军队、战争武器、金钱、众多臣民，也不是任何其他此类事物，而是灵魂的智慧与温和。据说，有福的君士坦丁有一次，当他的一座雕像被石头投掷时，许多人怂恿他对肇事者采取行动，说他们用石头砸毁了他的整个面容。他用手抚摸自己的脸，温和地微笑着说：\Quote{我完全看不出我的脸上有任何伤痕。头看起来完好，脸也完全完好。}于是这些人羞愧难当，放弃了他们不义的计谋。\par

这句话直至今日，人人传颂；时间的流逝既未削弱也未抹去对如此崇高智慧的记忆。这样的行为比任何战争胜利的光荣更加显赫！他建立了许多宏伟的称号，征服了许多野蛮部落；但我们如今一个也不记得；而这句话却被反复传颂，直至今日；我们的后代以及他们之后的人，都将听到。这并非唯一可赞叹的事：他们将听到它；而且人们提及它时，带着赞许和掌声；听到它的人也以同样的心情接受；没有人听了之后能保持沉默，而是每个人都立刻欢呼，赞扬说出这话的人，并祈求即使他已去世，也赐予他无数的祝福。如果在人中间，这句话为他赢得了如此多的荣誉，那么在慈悲的天主面前，他将获得多少冠冕呢！\par

12. 我为何需要提及君士坦丁和其他人的榜样呢？我本应用更贴近你自身、来自你值得称赞的行为的考量来劝诫你。你还记得，不久前，当这节日临近时，你向全世界发出一封诏书，命令释放监狱中的囚犯，赦免他们的罪行。仿佛这还不足以证明你的慷慨，你在信中说：\Quote{但愿我能召回并恢复死者，使他们回到先前的生命状态！}现在请记住这些话。请看，这就是召回和恢复死者、使他们重获生命的时节！因为这些人确实已经死了，甚至在判决宣布之前；这座城市如今已在哈德斯的门口搭起了帐篷！因此，请使它重新站立起来，这你可以做到，无需金钱、费用、时间或劳力的损失！你只需开口说话，就能使如今处于黑暗中的城市复活。现在，请让它从此因你的仁爱而得名；因为它所受于你判决的恩惠，将不亚于最初建立它的人的恩惠。这是极为合理的；因为那人只给了它开始，便离去了；而你，当它成长壮大，并在如此繁荣之后跌倒时，将成为它的重建者。如果城市被敌人攻陷、被蛮族蹂躏，你从危险中解救它，那没有什么神奇；而你现在赦免它，则更为神奇。许多皇帝经常做前一件事；但若你独自完成后一件事，你将是在那件事上的第一人，而且超乎所有期望。前一种善行，即保护臣民，一点也不神奇或非凡；它是经常发生的事件之一；但后一种，即忍受如此挑衅后消除愤怒，是超越人性的事。\par

13. 请思考，你现在所考虑的事不仅关乎这座城市，而且关乎你自己的荣耀；或者更准确地说，它关乎整个基督教事业。如今，外邦人、犹太人、整个帝国以及蛮族（因为后者也听说了这些事件）都急切地注视着你，等待着看你将对这些事情作出什么判决。如果你颁布一项仁慈而宽恕的判决，所有人都会赞扬这一决定，光荣天主，并彼此说：\Quote{天哪！基督信仰的力量何其伟大！它约束并抑制了一个在地球上无人能比的人、一个足以摧毁和蹂躏万物的强大君主，并教导他实践一种连一个平民都未必能表现出来的哲学！基督徒的天主必定是伟大的，他从人中造出天使，并使他们超越我们本性的一切强制力！}\par

14. 诚然，你绝不应心存那无谓的恐惧；也不应容忍那些说\Quote{其他城市会变得更糟、更加藐视权威，如果这座城市不受惩罚} 的人。因为，如果你无法报复；而他们在做了这些事之后，又强行向你挑战；双方的力量势均力敌；那么，产生这样的猜疑倒是合情合理的。但是，如果他们惊恐万分、半死不活，通过我跑来俯伏在你脚下；每天所期待的除了屠杀的深渊之外别无他物，并共同从事恳求；仰天呼求天主援助，并惠顾我们这个使团；每个人都像临终一样安排了自己的私事；那么，这样的恐惧怎能不是多余的？即使他们被判处死刑，所受的痛苦也未必有现在这样多，他们已活在恐惧战栗中这么多天；每当傍晚来临，不指望看到早晨；每当白昼到来，不希望能挨到黄昏！许多人在穿越荒野、前往人迹罕至之处时，还遭遇了野兽；不仅是男人，还有小孩和妇女；出身自由、家境良好；许多天许多夜躲藏在山洞、沟壑和荒野的洞穴中！这座城市确实遭受了一种新的俘虏方式。当房屋和城墙屹立时，他们所遭受的灾难比城市被烧毁时更重！虽然没有野蛮的敌人出现，没有敌人露面，他们的处境比被攻陷更为悲惨；每天一片树叶的沙沙声就足以使他们惊恐万分！这些事情是众所周知的；所以，即使所有人看到这座城市夷为平地，也不会像听到现在降临在它身上的灾难那样，学到如此节制的教训。因此，不要以为其他城市将来会变得更糟！即使你颠覆了其他城市，也未必能如此有效地纠正它们，像现在这样，通过悬而未决的命运，比任何惩罚都更严厉地惩戒了它们！\par

15. 那么，不要再把这灾祸推得更远；而是容许他们从此喘口气。因为惩罚有罪之人，并为这些行为索取刑罚，对任何人都是轻而易举和公开的；但是宽恕那侮辱了你的人，并饶恕那犯下不可饶恕之罪的人，却只有少数一两个人才能做到；尤其是当受辱的人是皇帝时，更是如此。使城市处于恐惧的统治之下是容易的；但要使人人都成为充满爱心的臣民，并说服他们对你政府保持好感，不仅为你的帝国献上共同的祈祷，还有个人的祈祷，这是一项艰难的工作。一位君主可以耗尽他的财富，或调动无数的军队，或随意做其他事，但他仍然无法像现在这样轻而易举地赢得如此多人的爱戴。因为那些受到善待的人，以及那些听闻此事的人，都会像受惠者一样对你有好感。为了在短时间内赢得整个世界，并能够说服所有现在活着的人以及未来的世代，像为自己的儿女祈求一样，为你祈求同样的祝福，你会付出多少金钱、多少辛劳！而如果你从人那里得到这样的回报，从天主那里你将得到何等更大的回报！这不仅是从正在发生的事件中，也是从将来他人所行的善事中。因为如果将来发生类似现在的事件，（但愿天主不允许！）而任何受辱的人当时正在考虑追究闹事者的责任，那么你的温和与道德智慧将代替一切其他教导和劝诫，成为他们的榜样；他们将因有这样的智慧典范而脸红羞愧，不甘落后。这样，你就成为所有后代的导师；即使他们达到了道德智慧的最高点，你仍将在他们中间获得棕榈枝！因为一个人自己首先树立这样温和的榜样，与通过观察他人来效法他们所行的善事，是不同的。因此，无论你的后人展现出怎样的仁爱或温和，你都将与他们一同获得赏报；因为提供根的人必须被视为果实的源头。为此，现在没有人能与你分享你慷慨所带来的赏报，因为善行完全是属于你的。但所有将来出现的人（如果有这样的人出现）的赏报，你都将有份；你将能够与他们一同平等地分享善工的功劳，并像教师与学生一样获得一份。即使没有这样的人出现，赞美和掌声的贡品将世世代代为你积累。\par

16. 因为请想一想，后代听到这样的事意味着什么：当一座如此大的城市应受惩罚和报复时，当所有人都惊恐时，当它的将军、地方官和法官都惊慌失措，不敢替可怜的人民说一句话时；一位单独的老人，拥有天主的司祭职，来到并仅凭他的神情和交往就打动了君主的心；而他未曾赐予任何其他臣民的恩惠，却因对天主的法律的敬畏而赐给了这位老人！因为，皇帝啊，正是这件事，即我被派来担任这个使团，城市给了你不小的荣誉；因为他们对你做出了最好、最可敬的评断，即你尊重天主的司祭，无论他们多么卑微，胜过你权下任何职位！\par

17. 但我现今前来，不仅是从这些人那里，而是从那位天使与人类的共同主宰那里，向您最仁慈、最温和的灵魂说这番话：\BibleQuote{如果你们宽恕别人的过犯，你们的天父也必宽恕你们的过犯。}\BibleRef{玛 6:12} 请记住那日，我们将为我们的一切行为交账！请思考，如果你在任何方面犯了罪，你便能通过这一判决和这一决定，毫不费力、不劳苦地洗刷所有过犯。有些人出使时，携带黄金、白银及其他此类礼物。但我带着神圣的法律来到你的御前；我献上这些，代替所有其他礼物；我劝你效法你的主，祂每日受我们侮辱，却仍不断向众人施恩！请不要辜负我们的希望，也不要使我们的承诺落空。我还希望你知道，如果你决心和解，恢复对这座城市从前的善意，平息这公义的愤怒，我将满怀信心地回去。但如果你决定弃绝这座城市，我将不仅永不返回，不再见到它的土地，而且今后将彻底不认它，并加入另一座城市成为其成员；因为天主决不允许我属于那个国家，而您，所有人的最温和、最仁慈者，却拒绝让它获得和平与和解！\par

18. 说了这些话，以及许多类似的话之后，他如此说服了皇帝，以至于发生了曾发生在\Person{若瑟}身上的事。因为正如\Person{若瑟}见到他的兄弟们时，渴望流泪，却克制了自己的情感，以免破坏他所扮演的角色；同样，皇帝在心里哭泣，却没有让人看见，为了在场的人。然而，他无法在会谈结束时掩饰这一情感，而是违背自己的意愿泄露了出来。因为这篇讲话结束后，不再需要任何言语，他只说出了一句情感，这比王冠更给他带来荣誉。那是什么？他说，我们这些不过是人的人，向那些冒犯我们的人息怒，这有什么奇妙的或伟大的呢？当宇宙的主宰，降临到世上，成为我们的仆人，被那些曾受过祂恩惠的人钉死，却为钉祂的人恳求圣父，说：\BibleQuote{宽恕他们吧，因为他们不知道自己在做什么？}\BibleRef{路 23:34} 那么我们宽恕我们的同仆，又有什么稀奇呢？这些话不是虚饰，随后的一切证明了这一点。尤其是我即将提到的那一特殊情况；因为这位司铎本应留在那里，与皇帝一同庆祝节日，他却催促他——尽管违背他的意愿——尽快、尽力赶去见他的同胞。\Quote{我知道，}他说，\Quote{他们的灵魂仍在动荡；还有许多灾难的遗存。去吧，给他们安慰！如果他们看到舵手，他们将不再记得已过的风暴；但所有这些悲伤事件的记忆都将被抹去！}当司铎急切地恳求他派遣自己的儿子时，他渴望提供最令人满意的证据，表明他已完全从灵魂中抹去一切愤怒的情感，回答说：\Quote{祈祷这些阻碍被清除，这些战争被终结；然后我必亲自前来。}\par

19. 还有什么比这样的灵魂更温和的呢？让外邦人从此羞愧吧；或者，与其羞愧，不如接受教导；离开他们与生俱来的错误，回归基督信仰的力量，从皇帝和司铎的榜样中学习我们的道理！因为我们最虔诚的皇帝并未止步于此；当主教离开城市，渡海之后，他也派遣了一些人，极其谨慎和努力，以免浪费时间，防止城市因主教在城墙外庆祝节日而失去一半的喜悦。哪里有如此慈父，会为了那些冒犯过他的人如此忙碌？但我必须提及另一个彰显这位义人之德的情况。当他完成此事后，他并没有像任何其他爱慕虚荣的人那样，竭力亲自递送那些将我们从沮丧中释放的信件；而是因为他的行程太慢，他认为最好另派一人代替他；一位擅长骑马的人，作为好消息的携带者前往城市；以免城中的悲伤因他迟迟未到而延长。因为他唯一热切渴望的，不是亲自带来这些充满喜悦的佳音，而是让我们的国家尽可能迅速地重新自由呼吸。\par

20. 所以，你们当时在装饰广场、挂灯、在商店前铺设绿叶床榻、并大事庆祝，仿佛城市刚刚诞生时所做的一切，如今你们要以另一种方式，在一切时间内去做：戴冠，不是鲜花，而是德行；点燃灵魂之光，来自善工；以属灵的喜乐欢欣。让我们永远不止息地为这一切感谢天主，不仅因为他使我们脱离这些灾祸，也因为他容许它们发生；让我们承认他丰厚的恩慈！因为藉着这两者，他装饰了我们的城市。如今这一切都按照先知的话：\BibleQuote{你们应将这事传报给你们的子孙，你们的子孙也应传报给他们的子孙，他们的子孙又应传报给下一代。}\BibleRef{岳 1:3} 好使将来所有的人，直到终结，获悉天主对这城市的慈爱行为，称我们有福，因获得了如此恩惠——并惊讶于我们的主，他竟在如此严重倾覆时重建了城市——他们自己也从中获益，因所发生的一切而被激发虔敬！因为最近发生在我们身上的历史，不仅有助于我们，若我们常记住它，也有助于后来的人。既然如此，让我们始终感谢爱人的天主，不仅因我们脱离了这些可怕的灾祸，也因它们被容许临到我们——从圣经以及最近降临我们的事件中学习：他永远以他那本有的慈爱，为我们的利益安排一切。愿天主恩赐我们不断享有这慈爱，从而获得天国，在耶稣基督我们的主内；愿光荣与权能归于他，至于无穷之世。阿们。\par

\SourceRecord{Translated by W.R.W. Stephens.；Nicene and Post-Nicene Fathers, First Series；Vol. 9.；Edited by Philip Schaff.；Buffalo, NY: Christian Literature Publishing Co.,；1889.；<http://www.newadvent.org/fathers/190121.htm>.}
